
\medterm{C And T Cells} B cells make antibodies, while T cells coordinate immune responses or destroy infected and abnormal cells. Both are types of lymphocytes, a group of white blood cells that protect against infection and disease.
\medterm{C Fiber} A small, slow nerve fiber without an insulating covering that carries dull pain, warmth, and itch from the body to the spinal cord and brain.
\medterm{C-Peptide} A small protein fragment released when the pancreas makes insulin. Measuring it helps show how much insulin a person's body is making, rather than how much injected insulin is present. Kidney function and blood-glucose levels can affect the result.

\medterm{C-Reactive Protein} A protein made by the liver that rises when there is inflammation or infection. It can help monitor inflammatory diseases and possible infection, but it does not identify the cause by itself. A high-sensitivity CRP test may be used with other information to assess heart-disease risk.

\medterm{C-Reactive Protein Test} A blood test that measures C-reactive protein. The level may rise with inflammation or infection, but the test cannot identify the cause by itself and must be interpreted with symptoms and other findings.

\medterm{C-Section} A cesarean birth in which a baby is delivered through surgical openings in the abdomen and uterus. It may be planned or performed urgently when vaginal birth could endanger the mother or baby.
\medterm{C. Diff} Alternate index label; see \textit{C. difficile infection}.

\medterm{C. Difficile} Clostridioides difficile is a spore-forming bacterium that can multiply in the bowel, especially after antibiotics disrupt normal bacteria. Its toxins cause watery diarrhea and colon inflammation, which can become severe or recur.

\medterm{C. Difficile Colitis} Alternate index label; see \textit{Pseudomembranous colitis}.

\medterm{C. Difficile Infection} An infection caused by Clostridioides difficile, a bacterium that can grow excessively in the bowel after antibiotic use. It commonly causes watery diarrhea and inflammation of the colon and may lead to dehydration, severe colitis, or recurrence.

\synonyms
C. Diff

\medterm{C1Q Nephropathy} A kidney disorder in which the immune protein C1q is deposited in the glomeruli, the kidney's filtering units. It may cause protein or blood in the urine, swelling, high blood pressure, or reduced kidney function.

\medterm{C282Y Mutation} A change in the HFE gene that is the most common genetic cause of hereditary hemochromatosis, a condition in which the body stores too much iron. Having two copies of this mutation increases the risk, although not everyone with two copies develops the disease.

\medterm{CA 19-9} A substance measured in blood that may be increased in pancreatic and other digestive-system cancers. It can also rise with blocked bile ducts or inflammation, so it cannot screen for or diagnose cancer by itself.

\medterm{CA 27-29} A blood test that measures a protein sometimes released by breast-cancer cells. It may help monitor selected people with known breast cancer but cannot by itself screen for or diagnose cancer.

\medterm{CA-125} A blood protein that may rise in ovarian and some other cancers. It can also increase with menstruation, pregnancy, endometriosis, infection, liver disease, or other noncancerous conditions, so it cannot diagnose ovarian cancer alone.

\synonyms
CA125

\medterm{CA-125 Blood Test} See \textit{CA-125}.

\medterm{Cabazitaxel} A taxane chemotherapy medicine used with prednisone for metastatic castration-resistant prostate cancer after docetaxel. It can cause low white-cell counts, infection, diarrhea, and serious allergic reactions.

\medterm{Cabergoline} A dopamine D2-receptor agonist used to suppress prolactin secretion and treat hyperprolactinemia; it is also used in selected Parkinson disease settings and requires monitoring for adverse effects.

\medterm{CABG} Coronary artery bypass grafting, surgery that creates a new path around a narrowed or blocked heart artery to improve blood flow to the heart muscle.

\synonyms
Coronary artery bypass grafting; coronary bypass surgery.

\medterm{Cabot Ring} A threadlike or ring-shaped remnant seen inside a red blood cell on a blood smear. It may occur in severe anemia or when the bone marrow is producing abnormal red blood cells.

\medterm{Cachectic} Affected by cachexia, a severe wasting syndrome involving involuntary weight loss, muscle loss, weakness, and often reduced appetite in chronic disease.

\medterm{Cachetic} An alternate spelling of cachectic, describing severe unintentional loss of body weight and muscle caused by a serious long-term illness.

\medterm{Cachexia} A wasting syndrome caused by serious long-term illness, marked by ongoing unintentional loss of weight, muscle, and body fat. It may cause poor appetite, weakness, and inflammation and cannot be corrected by nutrition alone.

\medterm{Cacosmia} A distorted sense of smell in which ordinary odors seem unusually foul or unpleasant. It may occur with nose or sinus disease, infection, seizures, head injury, or another nervous-system disorder.

\medterm{CAD} Coronary artery disease, narrowing or blockage of the arteries supplying the heart, usually caused by fatty plaque. It can reduce oxygen to heart muscle and cause chest pain, heart attack, heart failure, or dangerous rhythms.

\medterm{CADASIL Syndrome} An inherited disorder of small arteries in the brain caused by changes in the NOTCH3 gene. It may cause migraine with aura, repeated small strokes, mood problems, and progressive memory or thinking difficulties.

\medterm{Cadaver} A dead human body used for anatomical study, surgical training, research, or identification. A donated cadaver may be preserved so tissues remain suitable for examination and teaching.
\medterm{Cadaveric Spasm} A rare form of muscle stiffening that occurs immediately at death in a specific group of voluntary muscles, without the usual period of relaxation. It may preserve the position held at the moment of death.

\medterm{Cadherin} A cell-surface protein that helps neighboring cells attach to one another and keeps tissues organized. Loss or alteration of cadherins can allow some cancer cells to invade nearby tissue and spread.

\medterm{Cadmium} A toxic metal found in tobacco smoke, some batteries and pigments, workplace dust, and contaminated food or soil. Long-term exposure can damage the kidneys, lungs, and bones and increases the risk of some cancers.

\medterm{Cadmium Chloride} A toxic compound containing cadmium, used mainly in industry and research. Inhalation or swallowing can injure the lungs, kidneys, and digestive tract; repeated exposure can damage bones and increase cancer risk.

\medterm{Caduceus} A winged staff with two entwined snakes, traditionally associated with Hermes. It is often used incorrectly as a medical symbol and differs from the single-snake Rod of Asclepius.

\medterm{Cadusafos} An organophosphate pesticide that blocks acetylcholinesterase, causing excess nerve signaling after exposure. Poisoning may produce sweating, salivation, vomiting, diarrhea, small pupils, muscle weakness, breathing difficulty, or seizures.

\medterm{Caecal} Relating to the caecum, the pouch at the beginning of the large intestine where the small intestine joins the colon and the appendix begins.

\medterm{Caecitis} Inflammation of the caecum, the first part of the large intestine. It may result from infection, inflammatory bowel disease, reduced blood flow, neutropenia, or spread of nearby inflammation and can cause right-sided abdominal pain and fever.

\medterm{Caecopexy} An operation that attaches the caecum, the first pouch of the large intestine, to nearby tissue to prevent or treat abnormal movement or twisting.

\medterm{Caecostomy} An operation that creates an opening from the caecum to the abdominal wall. It may temporarily release pressure, drain the bowel, or divert stool when the normal route is blocked or unsafe.

\medterm{Caecotomy} An operation that opens the caecum, usually to remove an object, access the bowel, or assist another abdominal procedure. The opening may be closed during the same operation.

\medterm{Caecum} The first pouch-like part of the large intestine, located where the ileum joins the colon and attached to the appendix.

\medterm{Caenorhabditis Elegans} A small, free-living roundworm widely used in laboratory research. Its simple nervous system, short life cycle, and well-studied genes help scientists investigate development, aging, nerve function, and disease mechanisms.

\medterm{Caesarean Birth} Delivery of a baby through surgical openings in the abdomen and uterus. It may be planned or urgent when problems such as fetal distress, labor complications, abnormal presentation, placenta previa, or multiple pregnancy make vaginal birth unsafe.

\medterm{Caesarean Section} A caesarean section is surgical delivery of a baby through incisions in the abdomen and uterus when vaginal birth is unsafe or not feasible.
\medterm{Cafe-Au-Lait Macules} Flat, evenly colored light-brown skin patches, named for their coffee-with-milk color. A few may be harmless; numerous or unusually large patches can be a sign of an inherited condition and should be assessed.

\medterm{Cafe-Au-Lait Spot} A flat, sharply outlined, evenly pigmented light-brown patch of skin. One or a few are often harmless, while many lesions or a characteristic pattern may indicate neurofibromatosis type 1 or another inherited disorder.

\medterm{Café Au Lait Spot} A flat, sharply demarcated, light-brown skin macule caused by increased melanin.
Multiple lesions may be a feature of neurofibromatosis type 1 or other genetic
disorders.

\medterm{Café Coronary} Sudden blockage of the airway by food, usually during a meal. Complete blockage prevents breathing and can rapidly cause collapse or death unless emergency choking first aid and medical treatment restore the airway.

\medterm{Café-Au-Lait Macules} Flat, light-brown skin patches with relatively uniform pigmentation. Multiple lesions may be associated with genetic disorders such as neurofibromatosis type 1.

\medterm{Caffeic Acid} A naturally occurring plant compound found in coffee, fruit, vegetables, and some herbs. It is studied in laboratory research for antioxidant and other biological effects but is not an established treatment for disease.

\medterm{Caffeine} A stimulant found in coffee, tea, chocolate, soft drinks, and energy drinks. It increases alertness and can raise heart rate, cause shakiness, or disturb sleep.

\medterm{Caffeine Intoxication} A toxic reaction to excessive caffeine, causing restlessness, anxiety, tremor, insomnia, nausea, vomiting, rapid heartbeat, or agitation. Severe poisoning can cause abnormal heart rhythms, seizures, or dangerous metabolic changes.

\medterm{Caffeine Overdose} A dangerous amount of caffeine can cause severe agitation, tremor, vomiting, rapid heartbeat, abnormal heart rhythms, low blood pressure, seizures, or disturbed body chemistry. Severe symptoms require urgent medical care.

\medterm{Caffeine Withdrawal} A group of symptoms after suddenly reducing or stopping regular caffeine use. Headache, tiredness, low mood, poor concentration, irritability, nausea, or flu-like feelings usually begin within a day and improve within several days.

\medterm{Caisson Disease} Decompression sickness caused by gas bubbles forming in tissues or blood after a rapid fall in surrounding pressure, such as after diving. It can cause joint pain, skin changes, breathing problems, brain or spinal-cord injury, or collapse.

\medterm{Cal (Calorie)} A unit of energy. In nutrition, a food Calorie is a kilocalorie, equal to 1,000 small calories, and describes the energy supplied by food or used by the body.

\medterm{Caladium Plant Poisoning} Poisoning from chewing or swallowing caladium, whose needle-like calcium oxalate crystals irritate the mouth and throat. Burning, pain, drooling, and swelling may occur; trouble breathing or swallowing is an emergency.

\medterm{Calamine} A skin preparation containing zinc compounds used to soothe itching and irritation from minor rashes, insect bites, and sunburn. It should not be applied to deep wounds, infected skin, or the eyes.
\medterm{Calbindin} A protein that binds calcium inside cells and helps control calcium levels and cell signaling, especially in nerve cells and intestinal cells. Laboratories may also use it as a marker to identify certain tissues.

\medterm{Calcaneal Spur} A bony growth from the heel bone, usually where the plantar fascia or Achilles tendon attaches. It may cause no symptoms; when heel pain occurs, the associated soft-tissue inflammation is often more important than the spur itself.

\medterm{Calcaneum} The calcaneum, or calcaneus, is the large bone forming the heel. It supports body weight, connects the ankle to the foot, and transmits force from the calf muscles through the Achilles tendon.
\medterm{Calcaneus} The largest bone in the ankle and foot, forming the heel. It supports body weight and connects with the ankle bone and cuboid; fractures commonly follow falls from height and may require imaging and specialist care.

\medterm{Calcifediol} The main circulating storage form of vitamin D, also called 25-hydroxyvitamin D. A blood test for calcifediol is used to assess vitamin-D status before the body converts it to the active hormone calcitriol.

\medterm{Calciferol} A fat-soluble form of vitamin D that helps the body absorb calcium and phosphate and maintain healthy bones, muscles, and immune function. The body changes it into active vitamin-D hormones.

\medterm{Calcific Bursitis} A condition in which calcium deposits form in or around a bursa, often near the shoulder, causing pain, tenderness, swelling, and restricted movement. Treatment depends on severity and may include rest, physical therapy, medicines, drainage, or deposit removal.

\synonyms
Bursitis Calcific (Calcific Bursitis)

\medterm{Calcific Tendinitis} A condition in which calcium crystals build up inside a tendon, most often a rotator-cuff tendon in the shoulder. It can cause sudden severe pain, stiffness, and restricted movement, sometimes followed by gradual improvement as deposits break down.

\medterm{Calcific Uremic Arteriolopathy} A rare, life-threatening disorder, usually in advanced kidney disease, in which calcium deposits and clotting block small blood vessels supplying the skin and fat. It causes extremely painful wounds, tissue death, and a high risk of infection.

\medterm{Calcification} Deposition of calcium salts in body tissue. It may be a normal part of bone formation, a response to injury or aging, or a sign of abnormal calcium or phosphate balance; its importance depends on the site and pattern.

\medterm{Calcified Granuloma} A small area of past inflammation that has healed and become hardened by calcium deposits. It is often found in the lung after an old infection and usually does not represent active disease.

\medterm{Calcineurin} An enzyme that helps activate T cells, a type of immune cell, after they receive signals. Medicines that block calcineurin reduce immune activity and are used to prevent transplant rejection and treat some immune diseases.

\medterm{Calcinos} An older or shortened term for calcinosis, in which calcium salts build up in the skin or tissues beneath it. Firm white or yellow lumps may form, sometimes ulcerate, and release chalky material.

\medterm{Calcinosis} Abnormal buildup of calcium salts in the skin or tissues beneath it. It may cause firm white or yellow lumps or plaques that are painful, limit movement, ulcerate, or release chalky material.

\medterm{Calciphylaxis} A rare, life-threatening disorder, usually associated with advanced kidney disease, in which small blood vessels supplying the skin become blocked by clotting and calcium deposits. It causes severe pain, skin wounds, tissue death, and serious infection risk.

\medterm{Calcitonin} A hormone made by specialized cells in the thyroid gland. It helps lower blood calcium by reducing calcium release from bone and increasing calcium loss through the kidneys, although it has a smaller role than parathyroid hormone in adults.

\medterm{Calcitonin Blood Test} A blood test that measures calcitonin, a hormone produced by thyroid C cells. It is used mainly to evaluate or monitor medullary thyroid carcinoma and may be affected by other thyroid, kidney, gastrointestinal, or medicine-related factors.

\medterm{Calcitriol} The active hormone made from vitamin D, mainly by the kidneys. It increases intestinal absorption of calcium and phosphate and helps regulate parathyroid hormone, bone formation, and blood-mineral levels.
\medterm{Calcitriol And Analogues} Medicines related to active vitamin D that reduce excessive parathyroid-hormone release and help manage mineral imbalance in some kidney disorders. They can raise calcium or phosphate, so blood levels require monitoring.

\medterm{Calcium} A mineral needed for strong bones and teeth, muscle movement, nerve signals, blood clotting, and normal cell function. Blood calcium exists in an active free form and forms attached to proteins or other substances.

\medterm{Calcium - Ionized} A blood test measuring ionized calcium, the biologically active fraction of circulating calcium. It is useful when total calcium may be misleading, such as with abnormal albumin, critical illness, or acid-base disturbances.

\medterm{Calcium - Urine} A urine test measuring calcium excretion over a specified collection period or in a spot specimen. It helps evaluate kidney stones, parathyroid disorders, bone disease, and abnormal calcium handling, and must be interpreted with kidney function and diet.


\medterm{Calcium Blood Test} A blood test that measures calcium, usually as total calcium. The result may be interpreted with albumin, a blood protein, or measured as ionized calcium when the active form is needed. It helps assess parathyroid, bone, and kidney disorders.

\medterm{Calcium Carbonate} A calcium supplement and antacid that neutralizes stomach acid. Excessive use can cause constipation, high blood calcium, or kidney problems and may interfere with medicine absorption.

\medterm{Calcium Carbonate Overdose} Overdose of calcium carbonate that may cause gastrointestinal symptoms, hypercalcemia, alkalosis, kidney injury, or kidney stones.

\medterm{Calcium Carbonate With Magnesium Overdose} Taking too much of these antacid ingredients can cause nausea, vomiting, constipation or diarrhea, weakness, confusion, abnormal calcium or magnesium levels, and kidney injury. Severe symptoms or reduced kidney function require urgent medical advice.
\medterm{Calcium Channel} A protein in the cell membrane that controls calcium entry into cells. Calcium channels help regulate muscle contraction, the release of substances, nerve signals, and heart rhythm.

\medterm{Calcium Chloride} A calcium salt given by injection for severe low blood calcium, dangerously high potassium, or certain toxic exposures. It can cause serious tissue injury if it leaks outside the vein and therefore requires careful administration.

\medterm{Calcium Citrate} A calcium supplement that is absorbed even when stomach acid is low. It may prevent or treat calcium deficiency but can cause constipation and, in some people, contribute to kidney stones.

\medterm{Calcium Deficiency} Too little calcium intake or availability for the body's needs. Long-term deficiency can weaken bones; severe deficiency can lower blood calcium and cause tingling, muscle cramps, spasms, or abnormal heart rhythms.

\medterm{Calcium Elevated (Hypercalcemia)} Alternate index label for Hypercalcemia.

\medterm{Calcium Gluconate} A calcium salt taken by mouth or given by injection to treat calcium deficiency and some emergencies, including dangerously high potassium or magnesium. Intravenous treatment requires medical monitoring because it can affect the heart and veins.

\medterm{Calcium Hydroxide Poisoning} Exposure to a strongly alkaline chemical that can burn the mouth, throat, esophagus, skin, or eyes. Do not induce vomiting; eye, breathing, or swallowing problems require immediate emergency assessment.

\medterm{Calcium Phosphate} A mineral that forms much of the hard structure of bones and teeth. Excess calcium phosphate can also deposit in soft tissues, especially after tissue damage or when blood calcium or phosphate levels are abnormal.

\medterm{Calcium Pyrophosphate Arthritis} Inflammation caused by calcium pyrophosphate crystals in cartilage or a joint. Sudden attacks, often called pseudogout, commonly affect the knee or wrist; chronic deposits can cause ongoing arthritis and joint damage.

\medterm{Calcium Pyrophosphate Deposition} Alternate index label; see \textit{Calcium Pyrophosphate Deposition Disease}.

\medterm{Calcium Pyrophosphate Deposition Disease} A disorder in which calcium pyrophosphate crystals collect in cartilage and joints. It may cause sudden painful swelling of one joint, usually the knee or wrist, or a chronic arthritis that resembles osteoarthritis or rheumatoid arthritis.

\synonyms
Calcium Pyrophosphate Deposition

\medterm{Calcium Sensing Receptor} A protein, especially in the parathyroid glands and kidneys, that detects the calcium level in blood. It helps control parathyroid hormone and determines how much calcium and phosphate the kidneys keep or remove.

\medterm{Calcium Signaling} A way cells communicate using changes in calcium concentration. These signals help control muscle contraction, secretion of substances, energy use, movement, and activity of genes.

\medterm{Calcium Stones} The most common type of kidney stone, made mainly of calcium oxalate or calcium phosphate. Low urine volume and certain mineral imbalances increase risk; stones may cause severe side pain, blood in the urine, blockage, or infection.

\medterm{Calculi} Hardened mineral concretions, such as gallstones, kidney stones, or salivary stones, formed when dissolved substances precipitate and aggregate within a body organ or duct.

\medterm{Calculus} A calculus is a hard deposit formed from mineral salts, such as a kidney stone, gallstone, or dental tartar.
\medterm{Calculus (Tartar)} Hardened dental plaque formed when minerals from saliva collect on teeth. Brushing cannot remove it; professional cleaning is needed because it traps bacteria and can cause gum inflammation and periodontal disease.
\synonyms
Tartar

\medterm{Calculus Removal} Professional dental treatment that removes hardened plaque from teeth and below the gumline using hand or ultrasonic instruments. It helps control gum inflammation and may include deep cleaning of affected tooth roots.

\medterm{Calculus, Renal} A kidney stone, a hard mass of mineral crystals formed in the urinary tract. It may pass without symptoms or cause severe wave-like pain, blood in the urine, blockage, vomiting, or urinary infection.

\medterm{Caldwell-Luc Operation} An operation that reaches the maxillary sinus through the upper gum and jawbone. It is used selectively to treat or examine certain sinus lesions, chronic disease, or foreign bodies when less invasive routes are unsuitable.

\medterm{Calf} The back part of the lower leg between the knee and ankle. It contains the gastrocnemius, soleus, blood vessels, nerves, and other tissues and helps point the foot downward and propel walking.

\medterm{Calf Pseudohypertrophy} An enlarged, firm appearance of the calf muscles caused mainly by replacement of damaged muscle with fat and fibrous tissue rather than increased muscle strength. It is seen especially in Duchenne and Becker muscular dystrophies and some related disorders.

\synonyms
Pseudohypertrophy of the Calves; Muscular Pseudohypertrophy

\medterm{Calibration} Checking and, when necessary, adjusting a medical or laboratory instrument against a known standard so its measurements are accurate and reliable. Regular calibration supports correct diagnosis, treatment, and patient safety.

\medterm{Caliciviridae} A family of viruses that includes noroviruses and sapoviruses, which commonly cause vomiting and diarrhea. Other members infect animals; disease depends on the virus, host, and route of exposure.

\medterm{Calipers} An instrument with two arms used to measure distance, thickness, or body dimensions. In health care, calipers may measure skinfold thickness, joint movement, wounds, or the size of body structures.

\medterm{Calla Lily Poisoning} Poisoning from chewing or swallowing calla lily, whose needle-like calcium oxalate crystals irritate the mouth and throat. Burning, pain, drooling, and swelling may occur; breathing or swallowing difficulty requires emergency care.

\medterm{Callosal} Relating to the corpus callosum, the broad bundle of nerve fibers connecting the left and right halves of the brain. Callosal abnormalities can affect communication between the two sides.
\medterm{Callosity} A callosity is a localized area of thickened skin caused by repeated pressure or friction.
\medterm{Callosotomy} Surgical division of the corpus callosum to reduce the spread of certain severe, treatment-resistant seizures between the cerebral hemispheres.

\medterm{Callus} A thick, hardened area of skin formed by repeated pressure or rubbing. Calluses commonly occur on hands or feet, usually cause no harm, and may become painful or crack when extensive.

\medterm{Calmodulin} A protein inside cells that senses calcium and passes its signals to other proteins. It helps control muscle contraction, movement of substances through cell membranes, chemical reactions, and activity of genes.

\medterm{Calnexin} A protein inside cells that helps newly made proteins fold into the correct shape. It checks proteins before they are sent to other parts of the cell or broken down.

\medterm{Calor} Heat or increased warmth of a body part, one of the classic local signs of inflammation.

\medterm{Calor, Dolor, Rubor, And Tumor} The traditional four signs of inflammation are heat, pain, redness, and swelling. Inflamed tissue may also lose some of its normal function.
\medterm{Caloric Restriction} Deliberately eating fewer calories without causing malnutrition. It may help with weight loss and some metabolic risks, but it should provide enough protein, vitamins, minerals, and other nutrients.
\medterm{Caloric Stimulation} A balance test in which warm or cool air or water is placed in the ear canal. The temperature briefly stimulates the inner-ear balance organs and helps compare the balance nerve on each side.

\medterm{Calorie} A unit used to describe the energy in food and the energy used by the body. On food labels, one Calorie means one kilocalorie, the energy needed to raise one kilogram of water by one degree Celsius.

\synonyms
a Calorie.




\medterm{Calories} Units used to express the energy supplied by food or expended by the body; food labels generally report kilocalories, commonly called calories.

\medterm{Calorimetry} Measurement of heat production or gas exchange to estimate how much energy the body uses. Indirect calorimetry uses oxygen consumption and carbon-dioxide production to help assess metabolism and nutritional needs.

\medterm{Calpain} A group of enzymes activated by calcium that cut selected proteins inside cells. They help regulate cell movement, signaling, programmed cell death, and platelet activity; abnormal calpain activity contributes to some muscle and nervous-system diseases.

\medterm{Calpainopathy} An inherited muscle disease caused by harmful changes in both copies of the CAPN3 gene. It slowly weakens hip, pelvic, shoulder, and trunk muscles, often causing shoulder-blade winging, back curvature, difficulty climbing, and walking problems while usually sparing facial and heart muscles.

\synonyms
Limb-Girdle Muscular Dystrophy R1 (LGMD R1); LGMD 2A; Calpain-3 Deficiency

\medterm{Calponin} A protein that binds to actin, a component of the cell's internal framework, and helps regulate contraction in smooth muscle. Laboratory staining for calponin can help identify muscle-related cells in tissue samples.

\medterm{Calprotectin} A protein released mainly by activated white blood cells. Measuring it in stool helps detect inflammation in the intestines, particularly in inflammatory bowel disease, and can help distinguish it from some noninflammatory conditions.
\medterm{Calreticulin} A protein inside cells that helps newly made proteins fold correctly and helps regulate calcium storage and immune signaling. Harmful changes in its gene are found in some blood cancers involving excessive production of blood cells.

\medterm{Calretinin} A calcium-binding protein found in certain nerve cells and cells lining body cavities. Pathologists may use calretinin staining to help distinguish mesothelioma from some other cancers in a tissue sample.

\medterm{Calvaria} The dome-shaped upper part of the skull that covers and protects the brain. It is formed mainly by the flat bones of the forehead, top and sides of the head, and back of the skull.

\medterm{Calvarium} The dome-like portion of the skull formed mainly by the frontal, parietal, and occipital bones. It encloses and protects the upper parts of the brain.

\medterm{Calyx} A cup-shaped body structure that supports tissue or collects fluid. In the kidney, renal calyces collect urine from the pyramids and carry it toward the renal pelvis.

\medterm{Cambridge Classification} A grading system used with endoscopic retrograde cholangiopancreatography to describe chronic pancreatitis. It rates abnormalities of the main pancreatic duct and its branches, from normal or uncertain findings to mild, moderate, or severe disease.

\medterm{Cameron Lesions} Long, shallow erosions or ulcers in the stomach where it is squeezed by the diaphragm, usually in people with a large hiatal hernia. They may cause chronic blood loss, iron-deficiency anemia, or visible gastrointestinal bleeding.

\medterm{Campesterol} A plant sterol chemically similar to cholesterol and found in many plant foods. Most people absorb only small amounts, but excessive absorption in the rare disorder sitosterolemia can contribute to abnormal sterol levels and vascular problems.

\medterm{Campho-Phenique Overdose} Poisoning after swallowing or excessive exposure to a topical product containing camphor and other ingredients. It may cause nausea, vomiting, drowsiness, agitation, or seizures and requires poison-control or emergency advice.

\medterm{Camphor} A strong-smelling substance used in some topical products and traditional remedies. Swallowing even a relatively small amount can cause nausea, agitation, drowsiness, or rapidly developing seizures and may be life-threatening.

\medterm{Camphor Overdose} Toxic exposure to too much camphor, causing nausea, vomiting, agitation, drowsiness, seizures, coma, or slowed breathing. Seizures may begin quickly; suspected ingestion requires immediate poison-control or emergency assessment.

\medterm{Campimetry} Measurement and mapping of the area a person can see while looking straight ahead. It helps detect and describe visual-field loss caused by disease of the retina, optic nerve, or brain pathways.

\medterm{Campomelic Dysplasia} A rare inherited disorder caused by changes affecting the SOX9 gene. It causes bowed long bones, poorly developed shoulder blades, characteristic facial features, and often serious problems with breathing or other organs.
\medterm{Camptocormia} Severe forward bending of the trunk that worsens while standing or walking and improves when lying down. It may result from muscle disease, Parkinson disease, dystonia, or another nerve or muscle disorder.

\medterm{Camptodactyly} A fixed bending deformity of one or more fingers, usually at the middle finger joint. It may be present from birth or develop later because of shortened tendons, muscles, connective tissue, or joint structures.

\medterm{Camptothecin} A plant-derived compound that blocks topoisomerase I, an enzyme needed to copy and repair DNA. It is mainly a research compound; related medicines are used in some cancer treatments and can damage rapidly dividing cells.

\medterm{Campylobacter} A group of bacteria that commonly cause foodborne intestinal infection. Illness may include fever, abdominal cramps, diarrhea, or bloody stools; some species can invade the bloodstream or cause later nerve or joint problems.

\medterm{Campylobacter Coli} A Campylobacter species that can cause foodborne intestinal infection, usually after contaminated meat, water, or other food. Symptoms include fever, abdominal cramps, diarrhea, and sometimes bloody stools.
\medterm{Campylobacter Concisus} A Campylobacter species found mainly in the mouth and digestive tract. Its role in disease is uncertain, but it has been associated with occasional intestinal or bloodstream infection, especially in people with weakened immunity.

\medterm{Campylobacter Curvus} A Campylobacter species found in the mouth and digestive tract. It rarely causes human disease but has been reported in infections outside the intestine, particularly when normal body defenses are impaired.

\medterm{Campylobacter Fetus} A Campylobacter species that can cause bloodstream infection and infections of blood vessels, joints, the brain, or other tissues. Serious disease is more likely in older adults and people with weakened immunity.

\medterm{Campylobacter Gracilis} A Campylobacter species found mainly in the mouth and upper airways. It is usually part of the normal microbial community but may contribute to gum disease or rarely cause infection elsewhere.

\medterm{Campylobacter Infection} A bacterial infection, usually caused by C. jejuni or C. coli, acquired from contaminated poultry, food, water, or unpasteurized milk. It causes fever, cramps, and diarrhea that may be bloody; rare complications include bloodstream infection, arthritis, or Guillain--Barré syndrome.

\medterm{Campylobacter Jejuni} A major cause of foodborne intestinal infection, usually acquired from undercooked poultry, contaminated water, or unpasteurized milk. It commonly causes fever, abdominal cramps, and diarrhea that may contain blood.
\medterm{Campylobacter Lari} A Campylobacter species associated with marine environments and poultry. It can rarely cause diarrhea or infection outside the intestine, especially in people with weakened immunity.

\medterm{Campylobacter Rectus} A bacterium found mainly in the mouth and linked with gum disease. It can rarely cause infections of the bloodstream, heart, or other tissues, usually when normal defenses are impaired.

\medterm{Campylobacter Serology Test} A blood test that detects antibodies produced after Campylobacter infection. It may support evaluation of selected complications after infection, but stool molecular testing or culture is generally used to diagnose active intestinal disease.

\medterm{Campylobacter Upsaliensis} A Campylobacter species carried by dogs, cats, and other animals that can occasionally cause human diarrhea or bloodstream infection. Risk is higher in people with weakened immunity.

\medterm{Campylobacteraceae} A family of curved or spiral bacteria that includes Campylobacter and related species. Some members cause intestinal infection, while others live harmlessly in animals or the environment and only rarely infect people.

\medterm{Campylobacteriosis} Infection caused by Campylobacter bacteria, usually producing fever, abdominal pain, and diarrhea that may be bloody. Dehydration is common; rare complications include bloodstream infection, reactive arthritis, and Guillain--Barré syndrome.



\medterm{Canagliflozin} A medicine that lowers blood glucose by causing the kidneys to remove more glucose in urine. It may also reduce some heart-failure and kidney risks, but can cause genital infections, dehydration, low blood pressure, or ketoacidosis.

\medterm{Canal Of Schlemm} A circular drainage channel at the edge of the cornea that carries aqueous fluid from the front of the eye into the bloodstream. Poor drainage can raise eye pressure and contribute to glaucoma.

\medterm{Canaliculus} A small anatomical channel. The tear canaliculi are tiny passages in the eyelids that carry tears from the eye surface into the lacrimal sac and then the nose.
\medterm{Canalith Repositioning Maneuver} A series of guided head and body movements that moves displaced inner-ear crystals out of a semicircular canal. It is commonly used to treat posterior-canal benign paroxysmal positional vertigo.

\medterm{Canalith Repositioning Maneuvers} Head and body movements that move loose crystals from an inner-ear semicircular canal back to the utricle, where they no longer trigger false movement signals. The maneuver is selected according to the affected canal and side.

\medterm{Canalization} Formation or restoration of a channel through tissue, a blood vessel, or an embryonic structure. In a blood clot, canalization means that small channels form through the clot and partly restore blood flow.
\medterm{Canarypox Virus} A poxvirus that mainly infects birds. It does not normally cause human disease but is used as a research and vaccine-delivery vector because it can enter human cells without usually multiplying in them.

\medterm{Canavan Disease} A rare inherited disorder in which deficiency of the ASPA enzyme damages the brain's white matter. It usually begins in infancy with poor muscle tone, delayed development, increasing head size, feeding problems, and later severe movement and communication impairment.

\medterm{Cancell} A shortened or misspelled form of cancellous, meaning the porous inner part of bone that contains marrow and helps spread mechanical forces.

\medterm{Cancellous Bone} Porous bone made of thin supporting plates or struts. It is found inside many bones, especially near joints and in the vertebrae, and contains marrow while helping distribute forces.

\medterm{Cancer} A disease in which abnormal cells grow without normal control and may invade nearby tissues or spread to distant organs. Cancers are classified by where they begin and the cell changes they contain; treatment depends on these features and the extent of disease.


\medterm{Cancer And Lymph Nodes} Cancer cells may travel from a primary tumor to nearby lymph nodes through lymphatic vessels. Examination, imaging, or biopsy of these nodes helps determine the cancer's extent and guide treatment.

\medterm{Cancer Care Team} The health professionals who work together to diagnose, treat, and support a person with cancer. Depending on the illness, the team may include medical and radiation oncologists, surgeons, pathologists, nurses, pharmacists, therapists, and palliative-care specialists.

\medterm{Cancer Causes} Cancer develops after changes in cells allow abnormal growth and survival. Risk can be increased by inherited variants, tobacco, alcohol, certain infections, ultraviolet or other radiation, harmful chemicals, chronic inflammation, and random changes during cell division.

\synonyms
Cancer Risk (Cancer Causes)

\medterm{Cancer Cell} A cell with changes that allow it to grow and survive without normal control. Cancer cells may lose their specialized function, resist cell death, invade nearby tissues, enter blood or lymph, and form tumors elsewhere.

\medterm{Cancer Cluster} A group or area with more cancer cases than expected during a specified time. Investigation compares the number and types of cancers with reliable background rates and examines possible shared exposures while considering chance variation.


\medterm{Cancer Genetics} The study of inherited and acquired genetic changes that affect cancer risk, behavior, or treatment. Testing may identify inherited risk syndromes or changes in a tumor that help guide treatment and family counseling.

\medterm{Cancer Grading} Assessment of how abnormal cancer cells and their tissue pattern look under a microscope. The grade gives information about likely growth behavior and is different from staging, which describes how far the cancer has spread.

\medterm{Cancer Immunotherapy} Cancer treatment that helps the immune system recognize or attack cancer cells. Methods include checkpoint inhibitors, engineered immune cells, immune-stimulating proteins, and some vaccines; use and side effects depend on the cancer and treatment.

\medterm{Cancer Mutations} Genetic changes in cancer cells that can alter growth, survival, or spread. Most develop during life and are not inherited. Some drive cancer development, while others have little effect; testing selected changes can help guide prognosis or targeted treatment.

\medterm{Cancer In Lung (Lung Cancer)} Alternate index label for Lung Cancer.

\medterm{Cancer Nutrition} Nutrition care for people with cancer, who may lose appetite or weight because of the disease or treatment. Individual plans may address calories, protein, fluids, swallowing, nausea, diarrhea, and cachexia while respecting the person's goals and treatment.

\medterm{Cancer Of Unknown Origin} Alternate index label; see \textit{Carcinoma of unknown primary}.

\medterm{Cancer Of Unknown Primary} Cancer that has spread to another site but whose original starting point cannot be identified after appropriate examination and testing. Tissue studies, scans, and selected molecular tests may help guide treatment.

\medterm{Cancer Pain} Pain caused by a tumor, its spread, or its treatment. It may be constant or intermittent and can arise from bone, nerves, organs, surgery, radiation, or medicines; treatment combines control of the cause with suitable pain relief and supportive care.

\synonyms
Pain Cancer (Cancer Pain)

\medterm{Cancer Prevention} Actions that lower the chance of developing cancer or find precancerous changes early. They include avoiding tobacco, limiting alcohol, vaccination against cancer-causing infections, protecting skin from ultraviolet light, reducing harmful exposures, and recommended screening.

\medterm{Cancer Registry} An organized record of cancer diagnoses, including cancer type, location, stage, treatment, and outcome. Registries help measure cancer rates, plan services, improve care, and support research while protecting personal information.

\medterm{Cancer Risk} The likelihood that a person will develop cancer. It is influenced by age, inherited variants, family history, tobacco, alcohol, certain infections, radiation, environmental exposures, body weight, and chance; having a risk factor does not guarantee cancer.

\medterm{Cancer Risk (Cancer Causes)} Alternate index label for Cancer Causes.

\medterm{Cancer Screening} Testing people who have no symptoms to find selected cancers or precancerous changes early, when treatment may work better. Each test has specific age, risk, interval, benefits, and possible harms.

\medterm{Cancer Staging} Describing how far a cancer has grown or spread. Assessment may include the primary tumor, nearby lymph nodes, distant metastases, scans, and tissue examination; stage helps guide treatment and estimate prognosis.

\medterm{Cancer Surveillance} Ongoing monitoring for cancer recurrence, progression, treatment effects, or new cancers after diagnosis or treatment. It may include examinations, imaging, laboratory tests, and symptom review according to the cancer and treatment.

\medterm{Cancer Survivor} A person living with cancer from diagnosis through treatment and afterward. The term includes people whose cancer is controlled or cured and those with ongoing disease; follow-up may address recurrence, late effects, emotional health, and quality of life.

\medterm{Cancer Susceptibility Genes} Genes in which inherited harmful variants increase the chance of developing one or more cancers. Examples include BRCA1, BRCA2, APC, and mismatch-repair genes; carrying a risk variant does not guarantee that cancer will develop.

\medterm{Cancer Treatment - Dealing With Pain} Pain during or after cancer treatment may come from the tumor, surgery, radiation, chemotherapy, or nerve injury. Care may combine treatment of the cause, pain medicines, physical rehabilitation, psychological support, and other supportive measures.

\medterm{Cancer Treatment - Early Menopause} Cancer treatment may stop menstrual periods earlier than expected, especially after chemotherapy, pelvic radiation, ovarian surgery, or hormone-blocking treatment. It can cause hot flashes, vaginal symptoms, infertility, and bone loss and requires individualized follow-up.


\medterm{Cancer Treatment: Dealing With Hot Flashes And Night Sweats} Cancer treatment or hormone changes can cause sudden heat, sweating, and disturbed sleep. Cooling measures, layered clothing, sleep habits, behavioral approaches, and selected medicines may help, depending on the cancer and other health risks.


\medterm{Cancer Treatments} Cancer treatment may include surgery, radiation, chemotherapy, immunotherapy, targeted medicines, hormone therapy, stem-cell transplantation, or supportive and palliative care. The choice depends on cancer type, location, stage, molecular features, overall health, and patient goals.

\medterm{Cancer, Carcinoid Tumors} Alternate index label; see \textit{Carcinoid tumors}.

\medterm{Cancer, Chronic Lymphocytic Leukemia} Alternate index label; see \textit{Chronic lymphocytic leukemia}.

\medterm{Cancer, Chronic Myelogenous Leukemia} Alternate index label; see \textit{Chronic myelogenous leukemia}.

\medterm{Cancer, Colon} Alternate index label; see \textit{Colon cancer}.

\medterm{Cancer, Endometrial} Alternate index label; see \textit{Endometrial cancer}.

\medterm{Cancer, Esophageal} Alternate index label; see \textit{Esophageal cancer}.

\medterm{Cancer, Eye Melanoma} Alternate index label; see \textit{Eye melanoma}.

\medterm{Cancer, Floor Of The Mouth} Alternate index label; see \textit{Floor of the mouth cancer}.

\medterm{Cancer, Gallbladder} Alternate index label; see \textit{Gallbladder cancer}.

\medterm{Cancer, Gastric} Alternate index label; see \textit{Stomach cancer}.

\medterm{Cancer, Hairy Cell Leukemia} Alternate index label; see \textit{Hairy cell leukemia}.

\medterm{Cancer, Hodgkin Lymphoma} Alternate index label; see \textit{Hodgkin lymphoma (Hodgkin disease)}.

\medterm{Cancer, Hurthle Cell} Alternate index label; see \textit{Hurthle cell cancer}.

\medterm{Cancer, Inflammatory Breast} Alternate index label; see \textit{Inflammatory breast cancer}.

\medterm{Cancer, Kidney} Alternate index label; see \textit{Kidney cancer}.

\medterm{Cancer, Liver} Alternate index label; see \textit{Liver cancer}.

\medterm{Cancer, Lung} Alternate index label; see \textit{Lung cancer}.

\medterm{Cancer, Merkel Cell} Alternate index label; see \textit{Merkel cell carcinoma}.

\medterm{Cancer, Mouth} Alternate index label; see \textit{Mouth cancer}.

\medterm{Cancer, Multiple Myeloma} Alternate index label; see \textit{Multiple myeloma}.

\medterm{Cancer, Nasopharyngeal} Alternate index label; see \textit{Nasopharyngeal carcinoma}.

\medterm{Cancer, Neuroblastoma} Alternate index label; see \textit{Neuroblastoma}.

\medterm{Cancer, Neuroendocrine} Alternate index label; see \textit{Neuroendocrine tumors}.

\medterm{Cancer, Non-Hodgkin Lymphoma} Alternate index label; see \textit{Non-Hodgkin lymphoma}.

\medterm{Cancer, Ovarian} Alternate index label; see \textit{Ovarian cancer}.

\medterm{Cancer, Pancreatic} Alternate index label; see \textit{Pancreatic cancer}.

\medterm{Cancer, Paraneoplastic Syndromes} Alternate index label; see \textit{Paraneoplastic syndromes of the nervous system}.

\medterm{Cancer, Prostate} Alternate index label; see \textit{Prostate cancer}.

\medterm{Cancer, Rectal} Alternate index label; see \textit{Rectal cancer}.

\medterm{Cancer, Retinoblastoma} Alternate index label; see \textit{Retinoblastoma}.

\medterm{Cancer, Skin} Alternate index label; see \textit{Skin cancer}.

\medterm{Cancer, Soft Tissue Sarcoma} Alternate index label; see \textit{Soft tissue sarcoma}.

\medterm{Cancer, Squamous Cell} Alternate index label; see \textit{Squamous cell carcinoma of the skin}.

\medterm{Cancer, Stomach} Alternate index label; see \textit{Stomach cancer}.

\medterm{Cancer, Testicular} Alternate index label; see \textit{Testicular cancer}.

\medterm{Cancer, Throat} Alternate index label; see \textit{Throat cancer}.

\medterm{Cancer, Thyroid} Alternate index label; see \textit{Thyroid cancer}.

\medterm{Cancer, Uterine} Alternate index label; see \textit{Endometrial cancer}.

\medterm{Cancer, Vagina} Alternate index label; see \textit{Vaginal cancer}.

\medterm{Cancer, Vulvar} Alternate index label; see \textit{Vulvar cancer}.

\medterm{Cancer, Wilms' Tumor} Alternate index label; see \textit{Wilms tumor}.

\medterm{Cancer-Associated Thrombosis} A blood clot in a vein or artery occurring with active cancer or its treatment. Cancer, surgery, immobility, central lines, and some treatments increase risk; anticoagulant treatment is chosen according to clot location, bleeding risk, and cancer status.

\medterm{Cancrum Oris} Noma, a rapidly destructive gangrenous infection of the mouth and face that occurs
mainly in severely malnourished or immunocompromised children. It begins with gingival
or mucosal necrosis and can extend through facial tissues.

\medterm{Candela} The international unit used to measure the brightness of a light source in a specified direction. It is a physical unit, not a measure of a medicine's light dose or the amount of light reaching tissue.

\medterm{Candida} A group of yeasts that may normally live on the skin, in the mouth, or in the digestive and genital tracts. Overgrowth can cause mouth, genital, skin, or, in vulnerable people, bloodstream and organ infections.

\medterm{Candida Albicans} The Candida species most often causing human infection. It can cause oral thrush, vaginal yeast infection, moist skin-fold rash, or serious bloodstream and organ infection in people with weakened defenses.

\medterm{Candida Albicans Infection} Infection caused by Candida albicans. It may be limited to the mouth, vagina, or skin, or may invade the bloodstream and organs, especially in people who are hospitalized or have weakened immunity.

\medterm{Candida Auris} A yeast that can live on skin without symptoms or cause bloodstream, wound, or ear infection. It often resists several antifungal medicines, spreads readily in health-care settings, and requires specialized testing and strict infection-control measures.

\medterm{Candida Auris Infection} Infection or symptom-free colonization caused by Candida auris. Invasive infection may affect the blood or organs, especially in hospitalized or immunocompromised people; treatment is guided by susceptibility testing because resistance is common.

\medterm{Candida Boidinii} A yeast found mainly in the environment and occasionally in people. Human infection is uncommon but may occur opportunistically, particularly in people with severe illness, invasive devices, or weakened immunity.

\medterm{Candida Dubliniensis} A yeast related to Candida albicans that can live in the mouth and cause oral thrush or, rarely, invasive infection. It is more likely to cause disease when immune defenses are impaired.

\medterm{Candida Esophagitis} Yeast infection of the esophagus, usually caused by Candida. It commonly causes painful or difficult swallowing and chest discomfort, especially in people with weakened immunity or after certain medicines.

\medterm{Candida Famata} A yeast found in the environment and occasionally on people. It rarely causes bloodstream or other invasive infection, mainly in hospitalized people with weakened immunity or indwelling medical devices.

\medterm{Candida Glabrata} A Candida species that can cause genital, urinary, bloodstream, or other invasive infection. It may resist some commonly used antifungal medicines, so laboratory susceptibility testing can help guide treatment.

\medterm{Candida Inconspicua} A yeast that is an uncommon cause of human infection. It may cause bloodstream or other invasive disease mainly in hospitalized people with weakened immunity, and laboratory identification may be needed for treatment selection.


\medterm{Candida Infection Of The Skin} A superficial yeast infection, often in warm, moist skin folds. It causes a red, itchy, sore rash that may have small separate spots around its edge and is more common with moisture, diabetes, obesity, or immune suppression.

\medterm{Candida Intermedia} A yeast found in the environment and occasionally in people. It rarely causes infection but can produce bloodstream or other invasive disease in hospitalized people with weakened immunity.

\medterm{Candida Orthopsilosis} A yeast related to Candida parapsilosis that can rarely cause bloodstream or device-associated infection, particularly in hospitalized people, newborns, or those with weakened immunity.

\medterm{Candida Parapsilosis} A Candida species that can cause bloodstream and catheter-related infection, especially in newborns, hospitalized people, and those receiving intravenous treatment. It may spread between patients through hands or contaminated equipment.
\medterm{Candida Sake} A yeast found mainly in the environment and food. Human infection is rare but may occur opportunistically, especially in people with severe illness or weakened immune defenses.

\medterm{Candida Tropicalis} A Candida species that can cause mouth or genital infection and may invade the bloodstream or organs. Serious infection is more likely in people with neutropenia, cancer, or other immune impairment.

\medterm{Candida Zeylanoides} A yeast found mainly in the environment. It is an uncommon cause of human disease but may rarely cause bloodstream or other infection in people with severe illness or weakened immunity.

\medterm{Candidal Intertrigo} A Candida yeast infection in warm, moist skin folds such as the armpits, groin, under the breasts, or beneath an abdominal fold. It causes a red, moist, itchy or burning rash, sometimes with small spots at the edge.

\medterm{Candidal Vaginitis} Inflammation of the vagina and vulva caused by overgrowth of Candida yeast. It usually causes intense itching, soreness, redness, pain with sex or urination, and a thick white discharge without a strong odor.

\medterm{Candidal Vulvovaginitis} Alternate index label; see \textit{Vulvovaginal Candidiasis}.

\medterm{Candidemia} Candida yeast infection of the bloodstream. It can spread to organs and is more common with central lines, surgery, intensive care, recent antibiotics, or weakened immunity; it requires prompt medical treatment and evaluation for spread.

\medterm{Candidiasis} Infection caused by Candida yeast. It may affect the mouth, esophagus, vagina, or moist skin folds, or may invade the blood and organs in seriously ill people. Symptoms and treatment depend on the site and whether the infection is invasive.

\medterm{Candidiasis, Oral} Alternate index label; see \textit{Oral thrush}.

\medterm{Candidiasis, Vaginal} Alternate index label; see \textit{Yeast infection (vaginal)}.

\medterm{Candles Poisoning} Illness after swallowing candle wax or exposure to candle products. Small amounts of plain wax usually cause mild stomach upset, but scented or decorative products may contain harmful ingredients; choking or breathing problems require urgent care.

\medterm{Cane} A walking aid used to improve balance and reduce some weight on a painful or weak leg. It should be fitted to the person's height and used on the side opposite the weaker or painful leg unless a clinician advises otherwise.

\medterm{Canertinib Dihydrochloride} An experimental anticancer compound that blocks several epidermal growth-factor receptors, proteins that can promote tumor-cell growth. It was studied in clinical research but is not a routine approved cancer treatment.

\medterm{Canfosfamide Hcl} An experimental anticancer prodrug designed to become active after processing by an enzyme found at increased levels in some tumors. It was studied in cancer research but is not a routine approved treatment.

\medterm{Canicola Fever} A form of leptospirosis caused by a strain of Leptospira bacteria commonly carried by dogs. It may cause fever, headache, muscle aches, red eyes, kidney injury, or liver involvement after contact with infected urine or contaminated water.

\medterm{Canine Teeth} The pointed teeth between the incisors and premolars. They help tear food and guide the jaws during biting and chewing.
\medterm{Canker} A nontechnical word often used for an aphthous mouth ulcer, a painful shallow sore on the moist lining of the mouth. The word can have other meanings, so the site and appearance should be clarified.
\medterm{Canker Sore} A painful, shallow ulcer on the moist lining of the mouth, usually round or oval with a red border. It is not contagious and commonly heals within one to two weeks; persistent, severe, or unusual sores need assessment.

\synonyms
Recurrent Oral Ulcers
Ulcer, Aphthous

\medterm{Canker Sores} Painful, shallow, noncontagious ulcers on the moist lining of the mouth. Minor injury, stress, nutritional deficiency, immune disease, medicines, or other triggers may contribute; frequent, severe, or nonhealing sores need medical or dental evaluation.

\synonyms
Sores Canker (Canker Sores)

\medterm{Cannabidiol} A cannabinoid from the cannabis plant that usually does not produce the “high” caused by THC. A purified form treats selected severe seizure disorders; other products vary in strength, purity, side effects, and medicine interactions.

\medterm{Cannabinoid} A substance that acts on cannabinoid receptors or related signaling systems. Cannabinoids may come from plants, the body, or laboratories; effects can include pain relief, appetite changes, sedation, impaired memory, or intoxication.

\medterm{Cannabinoid Hyperemesis Syndrome} Repeated episodes of severe nausea, vomiting, and abdominal pain associated with long-term frequent cannabis use. Hot showers may temporarily ease symptoms, but lasting improvement usually requires stopping cannabis.

\medterm{Cannabinoid Receptor} A cell-surface protein, mainly CB1 or CB2, that responds to cannabinoids made by the body or found in cannabis and medicines. It helps regulate pain, appetite, mood, memory, immune activity, and inflammation.

\medterm{Cannabinol} A cannabinoid formed as THC breaks down with age or exposure to air. It may have mild psychoactive effects, but its medical benefits and safety are less well established than those of some other cannabinoids.

\medterm{Cannabis} A plant containing cannabinoids such as THC and CBD. Use can alter perception, memory, coordination, mood, and heart rate; it may have selected medical uses but can also cause dependence, anxiety, psychosis, or other harms.
\medterm{Cannabis Use Disorder} A pattern of cannabis use that causes significant difficulty controlling use or problems at home, work, school, relationships, or health. Tolerance and withdrawal may occur, with irritability, anxiety, sleep problems, or reduced appetite after stopping.

\medterm{Cannot Intubate Cannot Oxygenate} A life-threatening airway emergency in which clinicians cannot place a breathing tube and cannot provide enough oxygen by mask or another airway device. It requires immediate emergency action, often including a surgical airway.

\medterm{Cannula} A thin hollow tube placed in a blood vessel, body cavity, or tissue to deliver medicines or fluids, remove material, measure pressure, provide oxygen, or allow access for another device.

\medterm{Cannula Hub} The connector at the outer end of a cannula that attaches it to tubing, a syringe, or another device. It may include a port or valve for giving fluids, taking samples, or flushing the cannula.

\medterm{Cannulation} Placement of a cannula into a blood vessel, body cavity, duct, or other passage. It provides access for giving fluids or medicines, removing material, monitoring, collecting samples, or performing a procedure.

\medterm{Canola Oil} A vegetable oil relatively low in saturated fat and containing mostly monounsaturated and polyunsaturated fats. Replacing solid fats with canola oil can support a heart-healthy eating pattern when total calories remain appropriate.

\medterm{Canrenone} A metabolite related to spironolactone that blocks aldosterone's action. It promotes loss of sodium and water while reducing potassium loss and has been used in some countries for conditions involving fluid retention or excess aldosterone.

\medterm{Cantharides} Preparations made from blister beetles that contain cantharidin, a powerful irritant. Swallowing or applying it can cause severe mouth, stomach, kidney, and urinary-tract injury and may be life-threatening.
\medterm{Canthaxanthin} A red-orange pigment used as a food color and formerly in tanning products. Excessive intake can cause yellow skin, liver injury, and crystal deposits in the retina that impair vision.

\medterm{Canthotomy} A surgical cut through the outer corner of the eyelids, usually performed urgently to relieve dangerous pressure around the eye and protect vision after severe orbital swelling or injury.

\medterm{Canthus} The corner where the upper and lower eyelids meet. The inner canthus is beside the nose, and the outer canthus is toward the temple.

\medterm{Cantlie Line} An imaginary line across the liver from the gallbladder to the large vein behind the liver. It follows the middle hepatic vein and marks the functional division between the right and left liver, including their blood supply and bile drainage.

\medterm{Cantuzumab Ravtansine} An experimental antibody-drug conjugate designed to attach to MUC1, a protein found on some cancer cells, and deliver a cell-killing drug. It has been studied in trials but is not a routine approved treatment.

\medterm{Capacity} A person's ability, at a particular time, to understand information about a medical decision, appreciate how it applies to them, weigh the options, and communicate a choice.

\medterm{Capacity And Competence} Capacity is the clinical ability to make a specific decision by understanding, using, and communicating relevant information. Competence is a legal determination; the two terms are related but not interchangeable.

\medterm{Capecitabine} An oral chemotherapy medicine that the body converts to fluorouracil, which interferes with cancer-cell DNA production. It is used for selected cancers and can cause diarrhea, mouth sores, hand-foot syndrome, low blood counts, or heart problems.

\medterm{Capgras Syndrome} A fixed false belief that someone familiar has been replaced by an identical impostor. It may occur with psychotic illness, dementia, brain injury, seizures, or other neurologic disorders and needs clinical assessment.

\medterm{Capillaria} A group of parasitic roundworms. Different species can cause intestinal infection, liver disease, or infection of other tissues; symptoms and treatment depend on the species and affected organ.

\medterm{Capillaria Philippinensis} A parasitic roundworm acquired by eating raw or undercooked freshwater fish. It infects the intestine and can cause persistent watery diarrhea, abdominal pain, weight loss, poor absorption of nutrients, and loss of blood proteins.

\medterm{Capillaries} The body's smallest blood vessels. Their thin walls allow oxygen, nutrients, hormones, fluid, and waste products to move between the blood and nearby tissues.
\medterm{Capillaroscopy} Examination of tiny blood vessels, usually in the skin beside a fingernail, with magnification. It can show widened or missing vessels, bleeding, and abnormal patterns that support evaluation of connective-tissue and small-vessel disease.

\medterm{Capillary} The smallest type of blood vessel, forming networks that connect arterioles with venules. Its thin wall allows exchange of oxygen, nutrients, fluid, and waste between blood and tissues.

\medterm{Capillary Bed} A network of capillaries within a tissue that connects small arteries to small veins. It provides a large surface for exchange of oxygen, carbon dioxide, nutrients, hormones, fluid, and waste.

\medterm{Capillary Electrophoresis} A laboratory method that separates charged molecules in a very thin tube using an electric field. It can analyze proteins, medicines, DNA, or other substances in a sample, depending on the test.
\medterm{Capillary Leak Syndrome} A condition in which fluid and proteins escape from small blood vessels into surrounding tissues. It can cause sudden swelling, low blood pressure, thickened blood, shock, and organ failure; severe infection, medicines, immune disease, or rare unexplained episodes may trigger it.

\medterm{Capillary Microscope} A microscope used to magnify tiny blood vessels, especially those beside the fingernails. It helps assess vessel shape, size, bleeding, and areas without vessels in some circulation and connective-tissue disorders.
\medterm{Capillary Nail Refill Test} A bedside check of circulation in which a nail bed is pressed until it turns pale and the time for color to return is observed. Temperature, lighting, age, skin color, shock, and vascular disease affect the result, so it is not used alone.

\medterm{Capillary Plasma} The liquid component of blood obtained from small vessels in skin or tissue, containing water, electrolytes, proteins, hormones, nutrients, and waste products.

\medterm{Capillary Refill} A bedside assessment of peripheral perfusion performed by briefly pressing a nail bed or
skin and measuring the time for normal color to return. Delayed refill may indicate poor
circulation, shock, hypothermia, or vascular disease.

\medterm{Capillary Refill Time} The time required for color to return after pressure briefly blanches a nail bed or area of skin. A delayed
result may suggest impaired peripheral perfusion, but interpretation depends on age,
temperature, technique, and the clinical setting.

\medterm{Capillary Sample} A small blood specimen obtained by puncturing the skin, usually a finger or heel, for bedside glucose, hemoglobin, newborn screening, or other point-of-care tests.

\medterm{Capillary Vessel} A microscopic vessel with a thin endothelial wall where blood exchanges oxygen, carbon dioxide, nutrients, fluid, and waste with surrounding tissue.

\medterm{Capitate} The largest wrist bone, located centrally in the distal row of carpal bones. It connects with several other wrist bones and the third metacarpal, helps transmit force through the wrist, and can rarely develop blood-supply loss after fracture.

\medterm{Capitation} A health-care payment arrangement in which a clinician or organization receives a fixed amount per enrolled person for specified services over a set period, whether or not that person uses the services.
\medterm{Caplan'S Syndrome} A lung disorder combining rheumatoid arthritis with pneumoconiosis caused by long-term inhalation of coal or mineral dust. It produces multiple lung nodules and may cause cough or breathing difficulty.
\medterm{Capnocytophaga} A group of bacteria found mainly in the mouths of dogs, cats, and people. Some species can cause serious human infection after animal bites or scratches, especially in people without a functioning spleen or with weakened immunity.

\medterm{Capnocytophaga Canimorsus} A bacterium commonly found in the mouths of dogs and cats. It can cause rapidly severe bloodstream infection, meningitis, or skin infection after a bite or scratch, especially in older adults or people without a functioning spleen or with weakened immunity.

\medterm{Capnocytophaga Gingivalis} A bacterium that can live in the human mouth and has been associated with gum disease. It rarely causes infection outside the mouth, mainly in people whose immune defenses are impaired.

\medterm{Capnocytophaga Infection} Infection caused by Capnocytophaga bacteria, often after a dog or cat bite, scratch, or contact with animal saliva. It may cause a wound infection, fever, bloodstream infection, or meningitis, with greatest risk in people with impaired immunity or no spleen.
\medterm{Capnocytophaga Ochracea} A bacterium found mainly in the human mouth and associated with dental plaque and gum disease. Infection outside the mouth is uncommon but can occur when immune defenses are severely weakened.

\medterm{Capnocytophaga Sputigena} A bacterium that may live in the mouth and contribute to gum disease. It can rarely cause infection of the bloodstream or other tissues, usually in people with serious illness or weakened immunity.

\medterm{Capnography} Continuous measurement and display of carbon dioxide in exhaled air. The waveform and end-tidal value help assess breathing, airway blockage, circulation, and correct placement of a breathing tube during anesthesia, sedation, or emergency care.

\medterm{Capreomycin} An injectable antibiotic formerly used for drug-resistant tuberculosis. It can damage the kidneys, hearing, or balance and is now used only in selected specialist-directed situations because safer options may be available.

\medterm{Capromab Pendetide} A radioactive antibody formerly used to image prostate-cancer tissue by binding a protein on prostate cells. Its clinical use is limited and depends on local approval and availability of more accurate imaging methods.

\medterm{Caps} A prescription abbreviation for capsules, oral medicines enclosed in a shell that dissolves in the digestive tract. Abbreviations can be misread, so complete dosage-form wording is safer on prescriptions.

\medterm{Capsaicin} The substance that makes chili peppers feel hot. Applied in controlled amounts to the skin, it can reduce some nerve and joint pains by temporarily lowering pain signaling; it commonly causes initial burning or redness.

\medterm{Capsid} The protein shell surrounding a virus's genetic material. It protects the genome and helps the virus attach to and enter host cells; its structure influences which cells can be infected and how the immune system responds.

\medterm{Capsid Proteins} Proteins that assemble around a virus's genetic material to form its protective shell. Some also help the virus attach to and enter cells and are recognized by the immune system.

\medterm{Capsular Contracture} Abnormal tightening and hardening of scar tissue around an implanted device, especially a breast implant. It may make the area firm, painful, distorted, or displaced and sometimes requires corrective surgery.

\medterm{Capsule} A protective covering around an organ, joint, or other structure, or an oral medicine enclosed in a dissolvable shell. The exact meaning depends on whether the term is used in anatomy or pharmacology.

\medterm{Capsule Endoscope} A swallowable capsule containing a tiny camera and transmitter. As it travels through the digestive tract, it sends images, especially of the small intestine, where conventional endoscopes may not reach easily.

\medterm{Capsule Endoscopy} A test in which a person swallows a small wireless camera that takes pictures of the digestive tract, especially the small intestine. The capsule usually passes naturally but can become stuck where the bowel is narrowed.

\medterm{Capsule Of Glisson} The layer of connective tissue surrounding the liver. It extends inward around the liver's blood vessels, bile ducts, and nerves and helps support and protect these structures.

\medterm{Capsule Of Tenon} A thin connective-tissue sheath surrounding the eyeball and separating it from the fat of the eye socket. It allows smooth eye movement and creates a space used for some eye injections and operations.

\medterm{Capsulectomy} Surgical removal of some or all of the scar-tissue capsule around an organ, implant, or other structure. It may be performed for painful tightening, infection, abnormal tissue, or revision of a reconstruction.

\medterm{Captopril} An angiotensin-converting enzyme inhibitor used for high blood pressure, heart failure, and some kidney problems. It can cause cough, high potassium, low blood pressure, or swelling.

\medterm{Captopril Renal Scan} A kidney imaging test using a radioactive tracer before and after captopril. It compares kidney blood flow and function and may help evaluate high blood pressure caused by narrowing of a kidney artery.
\medterm{Caput Medusae} Enlarged veins spreading out from the navel, caused by increased pressure in the portal vein system, usually from liver disease. They are a physical sign of abdominal collateral circulation rather than a disease by themselves.

\medterm{Caput Succedaneum} Soft swelling of a newborn's scalp caused by pressure during labor. It can cross the skull's suture lines and usually disappears within hours to a few days; rapidly enlarging swelling or pallor requires urgent assessment.

\medterm{CAR T-Cell} A T lymphocyte genetically changed to carry a receptor that recognizes a specific marker on cancer cells. These cells can be used as treatment for selected blood cancers but may cause severe immune reactions or nervous-system effects.

\medterm{Chimeric Antigen Receptor T-Cell Therapy} A cancer treatment in which a person's T cells are collected, genetically changed in a laboratory to recognize a cancer marker, multiplied, and returned to the body. Treatment can cause serious inflammation, infection risk, or neurologic effects.

\medterm{Carbachol} A medicine that mimics acetylcholine and stimulates muscarinic receptors. In selected eye procedures, it constricts the pupil and can lower pressure inside the eye; effects may include sweating, slow heartbeat, or breathing problems.

\medterm{Carbadox} An antimicrobial formerly used to promote growth and prevent disease in food-producing animals. Concerns about toxic breakdown products and possible cancer risk have led to restrictions or bans in many places.

\medterm{Carbamate} A chemical group found in some medicines and pesticides. Carbamate insecticides can temporarily block acetylcholinesterase, causing excess nerve activity with sweating, salivation, vomiting, weakness, wheezing, or seizures.

\medterm{Carbamazepine} A medicine used for focal seizures, trigeminal neuralgia, and selected bipolar disorder. Important risks include severe skin reactions, low blood counts, liver injury, low sodium, dizziness, and interactions with many other medicines.

\medterm{Carbamide} Another name for urea, a waste product made when the body breaks down protein and removed mainly by the kidneys. Urea is also used in some skin creams to soften and hydrate thick or dry skin.
\medterm{Carbapenem} A group of powerful beta-lactam antibiotics that kill bacteria by blocking construction of their cell wall. They are reserved mainly for serious infections because resistance and important side effects can limit treatment choices.

\medterm{Carbapenem Resistant Enterobacteriaceae (CRE Infection)} Alternate index label for CRE Infection.

\medterm{Carbapenem-Resistant Enterobacteriaceae Infection} Infection caused by Enterobacterales bacteria that resist carbapenem antibiotics. It may affect the urinary tract, lungs, blood, or wounds, can spread in health-care settings, and requires culture and susceptibility testing to select treatment.

\medterm{Carbapenem-Resistant Organisms} Bacteria that cannot be reliably treated with carbapenem antibiotics, often because they produce enzymes that destroy these medicines. They can cause difficult infections and spread in health-care settings, where strict infection control is important.

\medterm{Carbaryl} A carbamate insecticide that temporarily blocks acetylcholinesterase. Substantial exposure can cause sweating, excess saliva, vomiting, diarrhea, small pupils, muscle weakness, wheezing, seizures, or breathing failure.

\medterm{Carbazilquinone} An experimental quinone-containing compound studied for its ability to damage cancer cells, especially under low-oxygen conditions. It is a research substance and is not a routine clinical medicine.

\medterm{Carbenicillin} A penicillin antibiotic formerly used for selected serious infections caused by gram-negative bacteria. Resistance, allergic reactions, diarrhea, and kidney or electrolyte problems can limit its use.

\medterm{Carbidopa} A medicine given with levodopa for Parkinson disease. It blocks conversion of levodopa to dopamine outside the brain, allowing more levodopa to reach the brain and reducing nausea, vomiting, and low blood pressure.

\medterm{Carbimazole} An antithyroid medicine that the body converts to methimazole. It reduces production of thyroid hormones and can rarely cause dangerously low white-cell counts or liver injury.
\medterm{Carbofuran} A highly toxic carbamate pesticide that blocks acetylcholinesterase. Ingestion, inhalation, or skin exposure can rapidly cause sweating, vomiting, muscle weakness, seizures, breathing failure, and death.

\medterm{Carbogen} A controlled breathing mixture containing oxygen and a small amount of carbon dioxide. It has limited specialized or experimental uses to stimulate breathing or alter tissue oxygenation and is not for unsupervised use.

\medterm{Carbohydrate} A nutrient found in sugars, starches, fruits, grains, vegetables, and milk. The body breaks digestible carbohydrates into glucose for energy; fiber is a carbohydrate that is not fully digested.

\medterm{Carbohydrate Counting} A way to plan meals by estimating the grams of carbohydrate in food. People using insulin may match their insulin dose to the amount eaten; accurate reading of food labels and advice from a diabetes team are important.

\medterm{Carbohydrate Metabolism} The processes cells use to break down, store, and make sugars. These processes provide energy and help keep blood glucose stable; problems can occur in diabetes and inherited disorders of sugar storage.

\medterm{Carbohydrates} Sugars, starches, and fiber that form one of the main types of nutrients in food. Digestible carbohydrates are broken down into glucose for energy, while fiber is not fully digested and helps support bowel function and blood-sugar control.

\medterm{Carbolic Acid (Phenol)} Phenol, a chemical once used as a surgical antiseptic and still found in some industrial products. It kills microorganisms by damaging proteins and cell membranes but can severely burn skin, eyes, and internal organs.

\synonyms
Phenol

\medterm{Carbolic Acid Poisoning} Exposure to phenol, a corrosive and systemically toxic chemical that can cause burns, altered consciousness, metabolic acidosis, and organ injury.

\medterm{Carbon} A chemical element forming the backbone of organic molecules such as carbohydrates, fats, proteins, and nucleic acids.

\medterm{Carbon Dioxide} A gas produced when cells use energy and removed mainly by the lungs during exhalation. The body controls its blood level through breathing; excess can cause headache, confusion, drowsiness, or respiratory failure.
\medterm{Carbon Dioxide Absorber} A device in some anaesthesia breathing systems that removes carbon dioxide from exhaled air. It allows selected gases to be reused while helping prevent carbon-dioxide build-up during controlled breathing.

\medterm{Carbon Dioxide Transport} The movement of carbon dioxide from body tissues to the lungs. Blood carries it dissolved in plasma, attached to haemoglobin, and mainly as bicarbonate; the lungs remove it during exhalation.

\medterm{Carbon Ion Therapy} A form of particle radiation treatment that uses carbon ions to damage tumor DNA. It may deliver concentrated energy and has higher biological effectiveness than conventional X-rays for some tumors, but availability is limited.

\medterm{Carbon Monoxide} A colorless, odorless gas that binds hemoglobin and prevents blood from delivering enough oxygen. Poisoning may cause headache, dizziness, nausea, confusion, collapse, or death and requires immediate fresh air and emergency care.
\medterm{Carbon Monoxide Poisoning} Poisoning from breathing carbon monoxide, which prevents blood from carrying and delivering oxygen. It may cause headache, dizziness, nausea, confusion, fainting, seizures, or death; leave the area and obtain emergency care immediately.

\medterm{Carbon Tetrachloride} A toxic industrial solvent that can damage the liver and kidneys when swallowed or inhaled. Suspected exposure requires urgent contact with poison-control or emergency services.

\medterm{Carbonate} A chemical form of carbon dioxide found in minerals and body fluids. Carbonate salts are used in some antacids and supplements; their effects, absorption, and risks depend on the particular salt.

\medterm{Carbonic Anhydrase Inhibitors} Medicines that block carbonic anhydrase, an enzyme involved in fluid and acid balance. They reduce fluid production in the eye to lower eye pressure and have selected uses for swelling, seizures, or metabolic disorders.

\medterm{Carbophenothion} An organophosphate insecticide that blocks acetylcholinesterase. Significant exposure can cause sweating, excess saliva, vomiting, muscle weakness, confusion, seizures, breathing failure, and death.

\medterm{Carboplatin} A platinum chemotherapy medicine that damages cancer-cell DNA and prevents it from being copied. It can lower blood counts, especially platelets, and may cause nausea, anemia, kidney effects, nerve symptoms, or allergic reactions.

\medterm{Carboprost} A prostaglandin medicine that strongly contracts the uterus. It is used in selected cases of severe bleeding after childbirth and for some pregnancy-related procedures; it can cause diarrhea, vomiting, fever, and dangerous airway tightening in people with asthma.

\medterm{Carbosulfan} A carbamate insecticide that blocks acetylcholinesterase and overstimulates nerves. Exposure can cause sweating, salivation, vomiting, diarrhea, small pupils, muscle weakness, wheezing, seizures, or breathing failure.

\medterm{Carboxy-Lyases} Enzymes that break a carbon-carbon bond and release carbon dioxide. They participate in many biochemical reactions, including processes that make or break down amino acids and other molecules.

\medterm{Carboxyhaemoglobin} Hemoglobin bound to carbon monoxide. This reduces the blood's ability to carry and release oxygen; an increased level supports recent carbon-monoxide exposure but must be interpreted with symptoms and timing.
\medterm{Carboxyhemoglobin} A compound formed when carbon monoxide binds to hemoglobin. It prevents hemoglobin from carrying and releasing oxygen normally, causing headache, dizziness, confusion, collapse, or death in severe poisoning.

\medterm{Carboxyhemoglobinemia} An abnormally high level of carboxyhemoglobin in the blood after carbon-monoxide exposure. It reduces oxygen delivery and may cause headache, nausea, confusion, fainting, heart injury, brain injury, or death.

\medterm{Carboxyl Group} A chemical group containing carbon, oxygen, and hydrogen that gives many molecules acidic properties. It is found in amino acids, fatty acids, and numerous medicines and metabolic compounds.
\medterm{Carboxylesterase} An enzyme that breaks ester bonds in medicines, toxins, and natural fats. Differences in its activity can change how quickly some substances are broken down and can affect their effects, interactions, or toxicity.

\medterm{Carboxylic Acid} An organic compound containing a carboxyl group, which gives it acidic properties. Many important fatty acids, amino-acid products, and substances involved in metabolism are carboxylic acids.

\medterm{Carboxypeptidase H} An enzyme inside cells that helps convert precursor proteins into active peptide hormones and nerve-signaling molecules by removing amino acids from one end.

\medterm{Carbuncle} A group of connected boils, usually caused by Staphylococcus aureus, forming a painful swollen area with several pus-draining openings. Fever may occur, and risk is increased by diabetes or a weakened immune system.

\medterm{Carbuncles And Boils} Alternate index label; see \textit{Boils and carbuncles}.

\medterm{Carcinoembryonal Antigen} Another name for carcinoembryonic antigen, or CEA, a blood marker that may rise in colorectal and other cancers. It can also increase for noncancerous reasons and is not suitable for general cancer screening.

\medterm{Carcinoembryonic Antigen} A protein measured in blood that may be increased in colorectal and some other cancers. It can also rise with smoking, liver disease, bowel inflammation, or other noncancerous conditions, so it is mainly used to monitor selected known cancers rather than to screen for cancer.
\medterm{Carcinogen} A substance, form of radiation, or infection that can cause cancer by damaging cells or changing how they grow. Risk depends on the type and amount of exposure.
\medterm{Carcinogen Testing} Laboratory, animal, or population research used to determine whether a chemical, radiation exposure, infection, medicine, or other agent can cause or promote cancer.

\medterm{Carcinogenesis} The multistep process by which normal cells acquire changes that allow uncontrolled growth and become cancer. Changes may arise from inherited traits, random errors, infections, tobacco, radiation, chemicals, or other exposures and often accumulate over time.

\medterm{Carcinogenic} Capable of causing or promoting cancer by damaging genetic material or altering how cells grow, divide, repair themselves, or avoid cell death.
\medterm{Carcinoid} A usually slow-growing neuroendocrine tumor that can arise in the digestive tract, appendix, or lungs. Some release hormones or other substances that cause flushing, diarrhea, wheezing, or heart-valve disease.
\medterm{Carcinoid Heart Disease} Damage to the inner lining and valves of the heart caused by hormone-like substances released by a neuroendocrine tumor. It mainly affects the tricuspid and pulmonary valves and can lead to right-sided heart failure.

\medterm{Carcinoid Syndrome} A group of symptoms caused by hormone-like substances released by some neuroendocrine tumors, especially after spread to the liver. It may cause facial flushing, diarrhea, wheezing, rapid heartbeat, and eventually heart-valve disease.

\medterm{Carcinoid Tumor} A usually well-differentiated neuroendocrine tumor arising in hormone-producing cells, commonly in the appendix, small intestine, rectum, or lungs. Its behavior ranges from slow-growing to capable of spreading and causing hormone-related symptoms.

\medterm{Carcinoid Tumors} Neuroendocrine tumors that often grow slowly and may release hormones or other active substances. They commonly arise in the digestive tract or lungs and can sometimes invade or spread to distant organs.

\synonyms
Cancer, Carcinoid Tumors

\medterm{Carcinoma} A cancer that begins in epithelial cells, which cover body surfaces and line organs and glands. Common types include adenocarcinoma, squamous-cell carcinoma, basal-cell carcinoma, and urothelial carcinoma; behavior depends on site and cell features.

\medterm{Carcinoma En Cuirasse} An advanced form of breast cancer in which cancer cells spread through lymphatic vessels in the skin, causing widespread hardening and thickening that can make the chest wall resemble armor.

\medterm{Carcinoma Erysipeloides} Spread of an internal cancer through lymph vessels in the skin, producing a red, firm, sharply outlined area that can resemble the infection erysipelas. It most often occurs with cancers of the breast or digestive tract.

\medterm{Carcinoma In Situ} Abnormal cancer-like cells occupying the full thickness of an epithelial lining but not yet invading through its supporting boundary. It is an early, noninvasive stage that may progress if untreated.

\medterm{Carcinoma Of Unknown Primary} Cancer found outside its original site when the starting organ cannot be identified despite evaluation. Treatment uses the tumor’s appearance, genetic features, spread pattern, symptoms, and likely tissue of origin.

\synonyms
Cancer Of Unknown Origin
Occult Primary Cancer

\medterm{Carcinomatosis} Widespread spread of carcinoma to multiple surfaces or organs, often along the lining of the abdomen, chest, or other body cavities. It usually indicates advanced cancer and may cause fluid buildup, pain, bowel problems, or breathing difficulty.

\medterm{Carcinosarcoma} A malignant tumor containing both carcinomatous and sarcomatous elements, most commonly
of the uterus (malignant mixed Müllerian tumor), with an aggressive course.

\medterm{Cardia} The upper part of the stomach surrounding the opening where the esophagus enters. It helps receive swallowed food and can be affected by reflux, inflammation, hiatal hernia, or stomach cancer.

\medterm{Cardiac} Relating to the heart. The word may describe the heart muscle, chambers, valves, electrical system, blood supply, rhythm, or diseases such as coronary disease, cardiomyopathy, and heart failure; in “cardia,” it instead refers to the upper stomach.

\synonyms
Heart-related; cardiac-related.

\medterm{Cardiac Ablation} A procedure that uses heat, cold, or another form of energy to create a small scar in heart tissue causing an abnormal rhythm. The scar blocks faulty electrical signals and may restore a regular rhythm.

\medterm{Cardiac Ablation Procedures} Procedures that use heat, cold, or another energy source to treat heart-rhythm disorders by creating controlled scars that block abnormal electrical signals. Risks include bleeding, blood-vessel injury, recurrence, and damage to nearby heart structures.

\medterm{Cardiac Allograft} A donor heart transplanted into another person, usually for end-stage heart failure. Lifelong immune-suppressing treatment is needed to reduce rejection; complications include infection, rejection, coronary-vessel disease, cancer, and graft failure.

\medterm{Cardiac Allograft Vasculopathy} Progressive narrowing of the heart's blood vessels after transplantation caused by abnormal thickening of their inner walls. It may silently reduce blood flow and cause graft failure, rhythm problems, heart attack, or sudden death.

\medterm{Cardiac Amyloidosis} Buildup of abnormal amyloid protein in the heart, making the heart muscle stiff and impairing filling. It can cause breathlessness, swelling, rhythm problems, or heart failure; treatment depends on whether the protein is light-chain or transthyretin.

\medterm{Cardiac Amyloidosis Diagnosis} Evaluation combines symptoms, examination, blood and urine tests, heart imaging, and sometimes tissue biopsy. Specialized nuclear imaging may support transthyretin amyloidosis only after testing has excluded light-chain amyloidosis.

\medterm{Cardiac Amyloidosis Staging} A system for estimating the severity and outlook of transthyretin cardiac amyloidosis using markers of heart strain and injury, sometimes with other findings. Higher stages indicate more advanced disease; treatment and prognosis must be individualized.

\medterm{Cardiac Arrest} Sudden stopping of effective heart pumping, causing unresponsiveness and absent or abnormal breathing. It is an immediate emergency requiring cardiopulmonary resuscitation, defibrillation when appropriate, and advanced life support.

\medterm{Cardiac Arrest Survival} The chance of surviving cardiac arrest depends on its cause, where it occurs, the initial rhythm, how quickly CPR and defibrillation begin, and care afterward. Early recognition, effective CPR, rapid defibrillation, and post-arrest treatment improve outcomes.

\medterm{Cardiac Arrest, Sudden} Alternate index label; see \textit{Sudden cardiac arrest}.

\medterm{Cardiac Arrhythmias} Abnormal heart rates or rhythms caused by problems forming or conducting electrical impulses. They include rhythms that are too fast, too slow, or irregular and may cause palpitations, dizziness, fainting, or no symptoms.

\medterm{Cardiac Axis} The average direction in which the heart's lower chambers become electrically activated, estimated from an electrocardiogram. A leftward or rightward shift may reflect normal variation or heart, conduction, or lung disease.

\medterm{Cardiac Biomarkers} Substances measured in blood to detect or assess heart-muscle injury. Troponin is the preferred marker; results must be interpreted with symptoms, electrocardiography, imaging, and changes in repeated tests because many conditions can raise levels.

\medterm{Cardiac Cachexia} Severe unintentional loss of weight and muscle associated with advanced heart failure or other serious heart disease. It can cause weakness and reduced function and reflects inflammation, altered metabolism, poor appetite, and increased nutritional needs.

\medterm{Cardiac Care} Prevention, diagnosis, treatment, and follow-up for heart and circulation disorders. It may include risk reduction, medicines, monitoring, procedures, exercise rehabilitation, nutrition support, and care for symptoms and quality of life.

\medterm{Cardiac Catheter} A thin flexible tube inserted through a blood vessel and guided to the heart. It can measure pressures, take blood samples, inject contrast for imaging, or deliver treatment.

\medterm{Cardiac Catheterization} A procedure in which thin tubes are guided through a blood vessel to the heart to measure pressures, assess oxygen levels, inject contrast, image coronary arteries, or perform treatment. Possible risks include bleeding, rhythm problems, vessel injury, and kidney effects.


\medterm{Cardiac Conduction System} Specialized heart tissue that creates and carries electrical signals so the chambers contract in the correct order. It includes the sinoatrial node, atrioventricular node, bundle of His, bundle branches, and Purkinje fibers.

\medterm{Cardiac CT} Computed-tomography imaging of the heart and nearby blood vessels. Depending on the protocol, it can assess coronary calcium, coronary narrowing, heart structure, valves, and selected aspects of heart function.

\medterm{Cardiac CT Angiography} A noninvasive scan using contrast and computed tomography to show the coronary arteries and heart structures. It is useful for assessing selected people with chest pain and can show narrowing, plaque, and coronary calcium.

\medterm{Cardiac Cycle} The events in one complete heartbeat. During diastole the heart chambers relax and fill with blood; during systole the ventricles contract and pump blood into the lungs and the rest of the body.

\medterm{Cardiac Defibrillation} Delivery of an electrical shock to stop ventricular fibrillation or pulseless ventricular tachycardia during cardiac arrest. It must be combined with high-quality CPR and advanced resuscitation according to the rhythm and situation.

\medterm{Cardiac Device Infections} Infection involving a pacemaker, implantable defibrillator, its pocket, or a lead. It may cause redness, pain, drainage, fever, bloodstream infection, or infection of the heart lining and often requires antibiotics and removal of infected hardware.

\medterm{Cardiac Diet} An eating pattern emphasizing vegetables, fruits, whole grains, beans, fish, nuts, and unsaturated fats while limiting excess sodium, added sugar, saturated fat, and trans fat. It should be adapted for heart failure, kidney disease, diabetes, and nutrition needs.

\medterm{Cardiac Enzymes} An older term for blood tests used to detect heart-muscle injury, especially creatine kinase-MB. Troponin is now preferred; any result must be interpreted with symptoms, electrocardiography, imaging, and changes over time.

\medterm{Cardiac Event Monitors} Portable devices that record or send an electrocardiogram when symptoms occur or at planned times. They help detect intermittent abnormal rhythms that may be missed during a short clinic ECG.

\medterm{Cardiac Failure} Alternate index label; see \textit{Heart failure}.
\medterm{Cardiac Fibroma} A rare noncancerous tumor made of fibrous tissue in the heart, usually in a ventricle or septum. It may obstruct blood flow, interfere with electrical signals, or cause rhythm problems and heart failure.

\medterm{Cardiac Gated SPECT} A heart scan using a small radioactive tracer while images are synchronized with the heartbeat. It can assess blood flow to heart muscle, pumping strength, wall movement, and areas that are damaged or still viable.

\medterm{Cardiac Glycoside} A medicine or plant substance that can strengthen heart contraction and slow electrical conduction through the heart. Digoxin is the main clinical example; the safe dose is narrow and toxicity can cause dangerous rhythms.

\medterm{Cardiac Glycoside Overdose} Toxic exposure to a cardiac glycoside such as digoxin. It may cause nausea, vomiting, confusion, visual changes, slow or fast abnormal rhythms, high potassium, seizures, or cardiac arrest and requires urgent assessment.

\medterm{Cardiac Hemangioma} A rare noncancerous tumor made of blood-vessel tissue in or around the heart. Many cause no symptoms, but a large or strategically located tumor may obstruct flow, disturb rhythm, or cause fluid around the heart.

\medterm{Cardiac Imaging Tests} Tests that show the heart's structure, movement, blood flow, or tissue. They include echocardiography, CT, MRI, and nuclear scans; the appropriate test depends on symptoms, the suspected condition, kidney function, and the clinical question.

\medterm{Cardiac Index} The amount of blood the heart pumps each minute adjusted for body size. It helps assess pumping function, especially in seriously ill people, but a single value must be interpreted with blood pressure, symptoms, organ function, and the overall clinical situation.

\medterm{Cardiac Intravascular Ultrasound} An ultrasound probe on a catheter that creates cross-sectional images from inside a coronary artery. It helps measure plaque, vessel size, narrowing, tears, calcium, and stent expansion during selected procedures.

\medterm{Cardiac Ischemia} Alternate index label; see \textit{Myocardial ischemia}.

\medterm{Cardiac Massage} Manual compression of the heart, historically performed through the chest or during open surgery. Modern cardiac-arrest treatment uses external chest compressions and advanced resuscitation rather than routine direct heart massage.
\medterm{Cardiac Metastases} Cancer that has spread to the heart or its surrounding tissues from another organ. Lung and breast cancers, melanoma, lymphoma, and kidney cancer are common sources; spread may cause fluid around the heart, obstruction, rhythm problems, or heart failure.

\medterm{Cardiac MRI} Magnetic-resonance imaging that shows heart structure, pumping, blood flow, scar, swelling, infiltrative disease, and congenital abnormalities without x-rays. Contrast may be avoided or used cautiously in selected people with severe kidney disease.

\medterm{Cardiac Muscle} Involuntary muscle that forms the myocardium, the main muscular layer of the heart. Its cells are electrically connected so the chambers contract in a coordinated way to pump blood.

\medterm{Cardiac Myocyte} A specialized heart-muscle cell that contracts to pump blood. Electrical signals and movement of calcium inside the cell coordinate its contraction and relaxation.

\medterm{Cardiac Myocytes} The muscle cells that make up most of the heart and contract with each heartbeat. Damage to them can reduce pumping function and release proteins such as troponin into the blood.

\medterm{Cardiac Myxoma} The most common primary noncancerous tumor of the heart, usually arising in the left atrium. It can obstruct blood flow, cause breathlessness or constitutional symptoms, or release clots; surgical removal is usually recommended.

\medterm{Cardiac Output} The volume of blood pumped by the heart each minute. It depends mainly on heart rate and stroke volume, the amount ejected with each beat, and helps indicate whether circulation is adequate.

\medterm{Cardiac Output Measurement} Measurement of blood pumped by the heart each minute using methods such as echocardiography, a pulmonary-artery catheter, oxygen-consumption calculations, or noninvasive monitors. The method's accuracy depends on the patient and testing conditions.

\medterm{Cardiac Pacemaker} An implanted or temporary device that sends electrical impulses to maintain an adequate heartbeat when the heart's natural rhythm is too slow or unreliable. It may be placed inside the body or used externally during emergencies.
\medterm{Cardiac Pacemaker Types} Single-chamber pacemakers stimulate one heart chamber; dual-chamber devices coordinate the atrium and ventricle; biventricular or resynchronization devices stimulate both ventricles to improve timing in selected heart-failure patients.

\medterm{Cardiac Perforation} A hole or tear through the heart wall, sometimes caused by injury or a medical procedure. It can cause bleeding around the heart, cardiac tamponade, shock, and death and requires immediate treatment.

\medterm{Cardiac PET} A nuclear-imaging test using a positron-emitting tracer to assess heart-muscle blood flow, metabolism, and living or scarred tissue. Selected tracers can also help evaluate inflammation or infiltrative heart disease.

\medterm{Cardiac PET Imaging} Positron-emission tomography using a radioactive tracer to measure blood flow and metabolism in heart muscle. It can quantify blood flow and help distinguish living from scarred tissue, but availability, radiation exposure, and tracer choice affect its use.

\medterm{Cardiac Power Output} A measure combining blood pressure and cardiac output to estimate how effectively the heart is pumping. It is useful in severe heart failure or cardiogenic shock, but clinicians interpret it with examination findings, organ function, and other circulation measurements.

\synonyms
CPO; Cardiac Power Index (CPI)

\medterm{Cardiac Radiation Injury} Damage to the heart or its blood vessels that may appear months or years after radiation to the chest. It can cause inflammation or scarring around the heart, coronary disease, valve disease, heart-muscle weakness, or rhythm and conduction problems.

\medterm{Cardiac Radionuclide Imaging} Heart imaging performed after a small radioactive tracer is given, usually with SPECT or PET. It can assess blood flow, metabolism, living or scarred heart muscle, inflammation, and, when synchronized with the heartbeat, pumping function.

\medterm{Cardiac Rehab Outcomes} Cardiac rehabilitation can improve exercise ability, blood-pressure and cholesterol control, confidence, quality of life, and psychological well-being after selected heart conditions. Benefits vary with the program, health condition, and participation.

\medterm{Cardiac Rehabilitation} A medically supervised program after conditions such as heart attack, coronary procedures, bypass surgery, or heart failure. It combines gradual exercise, education, risk-factor reduction, medicine support, nutrition, and emotional care to improve recovery and reduce future risk.

\medterm{Cardiac Resynchronization Therapy} A pacing treatment that coordinates contraction of the heart's lower chambers. It is used in selected people with symptomatic heart failure and delayed electrical activation and can improve symptoms, pumping, and survival when appropriate.

\medterm{Cardiac Rhabdomyoma} A usually noncancerous tumor of heart muscle, most often found in children and associated with tuberous sclerosis. It may shrink on its own but can obstruct blood flow or disturb rhythm and pumping if large.

\medterm{Cardiac Sarcoidosis} Sarcoidosis affecting the heart, where clusters of inflammatory cells can form in heart muscle and the electrical system. It may cause heart block, dangerous ventricular rhythms, heart failure, or sudden cardiac arrest.

\medterm{Cardiac Sarcoma} A rare cancer arising from supporting or connective tissue of the heart. It may obstruct blood flow, disturb heart rhythm, impair pumping, or release clots that travel to other parts of the body.

\medterm{Cardiac Skeleton} A strong connective-tissue framework around the heart valves and parts of the wall between the chambers. It supports the valves, provides attachment for heart muscle, and helps electrically separate the upper and lower chambers.

\medterm{Cardiac SPECT} A nuclear scan that shows blood flow to the heart muscle and can help identify reduced flow, prior damage, or living muscle. When synchronized with the heartbeat, it can also assess chamber movement and pumping function.

\medterm{Cardiac Stress Test} A test that observes the heart while it works harder, usually during exercise or after a medicine increases its workload. It may include an ECG or imaging and can help detect reduced blood flow or abnormal rhythms.

\medterm{Cardiac Stress Testing} Assessment of the heart during exercise or medicine-induced stress to look for reduced blood flow, abnormal rhythms, or limited pumping reserve. The method is chosen according to the person's ability to exercise and the clinical question.

\medterm{Cardiac Syncope} Transient loss of consciousness caused by a cardiac disorder that reduces cerebral perfusion, usually through bradyarrhythmia, tachyarrhythmia, structural obstruction, or severe cardiac dysfunction. Syncope during exertion, while supine, or with palpitations or known heart disease requires urgent cardiovascular evaluation because it may indicate a risk of sudden cardiac death.

\medterm{Cardiac Tamponade} A medical emergency in which blood or fluid builds up under pressure around the heart and prevents the chambers from filling normally. It reduces blood flow to the body and can cause shock or cardiac arrest.

\medterm{Cardiac Tamponade Physiology} Pressure from blood or fluid around the heart prevents the chambers from relaxing and filling. This lowers cardiac output and may cause low blood pressure, a marked breathing-related change in pulse pressure, and shock.

\medterm{Cardiac Transplantation} Replacement of a failing heart with a donor heart in carefully selected people with end-stage heart failure. Lifelong immune-suppressing treatment and regular monitoring are needed for rejection, infection, coronary-vessel disease, and other complications.

\medterm{Cardiac Tumors} Tumors involving the heart may begin there or spread from another organ. They are rare and may obstruct blood flow, disturb rhythm, cause fluid around the heart, or release clots; treatment depends on type, location, size, and spread.

\medterm{Cardiac Tumors Management} Management depends on whether a tumor is noncancerous or cancerous, its size, location, symptoms, and spread. Observation may suit some harmless tumors; surgery, medicines, radiation, or chemotherapy may be needed for others.

\medterm{Cardiac Ventriculography} Imaging of the heart's lower chambers to assess their size, wall movement, and ability to pump blood. It may use x-rays with contrast or another imaging method and is now often replaced by echocardiography or MRI.

\medterm{Cardiac Volume} The amount of blood in a heart chamber at a particular time. The amount before contraction and after contraction helps assess chamber size, filling, and pumping function.

\medterm{Cardiac Wall} The three-layered wall of the heart: the inner endocardium, muscular myocardium, and outer epicardium. It supports the heart's structure, electrical activity, contraction, and interaction with surrounding tissues.
\medterm{Cardinal Body Part} A major body region or structure used in anatomical description or clinical examination, such as the head, neck, chest, abdomen, pelvis, limbs, or an important internal organ.

\medterm{Cardinal Cell Part} A main component of a cell, such as its membrane, cytoplasm, nucleus, or an organelle. Each part has a role in maintaining the cell, processing information, or performing specialized functions.

\medterm{Cardinal Ligament} A band of strong connective tissue extending from the cervix and upper vagina to the sidewall of the pelvis. It helps support the uterus and vagina and contains the uterine blood vessels, which cross over the ureter.

\synonyms
Mackenrodt's Ligament; Lateral Cervical Ligament; Transverse Cervical Ligament

\medterm{Cardio-Oncology} The medical field focused on preventing and treating heart and blood-vessel problems in people with cancer. It monitors effects of chemotherapy, targeted or immune treatment, and chest radiation while helping cancer therapy continue safely when possible.

\medterm{Cardiobacteriaceae} A family of bacteria that includes Cardiobacterium. Some members can enter the bloodstream and infect heart valves, causing infective endocarditis.

\medterm{Cardiobacterium} A slow-growing bacterium in the HACEK group that can enter the bloodstream and infect heart valves. It may cause prolonged fever, fatigue, weight loss, a new murmur, or embolic complications.

\medterm{Cardiobacterium Hominis} The Cardiobacterium species most often associated with human infection. It can cause slowly developing infective endocarditis with fever, fatigue, weight loss, a new heart murmur, valve damage, or clots traveling to other organs.

\medterm{Cardiofaciocutaneous Syndrome} A rare genetic disorder affecting the heart, facial appearance, skin, hair, growth, and development. It is usually caused by a change in a gene involved in cell-growth signals.
\synonyms
CFC Syndrome

\medterm{Cardiogenic Shock} A state of inadequate tissue perfusion caused by severe cardiac dysfunction. It may produce
low blood pressure, altered mental status, cool extremities, reduced urine output, metabolic acidosis, or other signs of
organ hypoperfusion.

\medterm{Cardiogenic Shock Management} Emergency management of shock caused by severe cardiac dysfunction. It includes
support of oxygenation and circulation, treatment of the underlying cause, and selected use of vasopressors, inotropes,
mechanical support, or revascularization.

\medterm{Cardiolipin} A fat-like molecule found mainly in the inner membrane of mitochondria, the cell's energy-producing structures. Antibodies against cardiolipin can be associated with antiphospholipid syndrome and abnormal blood clotting.

\medterm{Cardiologist} A doctor trained to diagnose and treat diseases of the heart and blood vessels, including coronary disease, rhythm disorders, valve disease, heart failure, and high blood pressure.

\medterm{Cardiology} The branch of medicine concerned with the heart and blood vessels. It covers prevention, diagnosis, and treatment of coronary disease, heart failure, rhythm and valve disorders, cardiomyopathy, congenital heart disease, and high blood pressure.

\medterm{Cardiomegaly} Enlargement of the heart seen on examination or imaging. It may result from enlarged chambers, thickened heart muscle, fluid around the heart, or a combination; causes include high blood pressure, valve disease, cardiomyopathy, and congenital disease.

\medterm{Cardiomyopathy} Disease of the heart muscle that can impair contraction, relaxation, electrical rhythm, or valve function and may lead to heart failure or sudden death. Major patterns include dilated, hypertrophic, and restrictive cardiomyopathy.

\medterm{Cardiomyopathy (Restrictive)} Alternate index label for Restrictive.

\medterm{Cardiomyopathy Classification} A system for describing cardiomyopathy by its heart structure and function, other organs involved, genetic pattern, cause, and functional severity. This information helps communicate diagnosis, guide testing, and estimate prognosis.

\medterm{Cardiomyopathy, Dilated} Alternate index label; see \textit{Dilated cardiomyopathy}.

\medterm{Cardiomyopathy, Hypertrophic} Alternate index label; see \textit{Hypertrophic cardiomyopathy}.

\medterm{Cardiopathy} Any disease or disorder of the heart, including problems of the heart muscle, valves, coronary arteries, electrical rhythm, or structure present from birth.

\medterm{Cardioplegia} Temporary stopping and protection of the heart during open-heart surgery, usually by delivering a cold solution that changes electrical and metabolic activity. It allows the surgeon to operate while limiting heart-muscle injury.

\medterm{Cardiopulmonary} Relating to the heart and lungs or to their combined function. For example, cardiopulmonary exercise testing assesses how the heart, lungs, blood, and muscles respond during physical activity.

\synonyms
Cardiorespiratory.

\medterm{Cardiopulmonary Bypass} A heart-lung machine temporarily takes over circulation and oxygen exchange during some operations. It pumps and oxygenates blood while the heart is stopped or opened; possible complications include bleeding, inflammation, clots, stroke, and organ injury.

\medterm{Cardiopulmonary Bypass Machine} A machine that temporarily pumps and oxygenates blood during selected heart or major-vessel operations. It allows surgeons to work while the heart is stopped or the lungs are not carrying out their usual gas exchange.

\synonyms
Heart-lung machine; cardiopulmonary bypass pump.

\medterm{Cardiopulmonary Exercise Testing} A monitored treadmill or bicycle test that measures how the heart, lungs, blood, and muscles respond as exercise becomes harder. It helps assess unexplained breathlessness, exercise limitation, and severity of heart or lung disease.

\medterm{Cardiopulmonary Resuscitation} Emergency chest compressions and, when appropriate, rescue breaths, defibrillation, and advanced life support for cardiac arrest. Its purpose is to maintain blood flow and oxygen delivery until the heartbeat returns or definitive treatment is provided.

\medterm{Cardiorenal Syndrome} A condition in which heart and kidney dysfunction worsen one another. Sudden or long-term heart disease can injure the kidneys, and sudden or long-term kidney disease can strain the heart; treatment must address both organs.

\medterm{Cardiorespiratory} Relating to the heart, blood vessels, and respiratory system considered together, as in cardiorespiratory fitness, monitoring, or failure.
\medterm{Cardiorespiratory Endurance} The ability of the heart, lungs, blood, and muscles to deliver and use oxygen during sustained physical activity. It is influenced by fitness, heart and lung disease, anemia, muscle condition, and other factors.
\synonyms
Cardiorespiratory fitness; aerobic endurance.

\medterm{Cardioselective Beta-Blockers} Beta-blocker medicines that mainly block beta-1 receptors in the heart, slowing the heart rate and reducing its workload. They may still narrow airways at higher doses and are used for selected heart, blood-pressure, and rhythm conditions.

\medterm{Cardiospasm} An older term for failure of the lower esophageal muscle to relax normally, as occurs in achalasia. It can cause difficulty swallowing, regurgitation of food or liquid, chest discomfort, and weight loss.

\medterm{Cardiospasm (Achalasia)} Alternate index label for Achalasia.

\medterm{Cardiotocography} Continuous recording of a fetus's heart rate and the mother's uterine contractions during pregnancy or labor. The pattern helps assess fetal well-being, but results must be interpreted with the stage of labor and the overall clinical situation.

\medterm{Cardiotoxicity} Injury or impaired function of the heart caused by a medicine, toxin, radiation, or disease. It may produce rhythm problems, reduced heart-muscle strength, chest pain from poor blood flow, or heart failure.

\medterm{Cardiotoxicity From Chemotherapy} Heart injury caused by some chemotherapy medicines. It may weaken heart muscle, reduce pumping, or cause rhythm problems during or after treatment; risk depends on the medicine, dose, other treatments, and existing heart disease.

\medterm{Cardiotoxin} A substance that directly injures heart muscle or disrupts the heart's electrical activity or blood supply. It may cause abnormal rhythms, reduced pumping, heart failure, or heart-muscle injury.

\medterm{Cardiovascular} Relating to the heart and blood vessels, including circulation, blood pressure, and diseases affecting the heart or vascular system.

\medterm{Cardiovascular Agents} Medicines that act on the heart or blood vessels. They include treatments for high blood pressure, abnormal rhythms, angina, heart failure, blood clots, and abnormal cholesterol levels.

\medterm{Cardiovascular Changes In Pregnancy} During pregnancy, blood volume and cardiac output increase to support the uterus and developing baby. The rise in plasma is greater than the rise in red cells, causing normal dilutional anemia; heart rate and blood-vessel resistance also change.

\medterm{Cardiovascular Disease} Disease of the heart or blood vessels, including coronary artery disease, heart failure, abnormal rhythms, valve disease, stroke-related vascular disease, and disorders of arteries or veins.

\medterm{Cardiovascular Health} The condition of the heart and blood vessels and the habits that protect them. Regular activity, not using tobacco, healthy eating, adequate sleep, and control of blood pressure, cholesterol, diabetes, and weight can lower cardiovascular risk.

\medterm{Cardiovascular Infection} An infection involving the bloodstream, heart, heart valves, or blood vessels. Examples include infective endocarditis, myocarditis, infected vascular grafts, and infected pacemakers or other devices.

\medterm{Cardiovascular Injury} Damage to the heart or blood vessels caused by trauma, reduced blood flow, inflammation, toxins, surgery, or another disease. Effects may include bleeding, impaired circulation, heart-muscle injury, or abnormal rhythms.

\medterm{Cardiovascular Neoplasm} A noncancerous or cancerous tumor arising in the heart, blood vessels, or related tissues. Its symptoms and treatment depend on the tissue type, location, size, effect on circulation, and whether it has spread.

\medterm{Cardiovascular Pharmacology} The study of how medicines affect the heart, blood vessels, blood pressure, rhythm, blood clotting, and cholesterol or lipid metabolism, including their uses, interactions, benefits, and risks.

\medterm{Cardiovascular Risk Reduction} Actions that lower the chance of heart attack, stroke, heart failure, and vascular death. They include avoiding tobacco, controlling blood pressure, cholesterol, and diabetes, staying active, eating well, maintaining a healthy weight, and taking indicated medicines.

\medterm{Cardiovascular Screening} Assessment for cardiovascular risk or disease, often before symptoms develop. It may include blood-pressure measurement, cholesterol and diabetes testing, review of family history and tobacco use, and targeted tests based on age and risk.

\medterm{Cardiovascular Surgery} Operations on the heart or major blood vessels, including bypass surgery, valve repair or replacement, aortic surgery, and heart transplantation. Many open procedures require a heart-lung machine and specialized care before and after surgery.

\medterm{Cardiovascular System} The heart and network of arteries, arterioles, capillaries, venules, and veins that circulate blood. It delivers oxygen and nutrients, removes waste, regulates temperature, and helps maintain blood pressure.

\medterm{Cardioversion} Restoration of a more regular heart rhythm using a synchronized electrical shock or medicine. Before cardioversion for some atrial rhythms, clinicians assess clot risk and may give anticoagulant treatment to reduce the risk of stroke.

\medterm{Cardioversion (Synchronized)} Delivery of an electrical shock timed with the heart's electrical cycle to restore a more organized rhythm. It treats selected fast rhythms; sedation, clot prevention, monitoring, and energy settings depend on the rhythm and patient.

\synonyms
Synchronized

\medterm{Cardiovirus} A group of viruses in the picornavirus family. Some infect people or animals and may cause fever, muscle inflammation, heart inflammation, or other illness; effects depend on the species and the person's defenses.

\medterm{Carditis} Inflammation of the heart muscle, inner lining, valves, or surrounding sac. It may result from infection, an immune reaction, medicines, or another inflammatory disease and can cause chest pain, breathlessness, rhythm problems, or heart failure.

\medterm{Care Coordination} Organizing communication and health-care activities among clinicians, services, patients, and caregivers so care remains safe, continuous, and appropriate when people move between settings or receive several types of treatment.

\medterm{Caregiver} A family member, friend, or paid worker who helps another person with health care, personal activities, daily tasks, transportation, medicines, or emotional support.

\medterm{Caregiver Burden} The emotional, physical, social, or financial strain of providing ongoing care for another person. It may lead to exhaustion, low mood, sleep problems, isolation, or difficulty meeting one's own health needs.
\medterm{Caregiver Stress} Strain caused by caring for another person, often leading to tiredness, worry, irritability, sleep problems, or reduced well-being. Support, respite, counseling, and practical help can reduce its effects.




\medterm{Caries} A biofilm-mediated, multifactorial disease in which acids from bacterial metabolism of fermentable carbohydrates cause net demineralization of enamel, dentin, or cementum. Early lesions may be arrested or remineralized, but progressive disease forms cavities and can reach the pulp, causing pain or infection.

\medterm{Caries (Dental)} Tooth decay caused when bacteria in plaque use sugars to produce acids that remove minerals from enamel, dentin, or root surface. Repeated attacks can form cavities, pain, infection, and tooth loss; fluoride and good cleaning lower risk.

\synonyms
Dental

\medterm{Caries Risk Assessment} Evaluation of a person's chance of developing tooth decay by reviewing past and current cavities, plaque, diet, fluoride exposure, saliva, oral hygiene, tooth structure, medicines, and medical conditions. Results guide prevention and follow-up.

\medterm{Carina} The ridge at the bottom of the windpipe where it divides into the right and left main bronchi. It is a landmark during bronchoscopy and breathing-tube placement; inhaled objects more often enter the wider, more vertical right bronchus.


\medterm{Cariogenic} Likely to promote tooth decay, especially by supplying sugars that plaque bacteria convert into acids or by supporting conditions that remove minerals from teeth.

\medterm{Cariogram} A computer-based tool that combines factors such as diet, plaque bacteria, fluoride exposure, saliva, previous cavities, and social or clinical history to estimate tooth-decay risk and show where prevention may help.

\medterm{Carisoprodol} A prescription muscle relaxant used short term for acute painful muscle conditions. It can cause drowsiness, dizziness, impaired coordination, misuse, dependence, and dangerous breathing suppression, especially with alcohol or sedatives.

\medterm{Carminative} A substance traditionally used to ease intestinal gas or bloating. It may relax intestinal muscle or help gas pass, but persistent pain, vomiting, swelling, or altered bowel habits need medical assessment.

\synonyms
Antiflatulent; gas-relieving agent.

\medterm{Carmine} A red coloring made from cochineal insects and used in some foods, cosmetics, and medicines. Rarely, it causes allergy, wheezing, hives, or anaphylaxis in sensitive people.

\medterm{Carmustine} A chemotherapy medicine that damages DNA and is used for selected brain tumors, blood cancers, and other cancers. It can lower blood counts and cause nausea, infertility, liver problems, or delayed lung injury.

\medterm{Carney Complex} A rare inherited disorder, usually caused by a PRKAR1A gene change, that can produce multiple noncancerous skin, heart, and endocrine tumors. It may also cause excess hormones and requires lifelong, organ-specific monitoring.

\medterm{Carney Triad} A rare condition involving combinations of a stomach gastrointestinal stromal tumor, a benign cartilage tumor in the lung, and a paraganglioma outside the adrenal gland. It usually occurs sporadically and more often affects young women.

\medterm{Carnitine} A substance that transports long-chain fats into mitochondria, the cell structures that make energy, so the fats can be used as fuel. The body usually makes enough; severe deficiency can cause weakness, low blood sugar, or heart and muscle disease.

\medterm{Carnitine Palmitoyltransferase II Deficiency} A rare inherited disorder in which muscle cells cannot use long-chain fats normally for energy. Exercise, fasting, cold, or infection can trigger muscle pain, stiffness, muscle breakdown, dark urine, and kidney injury; symptoms may be absent between attacks.

\synonyms
CPT II Deficiency; CPT2 Deficiency

\medterm{Carnobacteriaceae} A family of bacteria that produce lactic acid and live mainly in food, animals, or the environment. Most do not cause human disease, although rare opportunistic infections have been reported.

\medterm{Carnosine} A small molecule made from beta-alanine and histidine, found mainly in muscle and nervous tissue. It helps buffer acidity; claims that supplements prevent disease or improve performance remain uncertain.

\medterm{Carnosinemia} A rare inherited disorder in which deficient carnosinase causes excess carnosine in blood and urine. Children may have developmental delay, low muscle tone, seizures, or other neurologic problems of varying severity.

\medterm{Caroli Disease} A rare inherited disorder in which the larger bile ducts inside the liver are widened in segments. It can cause repeated bile-duct infections, stones, abdominal pain, jaundice, and, in some people, liver scarring or kidney disease.

\medterm{Carotene} An orange-yellow plant pigment found in foods such as carrots and leafy vegetables. Some forms, especially beta-carotene, can be converted to vitamin A; excess intake may harmlessly turn the skin yellow-orange.

\medterm{Carotenoid} A family of fat-soluble yellow, orange, or red pigments found mainly in plants and some microorganisms. Some, such as beta-carotene, can form vitamin A; others, such as lutein, support eye tissues.

\medterm{Carotenoids} Plant pigments that give many fruits and vegetables their yellow, orange, or red color. Some can be converted to vitamin A, while others have roles in eye and cellular function.

\synonyms
Carotenoid pigments; provitamin-A pigments.

\medterm{Carotid} Relating to either of the two major arteries in the neck that carry blood from the heart toward the brain, face, and other head and neck structures.

\medterm{Carotid Artery} One of the major arteries supplying the head and neck. Each common carotid divides into an internal branch that supplies the brain and an external branch that supplies much of the face, scalp, and neck.

\medterm{Carotid Artery Disease} Narrowing or blockage of a carotid artery, usually from fatty plaque. It can reduce blood flow to the brain or release a clot and cause a transient ischemic attack or ischemic stroke.

\medterm{Carotid Artery Stenting} A procedure that places a small mesh tube inside a narrowed carotid artery to keep it open. It is used selectively to reduce stroke risk; clot release, stroke, bleeding, and re-narrowing are possible complications.

\medterm{Carotid Body} A small group of sensing cells near the division of the common carotid artery. It detects changes in blood oxygen, carbon dioxide, and acidity and sends signals that help adjust breathing.

\medterm{Carotid Bruit} A blowing sound heard over a carotid artery when blood flow is turbulent. It can occur with plaque but does not reliably show how narrow the artery is; ultrasound or other vascular imaging is needed when disease is suspected.

\medterm{Carotid Cavernous Fistula} An abnormal connection between the carotid artery and the cavernous sinus, a large venous space behind the eye. It may cause a bulging or red eye, pulsation, a whooshing sound, double vision, or vision loss and needs specialist treatment.

\medterm{Carotid Duplex Ultrasound} An ultrasound test that shows plaque and measures blood flow through the carotid arteries. It helps estimate narrowing without radiation and may guide decisions about medicines, surgery, or stenting.

\medterm{Carotid Endarterectomy} Surgery to open a carotid artery and remove fatty plaque from its inner wall. It can reduce stroke risk in selected people with significant narrowing, especially after symptoms, but carries risks such as stroke, heart attack, and nerve injury.

\medterm{Carotid Stenosis} Narrowing of a carotid artery, usually from fatty plaque. It may reduce blood flow to the brain or release a clot, causing a transient ischemic attack or stroke; treatment depends on symptoms, severity, and overall risk.

\medterm{Carpal Bones} The eight small bones forming the wrist: scaphoid, lunate, triquetrum, pisiform, trapezium, trapezoid, capitate, and hamate. They connect the forearm to the hand and work together to allow wrist movement and transmit force.

\synonyms
Carpals

\medterm{Carpal Joints} The joints between the wrist bones and between the wrist, forearm, and hand bones. Together they allow the wrist to bend, straighten, move side to side, and rotate with the forearm.

\medterm{Carpal Tunnel} A narrow passage on the palm side of the wrist formed by the wrist bones and a strong ligament. It contains the median nerve and finger-flexing tendons; pressure on the nerve causes carpal tunnel syndrome.

\medterm{Carpal Tunnel Biopsy} Removal of a small tissue sample from the carpal tunnel or its contents when an unusual infiltrative, inflammatory, infectious, or tumor-related process is suspected. It is not a routine test for ordinary carpal tunnel syndrome.

\medterm{Carpal Tunnel Release} Surgery that cuts the ligament forming the roof of the wrist tunnel to relieve pressure on the median nerve. It may improve numbness, tingling, pain, and weakness when symptoms do not respond to suitable nonsurgical care.

\medterm{Carpal Tunnel Syndrome} Pressure on the median nerve as it passes through the narrow wrist tunnel. It commonly causes numbness, tingling, burning, or pain in the thumb, index, middle, and part of the ring finger, sometimes with thumb weakness.

\synonyms
Median Nerve Compression (Carpal Tunnel Syndrome)

\medterm{Carpal Tunnel Syndrome Tests} Examination and nerve tests used to assess suspected median-nerve compression. Wrist-position or tapping tests may reproduce symptoms, while nerve-conduction studies and ultrasound can support diagnosis and assess severity when needed.

\medterm{Carpometacarpal Joint} A joint between a wrist bone and the base of a hand bone. The thumb joint is a mobile saddle-shaped joint that allows opposition and is commonly affected by osteoarthritis.

\medterm{Carpus} The group of eight small bones forming the wrist between the forearm and the metacarpals of the hand. The carpus provides stability while allowing controlled wrist movement.
\medterm{Carrageenan} A substance extracted from red seaweed and used to thicken or stabilize some foods and medicines. Different forms have different properties; evidence about effects on human inflammation depends on the type and amount consumed.

\medterm{Carrier} A person who has one altered copy of a gene for a recessive condition and usually has no symptoms. A carrier can pass the altered gene to a child; risk depends on the other parent's genes and the inheritance pattern.

\medterm{Carrier-Mediated Transport} Movement of a substance across a cell membrane with help from a specific carrier protein. It may use no direct energy when moving down a concentration gradient or use energy to move against that gradient.

\medterm{Carrier Excipient} An inactive ingredient that supports an active medicine in a formulation. It may help control the medicine's dose, stability, absorption, appearance, or delivery and can occasionally cause allergy or intolerance.

\medterm{Carrier Screening} Genetic testing for people without symptoms to find whether they carry a harmful gene variant that could be passed to a child. Results can guide reproductive choices and may identify the need to test a partner or relatives.

\medterm{Carrier State} A condition in which a person carries and may spread a microorganism without symptoms. Some carriers need testing, treatment, vaccination of contacts, or other measures to prevent transmission, depending on the infection.

\medterm{Carry-Over Effect} Persistence of an earlier treatment's effect into a later measurement or treatment period. It can distort results in crossover studies or make a person's current clinical state difficult to interpret.

\medterm{Carteolol} A nonselective beta-blocker used mainly as eye drops to lower pressure inside the eye in open-angle glaucoma or ocular hypertension. It can sometimes slow the heartbeat, lower blood pressure, or worsen breathing problems.

\medterm{Cartilage} A tough, flexible tissue found in joints, the ear, nose, rib cage, and spinal discs. It cushions joints, supports body structures, and resists pressure, but it has limited blood supply and heals slowly after injury.

\medterm{Cartilage Disease} A disorder that damages cartilage in joints or supporting structures. It can cause pain, stiffness, swelling, deformity, impaired growth, or loss of smooth movement and may result from injury, aging, inflammation, or inherited disease.

\medterm{Cartilaginous Joint} A joint in which bones are connected by cartilage and allow little or limited movement. Examples include joints between spinal bones and the joint joining the two front pelvic bones.

\medterm{Caruncles} Small fleshy anatomical projections. Examples include the lacrimal caruncle at the inner corner of the eye and the urethral caruncle near the opening of the urethra.
\medterm{Carvedilol} A beta-blocker that also relaxes blood vessels. It slows the heart and reduces its workload and is used for selected cases of high blood pressure, heart failure, and after a heart attack; it can lower blood pressure and heart rate.

\medterm{CAS} Alternate index label; see \textit{Childhood apraxia of speech}.

\medterm{Casanthranol} A stimulant laxative that increases bowel movement and secretion of water into the intestine. Prolonged or excessive use can cause diarrhea, dehydration, electrolyte abnormalities, and reliance on laxatives.

\medterm{Cascade} A series of linked reactions in which one activated step triggers the next. Examples include the blood-clotting, complement, and inflammatory cascades.

\medterm{Cascara} A plant-derived stimulant laxative that increases bowel movement. Prolonged or excessive use can cause diarrhea, dehydration, electrolyte imbalance, and reliance on laxatives; its use is restricted in some places.
\medterm{Case Definition} Standardized clinical, laboratory, and exposure criteria used to classify a person as a suspected, probable, or confirmed case during surveillance or an outbreak investigation. It may differ from the criteria used to diagnose one patient.

\medterm{Case Fatality Rate} The proportion of people diagnosed with a disease who die from it during a stated period. It describes severity in a particular group, but depends on accurate diagnosis, follow-up, treatment, and the population studied.

\medterm{Case Management} Coordinated planning, referral, and follow-up of health and social services to help a person receive appropriate care, avoid gaps or duplication, and meet clinical, functional, and personal goals.

\medterm{Case Study} A detailed examination of one patient, disease episode, or clinical situation. Case studies support education, research, and quality improvement by illustrating diagnostic reasoning, treatment decisions, outcomes, or unusual presentations.

\medterm{Case-Control Studies} Observational studies that compare people with a disease or outcome with similar people without it and look backward for differences in earlier exposures, treatments, behaviors, or risk factors.

\medterm{Case-Control Study} An observational research study comparing people who have a disease or outcome with similar people who do not. Researchers examine earlier exposures or characteristics to identify possible causes or risk factors.

\medterm{Case-Fatality Rate} The proportion of diagnosed people who die from a disease during a stated period, calculated as deaths from that disease divided by diagnosed cases. It estimates severity in the group studied and is not the same as population mortality.

\medterm{Casein} The main family of proteins in cow's milk and other mammalian milk. It supplies amino acids and minerals and forms curds during digestion and cheese-making; it can trigger reactions in people with cow's-milk protein allergy.

\medterm{Casein Kinase} A family of enzymes that attach phosphate groups to proteins and thereby change their activity. Casein kinases help regulate cell signaling, growth, DNA repair, development, and the body's internal clock.
\medterm{Caseous} Having a soft, crumbly, cheese-like appearance. The term often describes dead tissue in granulomas caused by tuberculosis or certain other long-lasting infections.
\medterm{Caseous Necrosis} A form of tissue death producing soft, crumbly, white or yellow material that resembles cheese. It is classically seen in granulomas caused by tuberculosis, though similar changes can occur in some other chronic infections.

\medterm{Cashew Allergy} An immune reaction to cashew proteins that may cause itching, hives, swelling, vomiting, wheezing, or life-threatening anaphylaxis. People with this allergy should avoid cashews and follow an emergency action plan.

\medterm{Caspase} A group of enzymes that cut specific proteins inside cells. Some initiate orderly programmed cell death, while others activate inflammation and immune defense against infection.

\medterm{Caspase Activation} Conversion of inactive caspase enzymes into their active form inside a cell. Depending on the caspase and triggering signal, this can start programmed cell death or activate inflammation and immune defense.

\medterm{Caspofungin} An intravenous antifungal medicine that blocks construction of an important fungal cell-wall component. It is used for serious Candida infections and selected Aspergillus infections and may cause liver-test abnormalities, infusion reactions, or low potassium.
\medterm{Caspofungin Acetate} The acetate salt form of caspofungin, an intravenous antifungal medicine that blocks fungal cell-wall construction. It is used for selected serious Candida and Aspergillus infections under medical supervision.

\medterm{Cassette} A holder that contains recording or imaging material. In traditional radiography, it held the film or image receptor during exposure; the word may also describe holders used in laboratory or monitoring equipment.

\medterm{Cast} A rigid or semi-rigid support placed around an injured limb or joint to limit movement while a bone, ligament, tendon, or other tissue heals. It must be monitored for swelling, pressure injury, numbness, or poor circulation.

\medterm{Castleman Disease} A group of uncommon disorders causing abnormal enlargement and activity of lymph nodes. Disease may affect one lymph-node area or the whole body and can cause fever, fatigue, anemia, weight loss, and organ problems.

\medterm{Castor Oil} A vegetable oil from the seeds of Ricinus communis. Taken by mouth, it acts as a stimulant laxative; excessive use can cause cramping, diarrhea, dehydration, and electrolyte disturbances.

\medterm{Castor Oil Overdose} Excessive ingestion of castor oil, causing intestinal irritation, cramping, nausea, vomiting, and diarrhea. Severe exposure can cause dehydration, low blood pressure, or fluid and electrolyte imbalance.

\medterm{Castration} Removal or functional suppression of the testes or ovaries, reducing production of sex hormones. It may be surgical or achieved with medicines and can affect fertility, sexual function, mood, and bone health.

\medterm{CAT Scan} An older name for computed tomography, or CT, an imaging test that uses many x-ray measurements and computer processing to create cross-sectional pictures of the body.

\medterm{Cat'S Eye Reflex} A bright white reflection from the pupil, called leukocoria, rather than the usual red reflection in photographs. It can signal cataract, retinal disease, or retinoblastoma and requires urgent eye assessment, especially in a child.
\medterm{Cat-Nap} A short period of sleep, usually during the day. Frequent unintended sleep may indicate inadequate sleep, a medicine effect, or a sleep disorder and should be assessed if it affects safety or daily life.
\medterm{Cat-Scratch Disease} An infection caused by Bartonella henselae, usually after a scratch or bite from an infected cat. It commonly causes a small skin bump, nearby swollen lymph nodes, fever, and fatigue; rarely it affects the eye, liver, spleen, brain, or heart.

\synonyms
Cat-Scratch Fever

\medterm{Catabolism} The set of chemical processes that break down complex molecules such as carbohydrates, fats, and proteins into simpler substances, releasing energy that cells can use. Examples include digestion-related breakdown, glycolysis, and fat oxidation.

\medterm{Catagen} The short transition phase of the hair cycle between active growth and rest. The hair follicle shrinks, hair growth stops, and the hair remains attached until the later shedding phase.

\synonyms
Catagen phase; regression phase of hair growth.

\medterm{Catalase} An enzyme that breaks down hydrogen peroxide into water and oxygen. It protects cells from oxidative damage and helps some bacteria survive inside immune cells.

\medterm{Catalepsy} A state of greatly reduced movement and responsiveness in which a person may remain rigid or hold an imposed posture. It can occur with catatonia, narcolepsy, neurologic disease, or some psychiatric conditions.

\medterm{Catalyst} A substance that speeds a chemical reaction without being used up. It works by providing an easier reaction pathway; enzymes are biological catalysts, and catalysts change reaction speed without changing the final balance of products and reactants.

\medterm{Catamenial Pneumothorax} Repeated collapse of a lung occurring around the start of menstruation, usually because of thoracic endometriosis. It often affects the right side and causes sudden chest pain or breathlessness; recurrent cases need specialist evaluation and may require surgery or hormone treatment.

\synonyms
Endometriosis-Related Pneumothorax; Catamenial Thorax

\medterm{Cataplasm} A poultice made from a soft, moist substance spread on cloth and applied to the body. It was historically used to warm or soothe a local problem and is not a standard treatment for serious disease.

\medterm{Cataplexy} A sudden, brief loss of muscle strength triggered by strong emotion such as laughter, surprise, or anger. Awareness usually remains clear; it is strongly associated with narcolepsy and may cause the head, jaw, knees, or whole body to buckle.

\medterm{Cataract} Clouding of the eye's natural lens, which reduces the passage of light to the retina. It can cause blurred or faded vision, glare, poor contrast, and difficulty seeing at night; lens-replacement surgery can restore sight when daily activities are affected.

\medterm{Cataract - Adult} Clouding of the eye lens that develops during adulthood and gradually reduces sharpness, contrast, color perception, or tolerance of glare. Updating glasses may help temporarily, but lens-replacement surgery is the definitive treatment when vision limits daily activities.

\medterm{Cataract Removal} Surgery that removes a cloudy natural lens from the eye and usually replaces it with an artificial intraocular lens. It is performed when the cataract interferes with vision, work, driving, reading, or other daily activities.

\medterm{Cataract Surgery} Removal of a clouded eye lens and replacement with an artificial lens to improve vision. Possible complications include infection, bleeding, retinal problems, increased eye pressure, or a need for later treatment, though serious complications are uncommon.

\medterm{Cataracts} Clouding of one or both natural eye lenses, causing blurred or faded vision, glare, reduced contrast, or difficulty seeing in dim light. They usually develop with aging but can also follow injury, disease, medicines, or be present at birth.


\medterm{Catastrophic Aging} A rapid loss of physical or mental function in an older person after an illness, operation, fall, or other stress. Reduced physiologic reserve makes recovery harder; early assessment, treatment, rehabilitation, and support may limit decline.

\medterm{Catastrophic Reaction} An unusually intense emotional or behavioral response to a situation that may seem minor to others. It can occur with dementia, brain injury, or other neurologic disease and may involve shouting, fear, crying, or agitation.

\synonyms
Catastrophic response; catastrophic emotional reaction.

\medterm{Catatonia} A serious syndrome involving marked changes in movement, speech, and response to the environment. Features may include immobility, mutism, staring, unusual postures, rigidity, repetitive movements, or copying speech or actions; causes include psychiatric and medical disorders.

\medterm{Catatonia, Fatal} A severe, life-threatening form of catatonia with extreme motor disturbance, fever, unstable blood pressure or heart rate, dehydration, and risk of organ failure. It requires emergency treatment.
\medterm{Catbite} A puncture wound caused by a cat's teeth. Cat bites can push bacteria deep into tissue and may infect tendons, joints, or bone, so prompt cleaning and medical assessment are often needed.
\medterm{Catch-Up Immunization} Vaccination of a person who is behind the recommended schedule using age- and risk-appropriate intervals. A delayed series usually does not need to be restarted, but minimum ages, intervals, and medical precautions must be followed.

\medterm{Catechin} A group of plant compounds found in tea, cocoa, some fruits, and other foods. They show antioxidant effects in laboratory studies, but catechin supplements are not proven treatments and may cause side effects or interact with medicines.

\medterm{Catechol} A chemical compound with two neighboring hydroxyl groups, used in industry and involved in making some biologically active substances. Contact or inhalation can irritate tissues and substantial exposure may be toxic.

\medterm{Catecholamine} A group of hormones and nerve-signaling chemicals that includes epinephrine, norepinephrine, and dopamine. They help control the body's stress response, heart rate, blood pressure, alertness, and metabolism.

\medterm{Catecholamine Blood Test} A blood test measuring epinephrine, norepinephrine, and dopamine. It may help evaluate unusual hormone release from a pheochromocytoma or related tumor, although metanephrine testing is often preferred for initial screening.

\medterm{Catecholamine Crisis} A sudden, life-threatening surge of stress hormones causing very high blood pressure, rapid or irregular heartbeat, sweating, headache, and possible injury to the heart, brain, kidneys, or other organs. It can occur with certain adrenal or nerve tumors.

\medterm{Catecholaminergic Polymorphic Ventricular Tachycardia} A rare inherited rhythm disorder in which exercise or strong emotion triggers dangerous fast rhythms from the heart's lower chambers. The heart may look normal and the resting ECG may be normal; fainting or cardiac arrest can occur, especially in children and young adults.

\synonyms
CPVT; Bidirectional Ventricular Tachycardia; Stress-Induced Polymorphic VT

\medterm{Catecholamines} Adrenaline, noradrenaline, and dopamine are catecholamines. These hormones and nerve-signaling chemicals help regulate stress responses, blood pressure, heart function, alertness, and metabolism.
\medterm{Catecholamines (Urine/Plasma)} Blood or urine tests measuring catecholamines or their breakdown products. They may help investigate pheochromocytoma or paraganglioma, but medicines, stress, posture, diet, and collection conditions can affect results.

\synonyms
Urine/Plasma

\medterm{Catecholamines - Urine} A urine test measuring catecholamines or their breakdown products over a specified period. It may help investigate pheochromocytoma, paraganglioma, or unexplained episodes of high blood pressure, sweating, headache, or palpitations.

\medterm{Category Five Hurricane} A Saffir--Simpson hurricane with sustained winds of at least 157 miles per hour (252 km/h), capable of catastrophic damage and prolonged loss of habitability.

\medterm{Category I Fetal Heart Tracing} A reassuring fetal heart-rate pattern with a baseline of 110--160 beats per minute, moderate variability, and no late or variable decelerations. It is usually consistent with adequate fetal oxygenation at that time.

\medterm{Category II Fetal Heart Tracing} A fetal heart-rate pattern that is not clearly normal or abnormal. It requires continued monitoring and assessment of the mother, fetus, labor, and changes over time because it may become reassuring or concerning.

\medterm{Category III Fetal Heart Tracing} An abnormal fetal heart-rate pattern that may indicate inadequate fetal oxygenation or acid buildup. It includes absent variability with recurrent late or variable decelerations or a true sinusoidal pattern and requires prompt clinical evaluation and action.

\medterm{Category Three Hurricane} A Saffir--Simpson hurricane with sustained winds of 111--129 miles per hour (178--208 km/h), capable of devastating damage.

\medterm{Category Two Hurricane} A Saffir--Simpson hurricane with sustained winds of 96--110 miles per hour (154--177 km/h), capable of extensive damage to structures and vegetation.

\medterm{Catenin} A family of proteins that helps neighboring cells attach and communicate. Beta-catenin also carries growth signals inside cells; abnormal activity can allow some cancers to grow or spread.
\medterm{Caterpillar Dermatitis} An itchy or painful skin reaction after contact with irritating or venomous caterpillar hairs or spines. Washing the area and removing visible hairs may help; eye exposure, breathing difficulty, or widespread symptoms need medical care.
\medterm{Caterpillars} The larval stage of butterflies and moths. Hairs or spines on some species can cause painful skin inflammation, eye injury, or allergic reactions after contact or inhalation.

\medterm{Catharsis} Catharsis is emotional release or expression; it is a historical psychological concept and not a specific medical treatment or diagnosis.
\medterm{Cathartic} A substance that promotes bowel emptying, usually by increasing stool water, intestinal movement, or both. Excessive or prolonged use can cause diarrhea, dehydration, electrolyte imbalance, or dependence on laxatives.
\synonyms
Purgative; laxative.

\medterm{Cathepsin} A family of enzymes that break down proteins inside cells, especially in lysosomes, the cell's recycling compartments. They help with immune activity, tissue remodeling, and bone turnover; abnormal activity may contribute to inflammation, bone loss, or cancer.
\medterm{Cathepsin D} An enzyme in lysosomes that breaks down proteins in an acidic environment. It also participates in cell death and immune processes; abnormal activity has been studied in cancer and nervous-system disease.
\medterm{Cathepsin H} An enzyme in lysosomes that breaks down proteins and helps activate some biologically active peptides. It contributes to normal cell maintenance and may be altered in certain diseases.

\medterm{Cathepsin K} An enzyme made by osteoclasts, the cells that remove old bone. It breaks down the protein framework of bone; excessive activity can contribute to bone loss and osteoporosis.

\medterm{Cathepsin L} An enzyme in lysosomes that breaks down proteins and helps process antigens for immune recognition. It also participates in tissue remodeling and can assist entry or spread of some infectious agents.

\medterm{Cathepsin Z} An enzyme found mainly in lysosomes that cuts proteins into smaller pieces. It helps cells process materials and may contribute to immune responses and inflammation.

\medterm{Cathepsins B} A group of protein-cutting enzymes, including cathepsin B, that help recycle proteins inside cells. Abnormal activity can contribute to inflammation, tissue remodeling, invasion of nearby tissue by tumors, and spread of cancer.

\medterm{Catheter} A thin flexible tube placed in a blood vessel, body cavity, or duct to deliver medicines or fluids, drain urine or other material, measure pressure, collect samples, or guide a procedure.

\medterm{Catheter, Foley} A urinary catheter that remains in the bladder and uses a small inflatable balloon to stay in place. It continuously drains urine into a collection bag but increases the risk of infection and blockage.

\medterm{Catheter-Associated UTI} A urinary infection associated with an indwelling catheter. Risk rises with prolonged use, and symptoms may include fever, lower abdominal or flank pain, or confusion; avoiding unnecessary catheters and removing them early lowers risk.

\medterm{Catheter-Related Infections} Infections caused by microorganisms entering along or contaminating an indwelling catheter. They may cause local redness, pain, drainage, fever, or bloodstream infection; treatment may require antibiotics and catheter removal or replacement.

\medterm{Catheter-Related UTI} A urinary tract infection associated with an indwelling urinary catheter, usually caused by bacteria entering or traveling along it. Diagnosis requires compatible symptoms and appropriately collected urine testing because bacteria may be present without infection.

\medterm{Catheterization} Placement of a catheter into a blood vessel, body cavity, duct, or other passage to drain, infuse, measure, collect samples, or perform a procedure. Risks include pain, bleeding, infection, blockage, and injury to nearby tissue.

\medterm{Catheterization (Urinary)} Insertion of a flexible tube through the urethra or through the lower abdominal wall into the bladder to drain urine. It may be intermittent, indwelling with a balloon, or suprapubic; infection and injury are possible risks.

\synonyms
Urinary

\medterm{Cathexis} The investment of emotional or psychological energy in a person, idea, object, or activity, a term used mainly in psychology and psychoanalytic theory.

\medterm{Cation} An electrically positively charged atom or molecule, such as sodium, potassium, calcium, or hydrogen. Cations help regulate fluid balance, nerve signals, muscle contraction, and the body's acid-base balance.

\synonyms
Positively charged ion.

\medterm{Cauda Equina} The bundle of lumbar, sacral, and coccygeal nerve roots that continues below the end of the spinal cord. These nerves control sensation and movement in the legs and help control the bladder, bowel, and sexual organs.

\synonyms
Cauda equina nerve roots; horse-tail nerve bundle.

\medterm{Cauda Equina Syndrome} A medical emergency caused by pressure on the bundle of nerves below the spinal cord. Severe low-back pain, leg weakness or numbness, numbness around the genitals or anus, or new bladder or bowel problems require immediate assessment.

\synonyms
Syndrome Cauda Equina (Cauda Equina Syndrome)

\medterm{Caudad} Toward the tail or lower end of the body; in human anatomy, generally toward the feet or away from the head.

\medterm{Caudal} Positioned toward the tail or lower part of the body; in human anatomy, usually toward the feet or away from the head.

\medterm{Caudal Dysplasia} A group of birth defects in which the lower spine and nearby structures do not develop fully. It may cause abnormalities of the legs, pelvis, bowel, bladder, or genital organs and is more common with maternal diabetes.
\medterm{Caudal Vein} A vein near the tail end of an animal's body. In human anatomy, the term is used mainly in embryology or when describing structures in comparative anatomy.

\medterm{Caudate Nucleus} A curved group of nerve cells deep in the brain that helps control movement, learning, memory, motivation, and behavior. Disease affecting it may contribute to movement or cognitive problems.

\medterm{Cauliflower Ear} A thick, irregular deformity of the outer ear caused when blood or fluid collects between its cartilage and covering, usually after blunt injury or repeated friction. Prompt drainage and compression can protect the cartilage and prevent permanent change.

\synonyms
Perichondrial Hematoma (Cauliflower Ear)

\medterm{Caulking Compound Poisoning} Poisoning or injury after contact with or swallowing a caulking product. Symptoms depend on its ingredients and may include eye, skin, or airway irritation, vomiting, drowsiness, or aspiration into the lungs.

\medterm{Caulobacter} A group of bacteria found mainly in freshwater and soil. Human infection is extremely rare but may occur opportunistically in people with severe illness or weakened immunity.

\medterm{Caulobacteraceae} A family of mainly freshwater and soil bacteria that includes Caulobacter. Most members are environmental and are not recognized as common causes of human disease.

\medterm{Causalgia} Severe persistent burning pain, sensitivity, and other sensory changes after injury to a peripheral nerve. It is now usually considered complex regional pain syndrome type II and may include swelling, color or temperature changes, and movement difficulty.

\synonyms
Complex regional pain syndrome type II; causalgic pain.

\medterm{Cause Of Death} The disease, injury, or toxic exposure that starts the chain of events leading directly to death. It differs from the immediate mechanism, such as respiratory failure, and from contributing conditions.


\medterm{Causes Symptoms Of Hormonal Imbalances Women} Hormone-related symptoms in women can result from thyroid, ovarian, adrenal, or pituitary disease, pregnancy, menopause, medicines, or weight changes. Possible symptoms include irregular bleeding, acne, hair changes, infertility, hot flashes, tiredness, and mood changes.

\medterm{Caustic} Able to chemically burn or destroy living tissue, as strong acids and alkalis can. Swallowing or contact with a caustic substance can injure the mouth, esophagus, stomach, skin, eyes, or airway and may be an emergency.
\medterm{Cauterization} Intentional destruction or sealing of tissue with heat, electricity, chemicals, or another energy source. It may control bleeding, remove abnormal tissue, or close small vessels and wounds.

\medterm{Cautery} Use of heat, electricity, chemicals, or another energy source to destroy or seal tissue, commonly to control bleeding or remove a small abnormal area.
\medterm{Cautery Device} A medical instrument that uses heat, electricity, or another energy source to seal bleeding vessels or destroy unwanted tissue. Use can cause burns or damage surrounding structures if poorly controlled.

\medterm{Caveola} A small pocket-like fold in a cell membrane that helps organize signals, move materials into or out of the cell, and respond to stretching or mechanical stress.

\medterm{Caveolins} Proteins that form or regulate caveolae, small pockets in cell membranes. They help organize cell signals, handle fats, move materials across membranes, and respond to mechanical stress.

\medterm{Cavernous Hemangioma} A cluster of enlarged, thin-walled blood vessels that may occur in the skin, brain, spinal cord, or other tissues. Symptoms depend on location and may result from bleeding, pressure, pain, or altered function.

\medterm{Cavernous Malformations} Clusters of unusually wide, fragile blood vessels, usually in the brain or spinal cord. They may cause no symptoms, but bleeding or pressure on nerve tissue can produce headaches, seizures, weakness, numbness, vision problems, or other neurologic changes.

\medterm{Cavernous Sinus} A large venous channel beside the pituitary gland that drains blood from the eye and face. It contains important eye-movement nerves and the internal carotid artery; infection or clotting there can be life-threatening.
\medterm{Cavernous Sinus Thrombosis} A dangerous blood clot, often infected, in the venous channel behind the eyes. It usually follows infection of the face, nose, sinuses, teeth, or orbit and may cause severe headache, fever, bulging eyes, double vision, or vision loss.

\medterm{Cavities} Holes or damaged areas in teeth caused by acids made by plaque bacteria from sugars and other carbohydrates. Decay can reach the inner pulp and cause pain or infection; prevention and treatment may include fluoride, cleaning, fillings, or root-canal care.

\synonyms
Dental Caries (Cavities)

\medterm{Cavities And Tooth Decay} Progressive loss of tooth structure caused by acids produced when plaque bacteria use sugars. Decay can move from enamel into dentin and the pulp, causing sensitivity, pain, infection, abscess, or tooth loss.
\medterm{Cavity} A hole or weakened area in a tooth caused by decay. Without treatment, it may reach the inner pulp and cause sensitivity, pain, infection, abscess, fracture, or tooth loss.

\synonyms
Dental cavity; carious lesion.

\medterm{Cavum} A cavity or hollow space in the body. The precise meaning depends on the anatomical phrase, such as cavum septi pellucidi, a fluid-containing space in the brain.
\medterm{CBC} See \textit{Complete Blood Count}.

\synonyms
Complete blood count; full blood count.

\medterm{CBC Blood Test} See \textit{Complete Blood Count}.

\medterm{CCP} Cyclic citrullinated peptide, a target used in a blood test for antibodies associated with rheumatoid arthritis. A positive anti-CCP result supports the diagnosis and may indicate greater risk of joint damage, but it is not diagnostic alone.

\medterm{CCTA} Alternate name; see \textit{Coronary Computed Tomography Angiography}.

\medterm{CCU} A hospital unit providing continuous heart monitoring and specialized treatment for people with serious acute cardiac problems, such as heart attack, dangerous rhythm disturbance, or severe heart failure.

\synonyms
Coronary care unit; cardiac care unit.

\medterm{Cd Antigens} Proteins or other markers on the surface of cells that help identify cell types, especially immune cells. Laboratory antibodies against these markers can classify cells, diagnose blood disorders, and guide selected treatments.

\medterm{CD4 Count} The number of CD4 T lymphocytes in a blood sample. It helps assess immune function and the risk of opportunistic infection in people with HIV; results are interpreted with viral load, symptoms, treatment, and trends over time.

\medterm{CD4-Positive T Cell} A type of white blood cell that coordinates immune responses. Different CD4 cells help activate other immune cells, support antibody production, defend mucosal surfaces, or limit excessive inflammation.

\medterm{CD8-Positive T Cell} A type of white blood cell that recognizes abnormal signals on infected or altered cells. Cytotoxic CD8 cells can destroy virus-infected, cancerous, or otherwise abnormal cells.

\medterm{Cea} An abbreviation for carcinoembryonic antigen, a blood marker used mainly to monitor selected cancers rather than to screen people without symptoms.
\medterm{CEA Assay} A laboratory test measuring carcinoembryonic antigen in a specimen, usually blood. It is mainly used to monitor selected known cancers and cannot diagnose or exclude cancer by itself.

\medterm{CEA Blood Test} A blood test measuring carcinoembryonic antigen, used mainly to monitor selected cancers such as colorectal cancer and assess treatment response or recurrence. Smoking, inflammation, and other noncancerous conditions can also raise it.

\medterm{Cecal Hemorrhage} Bleeding from the caecum, the first part of the large intestine. Causes include inflammation, infection, reduced blood flow, polyps, diverticular disease, or cancer and may produce blood in the stool or anemia.
\medterm{Cecal Infection} Infection or inflammation of the caecum caused by microorganisms, reduced blood flow, inflammatory bowel disease, or nearby disease. It may cause right-sided abdominal pain, fever, diarrhea, or tenderness.

\medterm{Cecal Neoplasm} A noncancerous or cancerous growth arising in the caecum. It may cause bleeding, anemia, abdominal pain, altered bowel habits, or blockage and is usually evaluated with colonoscopy and tissue sampling.

\medterm{Cecal Volvulus} Twisting of the caecum and nearby colon around the tissue that supports them, causing large-bowel blockage. It can produce sudden abdominal pain, swelling, vomiting, and inability to pass stool or gas and usually requires urgent treatment.

\medterm{Cecostomy} A surgically created opening from the caecum to the abdominal wall to release pressure, drain bowel contents, or divert stool. It may be temporary or permanent depending on the underlying problem.

\medterm{Cecum} The first pouch-like part of the large intestine, in the lower right abdomen where the small intestine joins the colon. The appendix arises from it; disease can include inflammation, tumors, bleeding, or twisting that blocks the bowel.

\medterm{Cedar Leaf Oil Poisoning} Toxic exposure to cedar-leaf oil, usually after swallowing or inhaling it. It may irritate the mouth, stomach, lungs, or nervous system and can cause vomiting, dizziness, seizures, or breathing problems; urgent poison-control advice is needed.
\medterm{Cedecea} A group of gram-negative bacteria found in the environment that rarely cause human disease. When infection occurs, it is usually opportunistic and affects people with severe illness, weakened immunity, or medical devices.

\medterm{Cedecea Davisae} A rare environmental bacterium that has caused bloodstream, respiratory, wound, and urinary infections. Disease is uncommon and usually occurs in people with serious underlying illness or impaired immunity.

\medterm{Cedecea Lapagei} A rare environmental bacterium that can occasionally cause bloodstream, respiratory, or other opportunistic infection, mainly in people with severe illness or weakened immune defenses.
\medterm{Cedecea Neteri} An uncommon gram-negative bacterium that rarely causes opportunistic human infection. A culture result must be interpreted with symptoms, the body site, and the quality of the specimen.

\medterm{Cediranib Maleate} The maleate salt of cediranib, an experimental medicine that blocks vascular endothelial growth factor receptors and may limit a tumor's blood-vessel growth. It has been studied in cancer trials but is not a routine treatment.

\medterm{Cefaclor} A second-generation cephalosporin antibiotic used for selected susceptible bacterial infections. It can cause allergic reactions, nausea, diarrhea, or antibiotic-associated colitis and should be used only when appropriate.

\medterm{Cefadroxil} An oral first-generation cephalosporin antibiotic used for selected susceptible skin, throat, urinary, and other infections. Possible effects include allergy, nausea, diarrhea, and antibiotic-associated colitis.

\medterm{Cefamandole} A second-generation cephalosporin antibiotic used in some settings for susceptible bacterial infections. It can cause allergy or diarrhea and may increase bleeding risk, especially in people with poor nutrition or liver disease.

\medterm{Cefatrizine} An older oral cephalosporin antibiotic used in some countries for susceptible bacterial infections. Its usefulness depends on local availability, bacterial resistance, allergy risk, and current treatment guidance.

\medterm{Cefazolin} An injectable first-generation cephalosporin antibiotic used for selected susceptible infections and commonly given before surgery to prevent infection. Allergy, diarrhea, and dose adjustment in kidney disease may be important.

\medterm{Cefdinir} An oral third-generation cephalosporin antibiotic used for selected susceptible respiratory, skin, and other infections. It may cause diarrhea, nausea, allergic reactions, or antibiotic-associated colitis and can interact with iron-containing products.
\medterm{Cefepime} An injectable fourth-generation cephalosporin used for serious infections caused by susceptible bacteria, including some hospital-acquired infections. Kidney-dose adjustment is important because accumulation can cause confusion, seizures, or other neurologic toxicity.

\medterm{Cefixime} An oral third-generation cephalosporin used for selected susceptible urinary, respiratory, ear, and other bacterial infections. It can cause diarrhea or allergy, and unnecessary use promotes antibiotic resistance.

\medterm{Cefmenoxime} An injectable third-generation cephalosporin used in some countries for susceptible bacterial infections. Possible effects include allergic reactions, diarrhea, and antibiotic-associated colitis.

\medterm{Cefmetazole} An injectable cephamycin antibiotic used in some countries for selected abdominal, pelvic, and other infections caused by susceptible bacteria. Allergy, diarrhea, and antibiotic-associated colitis are possible.

\medterm{Cefonicid} A long-acting second-generation cephalosporin antibiotic used for selected susceptible infections. Its use is limited in many places; dose selection depends on kidney function, local resistance, and the infection site.

\medterm{Cefoperazone} A third-generation cephalosporin active against many gram-negative bacteria and sometimes combined with a beta-lactamase inhibitor. It can cause allergy, diarrhea, bleeding, or impaired bile flow and requires appropriate monitoring.

\medterm{Cefotaxime} An injectable third-generation cephalosporin used for serious susceptible infections, including meningitis, pneumonia, bloodstream infection, and sepsis. Allergy, diarrhea, and antibiotic-associated colitis are possible.
\medterm{Cefotetan} An injectable cephamycin antibiotic used for selected serious infections and surgical prevention. It covers some anaerobic bacteria but can cause allergy, diarrhea, bleeding, or a reaction if alcohol is consumed around treatment.

\medterm{Cefotiam} A second-generation cephalosporin used in some countries for susceptible respiratory, urinary, skin, and other bacterial infections. Allergy, gastrointestinal effects, and antibiotic-associated colitis are possible.

\medterm{Cefoxitin} An injectable cephamycin antibiotic used for selected abdominal, pelvic, skin, and other infections, including some involving anaerobic bacteria. Allergy, diarrhea, and antibiotic-associated colitis can occur.

\medterm{Cefsulodin} An older cephalosporin with activity against Pseudomonas aeruginosa. It is used mainly in laboratory or historical contexts because resistance, limited availability, and newer treatments restrict routine clinical use.

\medterm{Ceftazidime} An injectable third-generation cephalosporin with activity against Pseudomonas aeruginosa and other susceptible gram-negative bacteria. It is used for selected serious infections and can cause allergy, diarrhea, or neurologic effects when it accumulates in kidney disease.
\medterm{Ceftazidime-Avibactam} A combination antibiotic in which avibactam blocks some bacterial enzymes that destroy ceftazidime. It is used for selected difficult gram-negative infections, including some resistant organisms; kidney-dose adjustment and susceptibility testing are important.

\medterm{Ceftibuten} An oral cephalosporin antibiotic used for selected susceptible respiratory, urinary, and other bacterial infections. It may cause diarrhea or allergy, and inappropriate use promotes resistance and antibiotic-associated colitis.

\medterm{Ceftizoxime} An injectable third-generation cephalosporin used for selected serious infections caused by susceptible bacteria, including some abdominal, pelvic, urinary, and anaerobic infections. Allergy, diarrhea, and antibiotic-associated colitis are possible.

\medterm{Ceftriaxone} An injectable third-generation cephalosporin used for many serious susceptible bacterial infections, including meningitis, pneumonia, bloodstream infection, and gonorrhea. Allergy, diarrhea, biliary problems, and antibiotic-associated colitis are possible.

\medterm{Cefuroxime} A second-generation cephalosporin available by mouth or injection for selected respiratory, ear, urinary, skin, and other bacterial infections. It can cause allergy, nausea, diarrhea, or antibiotic-associated colitis.

\medterm{Celiac Artery} A major branch of the abdominal aorta that supplies the upper digestive organs, including the stomach, liver, gallbladder, spleen, pancreas, and first part of the small intestine. It divides into the left gastric, splenic, and common hepatic arteries.

\medterm{Celiac Crisis} A rare, life-threatening severe flare of celiac disease causing profuse diarrhea or vomiting, dehydration, major electrolyte disturbances, and sometimes shock. It requires urgent hospital treatment and strict control of gluten exposure.

\medterm{Celiac Disease} An immune disorder triggered by gluten in wheat, barley, and rye in genetically susceptible people. It damages the small intestine and can cause diarrhea, bloating, weight loss, anemia, poor absorption, bone disease, or symptoms outside the intestine.

\synonyms
Nontropical Sprue

\medterm{Celiac Disease (Gluten Enteropathy)} Alternate index label for Gluten Enteropathy.

\medterm{Celiac Disease - Nutritional Considerations} Nutrition care for celiac disease requires lifelong avoidance of wheat, barley, rye, and cross-contaminated foods. Dietitians can help maintain adequate fiber and nutrients and correct deficiencies of iron, folate, vitamin B12, calcium, or vitamin D.


\medterm{Celiac Disease - Sprue} Another description of celiac disease, an immune reaction to gluten that damages the small intestine. It can cause diarrhea, bloating, weight loss, anemia, poor nutrient absorption, bone disease, and other symptoms.

\medterm{Celiac Plexus} A network of nerves in the upper abdomen that carries signals to and from the stomach, liver, pancreas, intestines, and other organs. It can be targeted with an injection for selected severe abdominal pain.

\medterm{Celiac Trunk} The first major branch of the abdominal aorta, supplying the upper digestive organs. It divides into the left gastric, splenic, and common hepatic arteries and is surrounded by nerves of the celiac plexus.

\medterm{Celiac Trunk Stenosis} Narrowing of the celiac trunk, caused by plaque or, in some people, compression by the diaphragm's median arcuate ligament. It may cause upper abdominal pain after meals and weight loss, although imaging findings do not always explain symptoms.

\medterm{Cell} The smallest living unit that can perform life processes. Human cells usually contain DNA in a nucleus and specialized internal structures that produce energy, make proteins, communicate, and maintain the body's tissues.

\medterm{Cell Adhesion} The attachment of cells to one another or to the supporting material around them. It helps maintain tissue structure, transmit signals, guide movement, and organize healing and development.

\medterm{Cell Aggregation} Clumping of cells through interactions between their surfaces or surrounding proteins. It is important in development, immune responses, wound healing, and formation of blood clots.

\medterm{Cell Body} The main part of a cell containing the nucleus and most internal structures. In a nerve cell, it excludes long extensions such as the axon and dendrites.

\medterm{Cell Count} The number of cells in a sample of blood, body fluid, tissue, or laboratory culture. It helps assess blood-cell levels, inflammation, infection, immune function, or abnormal cell growth.

\medterm{Cell Culture} Growth of cells outside the body in a controlled nutrient environment. It is used for research, medicine testing, diagnosis, production of biological products, and study of how cells respond to disease or treatment.

\medterm{Cell Cycle} The series of stages through which a cell grows, copies its DNA, checks for damage, and divides into two cells. Disruption of cycle control can contribute to developmental disorders and cancer.

\medterm{Cell Death} The loss of a cell's ability to remain alive. It may be a controlled process that removes old or damaged cells, or an uncontrolled result of injury, infection, toxins, lack of oxygen, or severe disease.

\medterm{Cell Dedifferentiation} Loss of some specialized features by a mature cell, making it more like an immature cell. It can occur during repair or development and may contribute to abnormal growth in some cancers.

\medterm{Cell Density} The number of cells in a given area or volume of tissue or fluid. It can change with growth, inflammation, healing, infection, or cancer and is measured in laboratory samples and tissue studies.

\medterm{Cell Differentiation} The process by which an unspecialized cell develops the structure and function of a specialized cell type, such as a nerve, muscle, blood, or skin cell.

\medterm{Cell Division} The process by which one cell produces new cells. Mitosis makes most body cells for growth and repair, while meiosis makes reproductive cells with half the usual number of chromosomes.

\medterm{Cell Fusion} Joining of two cells or their outer membranes to form one larger cell, sometimes with several nuclei. It occurs in normal development and repair and can also occur during infection or disease.

\medterm{Cell Growth} Increase in a cell's size and contents through production of proteins, membranes, energy stores, and other components. Abnormal growth signals can contribute to tissue overgrowth or cancer.

\medterm{Cell Maturation} Development of an immature cell into a mature cell with a specialized structure and function. It is essential for producing healthy blood, nerve, muscle, immune, and other body cells.

\medterm{Cell Migration} Movement of cells through tissue or across a surface in response to chemical, mechanical, or environmental signals. It is important in development, immune defense, wound healing, and cancer spread.

\medterm{Cell Motility} The ability of cells to move through or across a surface using their internal framework, attachment structures, or specialized projections such as cilia and flagella. It supports development, immunity, wound healing, and cancer spread.

\synonyms
Cell Locomotion

\medterm{Cell Movement} Directed or random movement of a cell through tissue or across a surface. It is important in development, immune responses, wound repair, inflammation, and the spread of cancer.

\medterm{Cell Nucleus} The membrane-covered compartment containing most of a cell's DNA. It helps control gene activity, growth, division, repair, and production of proteins needed for normal cell function.

\medterm{Cell Pellet} The concentrated layer of cells collected at the bottom of a tube after spinning a sample in a centrifuge. It may be removed for microscopic, genetic, chemical, or other laboratory testing.


\medterm{Cell Physiology} The study of how cells obtain energy, transport substances, communicate, respond to signals, repair themselves, and perform specialized functions.

\medterm{Cell Polarity} The unequal arrangement of structures and activities within a cell, giving it distinct sides or directions. It enables organized transport, signaling, movement, development, and tissue structure.

\medterm{Cell Proliferation} Increase in cell number through repeated division. Normal proliferation supports growth and repair, while excessive or uncontrolled proliferation can cause tissue overgrowth, benign tumors, or cancer.

\medterm{Cell Secretion} Release of substances made by cells, such as hormones, enzymes, nerve chemicals, mucus, or immune signals. Secretion supports communication, digestion, protection, movement, and regulation of body processes.

\medterm{Cell Senescence} A state in which a cell remains active but permanently stops dividing, often because of aging, stress, or DNA damage. Senescent cells can help prevent cancer but may contribute to inflammation when they accumulate.

\synonyms
Cellular senescence; replicative senescence.

\medterm{Cell Sorting} Separation of cells into groups according to features such as size, density, surface markers, or fluorescence. It is used in research, diagnosis, blood-cell analysis, and preparation of selected cells for treatment.

\medterm{Cell Structure} The arrangement of a cell's membrane, cytoplasm, nucleus, and internal structures. These parts work together to produce energy, make proteins, communicate, divide, and carry out specialized functions.

\medterm{Cell Surface} The outer membrane of a cell and the molecules attached to it. It controls movement of substances, attachment, recognition, communication, and interactions with immune cells and surrounding tissue.

\medterm{Cell Survival} The ability of a cell to remain alive and functional despite injury, lack of nutrients, toxins, infection, or other stress. It depends on the cell type, the severity and duration of the stress, and the cell's repair systems.

\medterm{Cell Therapy} Treatment that uses living cells to replace damaged cells, restore a function, or change an immune response. Uses range from established blood-forming stem-cell transplantation to experimental treatments; benefits and risks depend on the cells and disease.

\medterm{Cell Transplantation} Transfer of living cells into a person to replace missing cells, restore tissue function, or modify immunity. Blood-forming stem-cell transplantation is established for selected diseases; rejection, infection, graft-versus-host disease, and other risks may occur.

\medterm{Cell Wall} A rigid outer layer surrounding plant, fungal, and bacterial cells. It provides shape and protection; human and animal cells have a flexible membrane but no cell wall.

\medterm{Cell-Cell Junction} A specialized connection between neighboring cells that provides attachment, communication, or controlled passage of substances. Examples include tight junctions, desmosomes, and gap junctions.

\medterm{Cell-Cell Signaling} The ways cells send and receive messages through direct contact, electrical activity, or chemicals such as hormones and neurotransmitters. Signaling coordinates growth, movement, immune responses, repair, and other body functions.

\medterm{Cell-Free Fetal DNA Screening} A prenatal blood test that analyzes small DNA fragments from the placenta to estimate the chance of selected chromosome conditions. It is a screening, not a diagnostic, test; a high-risk or unclear result needs counseling and confirmatory testing.

\medterm{Cell-Free System} A laboratory preparation containing selected cell components but no intact living cells. It is used to study biochemical reactions, make proteins, test medicines, or examine biological processes under controlled conditions.

\medterm{Cell-Matrix Junction} A specialized connection anchoring a cell to the supporting material around it. It also transmits signals that influence cell shape, movement, survival, and tissue organization.

\medterm{Cell-Mediated Immune Response} Alternate index label; see \textit{Cell-mediated immunity}.

\medterm{Cell-Mediated Immunity} Immune defense carried out mainly by T lymphocytes and other cells rather than by antibodies. It is important for controlling infected cells, cancer cells, and organisms that live inside cells.

\synonyms
Cell-Mediated Immune Response

\medterm{Cellophane} A thin transparent film made from cellulose. In medicine it may be used in some dressings, packaging, or laboratory techniques; its clinical meaning depends on the context.
\medterm{Cells Stem (Stem Cells)} Alternate index label; see \textit{Stem Cell}.

\medterm{Cellular} Relating to cells or to processes occurring within or between cells.

\medterm{Cellular Adaptation} A reversible change in a cell's size, number, structure, or function in response to altered demands or mild injury. Adaptation may help a cell survive but can become harmful if the stress continues.

\medterm{Cellular Immunity} Immune defense carried out mainly by T cells, natural killer cells, macrophages, and other cells rather than antibodies alone. It helps destroy infected or abnormal cells and control inflammation.

\medterm{Cellular Metabolic Process} A chemical reaction or linked pathway in a cell that produces energy, builds cell components, or breaks down nutrients and waste.

\medterm{Cellular Respiration} The process cells use to release energy from food. Most cells use oxygen to convert nutrients into ATP, the cell's usable energy, while limited energy can be made without oxygen.

\medterm{Cellular Swelling} Early reversible cell injury in which failure of membrane pumps causes sodium and water to collect inside the cell. It commonly follows reduced blood flow, low oxygen, toxins, or other stresses.

\synonyms
hydropic change.

\medterm{Cellular Vesicle} A small membrane-bound compartment that transports, stores, or releases substances inside or outside a cell. Vesicles help move proteins, fats, hormones, and waste.

\medterm{Cellulite} A common harmless dimpled or uneven appearance of skin, usually on the thighs, hips, or buttocks, caused by the arrangement of fat, connective tissue, and skin. It is different from cellulitis, a bacterial infection.

\medterm{Cellulitis} A bacterial infection of the deeper skin and tissue beneath it, causing spreading redness, warmth, swelling, tenderness, and sometimes fever or chills. It commonly follows a break in the skin and can spread rapidly or enter the bloodstream.

\medterm{Cellulomonas} A group of bacteria found mainly in soil, plants, and other environmental materials. They break down cellulose; human infection is rare and usually occurs opportunistically in people with severe illness or weakened immunity.

\medterm{Cellulomonas Hominis} A rare environmental bacterium that has occasionally caused human infection, usually in people with weakened immunity or serious underlying illness. Its significance depends on symptoms, the body site, and laboratory identification.

\medterm{Cellulose} A complex carbohydrate forming the structure of plant cell walls. Humans cannot digest it, but it is dietary fiber that adds bulk to stool and can support regular bowel movements when enough fluid is consumed.

\medterm{Cellulosimicrobium} A group of environmental bacteria found mainly in soil. Human infection is rare but may affect people with serious illness, implanted devices, or weakened immune defenses and can involve the blood, heart, or other tissues.

\medterm{Cellulosome} A group of enzymes joined together by some microorganisms to break down cellulose in plant cell walls. It helps the organism digest plant material and is studied for biotechnology and biofuel production.

\medterm{Celosomy} A rare birth defect involving incomplete closure of the body wall, allowing abdominal or chest organs to protrude outside the body. The term is used for severe omphalocele-like malformations and requires specialist newborn care.

\medterm{Cement} A material used to fill spaces or attach an artificial joint or other prosthesis to bone. Dental cement bonds or protects dental work; the appropriate type depends on the procedure.
\medterm{Cemento-Enamel Junction} The boundary where tooth enamel covering the crown meets cementum covering the root. Its shape varies, and an exposed junction may be sensitive or more vulnerable to decay.
\medterm{Cementoma} A usually noncancerous jaw lesion containing cementum-like mineralized tissue. Its behavior and treatment depend on the specific subtype, location, size, symptoms, and relationship to nearby teeth.

\medterm{Cementum} A thin mineralized layer covering a tooth root. Fibers of the periodontal ligament attach to it and help anchor the tooth in its socket; it can be lost with gum disease.

\medterm{Cenesthesia} The general awareness of one’s body, including its condition, position, and internal sensations. Disturbance may occur in some neurologic or psychiatric illnesses and can alter how bodily sensations are experienced.
\medterm{Centenarian} A person aged 100 years or older. Centenarians are a diverse group; some remain independent, while others need substantial support, and age alone does not determine health or cognitive ability.

\medterm{Center} A central point, region, or facility organized around a function. In medicine, the term may describe an anatomical location, a treatment facility, or the middle of a measured structure.
\medterm{Centigram} A metric unit of mass equal to one hundredth of a gram, or 10 milligrams, used for small quantities of medicines or specimens.

\medterm{Centimeter} A metric unit of length equal to one hundredth of a meter, commonly used in clinical measurements such as wound size, tumor dimensions, and anatomical distances.

\medterm{Centimeter (Cm)} A unit of length equal to one hundredth of a meter, or ten millimeters, commonly used to measure body structures and objects.

\medterm{Centipede} A many-legged arthropod whose bite can inject venom and cause local pain, swelling, and redness; severe allergic reactions are uncommon but possible.

\medterm{Centipoise} A unit of dynamic viscosity equal to one millipascal-second, used to describe the flow resistance of blood, fluids, medicines, and other materials.

\medterm{Central} Located near the middle of the body or an organ, rather than near its edge. The word may describe a body structure, nervous-system area, blood line, symptom, or pattern of disease.

\medterm{Central (Brain) Fatigue} A sense of reduced mental energy, concentration, alertness, or motivation arising primarily from central nervous-system processes.

\synonyms
Central fatigue; mental fatigue.

\medterm{Central Apnea} A pause in breathing caused by absent or reduced respiratory effort from impaired brainstem signaling, rather than upper-airway obstruction.

\medterm{Central Auditory Processing Disorder} Difficulty interpreting sounds in the brain despite adequate hearing. People may struggle to understand speech in noise, follow rapid instructions, identify sounds, or distinguish similar words.

\medterm{Central Core Disease} An inherited muscle disorder, often caused by an RYR1 gene change, that may cause low muscle tone, delayed motor development, and weakness. Some affected people can develop life-threatening overheating with certain anesthetics.

\synonyms
CCD; Core Myopathy; RYR1-Related Central Core Disease

\medterm{Central Deafness} Severe hearing loss caused by damage to sound-processing pathways or areas of the brain, rather than the outer, middle, or inner ear. Hearing tests may show poor understanding despite detected sounds.

\medterm{Central Diabetes Insipidus} Excessive urination and intense thirst caused by too little antidiuretic hormone from the brain. The kidneys cannot conserve enough water, so dehydration and high blood sodium may develop.

\medterm{Central Hypothyroidism} A condition with too little thyroid hormone because the brain's pituitary gland or hypothalamus does not provide enough stimulation to the thyroid.

\medterm{Central Line} A catheter inserted into a large vein with its tip in or near the heart, used to give medications or nutrition, draw blood, or monitor central venous pressures.

\medterm{Central Line Infections - Hospitals} An infection related to a central venous catheter, including a central-line-associated bloodstream infection, in which microorganisms enter the bloodstream from or through the catheter.

\medterm{Central Line Placement} Insertion of a catheter into a large vein, commonly in the neck, chest, or groin, with its tip near the heart. It allows long-term medicines, nutrition, fluids, blood sampling, or monitoring.

\medterm{Central Line-Associated Bloodstream Infection} A bloodstream infection linked to a central venous catheter. It may cause fever, chills, or low blood pressure and is evaluated with blood cultures; treatment usually requires antibiotics and sometimes catheter removal.

\medterm{Central Nervous System} The brain and spinal cord. It receives and interprets information, controls movement and many automatic body functions, and supports thought, memory, emotion, and regulation of hormones.
\medterm{Central Nervous System Lymphoma} A lymphoma involving the brain, spinal cord, cerebrospinal fluid, or eye that can cause neurologic or visual symptoms.

\medterm{Central Nervous System Vascular Malformations} Abnormal groups or connections of blood vessels in or around the brain or spinal cord. Some cause no symptoms, while others can cause headaches, seizures, weakness, or bleeding; treatment depends on the type, size, location, and risk.

\medterm{Central Neurofibromatosis} Alternate index label; see \textit{NF2-related schwannomatosis (neurofibromatosis type 2)}.

\medterm{Central Pontine Myelinolysis} Brain injury caused most often by correcting very low blood sodium too quickly. Damage in the pons may cause weakness, trouble speaking or swallowing, confusion, reduced consciousness, or locked-in syndrome.

\medterm{Central Precocious Puberty} Central precocious puberty begins when the hypothalamus and pituitary activate the normal puberty pathway unusually early. It causes early secondary sexual characteristics and accelerated bone maturation; evaluation identifies causes and treatment may delay progression when appropriate.

\medterm{Central Retinal Artery Occlusion} A sudden blockage of the main artery supplying the retina, causing severe, usually painless loss of vision in one eye. It is an eye stroke requiring emergency assessment for related vascular disease.

\medterm{Central Serous Chorioretinopathy} A condition in which fluid collects beneath the retina, usually near the macula, causing blurred or distorted central vision. It may be linked to corticosteroids or stress and often needs eye-specialist follow-up.

\medterm{Central Serous Choroidopathy} A disorder in which fluid accumulates beneath the neurosensory retina, usually near the macula, causing blurred or distorted central vision. It is associated with corticosteroid exposure and stress-related factors; persistent or recurrent disease requires ophthalmic evaluation.

\medterm{Central Sleep Apnea} Repeated pauses in breathing during sleep because the brain temporarily stops sending signals to breathe, without physical airway blockage. It may occur with heart failure, high altitude, or opioid use.

\synonyms
Sleep Apnea, Central

\medterm{Central Venous Catheter} A thin tube ending in a large vein near the heart. It provides reliable access for prolonged or irritating treatments and monitoring but can cause infection, blood clots, bleeding, or injury.

\medterm{Central Venous Catheter - Dressing Change} Replacement of the sterile dressing covering a central venous catheter using aseptic technique to keep the insertion site clean, dry, and protected from contamination and infection.

\medterm{Central Venous Catheter - Flushing} Instillation of a prescribed sterile solution through a central venous catheter to maintain patency and reduce blockage, using the device-specific technique and volume.

\medterm{Central Venous Catheters - Ports} Implantable venous access ports are central venous catheters placed beneath the skin, accessed with a special needle to deliver medication, fluids, blood products, or obtain blood samples over prolonged treatment.

\medterm{Central Venous Line} A vascular catheter whose tip lies in a major central vein, commonly the superior vena cava, for reliable long-term or high-concentration intravenous access.

\medterm{Central Venous Line - Infants} A catheter whose tip lies in a large central vein of an infant, used for irritant medicines, nutrition, fluids, blood sampling, or monitoring; infection and thrombosis are key risks.

\medterm{Central Venous Pressure} Pressure measured in or near the right atrium, usually through a central venous catheter. It reflects
right-sided filling conditions but is affected by ventilation, heart function, vascular tone,
and measurement technique; it should not be used alone to guide fluid treatment.

\medterm{Central Vision} Detailed vision directed at the point of gaze, provided primarily by the macula and fovea; it is needed for reading and recognizing faces.

\medterm{Centriacinar Emphysema} A subtype of emphysema in which destruction begins in the respiratory bronchioles and
extends peripherally. It predominantly affects the upper lobes and is strongly
associated with cigarette smoking. It is the most common form of emphysema and often
coexists with chronic bronchitis.

\medterm{Centric-Fusion Translocation} A Robertsonian translocation in which the long arms of two acrocentric chromosomes fuse near their centromeres, leaving a combined chromosome and usually losing the short arms.

\medterm{Centrifugation} Spinning a specimen at high speed to separate components by density, such as cells from plasma or precipitated material from a liquid.

\medterm{Centrifuge} Spins samples to separate components by density. Used for blood processing, urine
analysis, tissue preparation. Speed: RPM or RCF.

\medterm{Centrilobular Emphysema} A type of permanent lung damage in which the small airways and nearby air sacs are destroyed, usually from long-term smoking or inhaled irritants. It often affects the upper lungs and can cause breathlessness, cough, wheezing, and reduced ability to exercise.

\synonyms
Centriacinar Emphysema; Proximal Acinar Emphysema

\medterm{Centriole} A cylindrical microtubule structure in the centrosome that helps organize the mitotic spindle and, in some cells, forms the basal body of a cilium or flagellum.

\medterm{Centromere} The constricted chromosome region where sister chromatids remain attached and the kinetochore forms for accurate chromosome segregation during cell division.

\medterm{Centronuclear Myopathy} A group of inherited muscle diseases in which muscle fibers develop abnormally. Symptoms range from severe weakness and breathing problems at birth to milder adult-onset weakness, often affecting the shoulders, hips, eyelids, eye movements, face, or feet.

\synonyms
CNM; Myotubular Myopathy (General)

\medterm{Centrosome} The main microtubule-organizing center of an animal cell, containing a pair of centrioles and coordinating the spindle during mitosis.

\medterm{Centrum} Centrum means a central part, especially the main body of a vertebra around which the neural arch develops.
\medterm{Cephacetrile} An older first-generation cephalosporin antibiotic that inhibits bacterial cell-wall synthesis; it is not commonly used in current practice.

\medterm{Cephalalgia} Headache or pain in the head, including pain arising from cranial or related structures.
It is a symptom with many possible primary and secondary causes.

\medterm{Cephalexin} An oral first-generation cephalosporin that inhibits bacterial cell-wall synthesis and treats selected skin, respiratory, urinary, and other susceptible infections.

\medterm{Cephalgia} Pain in the head, commonly called a headache. It may arise from a primary headache disorder or from infection, injury, medication, high blood pressure, or another illness.

\medterm{Cephalhaematoma} A collection of blood beneath the membrane covering a newborn's skull bone, usually caused by birth-related pressure or injury. It is limited by skull sutures and usually resolves without treatment.

\medterm{Cephalic} Relating to the head or toward the head; in obstetrics, a cephalic presentation means the fetus presents head first.

\medterm{Cephalic Vein} A superficial vein on the lateral side of the upper limb that drains the hand and forearm and joins the axillary vein near the shoulder.

\medterm{Cephalohematoma} A collection of blood beneath the membrane covering a newborn's skull bone, usually caused by birth-related pressure or injury. It does not cross skull sutures and generally clears gradually.

\medterm{Cephalometry} Measurement of the skull and facial structures, usually on standardized radiographs, used in orthodontic planning and assessment of craniofacial growth.

\medterm{Cephalopelvic Disproportion} A situation in which the fetal head cannot pass easily through the mother's pelvis because of their relative size or shape. It may prevent vaginal delivery, but cannot always be confirmed before labor.

\medterm{Cephaloridine} An older injectable first-generation cephalosporin antibiotic whose use declined because of nephrotoxicity and the availability of safer alternatives.

\medterm{Cephalosporin} A beta-lactam antibiotic that kills susceptible bacteria by blocking construction of their protective cell wall. Different generations act against different bacteria; allergy, diarrhea, and resistance are possible.

\medterm{Cephalosporins} Cephalosporins are beta-lactam antibiotics that inhibit bacterial cell-wall synthesis; generations differ in spectrum, and allergy, diarrhea, and resistance are important considerations.

\medterm{Cephalothin} An older injectable first-generation cephalosporin that inhibits bacterial cell-wall synthesis and was used for susceptible serious infections.

\medterm{Cephapirin} A first-generation cephalosporin antibiotic formerly used for susceptible bacterial infections; current availability and indications vary by country.

\medterm{Cephradine} A first-generation cephalosporin antibiotic related to cephalexin that inhibits bacterial cell-wall synthesis and treats selected susceptible infections.

\medterm{Ceramide} A type of fat that forms part of cell membranes and helps the outer skin retain water and block irritants. Ceramides also participate in signals controlling cell growth and survival.

\medterm{Cercaria} A cercaria is a free-swimming larval stage of a trematode parasite; contact with some species can cause swimmer’s itch or other infection.
\medterm{Cercarial Dermatitis} Alternate index label; see \textit{Swimmer's itch}.

\medterm{Cerclage} A procedure in which a strong stitch is placed around the cervix to help prevent it opening too early during pregnancy. It may be placed through the vagina or abdomen and requires obstetric follow-up.

\medterm{Cerebellar} Cerebellar means relating to the cerebellum, which coordinates balance, posture, eye movements, and skilled movement.
\medterm{Cerebellar Ataxia} Poor coordination caused by disease or injury affecting the cerebellum or its connections. It may cause an unsteady walk, clumsy movements, slurred speech, tremor, or abnormal eye movements.

\medterm{Cerebellar Cortex} The folded outer layer of the cerebellum containing Purkinje and granular neurons that coordinates balance, posture, timing, and precision of movement.

\medterm{Cerebellar Disease} A disorder of the cerebellum causing impaired coordination, gait ataxia, dysmetria, tremor, abnormal eye movements, or speech changes.

\medterm{Cerebellar Dysmetria} Inaccurate control of the distance or force of a movement because of cerebellar dysfunction. A person may reach past a target, stop short, or have difficulty judging movement size.

\medterm{Cerebellar Hemangioblastoma} A usually benign, highly vascular tumor of the cerebellum that may cause headache, imbalance, vomiting, or hydrocephalus and can occur with von Hippel-Lindau disease.

\medterm{Cerebellar Hypoplasia} Incomplete development of the cerebellum, usually beginning before birth. It may cause poor coordination, abnormal eye movements, seizures, delayed development, or problems with balance and movement.

\medterm{Cerebellar Neoplasm} A benign or malignant tumor involving the cerebellum; symptoms may include poor coordination, imbalance, abnormal eye movements, headache, vomiting, or hydrocephalus.

\medterm{Cerebellar Nuclei} Deep clusters of cerebellar neurons that receive processed signals from the cerebellar cortex and send coordinated motor output to the brainstem and thalamus.

\medterm{Cerebellar Tumor} A growth in or near the cerebellum that may cause poor coordination, unsteady walking, abnormal eye movements, vomiting, or pressure buildup.

\medterm{Cerebellar Vermis} The midline portion of the cerebellum that helps control trunk posture, gait, balance, and coordination of eye and head movements.

\medterm{Cerebellitis} Inflammation of the cerebellum, often after an infection and sometimes from an immune reaction. It may cause sudden poor coordination, unsteady walking, slurred speech, abnormal eye movements, headache, or reduced alertness.

\medterm{Cerebellopontine Angle} The angle-shaped region at the junction of the cerebellum, pons, and petrous temporal
bone. Lesions in this space may affect cranial nerves and commonly include vestibular
schwannoma.

\medterm{Cerebellopontine Angle Tumor} Lesion at the cerebellopontine angle, an anatomic space between the cerebellum and
pons, medial to the temporal lobe and lateral to the brainstem, containing CN V
(trigeminal), VII (facial), VIII (vestibulocochlear), and the anterior inferior
cerebellar artery (AICA).

\medterm{Cerebellum} Largest brain portion (approximately 10 percent of volume but over 50 percent of
neurons) in the posterior cranial fossa behind the brainstem. Two hemispheres and a
vermis, divided into vestibulocerebellum (balance, eye movements), spinocerebellum
(posture, gait, coordination), and cerebrocerebellum (motor planning, cognition).

\medterm{Cerebral} Relating to the cerebrum or, more broadly, to the brain.

\medterm{Cerebral Amyloid Angiopathy} Deposition of amyloid beta in the walls of small and medium cerebral and leptomeningeal
arteries. It commonly causes recurrent lobar intracerebral hemorrhage in older adults
and may also cause transient focal neurological episodes or cognitive decline.

\medterm{Cerebral Aneurysm} An abnormal focal enlargement of an artery in or around the brain. It may remain
asymptomatic or rupture and cause subarachnoid hemorrhage. Risk factors include hypertension, smoking, certain inherited
disorders, and family history.

\medterm{Cerebral Angiography} An X-ray test in which contrast dye is put into an artery to make the blood vessels of the brain visible. It can show aneurysms, abnormal artery-to-vein connections, narrowing, or blockage and is usually done with a thin tube called a catheter.

\medterm{Cerebral Aqueduct} A narrow channel through the midbrain connecting the third and fourth brain ventricles. It allows cerebrospinal fluid to flow, and blockage can enlarge the ventricles and increase pressure.

\medterm{Cerebral Arteriovenous Malformation} An abnormal direct connection between cerebral arteries and veins without a normal
capillary bed. It may cause intracranial hemorrhage, seizures, headache, or focal
neurological deficits and is evaluated with brain imaging and angiography.

\medterm{Cerebral Artery} An artery that carries blood to part of the brain. A blockage can cause an ischemic stroke, while a break in the vessel can cause bleeding in or around the brain; the effects depend on which artery is involved.

\medterm{Cerebral Atrophy} Loss or shrinkage of brain tissue caused by aging, neurodegenerative disease, injury, infection, stroke, or another condition. Effects vary from none to memory loss, weakness, seizures, or impaired thinking.
\medterm{Cerebral Autoregulation} The intrinsic ability of cerebral arterioles to maintain relatively stable cerebral
blood flow across a range of perfusion pressures. It may be impaired by severe brain
injury, ischemia, hypercapnia, or anesthetic and vasoactive drugs.

\medterm{Cerebral Calcification} Abnormal calcium deposition in brain tissue, with causes including inherited, metabolic, infectious, vascular, and age-related disorders.

\medterm{Cerebral Cavernous Malformation} A low-flow vascular malformation composed of dilated thin-walled vascular spaces in the brain or spinal cord. It may cause seizures, headaches, or hemorrhage, although some lesions remain asymptomatic.

\medterm{Cerebral Cavernous Malformations} Low-flow vascular malformations composed of dilated, thin-walled capillaries without
intervening brain parenchyma, often multiple and familial (autosomal dominant with
mutations in CCM1, CCM2, or CCM3 genes). They may present with seizures, focal
neurological deficits, or intracerebral hemorrhage.

\medterm{Cerebral Cortex} The outermost layer of the cerebrum, composed of gray matter (neuronal cell bodies)
folded into gyri (ridges) and sulci (grooves), greatly increasing surface area.

\medterm{Cerebral Crus} A group of nerve fibers in the middle part of the brain that carries movement commands down from the brain's outer layer. Damage can affect strength, movement, coordination, or control on the opposite side of the body.

\medterm{Cerebral Diplegia} Cerebral diplegia is a form of cerebral palsy with predominant weakness and spasticity of both legs, often greater than in the arms.
\medterm{Cerebral Edema} Swelling caused by excess fluid in or around brain tissue. It can raise pressure inside the skull, reduce blood flow, and cause headache, vomiting, confusion, seizures, coma, or life-threatening brain displacement.

\medterm{Cerebral Embolism} Alternate index label; see \textit{Embolic stroke (brain embolism)}.

\medterm{Cerebral Haemorrhage} Cerebral haemorrhage is bleeding within brain tissue or an adjacent intracranial space, causing pressure, tissue injury, and neurologic deficits.
\medterm{Cerebral Hemisphere} One of the two large halves of the cerebrum, each containing specialized but interconnected cortical regions.

\medterm{Cerebral Hemorrhage} Bleeding into brain tissue, usually from a damaged blood vessel. It can injure nearby brain cells, cause swelling and pressure, and produce sudden weakness, speech trouble, headache, reduced alertness, or death.

\medterm{Cerebral Hypoxia} Too little oxygen reaching the brain, often because of cardiac arrest, severe breathing failure, low blood pressure, or choking. Prolonged oxygen shortage can permanently injure brain cells.

\medterm{Cerebral Infarction} Permanent brain injury caused by blocked blood flow in a brain artery. The blockage may result from a clot formed in a narrowed artery or a clot traveling from the heart or another vessel.

\medterm{Cerebral Ischemia} Reduced blood flow and oxygen reaching brain tissue. It may cause temporary symptoms, such as a transient ischemic attack, or permanent injury if the interruption lasts long enough.
\medterm{Cerebral Lobes} Major regions of the cerebrum with different but overlapping functions: the frontal lobe supports planning and movement, the parietal sensation and spatial awareness, the temporal hearing and memory, and the occipital vision.

\medterm{Cerebral Malaria} A life-threatening form of malaria in which parasites cause brain dysfunction, often with confusion, seizures, abnormal behavior, or coma. It is usually caused by Plasmodium falciparum and requires urgent treatment.
\medterm{Cerebral Palsy} A group of permanent disorders of movement and posture development causing activity
limitation, attributed to non-progressive disturbances occurring in the developing fetal
or infant brain. It is the most common motor disability in childhood, affecting
approximately 1 in 500 live births.

\synonyms
Palsy, Cerebral
Palsy Cerebral (Cerebral Palsy)


\medterm{Cerebral Peduncle} One of two bundles of nerve fibers in the midbrain that carry movement signals from the cerebrum to the brainstem and spinal cord. Damage can cause weakness or loss of coordinated movement on the opposite side.

\medterm{Cerebral Perfusion Pressure} The pressure pushing blood through the brain, estimated from average arterial pressure minus pressure inside the skull. It falls when blood pressure drops or brain swelling raises skull pressure.

\medterm{Cerebral Salt Wasting} Excessive loss of sodium and water in urine associated with brain injury or disease. It can cause low blood sodium and dehydration and requires careful replacement guided by clinical and laboratory monitoring.

\medterm{Cerebral Thrombosis} Formation of a blood clot within a cerebral artery, reducing blood flow and potentially causing an ischemic stroke.

\medterm{Cerebral Tumor} A neoplasm arising within or involving the brain, producing symptoms according to its
location, growth rate, and effect on intracranial pressure.

\medterm{Cerebral Tumours} Cerebral tumours are abnormal growths in the cerebrum, which may be benign or malignant and cause seizures, headache, cognitive change, or focal deficits.
\medterm{Cerebral Vasospasm} Abnormal narrowing of cerebral arteries that can reduce brain perfusion and cause delayed cerebral ischemia, especially after aneurysmal subarachnoid hemorrhage.

\medterm{Cerebral Venous Sinus Thrombosis} A blood clot in a large vein or drainage channel of the brain. It can cause headache, seizures, visual problems, weakness, brain bleeding, or increased pressure inside the skull.

\medterm{Cerebral Ventricle} One of four connected spaces inside the brain that contain cerebrospinal fluid. The fluid cushions the brain and spinal cord, carries nutrients and waste, and flows through the ventricles before being absorbed around the brain.

\medterm{Cerebral Ventriculography} Imaging of the brain’s fluid-filled ventricles to assess their size, shape, connections, or blockage. Modern scans usually use computed tomography or magnetic resonance imaging rather than older invasive methods.

\medterm{Cerebral Ventriculomegaly} Enlargement of one or more fluid-filled spaces inside the brain. It may result from blocked fluid drainage, loss of brain tissue, or developmental variation; its importance depends on the cause and severity.
\medterm{Cerebral Visual Impairment} Reduced vision caused by damage to the brain's visual pathways or visual-processing areas rather than the eyes. It may cause poor visual attention, difficulty recognizing objects, or inconsistent use of vision.

\medterm{Cerebriform Nevus Sebaceous} A variant of nevus sebaceous presenting as a soft, cerebriform (brain-like),
pedunculated or sessile nodule, often found on the scalp or face. It is a benign adnexal
hamartoma composed of sebaceous glands, but like other nevus sebaceous lesions, it
carries a low risk of secondary neoplasms and may warrant monitoring or excision.

\medterm{Cerebritis} Inflammation and infection of brain tissue before a definite abscess forms, usually caused by bacteria, fungi, or parasites. It can cause headache, fever, seizures, or neurologic deficits and requires urgent imaging and antimicrobial treatment.

\medterm{Cerebroplacental Ratio} A calculation from fetal Doppler ultrasound comparing blood-flow patterns in a brain artery and the umbilical artery. A low ratio may suggest reduced placental function and fetal adaptation, but it must be interpreted with gestational age and other findings.

\medterm{Cerebrosides} Glycosphingolipids composed of ceramide plus a single sugar, abundant in myelin and nervous tissue and involved in membrane structure.

\medterm{Cerebrospinal Fluid} A clear, colorless fluid that fills the ventricles of the brain, the central canal of
the spinal cord, and the subarachnoid space around the brain and spinal cord. Produced
primarily by the choroid plexus of the lateral, third, and fourth ventricles (~500
mL\/day, total volume ~150 mL in adults).

\synonyms
Spinal Fluid

\medterm{Cerebrospinal Fluid Collection} A procedure, usually lumbar puncture, that collects cerebrospinal fluid for laboratory analysis; it can help evaluate infection, bleeding, inflammation, malignancy, and other disorders of the central nervous system.

\medterm{Cerebrospinal Fluid Culture} A laboratory test that attempts to grow bacteria, fungi, or other microorganisms from cerebrospinal fluid to investigate meningitis or another central nervous system infection. Results depend on prior antimicrobial treatment, specimen volume, transport, and the clinical context.

\medterm{Cerebrospinal Fluid Analysis} Laboratory examination of cerebrospinal fluid, usually obtained by lumbar puncture. It may
assess infection, bleeding, inflammation, demyelinating disease, or malignant cells.
Analysis can include opening pressure, appearance, cell counts, protein, glucose,
microbiology, and other tests selected for the clinical question.

\medterm{Cerebrospinal Fluid Cytology} Microscopic examination of cells in cerebrospinal fluid to evaluate suspected malignancy, infection, inflammation, or hemorrhage. It is especially useful for detecting leptomeningeal tumor involvement when interpreted with clinical and imaging findings.

\medterm{Cerebrovascular} Relating to the brain's blood vessels and blood flow. Cerebrovascular problems can reduce oxygen delivery and cause strokes, brief warning episodes, bleeding, or lasting problems with movement, speech, vision, or thinking.

\synonyms
Brain-vascular; cerebral vascular.

\medterm{Cerebrovascular Accident} A cerebrovascular accident is a stroke caused by interrupted brain blood flow from ischemia or bleeding, producing sudden neurologic dysfunction.
\medterm{Cerebrovascular Accident Prevention} Measures to reduce the risk of stroke, including controlling blood pressure, diabetes, and lipids, avoiding tobacco, treating atrial fibrillation when indicated, and using appropriate antithrombotic therapy.


\medterm{Cerebrovascular Accident (Stroke)} Sudden brain dysfunction caused by interrupted blood flow or bleeding in the brain. See Stroke.
\synonyms
Stroke

\medterm{Cerebrovascular Disorder} A disorder of brain blood vessels or cerebral blood flow, including ischemia, hemorrhage, aneurysm, malformation, and venous thrombosis.

\medterm{Cerebrum} The largest part of the brain, comprising two cerebral hemispheres connected by the
corpus callosum. Each hemisphere has four lobes: frontal (executive function, motor,
Broca's area), parietal (somatosensory, spatial awareness), temporal (auditory, memory,
Wernicke's area), occipital (visual processing).

\medterm{Cerium} A rare-earth metal whose compounds have industrial and experimental medical uses; exposure to some forms can irritate tissues or cause toxicity.

\medterm{Certification} A formal process confirming that a person, service, laboratory, or facility meets specified training, quality, safety, or professional standards.

\medterm{Certified Nurse-Midwife} An advanced-practice nurse trained and credentialed to provide reproductive, pregnancy, birth, postpartum, and newborn care, including management of normal pregnancy and referral of complications.

\medterm{Certolizumab Pegol} A biologic medicine that blocks tumor necrosis factor and treats selected inflammatory diseases, including rheumatoid arthritis and Crohn disease. It can increase the risk of serious infections.

\medterm{Ceruloplasmin} The main protein that carries copper in the blood. It is made mostly by the liver; unusually low or high levels can occur in Wilson disease, inflammation, liver disease, or other conditions.
\synonyms
Caeruloplasmin

\medterm{Ceruloplasmin Blood Test} A blood test measuring ceruloplasmin, the main copper-carrying protein in plasma; it is used with other copper studies to evaluate Wilson disease and disorders of copper metabolism.

\medterm{Cerumen} A substance that helps keep dirt out of the ear and lubricates the skin in the ear. More
commonly known as earwax.

\medterm{Cerumen Impaction} Earwax packed tightly enough to block the ear canal or prevent examination of the eardrum. It may cause blocked hearing, pressure, ringing, itching, pain, or dizziness.

\medterm{Ceruminous Adenocarcinoma} A malignant tumor of the ceruminous glands in the external ear canal that may cause pain, discharge, hearing loss, or a persistent ear-canal mass.

\medterm{Cervical} Relating to the neck or to the cervix, the lower opening of the uterus. The meaning depends on context, such as cervical spine, cervical pain, or cervical cancer.

\medterm{Cervical Adenocarcinoma} A malignant gland-forming tumor arising in the uterine cervix; it may be HPV-associated or HPV-independent and commonly presents with abnormal bleeding or discharge.

\medterm{Cervical Arthritis} Degenerative or inflammatory disease of the neck joints that can cause neck pain, stiffness, reduced motion, and sometimes nerve or spinal-cord compression.

\medterm{Cervical Atlas} The first bone of the neck, called C1. It is a ring-shaped vertebra that supports the skull and moves around the second neck vertebra when the head turns.

\medterm{Cervical Biopsy} Removal of a small piece of tissue from the cervix for laboratory examination. It is commonly performed after an abnormal screening result to look for precancerous or cancerous changes.
\medterm{Cervical Cancer} Cancer arising in the cervix, most often after persistent infection with a high-risk human papillomavirus type. Early disease may cause no symptoms; later signs include unusual bleeding, discharge, or pelvic pain.

\medterm{Cervical Cancer - Screening And Prevention} Clinical measures used to prevent or detect cervical cancer early, including human papillomavirus vaccination and screening with HPV testing, cervical cytology, or both according to age and risk.

\medterm{Cervical Cap} A flexible contraceptive barrier placed over the cervix to keep sperm from entering the uterus. It is used with spermicide and must remain in place for the recommended time after sex.
\medterm{Cervical Carcinoma} A malignant tumor of the uterine cervix, most commonly associated with persistent
high-risk human papillomavirus infection.

\medterm{Cervical Cerclage} Surgical placement of a non-absorbable suture around the cervix to provide mechanical
support and prevent preterm delivery in patients with cervical insufficiency.

\medterm{Cervical Change} A change in the uterine cervix or cervical spine, depending on context; clinical significance is determined by the specific structural, inflammatory, or pregnancy-related finding.

\medterm{Cervical Dilatation} Opening of the cervix during labor, measured from closed to about 10 centimeters. The cervix also becomes thinner and shorter, allowing the baby's head to pass into the birth canal.

\medterm{Cervical Dilation} The opening of the cervical canal during labor, progressing from closed to about 10 centimeters. It occurs with thinning and shortening of the cervix and is needed for vaginal birth.

\medterm{Cervical Disc Herniation} A neck disc bulges or tears and presses on a nearby nerve root or the spinal cord. It may cause neck pain, arm pain, tingling, numbness, or weakness.

\medterm{Cervical Dislocation} Abnormal displacement of one cervical vertebra relative to another, usually from trauma, that can destabilize the neck and injure the spinal cord.

\medterm{Cervical Dysplasia} Abnormal precancerous cells on the surface of the cervix, usually caused by persistent high-risk human papillomavirus infection. Mild changes often clear; more severe changes may progress to cancer if untreated.

\synonyms
cervical intraepithelial.
Cervical Intraepithelial Neoplasia Cervical (Cervical Dysplasia)
Dysplasia Cervical (Cervical Dysplasia)
Squamous Intraepithelial Lesion Cervical (Cervical Dysplasia)

\medterm{Cervical Dystocia} Failure of the uterine cervix to dilate normally during labor, which can delay or prevent vaginal delivery.

\medterm{Cervical Dystonia} Involuntary tightening of neck muscles that pulls the head into an abnormal position. It may cause neck pain, twisting, jerky movements, or tremor.

\synonyms
Dystonia, Cervical
Torticollis, Spasmodic

\medterm{Cervical Effacement} Thinning and shortening of the cervix during labor, measured from 0\% to 100\%. Complete effacement leaves the cervix very thin and usually occurs with opening of the cervix.

\medterm{Cervical Endometriosis} Endometrial-like tissue in or on the cervix that may cause bleeding after sex, pelvic pain, or an abnormal cervical lesion.

\medterm{Cervical Exam In Labor} A digital vaginal examination used to assess cervical dilation, effacement, consistency,
and position, together with the station and position of the fetal head. Serial examinations help assess labor progress.

\medterm{Cervical Ganglion} A cluster of nerve-cell bodies in the sympathetic chain of the neck that relays autonomic signals to the head, neck, and upper chest.

\medterm{Cervical Incompetence} The inability of the cervix to maintain pregnancy in the second trimester due to
painless cervical dilatation. Results in recurrent second trimester pregnancy loss or
preterm delivery. Risk factors include prior cervical trauma, cone biopsy, and
congenital anomalies.

\medterm{Cervical Insufficiency} Painless shortening or opening of the cervix during pregnancy before term, increasing risk of second-trimester loss or spontaneous preterm birth.

\medterm{Cervical Intraepithelial Neoplasia} Abnormal precancerous changes in the squamous cells of the cervix, usually associated with persistent high-risk human papillomavirus infection. CIN is graded by the severity of dysplasia.

\synonyms
CIN

\medterm{Cervical Intraepithelial Neoplasia Cervical (Cervical Dysplasia)} Alternate index label for Cervical Dysplasia.

\medterm{Cervical Length Screening} Transvaginal ultrasound measurement of the cervix between 18-24 weeks to assess risk of
spontaneous preterm birth. Normal cervix measures greater than 25 mm. Short cervix (less
than 25 mm) increases preterm delivery risk. Cervical funneling and dilatation are
additional concerning findings.

\medterm{Cervical Lymphadenopathy} Enlarged lymph nodes in the neck, commonly caused by infection but sometimes reflecting malignancy, autoimmune disease, or other systemic illness.

\medterm{Cervical MRI Scan} Magnetic resonance imaging of the cervical spine that uses a magnetic field and radio waves to produce detailed images of the neck vertebrae, discs, spinal cord, nerves, and surrounding soft tissues.

\medterm{Cervical Mucus} Secretions from the cervical glands that change throughout the menstrual cycle under
hormonal influence. Estrogen produces thin, clear, stretchy mucus favorable for sperm
transport around ovulation (spinnbarkeit). Progesterone produces thick, sticky, opaque
mucus that forms a plug in the cervical canal during pregnancy.

\medterm{Cervical Myelopathy} Dysfunction of the spinal cord caused by compression in the cervical spine, often producing hand clumsiness, gait disturbance, weakness, sensory changes, or abnormal reflexes.

\medterm{Cervical Osteoarthritis} Alternate index label; see \textit{Cervical spondylosis}.

\medterm{Cervical Pain} Alternate index label; see \textit{Neck pain}.

\medterm{Cervical Pain (Neck Pain)} Alternate index label for Neck Pain.

\medterm{Cervical Polyp} A usually benign polypoid overgrowth of endocervical or ectocervical tissue that may cause intermenstrual or postcoital bleeding or may be found incidentally during examination.

\medterm{Cervical Polyps} Usually noncancerous growths arising from the cervix or its canal. They may cause bleeding between periods, bleeding after sex, or discharge and are often removed for laboratory examination.

\medterm{Cervical Pregnancy} A rare ectopic pregnancy implanted in the cervical canal rather than the main uterine cavity. It may cause painless early-pregnancy bleeding and can lead to severe hemorrhage, requiring urgent specialist care.

\medterm{Cervical Radiculopathy} Irritation or compression of a nerve leaving the neck, often from a disc problem or arthritis. It can cause neck pain spreading into the arm, with tingling, numbness, or weakness.

\medterm{Cervical Rib} An extra rib growing from a neck vertebra, usually the lowest one. It often causes no symptoms but can narrow the space for nerves or blood vessels and produce arm pain, numbness, or swelling.

\medterm{Cervical Ripening Agents} Medicines or mechanical devices used to soften and prepare the cervix before labor
induction. Options include prostaglandins and cervical balloons; selection and dosing depend on the cervix, pregnancy,
contraindications, and local obstetric protocol.

\medterm{Cervical Screening (Pap Smear)} Screening for cervical cancer by collecting cells from the transformation zone of the
cervix. Liquid-based cytology (ThinPrep) or conventional Pap smear. HPV co-testing
(high-risk HPV 16, 18, others) improves sensitivity.

\synonyms
Pap Smear

\medterm{Cervical Spine} The seven bones in the neck, numbered C1 through C7. They support the head, protect the upper spinal cord, allow neck movement, and contain openings for important blood vessels.

\medterm{Cervical Spine CT Scan} Computed tomography of the cervical spine that uses x-rays and computer reconstruction to show the bones and alignment of the neck; it is particularly useful for detecting fractures and other bony abnormalities.

\medterm{Cervical Spine Trauma} An injury to the neck bones, discs, ligaments, spinal cord, or exiting nerves. It may cause pain, weakness, numbness, or paralysis; suspected serious injury requires immobilization and urgent assessment.

\medterm{Cervical Spondyloarthropathy} Wear or inflammation affecting the joints and discs of the neck. It may cause neck stiffness and pain, or press on nerves or the spinal cord, producing arm symptoms, weakness, or walking problems.

\medterm{Cervical Spondylosis} Age-related wear of the neck, including changes in discs, bones, joints, and ligaments. It may cause neck pain and stiffness and, if nerves or the spinal cord are compressed, limb symptoms.

\synonyms
Cervical Osteoarthritis
Osteoarthritis, Cervical
Spondylosis, Cervical

\medterm{Cervical Squamous Cell Carcinoma} A malignant squamous tumor of the cervix, most often caused by persistent high-risk human papillomavirus infection; it may present with abnormal bleeding, discharge, or an abnormal screening result.

\medterm{Cervical Stenosis} Narrowing or closure of the canal through the cervix, sometimes after surgery, radiation, infection, or menopause. It may cause painful or absent periods, trapped blood in the uterus, infertility, or no symptoms.

\medterm{Cervical Vertebrae} The seven vertebrae of the neck. Typical cervical vertebrae have small bodies, bifid
spinous processes, and transverse foramina transmitting the vertebral vessels; C1, C2,
and C7 are specialized.

\medterm{Cervical Vertebrae (C1-C7)} The seven vertebrae of the neck, numbered C1 through C7; they support the head, protect the cervical spinal cord, and allow neck movement.
\medterm{Cervicitis} Inflammation of the cervix, commonly caused by a sexually transmitted infection, changes in normal vaginal bacteria, chemical irritation, or allergy. It may cause discharge, bleeding after sex, or pelvic discomfort, but can be symptomless.

\medterm{Cervicogenic Headache} A headache caused by a disorder of the cervical spine or its soft tissues, with pain referred from the neck to the head.

\synonyms
Cervical headache; neck-related headache.

\medterm{Cervix} The cervix is the lower part of the uterus projecting into the vagina and producing mucus, allowing menstrual flow out and sperm passage in.

\medterm{Cervix Carcinoma} Cancer arising from the cervix, most often related to persistent high-risk human papillomavirus infection; screening and HPV vaccination reduce risk, and treatment depends on stage.

\medterm{Cervix Cryosurgery} Freezing abnormal surface cells of the cervix with a very cold probe. It may treat selected precancerous changes after appropriate testing; watery discharge and cramping are common temporarily.

\medterm{Cervix Disease} Any disorder affecting the lower part of the uterus, including infection, polyps, precancer, cancer, scarring, or pregnancy-related changes. Symptoms and treatment depend on the specific condition.

\medterm{Cervix Neoplasm} A benign, precancerous, or malignant growth of the cervix; evaluation may include pelvic examination, HPV testing, cytology, colposcopy, and biopsy.

\medterm{Cervix Uteri} The cervix uteri is the lower portion of the uterus that projects into the vagina and forms the passage between the uterus and birth canal.
\medterm{Cesarean Delivery} Surgical delivery of the fetus through abdominal and uterine incisions. Performed when
vaginal delivery is contraindicated or fails. Indications include malpresentation,
placenta previa, previous cesarean, fetal distress, and failed induction. Most common
uterine incision is low transverse (classical for very preterm or transverse lie).

\medterm{Cesarean Section} Delivery of a baby through surgical openings in the abdomen and uterus. It may be planned or urgent when vaginal birth is unsafe or unsuccessful, such as with fetal distress, abnormal position, placenta covering the cervix, or labor problems. Risks include bleeding, infection, blood clots, and complications in later pregnancies.

\medterm{Cesarean Section (C-Section)} Delivery of a baby through surgical cuts in the abdomen and uterus. It may be planned or emergency and can cause bleeding, infection, blood clots, or complications in future pregnancies.
\synonyms
C-section

\medterm{Cesium} A chemical element whose radioactive isotope cesium-137 can cause serious external or internal radiation injury after contamination or exposure.

\medterm{Cestoda} The tapeworm class of flatworms; human infection follows ingestion of eggs or larval tissue and can cause intestinal disease or tissue cysts.

\medterm{Cestode} A cestode is a tapeworm, a segmented parasitic flatworm that can infect humans or animals through contaminated food, water, or intermediate hosts.
\medterm{Cestodiasis} Infection or infestation caused by cestodes, commonly called tapeworms. Disease may
involve the intestine or tissues, depending on the species and life-cycle stage.

\synonyms
tapeworms.

\medterm{Cetirizine} A second-generation H1 antihistamine used for allergic rhinitis and urticaria; it is less sedating than older antihistamines but can still cause drowsiness.

\medterm{Cetoleic Acid} A long-chain monounsaturated fatty acid found mainly in marine oils and studied for effects on lipid metabolism and cellular membranes.

\medterm{Cetrimide} Cetrimide is a quaternary ammonium antiseptic and detergent used in some topical or laboratory preparations; concentrated exposure irritates tissue.
\medterm{Cetrorelix} A hormone-blocking medicine used during assisted reproduction to prevent premature ovulation. It temporarily blocks gonadotropin-releasing hormone receptors; injection-site reactions, headache, nausea, or allergy may occur.
\medterm{Cetuximab} A monoclonal antibody that blocks the epidermal growth factor receptor. It is used for selected colorectal, head-and-neck, and other cancers and commonly causes acne-like rash.

\medterm{Cevimeline} A muscarinic receptor agonist used to stimulate salivary secretion and relieve dry mouth, especially in selected patients with Sjögren syndrome.

\medterm{Cevimeline Hydrochloride} A medicine that stimulates muscarinic receptors to increase saliva. It is used mainly for dry mouth in Sjögren syndrome and can cause sweating, nausea, or blurred vision.

\medterm{CF} Alternate index label; see \textit{Cystic fibrosis}.

\medterm{CFS} Chronic fatigue syndrome, a condition involving substantial persistent fatigue with post-exertional worsening and associated symptoms that are not explained by another disorder.

\synonyms
Chronic fatigue syndrome; myalgic encephalomyelitis/chronic fatigue syndrome.

\medterm{CFTR Modulator Therapy} Medicines that improve the amount or function of defective CFTR protein in people with cystic fibrosis. The appropriate combination depends on the person's CFTR variants; treatment can improve breathing, nutrition, and salt-water balance.

\medterm{CHA2DS2-VASc Score} A score estimating stroke risk in a person with atrial fibrillation. It considers heart failure, high blood pressure, age, diabetes, previous stroke, blood-vessel disease, and sex; the result helps clinicians decide whether a blood thinner may offer more benefit than harm.

\synonyms
CHA2DS2-VASc; Atrial Fibrillation Stroke Risk Score

\medterm{Chaetomium} A group of molds found in soil and decaying plant material. They rarely infect people but may cause skin, eye, sinus, or widespread infection, mainly when immune defenses are severely weakened.

\medterm{Chafing} Skin irritation caused by repeated friction, moisture, heat, or tight clothing, producing redness, burning, tenderness, and sometimes erosions in opposing or rubbing areas.

\medterm{Chagas Cardiomyopathy} Heart damage caused by long-term Chagas disease. It can cause an irregular heartbeat, an enlarged heart, heart failure, or blood clots that may lead to a stroke.

\medterm{Chagas Disease} An infection caused by the parasite \textit{Trypanosoma cruzi}, usually spread by infected kissing bugs. It can also pass from mother to baby or through contaminated blood, organs, or food. Early illness may be mild or absent; years later, some people develop heart disease or enlargement of the oesophagus or colon.

\synonyms
Trypanosomiasis, American
American Trypanosomiasis

\medterm{Chagoma} A red, swollen lump where the parasite causing Chagas disease enters the body. It may be painful and accompanied by swollen nearby lymph nodes during the early stage of infection.

\medterm{Chain Of Custody} Meticulous paper trail and chronological documentation tracking the seizure, custody,
control, transfer, analysis, and disposition of physical and electronic evidence in
medico-legal investigations (e.g., blood samples for toxicology, bullets, sexual assault
kits, biological fluids).

\medterm{Chain-Termination Codon} A codon that signals the end of protein translation; the three standard stop codons are UAA, UAG, and UGA in messenger RNA.

\medterm{Chalazion} A usually painless, firm eyelid lump caused by blockage and inflammation of a meibomian oil gland; persistent or atypical lesions require ophthalmic assessment.

\medterm{Challenge Agent} A substance, organism, allergen, or physiologic stimulus deliberately administered in a controlled test to provoke and measure a response.

\medterm{Challenge Testing} A controlled exposure to a suspected trigger, such as a food or medication, to assess whether it produces a reproducible reaction.

\synonyms
Provocation testing; challenge test.

\medterm{Chamomile Extract} An extract of chamomile flowers used in traditional and complementary products; allergic reactions and medicine interactions are possible, while evidence varies by indication.

\medterm{Chancre} A usually painless ulcer at the site of inoculation of an infection, especially the
primary lesion of syphilis. It may be accompanied by regional lymph-node enlargement and
heals without treatment, although infection persists.

\medterm{Chancroid} A sexually transmitted infection caused by Haemophilus ducreyi, characterized by painful ragged genital ulcers and tender suppurative inguinal lymph nodes.

\medterm{Changes In The Newborn At Birth} The newborn transition from placental support to independent breathing and circulation, including lung aeration, closure of fetal shunts, temperature control, and adaptation of glucose and feeding.



\medterm{Channelopathy} A disorder caused by abnormal proteins that control movement of charged particles through cell membranes. It can affect nerves, muscles, or the heart and may cause seizures, weakness, or abnormal rhythms.

\medterm{Chaplain} A trained spiritual-care professional who supports patients and families with religious, existential, emotional, or ethical concerns in healthcare settings.

\medterm{Chapped Hands} Dry, cracked, inflamed skin of the hands caused by irritants, cold, low humidity, frequent washing, eczema, or allergy; fissures can become painful or infected.

\medterm{Chapped Lips} Dry, scaling, fissured lips caused by weather, licking, dehydration, irritants, medicines, or cheilitis; persistent lesions need assessment for infection or another disorder.

\medterm{Charcoal} Fine black powder produced by controlled pyrolysis of organic materials, processed to
yield a massive micro-porous surface area (500-1500 m2/g). See activated charcoal.

\medterm{Charcot Arthropathy} Progressive joint damage caused by loss of protective sensation, so repeated injuries go unnoticed. The joint may become warm, swollen, unstable, and deformed, especially in people with diabetic nerve damage.

\medterm{Charcot Foot} Progressive weakening and deformity of the foot or ankle caused by repeated unnoticed injury in a person with reduced sensation, commonly from diabetes. Early warmth and swelling require urgent off-loading to prevent collapse and ulcers.

\medterm{Charcot Joint} Severe destruction of a joint caused by loss of pain and position sensation. Repeated unnoticed injury can fragment bone, loosen or dislocate the joint, and produce major deformity.

\medterm{Charcot Triad} The combination of fever, pain in the upper-right abdomen, and jaundice, traditionally associated with acute infection of blocked bile ducts. This pattern is important but may be absent even in serious disease.

\medterm{Charcot Triad (Cholangitis)} Classic diagnostic triad of acute ascending cholangitis: (1) Fever, (2) Right Upper
Quadrant abdominal pain, and (3) Jaundice. Reynolds Pentad adds Hypotension and Altered
Mental Status. See Acute Cholangitis.

\synonyms
Cholangitis

\medterm{Charcot'S Joints} Charcot’s joints are severely damaged, unstable joints caused by loss of protective sensation, classically in diabetic neuropathy or other neurologic disease.
\medterm{Charcot'S Triad} Charcot’s triad is fever, jaundice, and right-upper-quadrant pain, a classic but incomplete clinical pattern of acute ascending cholangitis.
\medterm{Charcot-Marie-Tooth Disease} A group of inherited disorders that damage nerves supplying the feet, legs, hands, and arms. Symptoms usually include gradually worsening weakness, muscle wasting, numbness, foot deformity, and difficulty walking.

\synonyms
Neuropathy, Hereditary Motor And Sensory
Peroneal Muscular Atrophy

\medterm{Charcot-Wilbrand Syndrome} A rare neuropsychiatric syndrome after brain injury characterized by loss of visual imagery and visual dreaming, sometimes with difficulty recognizing or interpreting visual information.

\medterm{Chargaff Rule} The observation that double-stranded DNA contains approximately equal amounts of adenine and thymine and of guanine and cytosine, reflecting complementary base pairing.

\medterm{CHARGE Syndrome} A genetic condition, usually involving a new change in the CHD7 gene, that can affect the eyes, ears, heart, breathing passages, growth, development, and reproductive system. Features vary widely.

\medterm{Charley Horse} A sudden painful muscle cramp, commonly in the calf or thigh. It may follow exercise, dehydration, prolonged positioning, nerve problems, or certain medicines and usually lasts from seconds to minutes.

\medterm{Chasing The Dragon} Inhalation of vapor produced by heating heroin, a method of drug use that can cause dependence, overdose, toxic exposures, and neurological injury.

\medterm{Chaulmoogra Oil} Chaulmoogra oil is a historical plant oil once used to treat leprosy and some skin diseases; it has largely been replaced by effective antimicrobial therapy.
\medterm{CHD} Alternate index label; see \textit{Hip dysplasia}.

\medterm{Checkpoints} Regulatory signals that help prevent immune cells from attacking healthy tissues. Some cancers use checkpoint pathways to evade immunity; checkpoint-inhibitor medicines block selected signals to restore antitumor immune activity.

\medterm{Chediak-Higashi Syndrome} Autosomal recessive (LYST mutation, lysosomal trafficking regulator). Giant lysosomal
granules in neutrophils, melanocytes, neurons. Partial oculocutaneous albinism, silvery
hair, photophobia, nystagmus.

\medterm{Cheek} The facial region forming the side wall of the mouth, containing skin, muscle, fat, salivary structures, nerves, and blood vessels.

\medterm{Cheek Pouch} A distensible pocket inside the cheek found in some rodents and other animals, used to temporarily store and transport food.

\medterm{Cheilectropion} An eversion deformity of the lip, usually involving the lower lip. It may result from
scarring, congenital malformation, or loss of tissue support.

\medterm{Cheilitis} Inflammation of one or both lips, causing dryness, scaling, fissuring, erythema,
swelling, pain, or ulceration. Causes include irritant or allergic contact exposure,
infection, nutritional deficiency, chronic sun exposure, and inflammatory skin disease;
persistent lesions require evaluation for premalignant or malignant change.

\medterm{Cheilodynia} Pain of the lips, which may result from inflammation, infection, trauma, allergy, dryness, nerve pain, or another oral or dermatologic disorder.

\medterm{Cheilognathopalatoschisis} A congenital cleft involving the lip, jaw, and palate, caused by incomplete fusion of facial processes during embryonic development.

\medterm{Cheilognathoschisis} A congenital cleft of the lip and alveolar part of the upper jaw, with or without associated palatal involvement.

\medterm{Cheiloplasty} Plastic or reconstructive surgery of the lip, performed to repair congenital clefting,
trauma, tumors, scars, or other deformity.

\medterm{Cheiloschisis} Cleft lip, a congenital fissure of the upper lip caused by incomplete fusion of facial
processes during embryonic development. It may occur alone or with cleft palate.

\medterm{Cheilosis} A disorder of the lips characterized by dryness, scaling, fissuring, or inflammation.
Causes include irritation, nutritional deficiency, infection, and systemic disease.

\medterm{Cheiromegaly} Abnormal enlargement of one or both hands, classically associated with excess growth hormone but also possible with local overgrowth or other disorders.

\medterm{Cheiropompholyx} A type of hand eczema that causes intensely itchy, small fluid-filled blisters on the palms or sides of the fingers. The skin may later peel, crack, or become painful.

\medterm{Chelating Agent} A compound that binds a metal ion through multiple chemical sites, allowing the metal to be measured, mobilized, or removed in selected toxic exposures.

\medterm{Chelating Agents} Chelating agents bind specific metal ions to form complexes that can facilitate their removal or reduce toxicity; examples include agents for lead, iron, or copper poisoning.
\medterm{Chelation} The formation of a stable complex between a metal ion and a molecule with multiple
binding sites. In medicine, chelation can be used under supervision to bind and promote
elimination of selected toxic metals.

\medterm{Chelation Study} A test that uses a metal-binding substance and imaging or blood measurements to assess how the body handles a mineral or to locate abnormal deposits. The exact method depends on the clinical question.

\medterm{Chelation Therapy} Treatment using a medicine that binds a specific harmful or excess metal so the body can remove it in urine or stool. It is used only for selected poisonings or disorders because it can harm the kidneys or alter essential minerals.

\medterm{Chemexfoliation} Application of a chemical solution to remove controlled layers of skin, improving selected scars, pigmentation, wrinkles, or superficial lesions.

\medterm{Chemical} A substance defined by its molecular composition and properties; medical effects depend on dose, route, exposure duration, metabolism, and biological target.

\medterm{Chemical Ablation} Destruction of unwanted tissue by injecting or applying a chemical irritant or sclerosant, used in selected lesions, vessels, cysts, or other conditions.

\medterm{Chemical Burn} Damage to skin, eyes, or internal tissues caused by a corrosive chemical. Injury may worsen while the chemical remains present; contaminated clothing should be removed and the area flushed promptly with plenty of water.

\medterm{Chemical Burn Or Reaction} Tissue injury caused by contact with a corrosive or irritating chemical, commonly affecting the skin or eyes; severity depends on the substance, concentration, contact time, and tissue involved.

\medterm{Chemical Carcinogenesis} Development of cancer after exposure to chemicals that cause DNA damage, abnormal cell signaling, chronic tissue injury, or altered gene regulation.

\medterm{Chemical Cleavage} Breaking a chemical bond with a reagent, used in laboratory analysis, molecular biology, synthesis, or controlled degradation of a larger molecule.

\medterm{Chemical Cofactor} A nonprotein chemical helper required for an enzyme or biologic reaction, such as a metal ion, vitamin-derived coenzyme, or other small molecule.


\medterm{Chemical Dependency} Legacy, nonpreferred label; see \textit{Substance Use Disorder}.

\medterm{Chemical Eye Injury} Ocular damage caused by exposure to an acidic or alkaline substance. Immediate copious
irrigation, removal of particulate material, measurement of ocular-surface pH, and
urgent ophthalmic assessment are required; alkali injuries may penetrate deeply and
continue to worsen after exposure.

\medterm{Chemical Fractionation} Separation of a mixture into fractions according to chemical or physical properties such as solubility, charge, size, or binding affinity.

\medterm{Chemical Impurity} An unintended chemical present in a medicine, specimen, or material, potentially affecting safety, stability, measurement accuracy, or biologic activity.

\medterm{Chemical Injury} Tissue damage caused by contact with a corrosive or irritating chemical. Acids and alkalis can destroy tissue differently, and immediate thorough rinsing, removal of contaminated clothing, and medical assessment may be needed.

\medterm{Chemical Peel} Application of a chemical solution to the skin to produce controlled exfoliation and regeneration of superficial or deeper skin layers.

\synonyms
Chemexfoliation; chemical resurfacing.

\medterm{Chemical Peels} Treatments that apply a chemical solution to remove a controlled amount of skin. Superficial peels affect the outer layer; deeper peels reach the underlying dermis and have greater risks of scarring, pigment changes, infection, and prolonged healing.

\medterm{Chemical Pneumonitis} Lung inflammation and injury after inhaling or breathing in irritating chemicals, stomach contents, fuels, or other toxic substances. It may cause cough, shortness of breath, wheezing, low oxygen, and chest imaging changes.

\medterm{Chemistry} The study of substances, their composition, properties, reactions, and transformations, forming the basis of biochemistry, pharmacology, toxicology, and laboratory medicine.

\medterm{Chemo} Informal shorthand for chemotherapy, treatment using drugs that destroy cancer cells, inhibit their growth, or suppress abnormal immune cells.

\medterm{Chemo Brain} Problems with memory, concentration, word-finding, or mental processing that may occur during or after cancer treatment. Symptoms can affect work and daily activities and may have several contributing causes.

\synonyms
Post-Chemotherapy Cognitive Impairment

\medterm{Chemoattractant} A chemical signal that guides cells, especially immune cells, toward its source. Such signals help direct cells to areas of infection, inflammation, tissue repair, or normal development.
\medterm{Chemodectoma} A tumor arising from paraganglionic tissue, especially a carotid-body tumor. It is
usually slow-growing and may compress nearby cranial nerves or vessels.

\medterm{Chemoembolization} A procedure that combines delivery of an antineoplastic agent with embolization of the
blood supply to a tumor, most commonly for selected liver cancers.

\medterm{Chemokine} A small signaling protein that directs immune cells to sites of infection, inflammation, or tissue repair. Different chemokines attract different cell types by binding to matching receptors.

\medterm{Chemokine Receptor} A cell-surface receptor that binds chemokines and directs immune-cell movement; some also serve as entry cofactors for viruses such as HIV.

\medterm{Chemokinesis} Increased or altered movement of a cell in response to a chemical, without movement specifically toward or away from the chemical source. It differs from chemotaxis, which is directional.

\medterm{Chemonucleolysis} An intradiscal treatment intended to dissolve part of a herniated intervertebral disk, historically using an enzyme preparation.

\synonyms
Enzymatic disk dissolution; chemonucleolysis procedure.

\medterm{Chemoprevention} The use of drugs or natural agents to prevent, delay, or reverse cancer development,
e.g., tamoxifen or raloxifene in high-risk breast cancer, HPV and hepatitis B
vaccination, and aspirin in colorectal cancer.

\medterm{Chemoradiation} The concurrent combination of chemotherapy and radiation therapy, used to treat
locoregionally advanced cancers such as anal, cervical, head and neck, esophageal, and
limited-stage small cell lung cancer.

\medterm{Chemoradiotherapy} Cancer treatment that combines chemotherapy with radiation treatment. The medicines may make cancer cells more sensitive to radiation or treat cells elsewhere, but combined treatment can increase tiredness, nausea, skin changes, and other side effects.

\medterm{Chemoreceptor} A sensory receptor that responds to changes in blood or CSF chemistry, involved in
respiratory and cardiovascular regulation.

\medterm{Chemoreceptor Cell} A sensory cell that detects chemical changes such as oxygen, carbon dioxide, pH, taste molecules, or odorants and converts them into neural or physiologic signals.

\medterm{Chemoreceptor Trigger Zone} A brain area that detects vomiting-provoking chemicals in the blood or cerebrospinal fluid and activates brain circuits that produce nausea and vomiting.

\medterm{Chemoreceptors} Chemoreceptors detect chemical changes such as oxygen, carbon dioxide, pH, or odor molecules and trigger neural or physiologic responses.
\medterm{Chemoresistance} Reduced response of cancer cells or other target cells to a medicine that previously inhibited or destroyed them. It may be present from the start or develop during treatment.

\medterm{Chemosensitivity} How strongly cells or a tumor respond to a chemical treatment. A sensitive tumor may shrink or stop growing, but laboratory sensitivity tests do not always predict a person's treatment response.

\medterm{Chemosensitivity Assay} A laboratory test exposing cells to medicines to assess response; results are mainly investigational and may not predict treatment benefit.

\medterm{Chemosensitization} Making cancer cells or microorganisms more vulnerable to treatment by adding an agent that improves delivery, blocks resistance, or increases treatment-related damage.

\medterm{Chemosis} Swelling of the clear membrane covering the white part of the eye. It may result from allergy, infection, irritation, injury, surgery, or fluid retention and can make the eye look blistered.

\medterm{Chemosurgery} Chemosurgery uses a chemical agent to destroy selected abnormal tissue, historically including chemical treatment of skin lesions; modern use is specialized.
\medterm{Chemotactic Factors} Chemical signals that attract or repel cells by forming a concentration gradient. During infection or tissue injury, they help direct immune cells to the affected area.

\medterm{Chemotaxis} Directed movement of a cell or organism toward or away from a chemical stimulus. It is
important in immune-cell recruitment and inflammatory responses.

\medterm{Chemotherapy} Treatment using medicines to kill cancer cells or slow their growth, often by disrupting cell division or damaging DNA. It may be used alone or with surgery or radiation. Healthy fast-growing cells can also be affected, causing nausea, hair loss, anemia, or reduced immunity.


\medterm{Chemotherapy Extravasation} Leakage of an intravenous cancer medicine into tissue around the vein. Some medicines can cause severe pain, blistering, or tissue death; the infusion must stop immediately and the affected area needs urgent protocol-based care.

\medterm{Chemotherapy Regimen} A specific combination of anticancer drugs, dosing schedule, and treatment cycles designed for a particular cancer type and stage. Regimens are standardized by evidence-based protocols and may be modified based on response, toxicity, and patient factors.

\medterm{Chemotherapy Treatment} Chemotherapy uses medicines that damage or inhibit rapidly dividing cancer cells. It may be given before or after local treatment, for advanced disease, or for symptom control; effects and monitoring depend on the drugs, cancer, treatment goal, and organ function.

\medterm{Cherry Angioma} A common benign proliferation of small dermal blood vessels that appears as a bright-red
to purple papule, usually on the trunk or limbs. Lesions may bleed after trauma but
generally require no treatment unless diagnosis is uncertain or removal is desired.

\medterm{Cherry Eye} Prolapse of the gland of the third eyelid, producing a rounded red mass at the inner
corner of the eye. The term is used especially in veterinary ophthalmology.

\medterm{Cherry-Red Spot} A bright-red area at the center of the retina surrounded by pale retina. It occurs when the surrounding retina becomes swollen or loses blood flow and may indicate a serious metabolic, vascular, or storage disorder.

\medterm{Cherubism} A rare inherited disorder characterized by bilateral enlargement of the jaws from
expansile giant-cell lesions, usually beginning in childhood and often regressing after
puberty.

\medterm{Chest} The thorax, the region between the neck and abdomen containing the heart, lungs, major vessels, and other structures protected by the ribs and sternum.

\medterm{Chest CT} Computed tomography of the chest using x-rays and computer reconstruction to produce cross-sectional images of the lungs, heart, blood vessels, mediastinum, pleura, and chest wall.

\medterm{Chest MRI} Magnetic resonance imaging of the chest using a magnetic field and radio waves to evaluate soft tissues, blood vessels, heart structures, mediastinum, chest wall, and selected lung or pleural abnormalities.

\medterm{Chest Pain} Pain, pressure, burning, or discomfort in the chest. Causes include heart, lung, digestive, muscle, bone, or anxiety-related problems. New severe pain, breathlessness, sweating, faintness, or spreading pain requires emergency care.

\synonyms
Pain, Chest

\medterm{Chest Pain On The Left Side Above A Female Breast} Left-sided chest or breast-area pain can arise from muscle or rib injury, breast tissue, reflux, lung disease, anxiety, or heart disease. Pressure, exertional pain, shortness of breath, sweating, faintness, or pain spreading to the arm or jaw requires emergency assessment.

\medterm{Chest Physiotherapy} Techniques to mobilize and clear secretions from the lungs, including postural drainage,
percussion, vibration, and huffing. It is used in bronchiectasis, CF, and post-operative
patients. Positive expiratory pressure devices and oscillating PEP devices are also
used.


\medterm{Chest Tightness} A sensation of pressure, constriction, or heaviness in the chest that may arise from heart, lung, gastrointestinal, musculoskeletal, or anxiety-related causes.

\medterm{Chest Tube} A flexible tube inserted into the pleural space to remove air, blood, or fluid and help the lung re-expand.

\medterm{Chest Tube Insertion} Placement of a flexible tube through the chest wall into the space around a lung to remove air, blood, pus, or fluid and help the lung expand. It is connected to a drainage system and requires monitoring for bleeding, infection, or persistent air leak.

\medterm{Chest Tube Thoracostomy} Insertion of a tube into the pleural space to drain air, fluid, or blood. Indications:
pneumothorax (large, tension, or on mechanical ventilation), hemothorax, pleural
effusion (empyema, malignant, transudative), and postoperative (cardiac, thoracic
surgery).

\medterm{Chest Wall} The musculoskeletal boundary of the thoracic cavity, formed by the skin, fascia,
muscles, ribs, costal cartilages, sternum, thoracic vertebrae, and associated
neurovascular structures. It protects the thoracic organs and participates in breathing
through coordinated movement of the ribs and respiratory muscles.

\medterm{Chest Wall Deformity} An abnormal shape of the ribs, breastbone, spine, or surrounding tissues. A sunken or protruding breastbone is often mainly cosmetic, but severe deformity or scoliosis can restrict breathing or affect the heart.

\medterm{Chest Wall Pain} Alternate index label; see \textit{Costochondritis}.

\medterm{Chest X-Ray} An imaging test using a small amount of radiation to show the lungs, heart, ribs, and chest structures. It can detect pneumonia, collapsed lung, fluid, some fractures, and changes in heart size.

\medterm{Chewing} The coordinated use of teeth, jaw muscles, and joints to break food into smaller pieces for swallowing and digestion; pain or missing teeth can interfere.

\medterm{Cheyne-Stokes Respiration} A repeating pattern in which breathing gradually becomes deeper and faster, then shallower, followed by a pause. It may occur during sleep with heart failure, stroke, brain disease, or at high altitude.

\medterm{CHF} Heart failure in which the heart cannot pump or fill well enough for the body's needs. Reduced circulation can cause breathlessness, tiredness, leg swelling, and fluid buildup in the lungs or elsewhere.

\medterm{Chiari I Malformation} A condition in which the lower part of the cerebellum extends through the opening at the skull base. It may cause headaches worsened by straining, neck pain, balance problems, or no symptoms.

\medterm{Chiari Malformation} A structural abnormality in which part of the cerebellum extends downward through the opening at the base of the skull. Different types vary in severity and may cause headache, balance problems, swallowing difficulty, or no symptoms.

\medterm{Chiari-Frommel Syndrome} Continued milk production and absent menstrual periods after pregnancy, usually caused by persistently high prolactin levels. Evaluation looks for pituitary disorders, medicines, thyroid disease, and other causes.

\medterm{Chick Embryo} The developing chicken embryo, widely used as an experimental model for embryology, organ development, virology, toxicology, and tissue growth.

\medterm{Chicken Soup And Sickness} A supportive-care topic concerning the use of warm fluids and nutrition during illness; chicken soup may ease hydration or congestion but does not replace diagnosis or treatment.

\medterm{Chickenpox} A highly contagious infection caused by varicella-zoster virus. It causes an itchy rash that appears as spots, fluid-filled blisters, and scabs at different stages, with fever and tiredness; complications are more likely in adults and people with weakened immunity.

\medterm{Chickenpox (Varicella)} Alternate index label for Varicella.


\medterm{Chigger Bites} Itchy skin reactions caused by larval mites. Small red bumps or blisters often develop where clothing is tight and may itch for days. Washing the skin, avoiding scratching, and itch treatment usually provide relief.

\medterm{Chiggers} Larval harvest mites whose bite causes intensely itchy grouped red papules or welts; symptoms result from skin irritation, not from the mite burrowing into human tissue.

\medterm{Chiggers (Bites)} Alternate index label for Bites.

\medterm{Chignon} Temporary swelling of a baby's scalp where a vacuum device was applied during assisted delivery. It results from the suction and usually disappears within a day or two without treatment.

\medterm{Chikungunya} A mosquito-borne viral disease caused by chikungunya virus, usually producing abrupt fever and severe, often symmetric joint pain. Headache, muscle pain, joint swelling, nausea, fatigue, and rash may occur; joint symptoms can persist for months or years.

\synonyms
Chikungunya fever; chikungunya virus disease; CHIKV infection

\medterm{Chikungunya Virus} An alphavirus transmitted primarily by infected female Aedes aegypti and Aedes albopictus mosquitoes. Infection can cause abrupt fever, severe polyarthralgia, joint swelling, rash, and prolonged inflammatory joint symptoms.

\medterm{Chilblain} An inflammatory skin lesion caused by an abnormal response to cold and damp conditions.
It produces itchy, red or violaceous swollen areas, usually on fingers, toes, ears, or
nose.

\medterm{Chilblains} Itchy, painful, red or purple swellings caused by an abnormal skin response after exposure to cold and damp conditions. They usually improve with warmth but can blister, ulcerate, or become infected.
\synonyms
Pernio

\medterm{Chilblains (Pernio)} Itchy, painful, red or purplish inflammatory lesions caused by an abnormal response of small blood vessels to cold and subsequent rewarming.

\synonyms
Pernio

\medterm{Child Abuse} Physical, sexual, or emotional harm or neglect of a child by a caregiver. Recognized by
inconsistent histories, injury patterns inconsistent with mechanism, and behavioral
indicators. Types: physical abuse, sexual abuse, neglect, and emotional abuse. Mandatory
reporting requirements for healthcare providers.

\synonyms
Child Maltreatment


\medterm{Child Adoption} Child adoption is the legal process by which parental rights and caregiving responsibility are transferred to adoptive parents; medical care includes records, consent, and family history.
\medterm{Child Behavior} The observable actions, responses, and conduct of children, influenced by developmental stage, temperament, environment, and neurobiological factors. Behavioral assessment helps identify developmental disorders, emotional difficulties, and parenting needs.


\medterm{Child Development} Progressive physical, cognitive, language, emotional, and social changes from infancy through adolescence, assessed against developmental milestones and individual variation.

\medterm{Child Health} The physical, developmental, mental, and social well-being of children, including prevention, nutrition, immunization, safety, and treatment of illness.

\medterm{Child Health Services} Child health services provide preventive care, immunization, growth and developmental monitoring, screening, treatment, and support for children and families.

\medterm{Child Neglect} The failure of a caregiver to provide adequate food, shelter, clothing, supervision,
education, or medical care. The most common form of child maltreatment. Types: physical
neglect, educational neglect, emotional neglect, and medical neglect. May result in
failure to thrive, developmental delay, and chronic illness.

\medterm{Child Neglect And Emotional Abuse} Failure to provide a child’s basic needs or repeated behavior that harms emotional development, safety, or wellbeing; suspected abuse requires safeguarding assessment and appropriate reporting under local law.

\medterm{Child Physical Abuse} Intentional physical injury or harm inflicted on a child by a caregiver or another person, including hitting, burning, shaking, or fractures; suspected abuse requires immediate safeguarding assessment.

\medterm{Child Poisoning} Exposure of a child to a medicine, chemical, or other substance that may cause harm. Suspected poisoning requires prompt expert
assessment.

\medterm{Child Psychiatry} The medical specialty concerned with assessment and treatment of emotional, behavioral, developmental, and psychiatric disorders in children and adolescents.

\medterm{Child Psychology} The study of children’s cognitive, emotional, social, and behavioral development and the effects of relationships, learning, biology, and environment.



\medterm{Child-Pugh Classification} Clinical scoring system assessing prognosis of chronic liver disease by evaluating five
parameters: total bilirubin, serum albumin, international normalized ratio, ascites, and
hepatic encephalopathy. Each parameter is scored one to three points.

\medterm{Childbed Fever} Puerperal fever, usually infection of the reproductive tract after childbirth or miscarriage, often involving postpartum endometritis and requiring prompt evaluation and treatment.

\medterm{Childbirth} The process of giving birth, encompassing three stages: (1) First stage: from the onset
of regular uterine contractions (cervical effacement and dilation) to full dilation (10
cm). See Parturition (Childbirth).

\medterm{Childhood} The developmental period between infancy and adolescence, during which physical growth, brain development, learning, and social skills progress.

\medterm{Childhood Absence Epilepsy} An epilepsy syndrome beginning in childhood with frequent brief absence seizures,
characterized by sudden impairment of awareness and a typical generalized spike-and-wave
EEG pattern.

\medterm{Childhood Angiosarcoma} A rare malignant tumor arising from blood-vessel or lymphatic-endothelial cells in a child; diagnosis and treatment require specialist pediatric oncology care.

\medterm{Childhood Apraxia Of Speech} A speech-sound disorder in which a child has difficulty planning and coordinating the movements needed for clear speech.

\synonyms
CAS

\medterm{Childhood Asthma} Asthma beginning during childhood, involving variable airway narrowing and inflammation that cause wheezing, cough, chest tightness, or shortness of breath.

\medterm{Childhood Ependymoma} A tumor arising from ependymal cells lining the brain ventricles or spinal canal in children; symptoms depend on location and obstruction of cerebrospinal-fluid flow.

\medterm{Childhood Fibrosarcoma} A malignant tumor of fibroblastic connective tissue in children, usually presenting as a growing mass; treatment depends on location, extent, and tumor biology.

\medterm{Childhood Ganglioglioma} A usually slow-growing mixed neuronal and glial brain tumor in children or young adults, often presenting with long-standing focal seizures.

\medterm{Childhood Leiomyosarcoma} A rare malignant tumor of smooth muscle in a child; it may arise in soft tissue or an organ and requires biopsy, staging, and specialist treatment.

\medterm{Childhood Leukemia} Cancer of blood-forming cells in a child, most often acute lymphoblastic or acute myeloid leukemia; symptoms may include pallor, bruising, fever, bone pain, or recurrent infection.

\medterm{Childhood Liposarcoma} A rare malignant tumor derived from fat-forming tissue in a child; diagnosis requires imaging and biopsy, and treatment is individualized by a sarcoma team.

\medterm{Childhood Obesity} Alternate index label; see \textit{Pediatric Obesity}.

\medterm{Childhood Osteosarcoma} A malignant bone tumor that commonly develops near the ends of long bones during adolescence; pain, swelling, or a fracture may occur and treatment combines chemotherapy and surgery.

\medterm{Childhood Rhabdomyosarcoma} A malignant tumor showing skeletal-muscle differentiation, occurring mainly in children and adolescents; it may arise in the head and neck, urinary or reproductive tract, or limbs.

\medterm{Childhood Schizophrenia} A rare psychotic disorder beginning in childhood, involving disturbances in thinking, perception, behavior, and emotional functioning.

\synonyms
Juvenile Schizophrenia
Schizophrenia, Adolescent
Schizophrenia, Childhood
Schizophrenia, Juvenile

\medterm{Children} Humans in the developmental period between infancy and adolescence; medical assessment must account for age-specific anatomy, physiology, growth, communication, and consent.

\medterm{Children And Grief} The emotional and behavioral response of children to death or major loss, which varies by developmental stage and may include sadness, regression, anger, worry, or physical symptoms.







\medterm{Children'S Cancer Centers} Specialized services that coordinate diagnosis, treatment, rehabilitation, psychosocial support, survivorship care, and research for children and adolescents with cancer.



\synonyms
Childrens Health (Children's Health)


\medterm{Chill} A chill is a sensation of cold, often with shivering, that can accompany fever, infection, exposure, anxiety, or other illness.
\medterm{Chills} A sensation of cold with involuntary shivering, commonly accompanying fever or infection; chills without measured fever may also occur with anxiety, hypoglycaemia, or environmental cold.

\medterm{Chilomastix Mesnili} A usually harmless single-celled organism that may be found in stool. Its presence generally does not cause illness or require treatment, but it can indicate exposure to contaminated food or water.

\medterm{Chimera} An organism or tissue containing genetically distinct cell populations derived from different zygotes or embryos, or a laboratory construct combining genetic material from different sources.

\medterm{Chimerism} The presence in one individual of cell populations derived from genetically different
zygotes or sources. It may occur naturally, after transplantation, or through exchange
of cells between a pregnant person and fetus.

\medterm{Chin} The lower facial prominence formed mainly by the mandible and overlying soft tissue, contributing to facial contour, speech, chewing, and airway anatomy.

\medterm{Chin Augmentation} A cosmetic procedure that increases chin projection with an implant, injectable filler, or bone surgery to change facial contour; risks include infection, asymmetry, nerve injury, and implant problems.

\medterm{Chinese Restaurant Syndrome} An older, controversial term for short-lived symptoms such as headache, flushing, or palpitations attributed to eating restaurant food; evidence does not show that monosodium glutamate causes these symptoms in most people at usual dietary exposures.

\medterm{Chionablepsia} Temporary eye pain and blurred vision caused by ultraviolet light reflected from snow or ice, commonly called snow blindness. Symptoms may appear hours after exposure and usually improve as the cornea heals.

\synonyms
snow blindness.

\medterm{Chironomidae} The nonbiting midge family of insects; their larvae and body parts can act as environmental allergens and rarely trigger respiratory or skin reactions.

\medterm{Chiropody} Chiropody is care of the feet, including assessment and treatment of skin, nail, pressure, and biomechanical problems; it is also called podiatry.
\medterm{Chiropractic} A healthcare practice that mainly uses hands-on treatment for muscle, joint, and spine problems, especially spinal manipulation. Benefits and risks vary by the condition, technique, and practitioner's training.

\medterm{Chiropractic Care For Back Pain} Manual therapy, primarily spinal manipulation and mobilization, provided for acute and chronic low back pain when indicated. Evidence supports its use for acute low back pain in particular, though it is not appropriate for all causes of back pain.

\medterm{Chiropractor} A healthcare practitioner who provides manual or manipulative treatment, especially for musculoskeletal conditions, according to the practitioner's training and local regulation.

\synonyms
Chiropractic practitioner; chiropractic clinician.


\medterm{Chitin} A structural polysaccharide in fungal cell walls and arthropod exoskeletons.

\medterm{Chitosan} A material made by chemically modifying chitin from shellfish or fungi. It is used in some wound dressings and research products; claims that oral chitosan causes major weight loss are not well supported.

\medterm{Chlamydia} A group of bacteria that must live inside host cells to multiply. Different species cause sexually transmitted infections, eye disease, respiratory illness, or infections acquired from birds.

\medterm{Chlamydia Infections In Women} Infections in women caused by Chlamydia trachomatis, commonly involving the cervix or urethra and often producing no symptoms; untreated infection can cause pelvic inflammatory disease, infertility, or ectopic pregnancy.

\medterm{Chlamydia Pneumoniae Infection} A respiratory infection caused by Chlamydia pneumoniae, formerly called Chlamydophila pneumoniae. It commonly causes gradual-onset atypical pneumonia, bronchitis, or upper-respiratory symptoms.

\medterm{Chlamydia Psittaci} A bacterium that causes psittacosis, a zoonotic infection usually acquired by inhaling organisms from infected birds.

\medterm{Chlamydia Trachomatis} A bacterium that commonly causes sexually transmitted infection of the urogenital tract and can also cause conjunctivitis or pneumonia.

\medterm{Chlamydia Trachomatis Infection} A bacterial infection spread mainly through sexual contact. It often causes no symptoms but may cause discharge or painful urination and can lead to pelvic inflammatory disease, infertility, eye infection, or newborn infection.

\medterm{Chlamydiaceae} A family of bacteria that includes Chlamydia species. Some cause sexually transmitted infections, eye disease, or respiratory infection; diagnosis depends on the species, specimen, symptoms, and laboratory testing.

\medterm{Chlamydial Infections - Male} Infections in men caused by Chlamydia trachomatis, most commonly urethritis and sometimes epididymitis or proctitis; many infections are asymptomatic and diagnosis is usually by nucleic acid amplification testing.

\medterm{Chlamydial Pneumonia} A lung infection caused by Chlamydia bacteria, usually with gradually worsening cough, fever, tiredness, headache, or sore throat. Diagnosis and treatment depend on the species, severity, and other medical conditions.

\medterm{Chlamydiosis} Infection caused by bacteria of the genus Chlamydia. Depending on the species and site,
it may cause sexually transmitted infection, conjunctivitis, pneumonia, or other
disease.

\medterm{Chloasma} Patchy brown or gray-brown darkening of the skin, usually on the face. It is often related to pregnancy, hormones, and sunlight and is commonly called melasma.

\synonyms
melasma.

\medterm{Chloasma (Melasma)} Alternate index label for Melasma.

\medterm{Chloracne} An acne-like eruption caused by exposure to certain halogenated aromatic compounds, such
as dioxins and related industrial chemicals.

\medterm{Chloral Hydrate} A sedative-hypnotic compound with limited modern use because of toxicity, dependence, respiratory depression, and cardiac or gastrointestinal adverse effects.

\medterm{Chlorambucil} An oral chemotherapy medicine that damages DNA and is used mainly for some chronic leukemias and lymphomas. It can lower blood counts and increase infection risk.

\medterm{Chloramphenicol} A broad-spectrum antibiotic with activity against gram-positive and gram-negative
organisms, as well as Rickettsia. Chloramphenicol inhibits bacterial protein synthesis
by binding to the 50S ribosomal subunit. It is indicated for serious infections when
other agents are not appropriate, including bacterial meningitis and typhoid fever.

\medterm{Chloranemia} An older term for anemia in which red blood cells contain too little hemoglobin and appear unusually pale. Iron deficiency is a common cause; modern evaluation uses blood counts and iron studies.

\medterm{Chlordane} A persistent organochlorine insecticide formerly used for termite control; exposure can affect the nervous system and liver and has environmental persistence.

\medterm{Chlordecone} A persistent organochlorine pesticide that can cause neurologic, reproductive, and liver effects after significant exposure and remains in the environment for long periods.

\medterm{Chlordiazepoxide} A benzodiazepine medicine that slows brain activity and may be used short term for alcohol-withdrawal symptoms or severe anxiety. It can cause drowsiness, dependence, impaired coordination, and dangerous breathing suppression with alcohol or opioids.

\medterm{Chlordiazepoxide Overdose} Overdose of chlordiazepoxide, a benzodiazepine, that can cause drowsiness, impaired coordination, respiratory depression, and coma, especially with other sedatives.

\medterm{Chlorethoxyfos} An organophosphate insecticide that can block a nerve enzyme, causing sweating, vomiting, weakness, breathing problems, seizures, or other poisoning symptoms after exposure.

\medterm{Chlorfenvinphos} An organophosphate insecticide that can cause cholinergic toxicity by inhibiting acetylcholinesterase; use is restricted or prohibited in many places.

\medterm{Chlorhexidine} A broad-spectrum antiseptic used for skin, mucosal, and oral preparation; it can cause irritation, allergy, or rare severe anaphylaxis.

\medterm{Chloride} The principal negatively charged electrolyte in extracellular fluid, important for fluid balance, electrical neutrality, and gastric acid formation.

\medterm{Chloride - Urine Test} A test that measures chloride excretion in urine; interpreted with serum electrolytes and clinical findings to evaluate fluid balance, acid-base disorders, and causes of abnormal chloride loss or retention.

\medterm{Chloride Channel} A protein forming a passage for chloride ions across a cell membrane. Chloride channels help control electrical signals, fluid movement, and muscle or nerve function.

\medterm{Chloride In Diet} Chloride consumed in food, usually as part of salt. It helps maintain fluid balance, nerve signals, and stomach acid; excessive intake often accompanies excess sodium.

\medterm{Chloride Ion} The negatively charged form of chlorine and a major extracellular electrolyte that helps maintain fluid balance, osmotic pressure, acid--base status, and electrical neutrality.

\medterm{Chloride Test - Blood} A blood test measuring serum chloride, an electrolyte that helps regulate extracellular fluid volume and acid-base balance; results are interpreted with sodium, bicarbonate, kidney function, and the clinical context.

\medterm{Chlorinated Lime} Chlorinated lime, or bleaching powder, releases chlorine and is used as a disinfectant; concentrated exposure can burn skin, eyes, and airways.
\medterm{Chlorinated Lime Poisoning} Exposure to chlorinated lime, a hypochlorite-containing corrosive that can injure the skin, eyes, airways, or gastrointestinal tract.

\medterm{Chlorine} A reactive halogen element and disinfectant; inhalation of chlorine gas can injure the eyes and airways and cause chemical pneumonitis or pulmonary edema.

\medterm{Chlorine Poisoning} Inhalation or contact exposure to chlorine gas, which irritates the eyes and airways and can cause cough, bronchospasm, chemical pneumonitis, or pulmonary edema.

\medterm{Chlormadinone Acetate} A synthetic progestogen used in some hormonal medicines; effects and risks include changes in bleeding, mood, metabolism, and thrombotic risk depending on the formulation.

\medterm{Chlormerodrin} A historical radioactive mercury compound once used in diagnostic or research settings; it is not a routine modern medicine and mercury exposure is toxic.

\medterm{Chlormethine} Chlormethine, or mechlorethamine, is an alkylating anticancer drug used systemically or topically for selected cancers and cutaneous T-cell lymphoma.
\medterm{Chlorodeoxyadenosine} Another name for cladribine, a purine-analog medicine used for selected blood cancers and multiple sclerosis; prolonged lymphocyte suppression and infection risk are important considerations.

\medterm{Chloroform} A volatile solvent and environmental contaminant that can irritate the lungs and nervous system and damage the liver or kidneys after significant exposure.

\medterm{Chlorogenic Acid} A natural plant compound found in coffee, fruits, and vegetables. It has antioxidant effects in laboratory studies, but evidence is insufficient to use supplements to prevent or treat disease.

\medterm{Chloroleukemia} A rare form of myeloid leukemia associated with greenish tumor masses composed of
myeloid cells, also called granulocytic sarcoma or chloroma.

\synonyms
granulocytic sarcoma or chloroma.

\medterm{Chloroma} A greenish extramedullary mass of immature myeloid cells, usually associated with acute
myeloid leukemia. It is also called granulocytic sarcoma or myeloid sarcoma.

\synonyms
granulocytic sarcoma or myeloid sarcoma.

\medterm{Chlorophyll} The green substance in plants and algae that absorbs light for photosynthesis, the process that makes food from light, water, and carbon dioxide. It is not a human nutrient or a proven treatment for disease.

\medterm{Chlorophyll Poisoning} Ingestion of chlorophyll-containing products, which usually causes mild gastrointestinal effects; product ingredients and dose determine clinical risk.

\medterm{Chlorophyllin} A water-soluble, copper-containing derivative of chlorophyll used in some products and studied for effects such as odor control; clinical uses depend on evidence and context.

\synonyms
Chlorophyllin copper complex; sodium copper chlorophyllin.

\medterm{Chlorophyta} The group of green algae; some species produce toxins or contribute to harmful algal blooms, while others are used as food or in research.

\medterm{Chloropicrin} A pungent soil fumigant and irritant that can injure the eyes, skin, and respiratory tract after exposure.

\medterm{Chloroplast} A plant and algal cell organelle that performs photosynthesis and contains its own DNA; it is not found in human cells.

\medterm{Chloroplast DNA} The genetic material inside chloroplasts, inherited mainly through the maternal line in many plants and used in studies of plant evolution and disease.

\medterm{Chloroprene} An industrial chemical used to make chloroprene rubber; repeated or high exposure can irritate tissues and has carcinogenic concern.

\medterm{Chloropsia} A visual disturbance in which objects appear green or greenish. It may result from retinal or visual-pathway disease, medicines, toxins, migraine, or other conditions and needs assessment if persistent or sudden.

\medterm{Chloroquine} An antimalarial medicine used for susceptible Plasmodium species and selected autoimmune conditions; retinal toxicity, cardiotoxicity, and drug interactions require attention.

\medterm{Chlorothiazide} A thiazide diuretic used to treat high blood pressure and fluid retention. It increases urine production and may cause low sodium, low potassium, dehydration, or high blood sugar.

\medterm{Chlorotoxin} A peptide from scorpion venom that binds some tumor and ion-channel targets; it is mainly studied experimentally for brain-tumor imaging or treatment.

\medterm{Chlorotrianisene} A synthetic estrogen formerly used for menopausal symptoms; it is rarely used now because of estrogen-related risks and safer alternatives.

\medterm{Chlorpheniramine} A first-generation antihistamine used for allergy symptoms; drowsiness, dry mouth, blurred vision, urinary retention, and other anticholinergic effects may occur.

\medterm{Chlorphentermine} An older stimulant appetite suppressant related to amphetamines; it is not routinely used because of cardiovascular, psychiatric, and dependence risks.

\medterm{Chlorpromazine} A first-generation antipsychotic and phenothiazine used for schizophrenia, severe nausea, and selected other conditions; adverse effects include sedation, hypotension, movement disorders, and QT prolongation.

\medterm{Chlorpromazine Overdose} Overdose of chlorpromazine that may cause sedation, hypotension, anticholinergic effects, seizures, cardiac conduction abnormalities, or neuroleptic malignant syndrome.

\medterm{Chlorpropamide} A first-generation sulfonylurea that lowers blood glucose by stimulating pancreatic insulin release; prolonged hypoglycemia and hyponatremia are important risks.

\medterm{Chlorpyrifos} An organophosphate insecticide that inhibits acetylcholinesterase; significant exposure can cause cholinergic toxicity with secretions, bronchospasm, bradycardia, weakness, or seizures.

\medterm{Chlortalidone} Chlortalidone, also called chlorthalidone, is a long-acting thiazide-like diuretic used mainly for hypertension and fluid retention.
\medterm{Chlorthalidone} A long-acting thiazide-like diuretic used for hypertension and edema; it can cause hyponatremia, hypokalemia, hyperuricemia, and kidney-related electrolyte changes.

\medterm{Chlorzoxazone} A centrally acting muscle relaxant used for painful muscle spasm in some countries; drowsiness and rare serious liver injury are important risks.

\medterm{CHO Cell} A Chinese hamster ovary cell line widely used to produce recombinant proteins, antibodies, vaccines, and other biologic medicines.

\medterm{Choanal Atresia} A birth defect in which tissue blocks one or both passages from the nose to the throat. Both sides can cause severe breathing difficulty in a newborn, especially during feeding, and require urgent treatment.

\medterm{Chocolate Cyst} A chocolate cyst is an ovarian endometrioma filled with old blood from endometriosis, which can cause pelvic pain, infertility, or an adnexal mass.
\medterm{Choice Behavior} The process of selecting among alternatives, studied in psychology and neuroscience to understand decision-making, impulsivity, and reward processing. Abnormal choice behavior may be relevant to addiction, eating disorders, and psychiatric conditions.

\medterm{Choke} Airway blockage by food or another object that prevents normal airflow and can rapidly cause inability to speak, cyanosis, loss of consciousness, or death.

\medterm{Choking} Blockage of the airway by food or another object, making breathing difficult or impossible. A person who cannot speak, cough effectively, or breathe needs immediate first aid and emergency medical help.

\medterm{Choking - Adult Or Child Over 1 Year} First aid for severe airway obstruction in a person older than one year, using back blows and abdominal thrusts when appropriate and CPR if the person becomes unresponsive.

\medterm{Choking - Infant Under 1 Year} First aid for severe airway obstruction in an infant, using alternating back slaps and chest thrusts and avoiding blind finger sweeps; unresponsiveness requires infant CPR.

\medterm{Choking - Unconscious Adult Or Child Over 1 Year} Resuscitation for an unresponsive person with suspected airway obstruction, combining emergency activation, airway inspection, CPR, and removal of a visible object without blind sweeping.

\medterm{Cholagogue} A substance that encourages the gallbladder to release bile into the small intestine. Bile helps the body digest fats; medicines or foods described as cholagogues should be used carefully when gallstones or blocked bile ducts are possible.

\synonyms
Gallbladder stimulant; biliary evacuant.

\medterm{Cholangiitis} Inflammation of the tubes that carry bile from the liver and gallbladder to the intestine. Infection, a blockage, autoimmune disease, or another bile-duct disorder may cause it.

\medterm{Cholangiocarcinoma} Cholangiocarcinoma is cancer arising from bile-duct lining cells and may cause jaundice, itching, abdominal symptoms, or bile-flow obstruction.

\synonyms
Bile Duct Cancer (Cholangiocarcinoma)

\medterm{Cholangiocarcinoma (Bile Duct Cancer (Cholangiocarcinoma))} Alternate index label for Bile Duct Cancer (Cholangiocarcinoma).

\medterm{Cholangiocarcinoma (Bile Duct Cancer)} Cancer arising in the bile ducts, which may cause jaundice, itching, abdominal discomfort, weight loss, or impaired bile flow.

\medterm{Cholangiography} Imaging of the bile ducts after contrast material makes them visible. It can show gallstones, narrowing, leaks, or blockage and may be performed during surgery, through the skin, or during endoscopy.
\medterm{Cholangiohepatitis} Inflammation involving both the bile ducts and liver, often associated with ascending bacterial infection, biliary obstruction, parasites, or recurrent cholangitis. It may cause fever, right-upper-quadrant pain, jaundice, and cholestatic liver-test abnormalities.

\medterm{Cholangitis} Cholangitis is inflammation or infection of the bile ducts, often from obstruction by a stone or stricture, and can rapidly cause sepsis.

\medterm{Cholangitis, Primary Sclerosing} Alternate index label; see \textit{Primary sclerosing cholangitis}.

\medterm{Cholate} The ionized form of cholic acid, a primary bile acid that helps emulsify dietary fat and is conjugated and secreted by the liver.

\medterm{Cholecalciferol} Vitamin D3, produced in skin by ultraviolet light and obtained from food or supplements, then converted to active metabolites that regulate calcium and phosphate.

\medterm{Cholecyst} The gallbladder, a small sac beneath the liver that stores and concentrates bile. When food enters the intestine, it contracts and releases bile to help digest dietary fats.

\medterm{Cholecystectomy} Surgical removal of the gallbladder, usually by a laparoscopic approach when appropriate. It may be performed for symptomatic gallstones, acute cholecystitis, selected gallstone pancreatitis, or other gallbladder disease.

\medterm{Cholecystitis} Inflammation of the gallbladder, usually because a gallstone blocks the tube through which bile leaves it. It commonly causes steady upper-right abdominal pain, fever, nausea, or vomiting and may require urgent treatment.

\synonyms
Inflamed Gallbladder

\medterm{Cholecystography} Radiographic examination of the gallbladder, usually after administration of a contrast
substance, to assess its filling, concentration, and emptying.

\medterm{Cholecystokinin} A hormone released by the small intestine after food, especially fat and protein, enters it. It makes the gallbladder release bile, stimulates digestive enzymes, and helps create a feeling of fullness.

\medterm{Cholecystoses} Non-inflammatory disorders causing deposits, wall thickening, or other structural changes in the gallbladder, including cholesterolosis and adenomyomatosis. Many are incidental, but symptoms and imaging findings determine whether further evaluation is needed.

\medterm{Choledochal Cyst} An abnormal widening of one or more bile ducts, usually present from birth. It may cause abdominal pain, jaundice, pancreatitis, or infection and can increase later cancer risk, so specialist evaluation is needed.

\medterm{Choledochal Cysts} Choledochal cysts are congenital or acquired dilations of the bile ducts that can cause abdominal pain, jaundice, pancreatitis, cholangitis, or stones. Imaging establishes the anatomy, and surgical removal is often recommended because of recurrent complications and cancer risk.

\synonyms
Cysts Choledochal (Choledochal Cysts)

\medterm{Choledocholithiasis} A gallstone within the common bile duct, which may cause biliary pain, jaundice, cholangitis, or pancreatitis and is usually evaluated with liver tests and imaging.

\medterm{Cholelithiasis} The presence of stones in the gallbladder or biliary tract. Stones may contain cholesterol, pigment, or mixed material and may cause no symptoms or lead to biliary pain and complications.

\medterm{Cholelithiasis (Gallstones)} Alternate index label for Gallstones.

\medterm{Cholera} An acute intestinal infection caused by toxigenic Vibrio cholerae, producing profuse watery diarrhoea and potentially life-threatening dehydration; prompt oral or intravenous rehydration is central treatment.


\medterm{Choleresis} Choleresis is production and secretion of bile by liver cells.
\medterm{Choleretic} A substance that encourages the liver to make or release bile. Bile helps digest fats and remove some waste products; a choleretic should be used carefully when gallstones or blocked bile ducts may be present.

\synonyms
Bile-secretory agent; bile-flow stimulant.

\medterm{Cholestasis} Reduced production or flow of bile. It may cause itching, yellow skin or eyes, dark urine, pale stools, and abnormal liver tests; a blockage and liver-cell causes must be distinguished.

\medterm{Cholestasis Of Pregnancy} A liver disorder of pregnancy in which impaired bile flow causes itching and elevated bile acids.

\synonyms
Obstetric Cholestasis

\medterm{Cholesteatoma} A keratin-filled sac of abnormal squamous epithelium in the middle ear or mastoid that can erode bone and cause chronic discharge, hearing loss, or intracranial complications.

\medterm{Cholesterol} A lipid required for cell membranes and synthesis of steroid hormones and bile acids. Excess atherogenic cholesterol can contribute to arterial plaque.

\synonyms
Serum cholesterol; blood cholesterol.

\medterm{Blood Cholesterol} Cholesterol carried in the blood. The body needs it to build cell membranes and make hormones and bile acids. Too much LDL cholesterol can form fatty buildup in the arteries and increase the risk of heart disease and stroke. Blood cholesterol is usually checked as part of a lipid profile.

\medterm{Cholesterol And Lifestyle} Diet, physical activity, weight, smoking, and alcohol use can affect blood cholesterol and cardiovascular risk. Lifestyle changes complement medicines when cholesterol remains high.

\medterm{Cholesterol And Lipids} Fatty substances in the blood, including LDL cholesterol, HDL cholesterol, and
triglycerides, that play critical roles in cellular function and energy metabolism.
Abnormal lipid levels, particularly high LDL and triglycerides with low HDL, promote
atherosclerosis and increase the risk of heart attack and stroke.

\medterm{Cholesterol Crystal Embolization} Small pieces of cholesterol plaque breaking loose from a large artery and blocking smaller arteries. It may damage the kidneys, skin, eyes, brain, or intestines, often after a vascular procedure.

\medterm{Cholesterol Crystals} Needle-shaped, cleft-like spaces seen in histological sections representing dissolved
cholesterol esters, commonly found in atherosclerotic plaques, atheroembolic disease,
and cholesterol-rich effusions. In atheroembolic disease, cholesterol crystals embolize
from ruptured plaques to distal organs causing tissue ischemia.

\medterm{Cholesterol Emboli} Cholesterol crystals released from an atherosclerotic plaque that lodge in small arteries and cause
ischemia or inflammation in downstream tissues.

\medterm{Cholesterol Embolism} Blockage of small arteries by cholesterol crystals released from an artery wall, often after vascular procedures or disruption of atherosclerotic plaque.

\medterm{Cholesterol Embolization} A condition in which tiny pieces of fatty material from an artery wall break loose and travel into smaller blood vessels. It can damage the kidneys, skin, eyes, brain, or intestines, sometimes after a vascular procedure or blood-thinning treatment.

\medterm{Cholesterol Lowering (Cholesterol Management)} Alternate index label for Cholesterol Management.

\medterm{Cholesterol Management} Cholesterol management combines cardiovascular-risk assessment with dietary pattern, physical activity, tobacco avoidance, weight care, and medicines when indicated. Statins and other agents lower atherogenic lipoproteins, with intensity chosen according to prior cardiovascular disease and overall risk.

\synonyms
Cholesterol Lowering (Cholesterol Management)

\medterm{Cholesterol Synthesis} The process by which cells, especially liver and intestinal cells, make cholesterol from smaller molecules. Cholesterol is needed for cell membranes, hormones, and bile acids, but excess production contributes to high blood cholesterol.

\medterm{Cholesterol Testing And Results} Laboratory assessment of blood lipids, usually including total cholesterol, LDL cholesterol, HDL cholesterol, and triglycerides, used to estimate cardiovascular risk and guide prevention or treatment.

\medterm{Cholesterolemia} The presence and concentration of cholesterol in the blood. Abnormal elevation is a risk
factor for atherosclerotic cardiovascular disease.

\medterm{Cholesterolosis} Deposition of cholesterol esters in the gallbladder mucosa, producing yellow speckling or a "strawberry gallbladder." It is often incidental but may coexist with gallstones or chronic gallbladder inflammation.

\medterm{Cholestyramine} Cholestyramine is a bile-acid sequestrant used to lower LDL cholesterol and sometimes to treat bile-acid diarrhea or itching from cholestasis.
\medterm{Choline} Choline is an essential nutrient used to make cell membranes and acetylcholine and to support liver and nervous-system function.
\medterm{Cholinergic} Cholinergic means related to acetylcholine or receptors and pathways activated by it, including muscarinic and nicotinic signaling.
\medterm{Cholinergic Crisis} A life-threatening excess of acetylcholine, often after organophosphate or carbamate poisoning. It can cause sweating, salivation, vomiting, diarrhea, pinpoint pupils, wheezing, muscle weakness, seizures, and respiratory failure.

\medterm{Cholinergic Neuron} A nerve cell that releases acetylcholine as its principal neurotransmitter.

\synonyms
Acetylcholine-releasing neuron; cholinergic nerve cell.

\medterm{Cholinergic Toxidrome} A poisoning pattern caused by excessive acetylcholine activity, often producing salivation, tearing, urination, diarrhea, vomiting, bronchial secretions, slow heart rate, muscle fasciculations, weakness, or respiratory failure.

\medterm{Cholinesterase} Cholinesterase enzymes break down acetylcholine; inhibition increases cholinergic activity and is used therapeutically but can also cause poisoning.
\medterm{Cholinesterase - Blood} A blood test measuring acetylcholinesterase or butyrylcholinesterase activity, useful in selected pesticide exposures, inherited enzyme variants, and some liver-function assessments.

\medterm{Chondritis} Inflammation of cartilage, such as the cartilage of the ear, nose, ribs, or larynx. It
may result from infection, trauma, or immune-mediated disease.

\medterm{Chondroblastic Osteosarcoma} A variant of osteosarcoma in which the malignant cells produce a cartilaginous matrix.
It typically affects the long bones of adolescents and young adults, with a predilection
for the metaphyseal regions.

\medterm{Chondroblastoma} A chondroblastoma is an uncommon usually benign cartilage-forming bone tumor, often arising near the end of a long bone in adolescents or young adults.
\medterm{Chondrocalcinosis} Calcification of cartilage visible on X-ray. Associated with CPPD disease. Found in
knees, wrists, pubic symphysis. May be incidental or symptomatic.

\medterm{Chondrocyte} A mature cell that makes and maintains cartilage, the firm flexible tissue found in joints, airways, and other places. Chondrocytes help keep cartilage smooth and strong, but cartilage usually heals slowly because it has little direct blood supply.

\synonyms
Cartilage cell.

\medterm{Chondrodysplasia} A disorder of cartilage and bone development that causes abnormal skeletal growth. Depending on the type, it may produce short limbs, short stature, enlarged joints, or other bone abnormalities.

\medterm{Chondroitin Sulfate} A sulfated glycosaminoglycan found in cartilage and other connective tissues; it is also marketed in supplements for osteoarthritis, with variable evidence of benefit.

\medterm{Chondroma} Benign cartilage-forming tumor composed of mature hyaline cartilage, most often
intramedullary (enchondroma) or subperiosteal (juxtacortical / periosteal chondroma).

\medterm{Chondromalacia} Softening and degeneration of cartilage, especially articular cartilage. The term is
often used for changes affecting the patella and causing anterior knee pain.

\medterm{Chondromyxoid Fibroma} A rare benign cartilage-forming bone tumor, usually arising in the metaphysis of a long bone and presenting with localized pain or an incidental radiographic lesion.

\medterm{Chondrosarcoma} A cancer that forms cartilage-producing tissue, usually in a bone or nearby soft tissue. It often causes persistent pain or a growing lump, and treatment depends on its grade, site, and spread.

\medterm{Chondrosarcoma Of The Breast} An exceptionally rare malignant mesenchymal breast tumor producing cartilage. It is distinguished from metaplastic carcinoma and other sarcomas by histology and is managed with sarcoma-oriented complete excision.






\medterm{Chopart'S Joint} The transverse tarsal joint, consisting of the talonavicular and calcaneocuboid joints,
between the hindfoot and the midfoot.

\medterm{Chorangioma} A benign vascular tumor of the placenta, usually small and incidental but occasionally associated with fetal anemia, hydrops, growth restriction, or other complications when large.

\medterm{Chordae Tendineae} Strong, thin cords inside the heart that connect the mitral and tricuspid valves to small muscles. They help keep the valves from turning inside out when the heart contracts.

\medterm{Chordee} Chordee is abnormal downward or other curvature of the penis, commonly associated with hypospadias or scar tissue and sometimes requiring surgical correction.
\medterm{Chordoma} A rare cancer arising from leftover cells of early spinal development. It usually begins near the tailbone, base of the skull, or spine and can cause pain, weakness, or nerve problems as it grows.

\medterm{Chorea} Uncontrolled, brief, irregular movements that flow unpredictably from one body part to another. It may result from inherited disease, autoimmune illness, medicines, infection, or damage to movement-control areas of the brain.

\medterm{Chorea, Huntington'S} Alternate index label; see \textit{Huntington's disease}.

\medterm{Choreoathetosis} Choreoathetosis combines irregular, dance-like chorea with slow twisting athetotic movements and may result from basal-ganglia injury or genetic disease.
\medterm{Chorioamnionitis} Infection or inflammation of the fluid, membranes, or tissues around a fetus, usually caused by bacteria moving upward from the vagina. Maternal fever, uterine tenderness, foul fluid, or a fast fetal heartbeat may occur.

\medterm{Chorioamnionitis (Intraamniotic Infection)} Infection or inflammation of the amniotic fluid, membranes, placenta, or decidua,
usually due to ascending polymicrobial infection. Maternal fever, uterine tenderness,
fetal tachycardia, and maternal leukocytosis may occur; treatment requires
broad-spectrum antibiotics and delivery planning.

\synonyms
Intraamniotic Infection

\medterm{Choriocarcinoma} Choriocarcinoma is an aggressive trophoblastic cancer producing human chorionic gonadotropin, usually arising after pregnancy or a molar pregnancy.

\medterm{Chorion} The outer membrane surrounding an early embryo. Its finger-like projections help form the fetal part of the placenta, where oxygen, nutrients, and waste are exchanged between fetal and maternal blood supplies.

\medterm{Chorion Villus Sampling} Chorion villus sampling removes placental tissue for genetic or chromosomal testing, usually early in pregnancy; it carries a small procedure-related miscarriage risk.
\medterm{Chorionic Gonadotrophin} Chorionic gonadotrophin, especially hCG, is a placental hormone that maintains early pregnancy and is measured in pregnancy tests and selected tumor evaluations.
\medterm{Chorionic Villus Sampling} A prenatal test usually performed around 10 to 13 weeks by removing a small placental sample through the abdomen or cervix. It can identify many chromosome and inherited disorders but carries a small miscarriage risk.

\medterm{Chorioretinitis} Inflammation involving the choroid and retina, often caused by infection, autoimmune
disease, or systemic inflammatory illness. It can produce floaters and impaired vision.

\medterm{Chorioretinitis (Uveitis)} Alternate index label for Uveitis.

\medterm{Chorioretinitis Toxoplasma (Toxoplasmosis)} Alternate index label for Toxoplasmosis.

\medterm{Choroid} The vascular layer between retina and sclera that supplies the outer retina and absorbs stray light.

\medterm{Choroid Plexus Carcinoma} A cancer arising from the choroid plexus, tissue that produces cerebrospinal fluid, potentially causing pressure buildup, headaches, vomiting, or neurologic problems.

\medterm{Choroid Plexus Cyst} A fluid-filled cyst within the choroid plexus of a fetal cerebral ventricle. An isolated cyst
with otherwise reassuring prenatal assessment is usually benign, but its significance
depends on other ultrasound findings and screening results.

\medterm{Choroid Plexus Tumor} A papillary neoplasm of the choroid plexus, occurring most often in children, and
frequently presenting with hydrocephalus from excessive cerebrospinal fluid production.

\medterm{Choroidal Dystrophies} Inherited disorders causing progressive degeneration of the choroid, retinal pigment epithelium, and retina, leading to night blindness, visual-field loss, and eventual vision impairment.

\medterm{Choroidal Neovascularization} Growth of fragile abnormal blood vessels beneath or through the retina. Leakage or bleeding can damage central vision, especially in wet age-related macular degeneration and some inflammatory or inherited eye disorders.

\medterm{Choroideremia} Choroideremia is an X-linked inherited retinal degeneration causing progressive night blindness and loss of peripheral then central vision, mainly in males.
\medterm{Choroiditis} Choroiditis is inflammation of the choroid, the vascular layer behind the retina, often causing blurred vision, floaters, or visual-field disturbance.
\medterm{Christmas Disease} An inherited bleeding disorder, also called hemophilia B, caused by low or abnormal clotting factor IX. It can cause prolonged bleeding or bleeding into joints and muscles and is treated with replacement or other specialist therapy.
\medterm{Chromaffin} Chromaffin describes cells that stain with chromium salts because they contain catecholamine-rich granules, especially adrenal medullary cells and related neuroendocrine cells.
\medterm{Chromaffinoma} A tumor of cells that can produce adrenaline-like hormones, usually in the adrenal gland or nearby nerve tissue. Hormone release may cause episodes of high blood pressure, headache, sweating, and palpitations.

\medterm{Chromate} Chromate is a hexavalent chromium compound that can cause irritant or allergic dermatitis, nasal injury, respiratory disease, and cancer after occupational exposure.
\medterm{Chromatic} Chromatic means relating to color or wavelengths of light; chromatic vision depends on cone photoreceptors and neural processing.
\medterm{Chromatids} Chromatids are the two copies of a replicated chromosome joined at the centromere before separating during cell division.
\medterm{Chromatin} The combination of DNA and proteins that packages genetic material inside the cell nucleus. Its arrangement controls which genes are available for use and helps regulate DNA copying and repair.

\medterm{Chromatography} A laboratory separation method in which substances partition between a stationary phase and a moving phase, allowing identification or measurement of components in a mixture.

\medterm{Chromatopsia} A visual disturbance in which objects appear abnormally colored. It may result from
retinal disease, toxic exposure, medication effects, or other disorders of the visual
system.

\medterm{Chromhidrosis} A rare condition in which sweat is abnormally colored, often yellow, green, blue, or
black, because of pigment in the sweat glands or on the skin surface.

\medterm{Chromium} Chromium is a metallic element; trivalent chromium is a trace nutrient, while hexavalent compounds can be toxic, corrosive, mutagenic, and carcinogenic.
\medterm{Chromium - Blood Test} A blood test measuring chromium exposure or body burden, generally used when occupational or medical exposure to chromium is suspected; results must be interpreted with the exposure history and clinical findings.

\medterm{Chromium In Diet} Reference information on chromium as a trace element involved in carbohydrate and fat metabolism. Chromium deficiency is rare; most people obtain adequate chromium from normal dietary intake, and routine chromium supplementation is generally not recommended.

\medterm{Chromoblastomycosis} A chronic fungal infection of skin and subcutaneous tissue that produces slowly enlarging warty or
scaly lesions, usually after traumatic implantation of the fungus.

\medterm{Chromopertubation} A fertility test performed during laparoscopy in which colored fluid is passed through the uterus while the fallopian tubes are viewed. Fluid spilling from the tube ends suggests that they are open; no spill may indicate blockage or surrounding scar tissue.

\medterm{Chromosomal Abnormality} Change in chromosome number or structure. Numerical: aneuploidy (abnormal number, e.g.,
trisomy 21), polyploidy (extra sets). Structural: deletion, duplication, inversion,
translocation.

\medterm{Chromosomal Microarray} A genetic test that detects gains and losses of chromosome material, including copy-number variants and some regions of homozygosity; it is commonly used for developmental delay, intellectual disability, congenital anomalies, or autism-spectrum presentations.

\medterm{Chromosome} A chromosome is a DNA-protein structure carrying genetic information; humans normally have 46 chromosomes in most body cells.

\medterm{Chromosome Complement} The complete set of chromosomes in a cell, including their number and sex-chromosome composition, usually described by a karyotype.

\medterm{Chromosome Duplication} A structural chromosome abnormality in which a segment of genetic material is copied one or more additional times, potentially altering gene dosage and development.

\medterm{Chronic} Persisting or recurring over a long period, in contrast to acute; the appropriate duration depends on the disease or clinical context.

\medterm{Chronic Adrenal Insufficiency} Alternate index label; see \textit{Addison's disease}.

\medterm{Chronic Allograft Nephropathy} Gradual loss of function in a transplanted kidney over months or years, caused by immune injury, high blood pressure, recurrent disease, medicines, or other damage. Testing and biopsy help identify the cause.

\medterm{Chronic Atrophic Gastritis} Long-term inflammation that destroys glands in the stomach lining. It may reduce stomach acid and intrinsic factor, impair iron or vitamin B12 absorption, and increase the risk of anemia or stomach tumors.

\medterm{Chronic Autoimmune Hashimoto Thyroiditis} A long-term immune-system attack on the thyroid gland that gradually damages thyroid tissue and often causes too little thyroid hormone. It may produce tiredness, cold intolerance, weight change, or an enlarged thyroid.

\medterm{Chronic Autoimmune Thyroiditis} Long-term inflammation caused by the immune system attacking the thyroid, usually called Hashimoto thyroiditis. It commonly leads to low thyroid hormone and may be associated with thyroid antibodies.

\medterm{Chronic Bacterial Prostatitis} A recurring bacterial infection of the prostate that may cause repeated urinary infections, burning or urgent urination, pelvic discomfort, or painful ejaculation. Diagnosis uses appropriate urine tests, and treatment often lasts several weeks.

\synonyms
NIH Category II Prostatitis; Recurrent Bacterial Prostatitis

\medterm{Chronic Bronchitis} A cough producing mucus on most days for at least three months in each of two consecutive years, after other causes are excluded. It reflects long-term airway irritation and commonly occurs with chronic obstructive pulmonary disease.

\synonyms
Bronchitis Chronic (Chronic Bronchitis)

\medterm{Chronic Caries} Tooth decay that develops slowly over time. It may cause a firm dark area, a cavity, sensitivity, or no symptoms; without dental treatment, decay can reach the inner tooth, causing pain, infection, or tooth loss.

\medterm{Chronic Cholecystitis} Repeated or persistent inflammation of the gallbladder, usually related to gallstones and impaired emptying. It may cause episodic right-upper-abdominal pain, nausea, or food intolerance and can lead to scarring or a contracted gallbladder.

\medterm{Chronic Compartment Syndrome} Alternate index label; see \textit{Chronic exertional compartment syndrome}.

\medterm{Chronic Cough} A cough lasting longer than eight weeks. Common causes include asthma, nasal drainage, acid reflux, smoking, and some blood-pressure medicines; persistent cough requires assessment for its cause.

\synonyms
Cough Chronic (Chronic Cough)

\medterm{Chronic Daily Headaches} Headaches occurring on at least fifteen days each month for more than three months, caused by migraine, tension, medication overuse, or other conditions.

\medterm{Chronic Disease} A disease or condition that persists for a prolonged period, often months or years, and may require continuing monitoring or management.

\medterm{Chronic Endometritis} Persistent low-grade inflammation of the endometrium, often associated with infection, retained products, intrauterine devices, or reproductive failure; plasma cells support the diagnosis histologically.
delivery. Risk factors include prolonged labor, multiple vaginal examinations, premature
rupture of membranes, and cesarean section. Presents with fever, uterine tenderness,
foul-smelling lochia, and leukocytosis.

\medterm{Chronic Eosinophilic Leukaemia} A clonal myeloid neoplasm characterized by persistent eosinophilia, increased or
dysplastic eosinophils, and evidence of an underlying clonal process or organ injury.
Symptoms may include fatigue, fever, weight loss, rash, respiratory symptoms,
splenomegaly, or cardiac and neurological complications.

\medterm{Chronic Eosinophilic Pneumonia} An idiopathic inflammatory lung disease characterized by eosinophilic infiltration of
the pulmonary parenchyma, presenting with fever, cough, dyspnea, and peripheral airspace
opacities on imaging.

\medterm{Chronic Exertional Compartment Syndrome} Exercise-related increased pressure within a muscle compartment that causes predictable pain, tightness, or neurologic symptoms during activity.

\synonyms
Chronic Compartment Syndrome

\medterm{Chronic Fatigue Syndrome} A long-term illness involving substantial reduced ability to function, worsening after physical or mental activity, unrefreshing sleep, and often difficulty thinking or standing. Symptoms are not explained by another condition and need careful evaluation.

\synonyms
CFS (Chronic Fatigue Syndrome)
SEID (Chronic Fatigue Syndrome)

\medterm{Chronic Gastritis} Chronic, progressive inflammation of the gastric mucosa leading over time to mucosal atrophy, loss of specialized parietal and chief cells, and intestinal metaplasia. Major causes include persistent \textit{Helicobacter pylori} infection and autoimmune gastritis targeting parietal cells and intrinsic factor.

\medterm{Chronic Glomerulonephritis} Long-standing inflammation and scarring of the kidney glomeruli that can progressively reduce filtration and lead to chronic kidney disease or kidney failure.

\medterm{Chronic Granulomatous Disease} An inherited immune disorder in which certain white blood cells cannot effectively kill some bacteria and fungi. It causes repeated serious infections, swollen lymph nodes, pneumonia, skin infections, and inflammation called granulomas.

\medterm{Chronic Hives} Recurrent itchy, raised welts occurring repeatedly for more than six weeks. Individual welts usually fade within a day, while causes may include immune activity, medicines, infections, or no identifiable trigger.

\medterm{Chronic Hypertension In Pregnancy} Hypertension that predates pregnancy or is diagnosed early in pregnancy.
It increases the risk of preeclampsia, fetal growth problems, preterm birth, and placental complications and requires
ongoing maternal and fetal monitoring.

\medterm{Chronic Inflammation} Inflammation that continues for weeks, months, or years. It can occur with persistent infection, autoimmune disease, long-term irritation, or obesity and may cause ongoing tissue damage while the body attempts repair.

\medterm{Chronic Inflammatory Demyelinating Polyneuropathy} An immune disorder that damages the insulating covering of nerves outside the brain and spinal cord. It causes gradually worsening or recurring weakness, numbness, poor balance, and reduced reflexes in the limbs.

\medterm{Chronic Inflammatory Response Syndrome} A proposed multisystem illness involving persistent inflammatory symptoms after an exposure
such as certain water-damaged building environments. Reported symptoms are nonspecific and may include fatigue,
headache, cognitive difficulty, and respiratory complaints. Symptoms require a careful medical evaluation for
established infectious, allergic, toxic, autoimmune, and other causes.

\medterm{Chronic Intestinal Pseudo-Obstruction} Repeated symptoms of bowel blockage without a physical obstruction, caused by poor nerve or muscle function in the intestine. It may cause abdominal swelling, pain, vomiting, constipation, and inability to absorb nutrition.

\medterm{Chronic Kidney Disease} Long-term damage or reduced function of the kidneys lasting at least three months. Common causes include diabetes and high blood pressure; advanced disease can cause anemia, bone problems, fluid buildup, and kidney failure.

\synonyms
Chronic Kidney Failure
Chronic Renal Failure
Kidney Disease, Chronic
Kidney Failure, Chronic

\medterm{Chronic Kidney Disease Staging} Classification of kidney disease by cause, eGFR category, and albuminuria category. G1
through G5 represent eGFR from normal or high to kidney failure; staging is applied when
abnormalities persist for at least 3 months and helps estimate progression and
complications.

\medterm{Chronic Kidney Failure} Alternate index label; see \textit{Chronic kidney disease}.

\medterm{Chronic Lead Poisoning} Illness caused by long-term or repeated exposure to lead. It may cause abdominal pain, constipation, tiredness, anemia, nerve damage, high blood pressure, or kidney injury. In children, it can impair learning, attention, growth, or development. Treatment begins with stopping the exposure; some people need specialist-directed chelation.

\medterm{Chronic Leukemia} A blood cancer that develops gradually as abnormal white blood cells multiply and remain in the blood or bone marrow. Some types progress slowly but still require regular medical monitoring.

\medterm{Chronic Limb-Threatening Ischemia} Severe, long-lasting reduction of blood flow to a limb, usually from narrowed or blocked arteries. It causes pain at rest, ulcers that do not heal, or gangrene and requires urgent vascular assessment to protect the limb.

\medterm{Chronic Liver Disease} Persistent hepatic injury and structural change lasting at least 6 months, caused by
viral, metabolic, alcohol-related, autoimmune, cholestatic, vascular, or drug-related
disease. Progressive fibrosis can lead to cirrhosis, portal hypertension, liver failure,
and hepatocellular carcinoma.

\medterm{Chronic Lymphocytic Leukemia} The most common leukemia in Western countries, characterized by the clonal proliferation
of mature-appearing but functionally incompetent B lymphocytes that accumulate in blood,
bone marrow, and lymphoid tissues.

\synonyms
Cancer, Chronic Lymphocytic Leukemia
CLL
Leukemia, Chronic Lymphocytic

\medterm{Chronic Lymphocytic Thyroiditis} Alternate index label; see \textit{Hashimoto's disease}.

\medterm{Chronic Mesenteric Ischemia} Reduced blood flow to the intestines, usually from hardening and narrowing of abdominal arteries. It often causes cramping after meals, food avoidance, and weight loss and needs vascular evaluation.

\medterm{Chronic Migraine} Headache on at least 15 days each month for more than three months, with migraine features on at least eight days. Pain may be one-sided or throbbing and may include nausea or sensitivity to light or sound.

\medterm{Chronic Motor Or Vocal Tic Disorder} A neurodevelopmental disorder with persistent motor or vocal tics lasting more than one year; symptoms often fluctuate and may be treated with education, behavioral therapy, or medicines when impairment is significant.

\medterm{Chronic Mucocutaneous Candidiasis} Repeated or persistent yeast infections of the skin, nails, or moist linings of the mouth, throat, or vagina. It may result from an inherited immune problem and can occur with autoimmune or hormone disorders.

\medterm{Chronic Myelogenous Leukemia} A blood cancer in which the bone marrow makes too many abnormal white blood cells, usually because of the BCR::ABL1 gene change. It may progress over time but is often controlled with targeted medicines.

\medterm{Chronic Myeloid Leukemia} A blood cancer usually caused by the Philadelphia chromosome, which creates the BCR::ABL1 gene. The bone marrow produces too many abnormal white blood cells, and targeted treatment can often control the disease.

\medterm{Chronic Myelomonocytic Leukemia} A rare blood cancer in which the bone marrow makes too many monocytes, a type of white blood cell, along with abnormal blood-forming cells. It can cause anemia, infections, bleeding, fatigue, or an enlarged spleen.

\medterm{Chronic Neutrophilic Leukaemia} A rare clonal myeloproliferative neoplasm characterized by sustained mature neutrophilic
leukocytosis without a primary infectious or inflammatory cause. It may cause fatigue,
fever, weight loss, splenomegaly, and progressive constitutional symptoms.

\medterm{Chronic Obstructive Lung Disease} An older abbreviation for chronic obstructive lung disease, now usually called chronic obstructive pulmonary disease, characterized by persistent airflow limitation.
\medterm{Chronic Obstructive Pulmonary Disease} A long-term lung disease in which damaged airways or air sacs limit airflow. It causes cough, mucus, wheezing, and breathlessness, and is most often related to smoking or prolonged exposure to irritating dusts or gases.

\synonyms
COPD (Chronic Obstructive Pulmonary Disease)

\medterm{Chronic Pain} Pain that persists or recurs for longer than three months. It may have an identifiable ongoing cause or persist after healing, and is influenced by
complex interactions among biological, psychological, and social factors.


\medterm{Chronic Pain Management} Care aimed at reducing long-lasting pain and improving sleep, movement, and daily function. It may combine exercise, physical therapy, psychological support, nonopioid medicines, procedures, and carefully monitored opioids when appropriate.

\medterm{Chronic Pain Syndrome} Long-lasting pain accompanied by substantial effects on movement, sleep, mood, relationships, or daily activities. Treatment addresses pain and these related effects rather than relying only on pain medicines.
\medterm{Chronic Pancreatitis} Repeated or long-term inflammation that permanently scars the pancreas. It may cause upper-abdominal pain, poor digestion, weight loss, diabetes, or oily stools; alcohol, smoking, inherited disorders, and blocked ducts are possible causes.

\medterm{Chronic Paroxysmal Hemicrania} A primary headache disorder causing frequent, short-lasting, severe one-sided headaches with autonomic symptoms and a characteristic response to indomethacin.

\synonyms
CPH; chronic paroxysmal hemicrania syndrome.

\medterm{Chronic Pelvic Pain} Pain in the lower abdomen or pelvis that persists or returns for at least six months. Causes may involve reproductive, urinary, digestive, muscle, nerve, or pain-processing disorders and require individualized evaluation.
\synonyms
Pelvic Pain, Chronic

\medterm{Chronic Pelvic Pain Syndrome} Persistent or recurring pain in the pelvis, area between the genitals and anus, lower abdomen, back, or scrotum without a proven ongoing bacterial infection. Urinary symptoms, painful ejaculation, muscle tension, and sexual problems may also occur.

\synonyms
CPPS; NIH Category III Prostatitis; Chronic Nonbacterial Prostatitis

\medterm{Chronic Progressive External Ophthalmoplegia} A slowly worsening muscle disorder causing drooping eyelids and limited movement of both eyes. It often begins in adulthood and may occur alone or with weakness of other muscles; inherited mitochondrial disease is a common cause.

\synonyms
CPEO; Progressive External Ophthalmoplegia (PEO)

\medterm{Chronic Recurrent Multifocal Osteomyelitis} A rare autoinflammatory bone disorder, often beginning in children or adolescents, causing recurrent sterile inflammation and pain at multiple bone sites; it is also called CRMO.

\medterm{Chronic Pyelonephritis} Long-term inflammation and scarring of the kidneys, often after repeated infections or backward flow of urine. It can gradually reduce kidney function and may cause high blood pressure, poor growth, or kidney failure.

\medterm{Chronic Renal Failure} Alternate index label; see \textit{Chronic kidney disease}.

\medterm{Chronic Rhinitis} Persistent inflammation or irritation of the nasal lining causing congestion, runny nose, sneezing, itching, or
postnasal drainage. Allergic, nonallergic, medication-related, and structural causes are possible. Management depends
on the cause and may include trigger avoidance, saline irrigation, and clinician-directed nasal treatment.

\synonyms
Post-Nasal Drip (Chronic Rhinitis)

\medterm{Chronic Rhinosinusitis} Persistent inflammation of the nose and paranasal sinuses lasting at least 12 weeks, usually
with nasal obstruction or discharge and sometimes facial pressure or reduced smell. It
may occur with or without nasal polyps.

\medterm{Chronic Sinusitis} Persistent inflammation of the paranasal sinuses for at least 12 weeks, causing nasal
obstruction, discharge, facial pressure, and reduced smell.

\synonyms
Sinusitis, Chronic

\medterm{Chronic Subdural Hematoma} A collection of blood and breakdown products between the dura and arachnoid that develops over weeks, often after minor trauma in older adults, causing headache, confusion, weakness, or gait change.

\medterm{Chronic Suppurative Otitis Media} Persistent inflammation and infection of the middle ear with recurrent or continuous ear discharge,
usually through a perforated tympanic membrane. It can cause conductive hearing loss and
may lead to cholesteatoma or other complications.

\synonyms
CSOM; chronic otitis media with discharge; chronic suppurative middle-ear infection.

\medterm{Chronic Tension-Type Headache} A recurring headache occurring on at least 15 days per month for more than three months, usually causing
bilateral pressing or tightening pain rather than throbbing. It may be associated with scalp or neck tenderness
but typically lacks prominent nausea; a new or changing pattern requires clinical assessment.

\medterm{Chronic Thromboembolic Pulmonary Hypertension} Pulmonary hypertension caused by organized, persistent thromboembolic obstruction and
secondary pulmonary-vessel remodeling after pulmonary embolism. It is potentially
curable in selected patients by pulmonary endarterectomy or balloon pulmonary
angioplasty.

\medterm{Chronic Thyroiditis (Hashimoto Disease)} An autoimmune disease in which immune attack gradually damages the thyroid, often causing hypothyroidism, fatigue, cold intolerance, constipation, dry skin, or goiter.

\medterm{Chronic Toxicity} Adverse effects from repeated or prolonged exposure to low doses. Cumulative damage may
be irreversible. Examples: heavy metals (lead, mercury), asbestos, alcohol, solvent
exposure. Monitoring: biomarkers, organ function tests.

\medterm{Chronic Traumatic Encephalopathy} A progressive brain disorder associated with repeated head impacts and characterized by cognitive, mood, or behavioral changes.

\medterm{Chronic Venous Disease} A spectrum of venous outflow disorders caused by reflux, obstruction, or
calf-muscle-pump dysfunction. Findings range from edema and varicose veins to skin
changes, lipodermatosclerosis, and venous ulceration.

\medterm{Chronic Venous Insufficiency} Impaired venous return from the lower extremities due to venous valve incompetence or
obstruction. Leads to venous hypertension, edema, skin changes (lipodermatosclerosis,
hyperpigmentation), and venous stasis ulcers (medial malleolus). Risk factors: varicose
veins, DVT, obesity, prolonged standing.

\medterm{Chronic Venous Stasis} Persistent pooling of blood in the lower limbs caused by impaired venous return. It can produce
edema, skin discoloration, dermatitis, and venous ulcers.

\medterm{Chronic Vulvar Pain} Alternate index label; see \textit{Vulvodynia}.

\medterm{Chronicity} How long a disease, symptom, or process has lasted or how long it tends to continue. A condition may be described as acute when it is recent and short-lasting, or chronic when it persists over a long period.

\medterm{Chronobiology} Chronobiology studies biological timing and rhythms, including circadian sleep-wake cycles, seasonal rhythms, and their effects on health.
\medterm{Chronotropic} Chronotropic describes an influence on heart rate; a positive chronotropic effect increases rate, while a negative effect decreases it.
\medterm{Chrysarobin} Chrysarobin is a plant-derived anthracene compound formerly used topically for psoriasis and other skin disorders; it can irritate and stain skin.
\medterm{Churg Strauss Syndrome (Churg-Strauss Syndrome)} Alternate index label for Churg-Strauss Syndrome.

\medterm{Churg-Strauss Syndrome} A rare disease in which inflamed small and medium blood vessels are associated with asthma and too many eosinophils, a type of white blood cell. It can damage the lungs, nerves, skin, heart, or kidneys.

\synonyms
Churg Strauss Syndrome (Churg-Strauss Syndrome)

\medterm{Chvostek Sign} Twitching of facial muscles after tapping over the facial nerve near the ear. It can occur with low blood calcium but is not reliable enough by itself to diagnose that condition.

\medterm{Chvostek'S Sign} Chvostek’s sign is twitching of facial muscles after tapping near the facial nerve, sometimes associated with hypocalcemia but neither sensitive nor specific.
\medterm{Chylangioma} A usually noncancerous abnormal collection of lymphatic vessels, often in the neck or chest. It may form a soft swelling and can press on nearby structures; the modern term is often lymphangioma.

\medterm{Chyle} A milky fluid formed in intestinal lymph vessels after a meal. It contains lymph and absorbed dietary fat and normally travels through the thoracic duct into the bloodstream.

\medterm{Chylomicron} A particle made by the intestine after eating fat. It carries dietary fat and cholesterol through the lymph and blood, then leaves remnants for the liver to remove.

\medterm{Chylomicronemia Syndrome} Severe elevation of triglyceride-rich chylomicrons in blood, usually from genetic or acquired impairment of triglyceride clearance; it can cause abdominal pain, eruptive xanthomas, and recurrent acute pancreatitis.

\medterm{Chylothorax} A buildup of chyle in the space around a lung, usually after injury or blockage of the thoracic duct. It can cause breathlessness and requires treatment of the leak, often with drainage and dietary or surgical measures.

\medterm{Chyluria} Passage of lymphatic fluid into the urine, giving it a milky appearance. It usually
results from an abnormal communication between lymphatic vessels and the urinary tract.

\medterm{Chyme} A semifluid mixture of partially digested food and gastric or intestinal secretions that passes from the stomach into the small intestine.

\synonyms
Digestive chyme; gastric chyme.

\medterm{Chymotrypsin} Chymotrypsin is a pancreatic digestive enzyme that cleaves proteins in the small intestine; abnormal pancreatic release can contribute to tissue injury.
\medterm{Ci (Curie)} A unit of radioactivity equal to 3.7 billion nuclear disintegrations per second; the SI unit is the becquerel.

\medterm{CIC} Commonly, clean intermittent catheterization: periodic insertion of a catheter to empty the bladder, followed by removal; other expansions are possible by context.

\medterm{Cicatricial Alopecia} Hair loss caused by destruction and scarring of hair follicles, usually leading to permanent loss in affected areas.

\synonyms
Scarring alopecia; cicatricial hair loss.

\medterm{Cicatrization} Cicatrization is healing by formation of a scar after tissue injury, with replacement of damaged tissue by fibrous connective tissue.
\medterm{Ciclopia} A rare severe congenital malformation in which the eyes are fused or a single median eye
develops, usually as part of holoprosencephaly.

\medterm{Cidofovir} Cidofovir is an antiviral nucleotide analogue used for selected serious DNA-virus infections; kidney toxicity requires careful dosing and monitoring.
\medterm{CIDP} Chronic inflammatory demyelinating polyneuropathy is a chronic, immune-mediated,
peripheral neuropathy characterized by progressive or relapsing symmetrical weakness and
sensory loss in all four limbs, with areflexia or hyporeflexia, lasting for more than
eight weeks.

\medterm{Cigarette} A small cylinder of tobacco or other material wrapped in paper for smoking; smoke exposure delivers nicotine and many toxic and carcinogenic substances.

\medterm{Ciguatera} Food poisoning caused by eating reef fish containing ciguatoxins. It may cause
gastrointestinal illness followed by neurologic symptoms such as paresthesias and
temperature-sensation disturbance.

\medterm{Ciguatera Poisoning} Ciguatera poisoning follows eating reef fish containing ciguatoxins. Gastrointestinal symptoms may be followed by neurologic sensations such as temperature reversal, tingling, itching, weakness, or abnormal heart rate; treatment is supportive and prevention depends on avoiding high-risk fish.

\synonyms
Poisoning Ciguatera (Ciguatera Poisoning)
Toxin Ciguatera (Ciguatera Poisoning)

\medterm{Ciguatoxin} A heat-stable toxin produced by marine microorganisms and accumulated in some reef fish, causing ciguatera poisoning with gastrointestinal, neurologic, and temperature-sensation symptoms.

\medterm{Cilia} Tiny hair-like structures on cells that move fluid, mucus, or particles across a tissue surface. In the airways, they help clear inhaled material; some cilia also sense signals and help move an egg cell.
\medterm{Ciliary Body} Circumferential structure between iris and choroid, with ciliary muscle (accommodation
control via zonular fibers) and ciliary processes (aqueous humor production). Part of
the uveal tract (iris, ciliary body, choroid).

\medterm{Ciliary Dyskinesia} Abnormal movement of cilia that impairs clearance of mucus from the airways and can cause recurrent respiratory infections.

\medterm{Ciliary Muscle} A ring of smooth muscle in the eye that changes the lens shape for focusing at different distances and helps regulate aqueous humor outflow through ciliary-body structures.

\medterm{Ciliary Neuralgia} Pain arising from or felt along the ciliary nerves around the eye, often described as deep orbital or eye pain; the cause requires clinical evaluation.

\medterm{Ciliopathy} A disorder caused by abnormal cell cilia, the tiny structures that move material or sense signals. It may affect development, vision, kidney function, breathing, or several organs.

\medterm{Cimetidine} Cimetidine is an H2-receptor antagonist that reduces gastric acid secretion; it has numerous drug interactions and is used less often than newer agents.
\medterm{Cinacalcet} A calcimimetic drug that activates calcium-sensing receptors and lowers parathyroid hormone, used for selected forms of hyperparathyroidism and hypercalcemia.

\medterm{Cineangiocardiography} Radiographic recording of the heart and blood vessels during passage of contrast
material, used to visualize cardiac chambers, valves, and coronary or great vessels.

\medterm{Cinefluorography} Motion-picture recording of fluoroscopic images, used to document movement of organs or
contrast material during a diagnostic examination.

\medterm{Cineplasty} Surgical creation of a muscle-controlled stump or related reconstruction, historically
used to improve control of a prosthetic limb.

\medterm{Cingulotomy} Surgical interruption of selected fibers of the cingulum, sometimes performed as a
psychosurgical treatment for severe, treatment-resistant psychiatric illness or chronic
pain.

\medterm{Cinnarizine} Cinnarizine is an antihistamine with calcium-channel-blocking properties used in some countries for motion sickness and vestibular symptoms.
\medterm{Ciprofloxacin} Ciprofloxacin is a fluoroquinolone antibiotic used for selected bacterial infections; tendon injury, nerve effects, QT effects, and resistance limit indiscriminate use.
\medterm{Circadian Rhythm} An internal cycle of about 24 hours that helps regulate sleep, alertness, hormones, body temperature, and metabolism. Light and darkness reset the cycle each day through signals from the eyes to the brain.

\medterm{Circadian Rhythms} Approximately 24-hour biological cycles that regulate sleep and wakefulness, hormone release, body temperature, metabolism, and other physiological functions in response chiefly to light and darkness.
\medterm{Circle Of Willis} A ring of arteries at the base of the brain that links the main front and back blood supplies. It can provide an alternate route for blood if one artery becomes narrowed or blocked.

\medterm{Circulating Tumor DNA} Fragments of tumor-derived DNA detectable in blood, enabling molecular profiling,
treatment monitoring, and detection of minimal residual disease without tissue biopsy
(liquid biopsy).

\medterm{Circulation} Continuous movement of blood through the heart and blood vessels, or of lymph through lymphatic vessels, to transport oxygen, nutrients, hormones, and waste.

\medterm{Circulatory} Relating to the movement of blood or lymph and to the organs and vessels that carry them.

\medterm{Circulatory System} The heart and blood-vessel network that pumps and distributes blood through pulmonary and systemic circuits; broadly, the term may include the lymphatic system.

\medterm{Circumcision} Surgical removal of the foreskin covering the head of the penis. It may be performed for cultural, religious, preventive, or medical reasons; possible complications include bleeding, infection, pain, and injury, although serious problems are uncommon with proper care.
medical reasons such as severe phimosis, recurrent balanoposthitis, paraphimosis, or
lichen sclerosus (balanitis xerotica obliterans); elective infant circumcision for
religious, cultural, or preventive health reasons.

\medterm{Circumflex} Curving around; the circumflex coronary artery is a branch of the left coronary artery that travels around the heart's lateral and posterior surface.

\medterm{Circumflex Artery} A branch of the left coronary artery that travels in the left atrioventricular groove.
It supplies the lateral wall of the left ventricle, the left atrium, and in
left-dominant systems, the inferior wall and posteromedial papillary muscle.

\medterm{Circumvallate Papillae} Circumvallate papillae are large taste papillae arranged near the back of the tongue and surrounded by a trench containing taste buds and serous glands.
\medterm{Circumvallate Placenta} An extrachorial placental configuration in which the chorionic plate is smaller than the basal plate, producing a raised circumferential edge; it may be associated with bleeding, abruption, or growth restriction.

\medterm{Cirrhosis} Advanced scarring of the liver that permanently changes its structure and reduces blood flow and function. Causes include viral hepatitis, alcohol-related disease, fatty-liver disease, and autoimmune disorders; complications include jaundice, fluid buildup, bleeding, and liver failure.


\medterm{Cirrhosis Primary Biliary (Primary Biliary Cirrhosis)} Alternate index label for Primary Biliary Cirrhosis.

\medterm{Cisplatin} A platinum-based chemotherapy drug that forms DNA cross-links and is used to treat several cancers; important toxicities include kidney injury, hearing loss, neuropathy, and nausea.

\medterm{Citalopram} Citalopram is a selective serotonin-reuptake inhibitor used to treat depression and some anxiety disorders; nausea, sexual effects, and QT prolongation are important considerations.
\medterm{Citrate} Citrate is a salt or ion of citric acid that participates in metabolism and binds calcium; it is used in anticoagulant blood collection and some stone-prevention treatments.
\medterm{Citrate Synthase} An enzyme that begins an important energy-producing process inside cells. It joins acetyl-CoA with oxaloacetate to make citrate, the first step of the citric acid cycle, which helps release energy from food.

\medterm{Citric Acid Cycle} The citric acid cycle is a mitochondrial pathway that oxidizes acetyl-CoA and generates reducing equivalents for ATP production.
\medterm{Citric Acid Cycle (Krebs Cycle)} A central metabolic pathway occurring in the mitochondrial matrix that oxidizes
acetyl-CoA (derived from carbohydrates, lipids, and proteins) into CO2, generating 3
NADH, 1 FADH2, and 1 GTP (ATP) per turn. Rate-limiting enzyme: isocitrate dehydrogenase
(inhibited by ATP and NADH; activated by ADP).

\synonyms
Krebs Cycle

\medterm{Citric Acid Urine Test} A urine test measuring citrate excretion, used mainly in the evaluation of recurrent kidney stones because low urinary citrate increases the risk of calcium stone formation.

\medterm{Citrullinaemia} An inherited disorder of the urea cycle causing citrulline and ammonia buildup, which can lead to vomiting, confusion, seizures, coma, or death.
\medterm{Citrullinemia} An inherited disorder in which the body cannot properly remove nitrogen through the urea cycle. Citrulline and ammonia build up, causing vomiting, poor feeding, sleepiness, confusion, seizures, coma, or death without urgent treatment.
\medterm{CJD} Alternate index label; see \textit{Creutzfeldt-Jakob disease}.

\medterm{CK} Creatine kinase, an enzyme found in skeletal muscle, heart muscle, and brain whose blood level may rise after tissue injury.

\synonyms
Creatine phosphokinase; creatine kinase enzyme.

\medterm{CKD-Mineral Bone Disorder} A complication of chronic kidney disease in which abnormal calcium, phosphate, vitamin D, and parathyroid hormone levels weaken bones and may cause calcium deposits in blood vessels and other tissues.

\medterm{Cl (Chloride)} The chemical symbol for chloride, the negatively charged ion of chlorine and a major extracellular electrolyte.

\medterm{Cladribine} Cladribine is a purine analogue used for selected blood cancers and multiple sclerosis; lymphocyte depletion and infection risk require monitoring.
\medterm{Clairvoyance} Clairvoyance is an alleged ability to perceive information beyond ordinary sensory means; it is not an established medical diagnostic finding.
\medterm{Clamp} A surgical instrument used to grasp, compress, occlude, or hold tissue, blood vessels, tubing, or other structures.

\medterm{Clang Association} Clang association is speech in which word choice is driven by sound, rhyme, or pun rather than meaning, seen particularly in mania or psychosis.
\medterm{Clarithromycin} Clarithromycin is a macrolide antibiotic used for selected bacterial and mycobacterial infections; interactions and QT prolongation are important risks.
\medterm{Clark Level Of Invasion} An older system describing how deeply a melanoma has grown through the skin layers. Modern melanoma assessment relies more on tumor thickness, ulceration, spread, and other features.

\medterm{Class IA Antiarrhythmics} Heart-rhythm medicines that slow electrical conduction and lengthen the time heart cells need before they can fire again. They treat selected abnormal rhythms but can themselves cause dangerous rhythm disturbances.

\medterm{Class IB Antiarrhythmics} Heart-rhythm medicines that reduce abnormal electrical activity, especially in injured heart muscle. They are used mainly for selected fast rhythms arising from the lower chambers and require monitoring.

\medterm{Class IC Antiarrhythmics} Strong heart-rhythm medicines that slow electrical conduction without greatly changing recovery time. They can treat selected upper-chamber rhythms but may be dangerous in people with structural heart disease.

\medterm{Class III Antiarrhythmics} Heart-rhythm medicines that lengthen the recovery period between electrical beats. They treat selected abnormal rhythms but may slow the heart, affect organs, or prolong the electrical interval and cause dangerous rhythms.

\medterm{Classic Migraine} A migraine attack preceded or accompanied by an aura, such as temporary visual, sensory, or speech changes. Headache may follow, with nausea and light or sound sensitivity, although patterns vary.

\medterm{Claudication} Muscle pain or cramping brought on by exertion and relieved by rest, usually because of reduced arterial blood flow to a limb.

\synonyms
Vascular Claudication

\medterm{Claudication Distance} The distance a person can walk before exertional limb pain begins. It is used to describe symptoms and monitor change.

\medterm{Claustrophobia} An intense fear of enclosed or confined spaces that may cause anxiety, panic, or avoidance.
Symptoms can occur in places such as elevators, tunnels, or MRI scanners.

\medterm{Clavicle} The collarbone, a curved bone running from the breastbone to the shoulder blade. It supports the shoulder and commonly breaks after a fall or direct blow, causing pain, swelling, or visible deformity.

\medterm{Clavicle Fracture} A break in the collarbone that may occur during birth or after trauma. A newborn may show
asymmetric arm movement or tenderness; older children and adults may have pain, swelling, or deformity. Diagnosis is
usually made by examination and imaging.

\medterm{Clavicle Fracture (Intentional)} Intentional fracture of the fetal clavicle as a last resort to reduce the bisacromial
diameter in shoulder dystocia. The operator applies upward pressure on the midportion of
the anterior fetal clavicle with a finger, fracturing it. The clavicle heals without
long-term sequelae.

\synonyms
Intentional

\medterm{Clavus} A clavus is a localized painful thickening of skin, or corn, caused by repeated pressure or friction.
\medterm{Claw Foot} A toe deformity in which the toes bend upward at their bases and downward at their middle and end joints. It may result from muscle imbalance, nerve damage, arthritis, or unsuitable footwear.

\medterm{Claw Hand} A deformity in which the large finger joints bend backward while the smaller joints bend inward. It is commonly caused by ulnar-nerve damage and may impair grip and hand function.

\medterm{Claw Toe} A toe deformity characterized by extension at the metatarsophalangeal joint and flexion
at the proximal and distal interphalangeal joints.

\medterm{Clean Catch Urine Sample} A midstream urine specimen collected after cleaning the genital area to reduce contamination and improve interpretation of urinalysis or culture.

\medterm{Clean-Catch Specimen} A urine sample collected after cleansing the surrounding area and beginning urination, with the middle portion collected in a sterile container to reduce contamination.

\medterm{Cleaning Supplies And Equipment} Materials and devices used to remove soil and microorganisms from surfaces or equipment; selection and disinfection level depend on the item and infection risk.

\medterm{Cleaning To Prevent The Spread Of Germs} Cleaning and disinfection practices that remove or inactivate infectious organisms on hands, surfaces, equipment, and shared objects to reduce transmission.

\medterm{Clear Cell Carcinoma} Clear cell carcinoma is a cancer whose cells appear clear because of glycogen or lipid; it can arise in organs such as kidney, ovary, or endometrium.

\medterm{Clear Cell Carcinoma Of The Ovary} A type of ovarian cancer whose cells look clear under a microscope. It is often associated with endometriosis and may be found at an earlier stage, but can resist some standard treatments and increase the risk of blood clots; treatment depends on stage and overall health.

\medterm{Clear Cell Renal Cell Carcinoma} The predominant histologic subtype of renal cell carcinoma, accounting for approximately 75--80\% of all adult primary kidney cancers. Arises from the proximal convoluted tubule and is histologically characterized by malignant cells with clear lipid- and glycogen-rich cytoplasm; strongly associated with loss or inactivation of the \textit{VHL} tumor suppressor gene.

\synonyms
ccRCC; Conventional renal cell carcinoma
\medterm{Clear Liquid Diet} A limited diet consisting of transparent liquids such as water, clear broth, apple juice, and black coffee that leave minimal residue. It is used for short periods before procedures, during acute illness, or after surgery to maintain hydration while resting the gastrointestinal tract.

\medterm{Clear Margins} A histological finding indicating that no abnormal cells (e.g., cervical intraepithelial
neoplasia or cancer cells) are present at the edge of a tissue specimen removed during a
procedure such as LLETZ, cone biopsy, or surgical excision.

\medterm{Clearance} The volume of plasma from which a drug is completely removed per unit time, typically
expressed in milliliters per minute or liters per hour. Total clearance is the sum of
all clearance mechanisms, primarily hepatic and renal clearance.

\medterm{Cleft Lip} A birth defect in which the upper lip does not join together fully
before birth. It may occur alone or with a cleft palate.

\medterm{Cleft Lip And Cleft Palate} Birth defects in which the upper lip, the roof of
the mouth, or both do not join together fully before birth. They can cause problems
with feeding, speech, hearing, or teeth and may need treatment by several specialists.

\medterm{Cleft Lip And Palate Repair} Surgery to close a cleft in the lip, palate, or both and improve feeding, speech, hearing, breathing, or appearance. Repair is usually staged during childhood and requires long-term dental and specialist follow-up.


\medterm{Cleft Palate} A birth defect in which there is an opening in the roof of the
mouth. It may occur alone or with a cleft lip and can cause problems with feeding,
speech, hearing, or teeth.



\medterm{Cleidocranial Dysostosis} An inherited bone disorder, usually caused by a RUNX2 gene change, in which the collarbones may be small or absent and skull bones close late. Extra teeth, delayed tooth eruption, short stature, and joint problems are common.

\medterm{Cleidocranial Dysplasia} An inherited disorder of bone development affecting the collarbones and skull. It can cause delayed closure of skull joints, unusual teeth, short stature, and other skeletal differences.

\medterm{Clemastine} Clemastine is a first-generation antihistamine used for allergic symptoms; sedation, anticholinergic effects, and impaired driving are concerns.
\medterm{Clendamycin} Clendamycin is a misspelling of clindamycin, a lincosamide antibiotic that inhibits bacterial protein synthesis and can cause severe antibiotic-associated colitis.
\medterm{Click-Murmur Syndrome} Alternate index label; see \textit{Mitral valve prolapse}.

\medterm{Climacteric} The period of physiologic change around the end of reproductive capacity, including the
menopausal transition. It may involve changes in menstruation, vasomotor symptoms, and
other endocrine effects.

\medterm{Climate} Climate is the long-term pattern of weather and environmental conditions; climate affects infectious disease, heat illness, allergies, and population health.
\medterm{Clindamycin} Clindamycin is a lincosamide antibiotic used for selected infections; diarrhea and Clostridioides difficile colitis are important adverse effects.
\medterm{Clinic Visit} A planned outpatient appointment for consultation, examination, treatment, follow-up, or monitoring. It may include reviewing symptoms, medicines, test results, prevention needs, and decisions about further care.

\medterm{Clinical} Relating to direct observation, examination, diagnosis, or treatment of a patient, rather than solely to laboratory or basic research findings.

\medterm{Clinical Communication} Exchange of information between healthcare professionals, patients, and caregivers to
support diagnosis, treatment, safety, and informed decisions.

\medterm{Clinical Counseling} Guidance provided in healthcare to help patients make informed decisions about health,
prevention, treatment, and self-management.

\medterm{Clinical Depression} A depressive illness causing persistent sadness or loss of interest, together with changes in sleep, appetite, energy, thinking, or self-worth that significantly interfere with daily life.

\medterm{Clinical Depression (Depression)} Alternate index label for Depression.

\medterm{Clinical Deterioration} Worsening health that increases the risk of organ failure, disability, or death. Warning signs may include new confusion, breathing difficulty, abnormal vital signs, reduced urine, worsening pain, or poor circulation.

\medterm{Clinical Disease} Disease that produces recognizable signs or symptoms or other findings detectable in a patient, as opposed to a silent or subclinical state.


\medterm{Clinical Frailty Scale} A scale used by a healthcare professional to describe how fit, independent, or frail an older person is. It considers activity, daily function, and illness burden and supports care planning.

\medterm{Clinical Guideline} Carefully developed recommendations helping healthcare professionals make decisions about prevention, diagnosis, treatment, monitoring, and follow-up for specific conditions.

\medterm{Clinical Handover} Transfer of responsibility and relevant patient information from one healthcare
professional or team to another. Effective handover communicates diagnosis, current
status, risks, pending tasks, contingency plans, and accountability.

\medterm{Clinical Practice Guideline} Evidence-informed recommendations intended to help healthcare professionals and patients make decisions about a specific
clinical condition or situation.

\medterm{Clinical Presentation} The symptoms, signs, examination findings, and test results seen in a person with a health problem. The presentation can differ with the cause, severity, age, other illnesses, and the time since the problem began.

\medterm{Clinical Reasoning} The cognitive process of collecting, weighting, and integrating clinical information to
formulate diagnostic and management decisions. It includes problem representation,
differential diagnosis, probability assessment, and recognition of cognitive bias.

\medterm{Clinical Research} Research involving people or human health data to study disease, diagnosis, prevention, treatment, or healthcare delivery.

\medterm{Clinical Research Bias} A systematic error in how a health study is designed, conducted, analyzed, or reported that can make results misleading. Examples include choosing unrepresentative participants, measuring inaccurately, or failing to account for other influences.
\medterm{Clinical Trial} A planned study in which people receive or are observed for an intervention to evaluate its safety, effects, or outcomes. Trials may
be conducted in phases and for many types of health interventions.

\medterm{Clinical Trial Phases} Stages for testing a medical treatment: Phase I mainly examines safety and dose, Phase II studies early effectiveness and safety, Phase III compares it with usual care in larger groups, and Phase IV monitors longer-term use.

\medterm{Clinical Trials} Research studies involving human participants that evaluate the safety, effectiveness, risks, or best use of medical interventions or health strategies.

\medterm{Clinically Proven} A description meaning that research in people supports a treatment or claim. The quality, number, participants, outcomes, and limitations of the studies should be considered before judging how strong the evidence is.
\medterm{Clinicopathologic Correlation} Combining a person's symptoms, examination, scans, laboratory tests, and tissue findings to reach the most accurate diagnosis, especially when a small sample may not show the whole disease.

\medterm{Clinistix} Clinistix is a reagent-strip method historically used to detect glucose in urine through an enzymatic color reaction.
\medterm{Clinitest} Clinitest is a tablet-based copper-reduction test historically used to detect reducing substances in urine, including glucose; it is less specific than modern methods.
\medterm{Clinodactyly} A finger or toe that curves sideways from birth or during development, most often the little finger. It is usually mild but may interfere with movement or occur with a genetic condition.

\medterm{Clinoid Process} A small bony projection at the base of the skull near the eye socket and major nerves and blood vessels. It helps form the support around the pituitary region and optic canal.

\medterm{CLIP} A device used to hold tissue or structures together; surgical clips can close blood vessels, mark tissue, or occlude structures such as the vas deferens.

\medterm{Clip (Vascular)} A small device placed during a procedure to close or control a blood vessel. It may stop bleeding, prevent blood flow to a selected area, or secure a vessel permanently.
\synonyms
Vascular

\medterm{Clitoral} Relating to the clitoris, an erectile organ with a dense concentration of sensory nerves at the front of the vulva.

\medterm{Clitorectomy} Surgical removal of all or part of the clitoris, sometimes with removal of adjacent vulvar tissue; nonconsensual forms constitute female genital mutilation and can cause serious harm.

\medterm{Clitoridectomy} Surgical removal of part or all of the clitoris. In medicine it may be performed for
selected tumors or severe disease; nonmedical removal is a harmful practice with no
health benefit.

\medterm{Clitoris} A small, sensitive erectile organ located at the anterior junction of the labia minora,
beneath the clitoral hood (prepuce). The clitoris is the primary source of female sexual
pleasure and is homologous to the male penis.

\medterm{Clitoromegaly} Clitoromegaly is enlargement of the clitoris, which may be congenital or acquired from androgen excess, endocrine disease, medication, or other causes.
\medterm{CLL} Alternate index label; see \textit{Chronic lymphocytic leukemia}.

\medterm{Cloacal Exstrophy} The most severe congenital anomaly of the ventral abdominal wall, also called the OEIS
complex (Omphalocele, Exstrophy of the bladder, Imperforate anus, Spinal defects).
Incidence 1 in 200,000-400,000 births.

\synonyms
the OEIS.

\medterm{Cloacitis} Inflammation of a cloaca or cloacal region, particularly in disorders affecting the
common terminal outlet of the urinary, genital, and intestinal tracts.

\medterm{Clodronate} Clodronate is a bisphosphonate that reduces bone resorption and is used in selected bone and cancer-related conditions.
\medterm{Clofazimine} Clofazimine is a riminophenazine antimicrobial used mainly in multidrug treatment of leprosy and selected mycobacterial infections; skin discoloration is common.
\medterm{Clomifene} Clomifene, or clomiphene, is a selective estrogen-receptor modulator used to induce ovulation in selected infertility settings.
\medterm{Clomipramine} Clomipramine is a tricyclic antidepressant particularly effective for obsessive-compulsive disorder; anticholinergic, cardiac, and overdose toxicity are important.
\medterm{Clonal Evolution} The selection and expansion of tumor-cell subclones with biologic or treatment-related
advantages over time. It contributes to progression, metastasis, and acquired resistance
to anticancer therapy.

\medterm{Clonal Expansion} Multiplication of cells descended from one original cell. It may occur when immune cells respond to a trigger or when an abnormal cell population grows in cancer or another disorder.
\synonyms
Clonal proliferation; clone expansion.

\medterm{Clonazepam} Clonazepam is a benzodiazepine anticonvulsant and anxiolytic that enhances GABA activity; sedation, dependence, and respiratory depression can occur.
\medterm{Clonic Seizure} A seizure characterized by rhythmic, repetitive jerking caused by recurrent involuntary
muscle contractions. It may be focal or generalized, with or without impaired awareness,
and is classified using the clinical pattern, electroencephalography, and the underlying
cause.

\medterm{Clonidine} Clonidine stimulates central alpha-2 receptors to reduce sympathetic outflow and is used for hypertension and selected other conditions; abrupt withdrawal can cause rebound hypertension.
\medterm{Clonorchiasis} Infection with the Chinese liver fluke, Clonorchis sinensis. It affects the biliary
tract and may cause inflammation, obstruction, and long-term hepatobiliary
complications.

\medterm{Clonus} Repeated, rhythmic muscle jerks triggered by suddenly stretching a muscle, often at the ankle. It commonly indicates overactive reflexes from a problem affecting the brain or spinal cord.

\medterm{Clopidogrel} Clopidogrel is an antiplatelet drug that irreversibly blocks the platelet P2Y12 receptor, reducing arterial thrombosis risk after selected cardiovascular events or procedures.
\medterm{Closed Comedo} A small skin lesion formed when a follicular opening is blocked by sebum and keratin but
remains covered by the epidermis. It is commonly called a whitehead.

\synonyms
a whitehead.

\medterm{Closed Fracture} A fracture in which the broken bone does not communicate through a wound with the
external environment.

\medterm{Closed Reduction} Nonoperative realignment of a fractured or dislocated bone or joint by manipulation
without surgically opening the site.

\medterm{Closed Reduction Of A Fractured Bone} A nonsurgical procedure that realigns a broken bone by manipulating it through the skin without opening the fracture site, usually followed by immobilization and imaging to confirm alignment.


\medterm{Closed Suction Drain With Bulb} A flexible drain connected to a compressible bulb that creates gentle negative pressure to remove blood or fluid from a surgical or injured area; the bulb is emptied and output recorded as directed.

\medterm{Closed-Angle Glaucoma} Glaucoma caused by blockage of the eye's drainage angle, so fluid cannot leave and pressure rises. Sudden attacks can cause severe eye pain, blurred vision, halos, headache, nausea, and permanent sight loss.

\medterm{Closed-Loop Insulin Delivery} Automated insulin delivery systems that combine continuous glucose monitoring with an
insulin pump and control algorithm. The algorithm adjusts basal insulin delivery based
on sensor glucose values in real time. Hybrid closed-loop systems require user-initiated
meal boluses but automatically manage basal rates.

\medterm{Clostridioides Difficile} A spore-forming bacterium that can overgrow after antibiotics disturb normal gut bacteria. Its toxins inflame the colon and cause watery diarrhea, abdominal pain, fever, and sometimes life-threatening colitis.

\medterm{Clostridioides Difficile Infection} Diarrhea and colon inflammation caused by overgrowth of Clostridioides difficile, often after antibiotic use. It may be mild or severe and can recur; diagnosis requires appropriate stool testing in a person with symptoms.

\medterm{Clostridioides Difficile Pathogenesis} The mechanisms by which C. difficile causes disease. Antibiotic disruption of the normal
gut flora allows the organism to colonize; its toxins injure the colonic mucosa and
cause inflammation and diarrhea. Spores persist in the environment and facilitate
transmission.

\medterm{Clostridium} A genus of anaerobic, usually spore-forming bacteria that includes species capable of
causing diseases such as tetanus, botulism, gas gangrene, and antibiotic-associated
colitis.

\medterm{Clostridium Botulinum} A bacterium that produces botulinum toxin, which blocks nerve signals to muscles and causes weakness or paralysis. Food, wounds, or swallowed spores can cause botulism, a medical emergency affecting breathing and swallowing.

\medterm{Clostridium Difficile} Clostridioides difficile infection: a toxin-producing anaerobic gram-positive rod
causing antibiotic-associated diarrhea and pseudomembranous colitis, most commonly after
broad-spectrum antibiotics (clindamycin, cephalosporins, fluoroquinolones).

\medterm{Clostridium Difficile Colitis} Clostridioides difficile colitis is antibiotic-associated inflammation of the colon caused by toxin-producing bacteria. It causes watery diarrhea, abdominal pain, fever, and sometimes severe colitis; diagnosis and treatment depend on symptoms, testing, recurrence, and severity.

\synonyms
Difficile Clostridium (Clostridium Difficile Colitis)

\medterm{Clostridium Difficile Infection} An anaerobic, spore-forming, toxin-producing gram-positive bacillus causing antibiotic-associated diarrhea and pseudomembranous colitis. See Clostridioides difficile.
\medterm{Clostridium Perfringens} A bacterium that can cause rapidly destructive muscle infection with gas in tissues or food poisoning with sudden diarrhea and abdominal cramps. Severe wound infection requires emergency surgery and antibiotics.

\medterm{Clostridium Tetani} A soil bacterium whose toxin causes tetanus after entering a wound. It produces painful muscle stiffness and spasms, often beginning in the jaw, and can interfere with swallowing and breathing; vaccination prevents it.

\medterm{Closure} Closure is the joining or sealing of a wound, opening, vessel, duct, or anatomical structure, either naturally or by treatment.
\medterm{Clot} A clot is a semisolid mass of coagulated blood; clots may stop bleeding or obstruct a vessel and cause thrombosis or embolism.
\medterm{Clot Blood (Blood Clots)} Alternate index label for Blood Clots.

\medterm{Clot Buster} A medicine that breaks down a blood clot, also called a thrombolytic medicine. It may be used urgently for selected strokes, heart attacks, or lung clots, but it can cause dangerous bleeding and is not safe for everyone.

\medterm{Cloth Dye Poisoning} Exposure to cloth dyes that may cause irritation, gastrointestinal symptoms, allergic reactions, or systemic toxicity depending on the chemical composition.

\medterm{Clotrimazole} Clotrimazole is an azole antifungal used topically or orally in selected formulations for fungal infections of skin, mouth, or vagina.
\medterm{Clotting} The body's process of stopping bleeding by narrowing a damaged vessel, forming a platelet plug, and producing a protein mesh that strengthens the plug. Excessive or inappropriate clotting can block blood vessels.

\medterm{Cloudy Cornea} Loss of corneal transparency from edema, scarring, infection, dystrophy, trauma, or deposits, potentially causing blurred vision, pain, photophobia, or glare.

\medterm{Clove Oil} Clove oil contains eugenol and is used in some dental and topical products; concentrated ingestion or exposure can cause irritation, liver injury, or seizures.
\medterm{Clozapine} Clozapine is an atypical antipsychotic for treatment-resistant schizophrenia and selected suicidality risk; neutropenia, myocarditis, seizures, constipation, and metabolic effects require monitoring.
\medterm{Club Drugs List And Side Effects} Club drugs include substances such as MDMA, ketamine, GHB, and others used recreationally in social settings. Effects may include agitation, hallucinations, seizures, dangerous temperature elevation, respiratory depression, cardiovascular toxicity, impaired consent, and overdose.

\medterm{Clubbing} Enlargement and rounding of the fingertips or toes with more curved nails and loss of the normal nail angle. It may accompany long-term lung, heart, bowel, liver, or other disease and should be medically assessed.

\medterm{Clubbing Of The Fingers Or Toes} Bulbous enlargement of fingertip or toe soft tissue with increased nail curvature caused by chronic changes in the nail bed; it can accompany lung, heart, gastrointestinal, or other disease.

\medterm{Clubfoot} A birth defect in which a baby's foot is turned inward and
downward. The foot may be stiff and difficult to move into a normal position.

\medterm{Clubfoot (Talipes Equinovarus)} Alternate name; see \textit{Clubfoot}.

\medterm{Clubfoot Repair} Treatment to gradually straighten a baby's turned-in foot, usually with gentle stretching and casts, followed by a brace. Some children need a small tendon procedure or later surgery if the deformity returns or is severe.

\medterm{Clubhand} A birth defect in which the hand and wrist curve toward the thumb or little-finger side because forearm bones or other tissues developed abnormally. Treatment depends on function and may include therapy or surgery.

\medterm{Clumping} The gathering of cells, particles, or proteins into groups. In a laboratory sample, clumped red blood cells or platelets can affect test results; in the body, clumping can sometimes block vessels.
\medterm{Cluster} A group of similar events, findings, cells, or cases occurring together in space, time, or clinical pattern.

\medterm{Cluster Headache} Repeated attacks of extremely severe pain around or behind one eye, lasting about 15 minutes to three hours. The same side may tear, redden, swell, or feel congested, and the person often becomes restless.

\medterm{Cluster Headaches} Cluster headache causes repeated attacks of severe one-sided orbital or temporal pain with tearing, nasal symptoms, eyelid drooping, or restlessness. Attacks occur in bouts and require targeted acute and preventive treatment; a first sudden severe headache needs urgent evaluation.

\medterm{Clutton'S Joints} Painless swelling of both knees or other large joints caused by long-term inflammation in children with congenital syphilis. Joint movement is usually preserved, but the underlying infection requires treatment.

\medterm{Clysis} An older method of giving fluid or medicine through the rectum, sometimes used for hydration or bowel treatment. It is rarely used today because safer and more effective routes are available.
\medterm{Cm (Centimeter)} The abbreviation for centimeter, a unit of length equal to one hundredth of a meter or 10 millimeters.

\medterm{CMV} Cytomegalovirus, a common herpesvirus that may cause mild or asymptomatic infection but can cause serious disease in fetuses and people with
weakened immunity.

\synonyms
Cytomegalovirus (CMV)

\medterm{CMV - Gastroenteritis/Colitis} Inflammation and ulceration of the stomach or intestine caused by cytomegalovirus, usually in people with substantial immune suppression. It may cause abdominal pain, fever, diarrhea, bleeding, or perforation and requires tissue or other targeted evaluation.

\medterm{CMV Blood Test} A blood test that looks for cytomegalovirus antibodies, viral proteins, or viral genetic material. Antibodies usually show past or recent exposure, while viral material can help detect active infection, especially in people with weakened immunity.

\medterm{CMV Colitis} Inflammation of the colon caused by cytomegalovirus, usually in someone with severely weakened immunity. It may cause bloody diarrhea, abdominal pain, fever, or weight loss and is confirmed with appropriate testing, often including biopsy.

\medterm{CMV Pneumonia} Pneumonia caused by cytomegalovirus, usually in people with substantial immune suppression; fever, cough, shortness of breath, and low oxygen may occur and require targeted evaluation and antiviral treatment.

\medterm{CMV Retinitis} Infection and inflammation of the retina caused by cytomegalovirus, usually in people with severely weakened immunity. It may cause floaters, blurred or missing areas of vision, and blindness unless treated urgently.

\medterm{CNS Infection} Infection of the brain, spinal cord, meninges, or cerebrospinal-fluid spaces caused by bacteria, viruses, fungi, parasites, or other organisms; presentation and treatment depend on the pathogen and anatomic site.

\medterm{CNS Lymphoma} A lymphoma involving the central nervous system, most often diffuse large B-cell lymphoma. It may involve brain parenchyma, meninges, eyes, or cerebrospinal fluid and requires urgent hematologic and neurologic evaluation.

\medterm{CO Poisoning} Toxicity from inhalation of carbon monoxide, a colorless, odorless gas binding
hemoglobin ~200-300 times more avidly than oxygen, forming carboxyhemoglobin (COHb) and
reducing oxygen delivery. See carbon monoxide poisoning.

\medterm{CO2 Blood Test} A blood chemistry measurement that usually estimates serum total carbon dioxide, which largely reflects bicarbonate. It helps assess acid-base and electrolyte disorders but must be interpreted with sodium, chloride, kidney function, and, when needed, blood-gas results.

\medterm{Coagulate} To change from a liquid to a semisolid or solid state, especially the formation of a blood clot through coagulation.

\synonyms
Clot; clot formation.

\medterm{Coagulation} The body's process for stopping bleeding by forming a stable blood clot. Platelets first gather at an injured vessel, then blood proteins create a mesh of fibrin that strengthens the plug; too little clotting causes bleeding, while too much can cause harmful clots.
\medterm{Coagulation Cascade} The coagulation cascade is a series of enzymatic reactions leading to fibrin clot
formation, involving serine protease activation cascades. The intrinsic pathway is
initiated by contact activation when factor XII binds exposed collagen, sequentially
activating factors XI, IX, and the cofactor VIII.

\medterm{Coagulation Defects} Coagulation defects are abnormalities of clotting proteins, platelets, vessels, or regulation that cause excessive bleeding or inappropriate thrombosis.
\medterm{Coagulation Factor} A blood protein required for the coagulation cascade and formation of a stable fibrin clot; deficiencies or inhibitors can cause bleeding, while excessive activation can contribute to thrombosis.

\medterm{Coagulation Factor Deficiencies} Inherited or acquired disorders in which one or more clotting factors are reduced or
function abnormally, causing excessive bleeding. The clinical pattern and laboratory
findings depend on the affected factor.

\medterm{Coagulation Profile} Panel of tests assessing the coagulation cascade. PT (prothrombin time, 11-14 sec):
measures extrinsic and common pathways; elevated in vitamin K deficiency, warfarin,
liver disease, DIC. INR standardizes PT for warfarin monitoring.

\medterm{Coagulation Time} Coagulation time is the time required for a blood sample to form a clot under specified test conditions; modern tests usually measure particular pathways more precisely.
\medterm{Coagulative Necrosis} The most common form of necrosis, typically resulting from ischemic injury or infarction
in most solid organs except the brain. The fundamental process involves denaturation of
both structural proteins and enzymatic proteins, which prevents the immediate
liquefaction of the dead tissue.

\medterm{Coagulopathy} A disorder of hemostasis that impairs normal clot formation or regulation and
predisposes to abnormal bleeding or thrombosis. It may result from coagulation-factor
deficiency, platelet dysfunction, anticoagulant exposure, liver disease, consumptive
coagulopathy, or systemic illness.

\medterm{Coal Worker'S Pneumoconiosis} Pneumoconiosis from coal dust. Simple CWP shows upper lobe macules. Complicated CWP
(progressive massive fibrosis) involves coalescent macules. May be associated with
Caplan syndrome.

\medterm{Coarctation} A localized narrowing of a tubular anatomical structure, especially coarctation of the
aorta, which obstructs blood flow and produces characteristic pressure differences.

\medterm{Coarctation Of The Aorta} A congenital narrowing of the aorta, typically near the ductus arteriosus. Presents with
upper extremity hypertension, diminished femoral pulses, and a blood pressure gradient
between arms and legs. Associated with bicuspid aortic valve and other congenital heart
defects.

\medterm{Coarse Crackles} Lower-pitched, longer-duration crackles heard throughout inspiration, caused by fluid or
secretions in larger airways. They are heard in COPD, bronchiectasis, pneumonia, and
pulmonary edema. They are less specific than fine crackles.

\medterm{Coats Disease} An uncommon retinal disorder characterized by abnormal retinal blood vessels, leakage,
exudation, and sometimes retinal detachment, usually affecting one eye in childhood.

\medterm{Coats Disease (Coats' Disease)} Alternate index label for Coats' Disease.

\medterm{Cobalamin} Vitamin B12, a nutrient needed for DNA production, healthy red blood cells, nerve function, and normal processing of certain fats and proteins.

\medterm{Cobalamin Deficiency} Too little vitamin B12, which is needed for healthy blood cells and normal nerve function. It may cause anemia, fatigue, sore tongue, numbness, balance problems, memory changes, or irreversible nerve injury if untreated.
\medterm{Cobalt} Cobalt is a metal present in vitamin B12 and some industrial materials; excess exposure can injure the heart, thyroid, lungs, or skin.
\medterm{Cobalt Poisoning} Toxic exposure to cobalt that may cause gastrointestinal, neurologic, thyroid, cardiac, or lung effects after significant or prolonged exposure.

\medterm{Cobalt-57} A radioactive form of cobalt used as a test source for checking and calibrating gamma cameras. It is mainly used for equipment quality control and is not normally given to patients as a diagnostic or treatment medicine.

\medterm{Cobalt-Chromium} A strong, corrosion-resistant metal mixture used for cast dental restorations, partial dentures, crowns, and selected medical devices. Some people may develop sensitivity, and prolonged exposure to cobalt can be harmful.

\medterm{Cocaine} A stimulant drug that blocks monoamine reuptake and can also act as a local anesthetic; use can cause dependence, cardiovascular complications, seizures, stroke, and overdose.

\medterm{Cocaine Use Disorder} A problematic pattern of cocaine use leading to clinically significant impairment or distress.
It may involve craving, impaired control, continued use despite harm, tolerance, or withdrawal.

\synonyms
Cocaine And Crack Abuse

\medterm{Cocaine Cardiotoxicity} Cocaine blocks norepinephrine and dopamine reuptake, causing sympathomimetic effects:
tachycardia, hypertension, and coronary vasospasm. Acute myocardial infarction can occur
even with normal coronary arteries due to vasospasm, accelerated atherosclerosis, and
prothrombotic effects. Treatment: benzodiazepines for sympatholysis.

\medterm{Cocaine Intoxication} A toxic reaction to cocaine that may cause agitation, overheating, high blood pressure, seizures, abnormal rhythms, heart attack, stroke, or sudden death.

\medterm{Cocaine Toxicity} A toxic state caused by cocaine, often producing agitation, high blood pressure, rapid heart
rate, hyperthermia, seizures, coronary vasospasm, myocardial infarction, or stroke. Treatment is supportive and may
include sedation, temperature control, and management of cardiovascular or neurologic complications.

\medterm{Cocaine Withdrawal} Physical and emotional symptoms after stopping or reducing cocaine, including fatigue, depressed mood, increased appetite, sleep changes, irritability, and cravings. Severe depression or suicidal thoughts require urgent help.

\medterm{Cocarboxylase} Cocarboxylase is thiamine pyrophosphate, the active coenzyme form of vitamin B1 required for several oxidative carbohydrate and amino-acid reactions.

\medterm{Coccidioides Complement Fixation} A blood test that detects complement-fixing antibodies to Coccidioides; antibody levels can support diagnosis and help assess the severity or progression of coccidioidomycosis.

\medterm{Coccidioides Immitis} A dimorphic soil fungus and a principal cause of coccidioidomycosis, or Valley fever, acquired by inhaling airborne arthroconidia. Infection is usually pulmonary but can disseminate, especially in people with impaired cellular immunity.

\medterm{Coccidioides Precipitin Test} A serologic test that detects antibodies to Coccidioides, typically early immunoglobulin M antibodies, and supports evaluation of suspected coccidioidomycosis when interpreted with symptoms and other testing.

\medterm{Coccidioidomycosis} A fungal infection caused by Coccidioides fungi, usually acquired by breathing spores from disturbed soil. It most often affects the lungs and causes cough, fever, tiredness, or chest pain; rarely, it spreads to the skin, bones, joints, or brain.
\commonname{Valley fever}

\medterm{Coccidiosis} Infection by coccidian protozoa. In humans it may cause intestinal disease or, in
immunocompromised patients, infection of other organs.

\medterm{Coccus} A round or nearly round bacterium. These bacteria may occur alone, in pairs, chains, or clusters; their shape helps describe them but does not identify the species or illness by itself.
\medterm{Coccydynia} Pain in the region of the coccyx, often worsened by sitting. Causes include trauma,
childbirth, degenerative change, and less commonly infection or tumor.

\synonyms
Pain Tailbone (Coccydynia)

\medterm{Coccygectomy} Surgical removal of the coccyx, reserved for selected cases of severe persistent
coccygeal pain or a lesion that has not responded to conservative treatment.

\medterm{Coccyx} The small pointed bone at the bottom of the spine, commonly called the tailbone. It connects to the sacrum, supports pelvic-floor tissues, and may become painful after a fall, childbirth, or prolonged sitting.

\medterm{Cochlea} The spiral-shaped part of the inner ear that changes sound vibrations into nerve signals. Fluid movement inside it stimulates sensory hair cells, allowing the brain to recognize pitch and loudness.

\medterm{Cochlear Implant} An electronic device for severe hearing loss that bypasses damaged inner-ear hair cells and stimulates the hearing nerve. It can provide useful sound perception but does not restore normal hearing and requires training.

\medterm{Cockayne Syndrome} A rare inherited disorder of DNA repair characterized by progressive neurologic
dysfunction, growth failure, developmental impairment, photosensitivity, and a
prematurely aged appearance.

\medterm{Code Black} A hospital emergency alert whose meaning varies by institution, often indicating a severe security threat, bomb threat, or other critical incident; local policy determines the response.

\medterm{Code Blue} A hospital emergency signal indicating that a person may be in cardiopulmonary arrest
and requires immediate resuscitation response.

\medterm{Code, Hospital} A standardized hospital announcement used to summon a specialized response to an emergency; the meaning of each code depends on the hospital's policy.

\medterm{Codeine Overdose} Overdose of codeine, an opioid, that can cause respiratory depression, impaired consciousness, pinpoint pupils, and coma. Naloxone may be required in emergency care.

\medterm{Codman Triangle} A triangular area of new bone seen on an X-ray when a rapidly growing bone lesion lifts the outer covering of bone. It suggests aggressive bone disease but does not identify a specific diagnosis.

\medterm{Codon} A group of three genetic building blocks in messenger RNA that specifies one amino acid or signals the end of protein production. Several different codons can specify the same amino acid.

\medterm{Coenzyme} A non-protein organic molecule that binds to an enzyme (apoenzyme) and is required for
its catalytic activity. Coenzymes often act as transient carriers of specific chemical
groups or electrons.

\medterm{Coenzyme Q} A fat-soluble molecule in mitochondrial membranes that transfers electrons during cellular energy production. It also acts as an antioxidant, and the body makes it naturally, although levels may change with age or illness.

\medterm{Coenzyme Q10} CoQ10, a lipid-soluble compound in mitochondrial electron transport and an antioxidant; it is also sold as a dietary supplement, although benefits vary by condition.

\medterm{Cofactor} A nonprotein helper, such as a mineral ion or vitamin-derived molecule, that an enzyme needs to perform a chemical reaction. Without its cofactor, the enzyme may not work normally.

\medterm{Coffin-Lowry Syndrome} A rare inherited developmental disorder, usually linked to an RPS6KA3 gene change. It can cause intellectual disability, distinctive facial features, short stature, spinal curvature, tapered fingers, and sudden falls triggered by surprise.

\synonyms
CLS

\medterm{Cogan’S Syndrome} A rare immune disorder causing inflammation of the cornea and inner ear. It may produce eye pain, reduced vision, dizziness, ringing, or hearing loss and can sometimes inflame large blood vessels.

\medterm{Cognition} The mental processes involved in acquiring, storing, retrieving, and using information, including attention, memory, language, reasoning, and judgment.

\medterm{Cognitive} Relating to mental processes such as attention, memory, learning, language, reasoning, perception, judgment, problem solving, and decision-making.

\medterm{Cognitive Behavior Therapy} A talking treatment that helps a person notice unhelpful thoughts and behaviors and develop more useful ways to respond. It is used for problems such as anxiety, depression, pain, insomnia, and some eating disorders.

\medterm{Cognitive Behavioral Therapy} A structured talking therapy that helps a person identify unhelpful thoughts and behaviors and practice more useful responses. It is used for depression, anxiety, insomnia, pain, and other conditions.

\medterm{Cognitive Behavioral Therapy For Back Pain} A structured psychological treatment that addresses pain-related thoughts, behaviors, sleep, activity, and coping to reduce disability and improve function in chronic back pain.

\medterm{Cognitive Behavioral Therapy For Insomnia} Cognitive behavioral therapy for insomnia is a structured, evidence-based program that
targets the thoughts and behaviors perpetuating poor sleep. Components include sleep
restriction, stimulus control, cognitive restructuring, and sleep hygiene education.

\medterm{Cognitive Decline} A reduction over time in one or more cognitive abilities, such as memory, attention, language, reasoning, or visuospatial function.

\medterm{Cognitive Distortions} Systematic errors in thinking that perpetuate negative emotional states and maladaptive
behaviors. Common distortions include all-or-nothing thinking, catastrophizing, mind
reading, overgeneralization, emotional reasoning, and should statements. CBT focuses on
identifying and modifying these patterns.

\medterm{Cognitive Fitness} The active maintenance of mental sharpness and brain health through stimulating
activities, physical exercise, social engagement, and proper nutrition. Regular
challenges such as learning new skills, puzzles, reading, and social interaction support
neuroplasticity and cognitive reserve.

\medterm{Cognitive Function} The mental abilities used for attention, memory, language, learning, reasoning, judgment, planning, problem solving, and understanding information.

\synonyms
Cognition; cognitive ability.

\medterm{Cognitive Health And Memory} Cognitive health encompasses the ability to think clearly, learn, remember, and make
decisions, with memory being a core component that involves encoding, storing, and
retrieving information. Age-related changes in processing speed and recall are normal,
but persistent decline may signal mild cognitive impairment or dementia.

\medterm{Cognitive Impairment} A decline in thinking abilities such as memory, attention, language, judgment, or problem solving. It may be mild with independence preserved or severe enough to interfere with everyday activities.

\medterm{Cognitive Rehabilitation} Structured therapy to improve cognitive function or help a person compensate for difficulties
with memory, attention, language, or executive skills after illness or injury.

\medterm{Cognitive Reserve} The brain's ability to use alternative strategies or networks to maintain function despite injury or disease.

\medterm{Cognitive Screening} A brief assessment used to identify possible problems with memory, attention, language, or other cognitive functions. It does not by
itself diagnose dementia or another disorder.

\medterm{Cognitive Therapy} A psychological treatment that helps a person identify and change unhelpful thoughts and beliefs and develop more adaptive behaviors; cognitive behavioral therapy combines it with behavioral methods.

\medterm{Cohort} A group of people who share a defined characteristic or experience during a specified period and are studied over time or compared with another group.

\medterm{Cohort Study} A research study following groups with different exposures over time to compare new disease rates, outcomes, or risks.

\medterm{Coil (Vascular)} A catheter-delivered wire device used to block selected blood vessels, often to treat aneurysms or bleeding.

\synonyms
Vascular

\medterm{Coincidence Detection} Principle of PET imaging. Two 511 keV photons detected simultaneously (within
nanoseconds) by opposing detectors. Line of response defined. Time-of-flight PET
improves localization.

\medterm{Coitus} Sexual intercourse, traditionally meaning penile-vaginal penetration; the term does not include every form of sexual activity or intimate contact.

\medterm{Coitus Interruptus} The withdrawal method of contraception, in which the penis is removed from the vagina before ejaculation; it is less reliable than many other contraceptive methods.

\medterm{Coke} Slang for cocaine, a stimulant drug; the word can also mean a carbon fuel product, so context is essential.

\medterm{Cold (Common Cold)} Alternate index label for Common Cold.

\medterm{Cold Abscess} A pocket of pus that develops slowly with little redness, warmth, or pain. It may occur with tuberculosis or other long-term infections and still requires medical evaluation and treatment.

\medterm{Cold Agglutinin} An antibody that causes red blood cells to agglutinate at relatively low temperatures.
Pathologic cold agglutinins can cause autoimmune hemolytic anemia.

\medterm{Cold Agglutinin Disease} An autoimmune blood disorder in which antibodies become active in cooler parts of the body and cause red blood cells to stick together and break down. It may cause anemia, fatigue, jaundice, or cold-related color changes.

\medterm{Cold Agglutinin Syndrome} An immune-mediated disorder in which antibodies react with red blood cells at low temperatures,
potentially causing hemolytic anemia or circulation-related symptoms.

\medterm{Cold Agglutinins} IgM autoantibodies that bind red blood cell antigens at temperatures below 37 degrees
Celsius, causing agglutination and potentially complement-mediated hemolysis in cooler
peripheral circulation.

\medterm{Cold Allergy} Alternate index label; see \textit{Cold urticaria}.

\medterm{Cold Exposure} Alternate index label; see \textit{Hypothermia}.

\medterm{Cold Flu Allergy Treatments (Cold, Flu, Allergy)} Alternate index label for Cold, Flu, Allergy.

\medterm{Cold Hands} Hands that feel unusually cold, which may result from environmental exposure, reduced blood flow, Raynaud phenomenon, anemia, thyroid disease, or other medical conditions.

\medterm{Cold Injury} Harm caused by low temperature, including hypothermia, frostbite, chilblains, and nonfreezing tissue injury. Risk increases with prolonged exposure, wetness, wind, or poor circulation and may affect the whole body or a local area.

\medterm{Cold Intolerance} Unusual sensitivity or inability to tolerate cold temperatures, which may occur with hypothyroidism, anemia, low body weight, vascular disease, or other conditions.

\medterm{Cold Medicines And Children} Safety guidance for giving medicines to children with cold symptoms; many combination products provide little benefit in young children and can cause dosing errors or serious adverse effects.

\medterm{Cold Sore} A contagious infection of the lips or skin around the mouth, usually caused by herpes simplex virus type 1 and occasionally by HSV-2. It often begins with tingling or burning, followed by grouped fluid-filled blisters that crust and heal in about 1--2 weeks. The virus remains in the body and can reactivate, sometimes after illness, stress, fever, or sunlight exposure.


\medterm{Cold Urticaria} Raised, itchy patches or swelling that appear after the skin is exposed to cold. A large area of cold exposure, such as swimming in cold water, can rarely cause fainting or a life-threatening allergic reaction.

\synonyms
Cold Allergy

\medterm{Cold Wave Lotion Poisoning} Harmful exposure to chemicals in a permanent-wave hair product. Swallowing or contact may burn the mouth, skin, eyes, or airways and can cause vomiting or breathing problems; urgent poison-control advice is needed.
\medterm{Cold, Common} Alternate index label; see \textit{Common cold}.

\medterm{Cold, Flu, Allergy} Colds, influenza, and allergies can all cause nasal symptoms and cough but differ in cause, fever pattern, duration, and contagiousness. Viral illness usually improves with supportive care, while allergies respond to trigger reduction and appropriate antihistamine or nasal treatment.

\synonyms
Cold Flu Allergy Treatments (Cold, Flu, Allergy)

\medterm{Colds And Flu - Antibiotics} An explanation that antibiotics treat susceptible bacterial infections, not uncomplicated viral colds or influenza; unnecessary use causes adverse effects and antimicrobial resistance.



\medterm{Colectomy} Surgery to remove part or all of the colon. It may treat cancer, severe inflammatory bowel disease, diverticular complications, blocked bowel, or inherited high cancer risk; the amount removed depends on the disease.

\medterm{Colibacillosis} Disease caused by pathogenic strains of Escherichia coli, which may produce intestinal
infection, urinary-tract infection, sepsis, or other illness.

\medterm{Colic} Severe wave-like pain caused by contraction or blockage of a hollow organ, such as the intestine, gallbladder, or urinary tract. The term also describes repeated crying in an otherwise healthy young infant.

\medterm{Colicky Pain} Intermittent, wave-like cramping pain caused by contraction, distension, or obstruction of a hollow organ; common examples include biliary, renal, and intestinal colic.


\medterm{Coliform} A group of bacteria used as an indicator of possible contamination in water or food. Finding them does not prove that disease-causing bacteria are present, because many coliforms are harmless environmental organisms.

\medterm{Colitis} Inflammation of the colon, which may cause diarrhea, abdominal pain, blood or mucus in stool, fever, or urgency. Causes include infection, inflammatory bowel disease, reduced blood flow, medicines, and radiation.

\medterm{Colitis Collagenous (Lymphocytic Colitis)} Alternate index label for Lymphocytic Colitis.


\medterm{Colitis Lymphocytic (Lymphocytic Colitis)} Alternate index label for Lymphocytic Colitis.

\medterm{Colitis Microscopic (Lymphocytic Colitis)} Alternate index label for Lymphocytic Colitis.

\medterm{Colitis Ulcerative (Ulcerative Colitis)} Alternate index label for Ulcerative Colitis.

\medterm{Colitis, Ischemic} Alternate index label; see \textit{Ischemic colitis}.

\medterm{Colitis, Microscopic} Alternate index label; see \textit{Microscopic colitis}.

\medterm{Colitis, Pseudomembranous} Alternate index label; see \textit{Pseudomembranous colitis}.

\medterm{Colitis, Ulcerative} Alternate index label; see \textit{Ulcerative colitis}.

\medterm{Collagen} A strong protein that forms much of the body's connective tissue. Different types support and strengthen skin, bone, cartilage, tendons, ligaments, blood vessels, and many internal structures.

\medterm{Collagen Disease} A systemic disorder involving connective tissue and often immune-mediated injury to
collagen-rich structures. The term has historically included diseases such as systemic
lupus erythematosus and systemic sclerosis.

\medterm{Collagen Implant} A medical material made from or based on collagen, a structural protein, used to support tissue repair or add soft-tissue volume. Sources may be animal-derived or manufactured. Its uses, durability, and risks depend on the product and the treated body area.

\medterm{Collagen Vascular Disease} An older term for systemic connective-tissue or autoimmune diseases such as lupus, systemic sclerosis, and inflammatory myopathies; modern diagnosis uses the specific disease name.

\medterm{Collagenase} An enzyme that breaks down collagen, a major structural protein in skin, bone, cartilage, and connective tissue. Natural collagenases help remodel tissue, while medicines or bacterial enzymes can be used in selected treatments or cause tissue damage.

\synonyms
Collagen-digesting enzyme; collagen-degrading enzyme.

\medterm{Collagenous Colitis} A type of microscopic colitis in which a thick layer of collagen develops beneath the colon lining. It usually causes persistent, watery, nonbloody diarrhea and may require colon biopsy for diagnosis.

\medterm{Collapsed Arch} Alternate index label; see \textit{Flatfeet}.

\medterm{Collapsed Lung} A common description of pneumothorax, in which air enters the pleural space and
partially or completely collapses the lung. It may cause sudden chest pain and shortness
of breath and can be life-threatening when tension develops. See also Pneumothorax.

\medterm{Collapsed Lung (Pneumothorax)} Air trapped between the lung and chest wall, causing part or all of the lung to collapse. It may cause sudden chest pain and breathlessness; pressure that shifts the heart and other structures is immediately life-threatening.

\medterm{Collateral} Providing an alternative or parallel route, such as collateral blood vessels that supply tissue around an obstruction; in other contexts it means accompanying or secondary.

\medterm{Collateral Circulation} Alternative vascular pathways that develop when coronary arteries become chronically
narrowed. Collaterals connect different coronary artery territories and can partially
compensate for stenosis. They develop gradually in response to ischemia and are more
prominent in patients with chronic coronary artery disease.


\medterm{Collateral Ligament Injury} MCL (medial) or LCL (lateral) knee ligament injury. MCL more common, often from valgus
stress. Treatment: braces, physical therapy, surgical repair if severe.

\medterm{Collecting Duct} Final segment of the nephron where fine-tuning of urine concentration and composition
occurs. Principal cells respond to antidiuretic hormone by inserting aquaporin-2 water
channels into the luminal membrane, enabling water reabsorption and urine concentration.

\medterm{College Students And The Flu} Influenza prevention and care for college students, emphasizing vaccination, hand and respiratory hygiene, staying home when ill, and prompt assessment for severe symptoms or high-risk conditions.

\medterm{Colles Fracture} A break near the wrist end of the larger forearm bone, with the broken hand-side fragment displaced backward. It usually follows a fall onto an outstretched hand and causes wrist pain, swelling, and deformity.


\medterm{Collimator} A device in an imaging camera that allows radiation traveling in selected directions to reach the detector. Its design affects the balance between image detail and detection sensitivity.

\medterm{Colloids} Fluids containing large particles that tend to remain in the bloodstream longer than simple salt solutions. Some are used for selected fluid replacement, but benefits and risks depend on the product and the patient's condition.

\medterm{Coloboma} A birth defect caused by incomplete formation of part of the eye, leaving a notch or gap in the iris, retina, choroid, or optic nerve. Vision problems vary with the structure and amount affected.

\medterm{Coloboma Of The Iris} A congenital keyhole- or notch-shaped defect in the iris caused by incomplete embryonic closure, sometimes associated with colobomas in the retina, optic nerve, or other organs.

\medterm{Colocolonic Intussusception} Telescoping of one part of the colon into the next section. It is uncommon in adults and may be caused by a polyp, tumor, or other lesion, producing intermittent pain, bloating, vomiting, or altered bowel movements.

\medterm{Cologne Poisoning} Ingestion of cologne containing alcohol and fragrances that may cause intoxication, hypoglycemia, respiratory depression, or gastrointestinal irritation.

\medterm{Cologuard} A stool-based colorectal-cancer screening test combining DNA-marker analysis with a fecal blood test; a positive result requires diagnostic colonoscopy.

\medterm{Colon} The main part of the large intestine between the small intestine and rectum. It absorbs water and salts and stores and moves stool through its ascending, transverse, descending, and sigmoid sections.

\medterm{Colon Cancer} Cancer arising in the lining of the colon, often developing gradually from a precancerous polyp. It may cause blood in stool, changed bowel habits, abdominal pain, or anemia, but early cancer may cause no symptoms.

\synonyms
Cancer, Colon
Colorectal Cancer (Colon Cancer)

\medterm{Colon Cancer Prevention} Colorectal-cancer risk can be reduced through recommended screening, physical activity, healthy weight, limiting alcohol, avoiding tobacco, and a diet rich in fiber-containing plant foods. Screening can detect precancerous polyps and early cancer before symptoms develop.

\medterm{Colon Cancer Screening} Testing intended to detect colorectal cancer or precancerous polyps before symptoms develop. Options include stool-based tests and direct visualization such as colonoscopy; the appropriate method and interval depend on age, risk, prior findings, and local guidelines.

\medterm{Colon Polyps} Growths arising from the inner lining of the colon; some types can develop into colorectal cancer.

\synonyms
Polyps, Colon
Polyps Colon (Colon Polyps)
Polyps Rectal (Colon Polyps)

\medterm{Colonic Diverticulitis (Diverticulosis)} Alternate index label for Diverticulosis.

\medterm{Colonic Inertia} Abnormally slow movement of stool through the colon, causing persistent constipation, bloating, abdominal discomfort, and infrequent bowel movements. It is diagnosed after other causes of constipation have been considered.

\medterm{Colonic Ischemia} Reduced blood flow to the colon causing tissue injury. It may produce sudden abdominal pain, bloody diarrhea, and tenderness and can result from low blood pressure, narrowed arteries, clots, medicines, or inflammation of blood vessels.

\medterm{Colonic Stenosis} Narrowing of part of the large intestine that makes it harder for stool and gas to pass. Causes include scar tissue, inflammation, cancer, or reduced blood supply; symptoms may include pain, bloating, constipation, vomiting, or bowel blockage.

\medterm{Colonization} Presence and multiplication of microorganisms on or in a body site without tissue
invasion or clinical disease. Colonization may precede infection but does not by itself
establish that symptoms are caused by the organism.

\medterm{Colonoscopy} An examination of the rectum and colon using a flexible camera. It can find inflammation, bleeding, tumors, and polyps and allows tissue sampling or removal of some polyps during the same procedure.


\medterm{Colony-Forming Unit} A laboratory estimate of the number of living microorganisms in a sample, based on the colonies that grow in culture. One colony may come from one organism or from a group of organisms.

\medterm{Color Blindness} Reduced or absent ability to distinguish certain colors under normal lighting. It is often
inherited, especially through X-linked variants, but may also result from eye disease, neurologic disease, aging, or
medicines.

\synonyms
Deficient Color Vision
Poor Color Vision

\medterm{Color Doppler} An ultrasound technique that uses colors to show the direction and relative speed of blood flow over a regular gray-scale image. It helps assess vessels, heart valves, and circulation.

\medterm{Color Vision Deficiency} Reduced ability to distinguish some colors, commonly red and green. It is often inherited because of differences in the eye’s color-sensing cells, but it can also develop from eye disease, optic-nerve or brain disorders, medicines, or toxins. The type and severity vary.

\medterm{Color Vision Test} An eye test that assesses the ability to distinguish colors or identify color patterns. It helps detect inherited color-vision deficiency and acquired abnormalities caused by eye, optic-nerve, brain, medicine, or toxic disorders.

\medterm{Colorado Tick Fever} An acute viral infection transmitted by ticks and characterized by fever, headache,
myalgia, and sometimes leukopenia or rash.

\medterm{Colorectal Adenocarcinoma} The most common cancer of the colon and rectum, arising from gland-forming cells in the bowel lining. It may cause bleeding, anemia, altered bowel habits, pain, or no early symptoms and can spread to other organs.

\medterm{Colorectal Adenoma} A usually noncancerous gland-forming growth in the colon or rectum that may gradually develop into colorectal cancer.

\medterm{Colorectal Cancer} A malignant tumor of the colon or rectum. It may cause bleeding, altered bowel habits, anemia, abdominal pain, or no early symptoms;
evaluation and treatment depend on the tumor's location and stage.

\medterm{Colorectal Cancer (Colon Cancer)} Alternate index label for Colon Cancer.


\medterm{Colorectal Cancer Genetics} The study of inherited variants that increase the risk of colorectal cancer.
Important inherited syndromes include Lynch syndrome, involving DNA mismatch-repair genes, and familial adenomatous
polyposis, involving the APC gene.

\medterm{Colorectal Cancer Screening} Public health strategy to detect colorectal cancer and precancerous polyps in
asymptomatic individuals. Current guidelines recommend screening beginning at age
forty-five.

\medterm{Colorectal Cancer Screening Tests} Tests used to detect colorectal cancer or precancerous polyps in people
without symptoms. They include stool-based tests and examinations that visualize the colon, such as colonoscopy, CT
colonography, and flexible sigmoidoscopy. The appropriate test and interval depend on age, risk, and local guidance.

\medterm{Colorectal Polyps} Growths arising from the lining of the colon or rectum. They may be non-neoplastic, adenomatous, or serrated; some can progress to colorectal cancer, so removal and microscopic examination guide surveillance.

\medterm{Colorectal Surgery} Surgical treatment of diseases of the colon, rectum, and anus. Procedures may include
segmental bowel resection, restorative reconstruction, stoma formation, or local
procedures, depending on the disease and its extent.

\medterm{Colostomy} A surgically created opening that brings part of the colon through the abdominal wall to divert stool into an external appliance.

\medterm{Colostrum} The thick, yellowish fluid produced by the mammary glands in late pregnancy and the
first few days postpartum. Rich in immunoglobulins (especially IgA), proteins,
leukocytes, and growth factors.

\medterm{Colpitis} Inflammation of the vagina caused by infection, altered vaginal bacteria or yeast, irritation, allergy, or hormonal changes. Symptoms may include discharge, itching, burning, pain, odor, or painful intercourse.
\medterm{Colporrhaphy} A surgical repair of vaginal wall defects performed to correct pelvic organ prolapse. Anterior colporrhaphy plicates and reinforces the attenuated pubocervical fascia to treat a cystocele (bladder prolapse), while posterior colporrhaphy plicates the rectovaginal fascia and perineal body to repair a rectocele (rectal prolapse) and restore normal vaginal caliber.

\medterm{Colposcope} A lighted magnifying instrument used to examine the cervix, vagina, or vulva closely and guide biopsy of suspicious areas.

\medterm{Colposcopy} Magnified visualization of the cervix, vagina, and vulva after application of acetic
acid. Indication: abnormal Pap smear (ASCUS with high-risk HPV, LSIL, HSIL, AGC). Acetic
acid causes acetowhite lesions in dysplastic epithelium. Lugol iodine (Schiller test):
normal glycogenated epithelium stains brown; abnormal epithelium does not.

\medterm{Colposcopy - Directed Biopsy} Removal of a small tissue sample from an unusual area of the cervix, vagina, or vulva during a magnified examination. Laboratory testing checks for infection, precancerous changes, or cancer.

\medterm{Columnar Epithelium Lined Lower Esophagus} Alternate index label; see \textit{Barrett's esophagus}.

\medterm{Coma} A state of profound unresponsiveness in which a person cannot be awakened and does not show purposeful
responses. It is a medical emergency with many possible causes.

\synonyms
Severe Brain Injury, Coma

\medterm{Coma Myxedema (Myxedema Coma)} Alternate index label for Myxedema Coma.

\medterm{Coma Scale} A standardized tool for describing level of consciousness using eye opening, verbal response, and motor
response.

\medterm{Combination Inhalers} Inhalers containing two or more medicines, such as an inhaled corticosteroid with a long-acting bronchodilator. They may simplify regular treatment for asthma or chronic obstructive pulmonary disease; the prescribed device and schedule must be followed.

\medterm{Combination Therapy} Use of two or more treatments together to address a disease or mechanism. The benefits and risks depend on the condition and
the specific combination.

\medterm{Combined Hormone Therapy} Hormone therapy containing both estrogen and a progestogen, generally used when endometrial protection is needed in a person with a uterus.

\synonyms
Estrogen-progestogen therapy; combined menopausal hormone therapy.

\medterm{Combined Hyperlipidemia} A blood-fat disorder with high cholesterol and triglycerides, often increasing harmful lipoproteins and raising the long-term risk of artery disease.

\synonyms
Mixed hyperlipidemia; combined hyperlipoproteinemia.

\medterm{Combined Spinal-Epidural Analgesia} A form of regional pain relief that gives medicine into the spinal fluid for rapid relief and places a catheter outside it for continued dosing. It is often used during labor or surgery and requires monitoring.

\medterm{Comedocarcinoma} A form of ductal carcinoma in situ of the breast characterized by central necrosis
within malignant ducts, which may produce a comedo-like appearance.

\medterm{Comedones} Plugs of sebum and keratin within hair follicles, appearing as open blackheads or closed whiteheads and representing a primary lesion of acne.

\medterm{Comfort Measures Only} A care plan focused on relief of pain, breathlessness, distress, and other symptoms rather than treatments intended to prolong life.

\medterm{Commando Procedure} Extensive heart surgery for severe infection destroying both the aortic and mitral valve areas and the tissue between them. The surgeon removes infected tissue, rebuilds the damaged area, and replaces both valves when necessary.

\synonyms
Aortomitral Fibrous Body Reconstruction; Intervalvular Fibrosa Reconstruction

\medterm{Commensal} Describing an organism that lives with another organism and benefits from the relationship without normally harming or helping the host; normal skin and intestinal flora may be commensal.

\medterm{Comminuted} Describing a fracture in which one bone breaks into several pieces. These fractures may be caused by high-energy injury and can be unstable or involve nearby tissues, so treatment depends on the bone, alignment, skin, nerves, and blood vessels.

\medterm{Comminuted Fracture} A broken bone divided into several pieces, usually by a high-energy injury. It may be unstable and damage nearby skin, muscles, nerves, or blood vessels, so treatment depends on the site and severity.

\medterm{Commissurotomy} Surgical incision or separation of a commissure, especially a fused cardiac valve
commissure, to relieve obstruction.

\medterm{Commode} A portable toilet or chair placed near a bed for a person who cannot easily reach a bathroom.

\medterm{Common Bile Duct} The tube formed when ducts from the liver and gallbladder join. It carries bile into the first part of the small intestine, where bile helps digest dietary fat.

\medterm{Common Cold} A short-term viral infection of the nose and throat causing runny or blocked nose, sneezing, sore throat, cough, and sometimes mild fever. Symptoms usually improve within one to two weeks.

\synonyms
Cold, Common
Cold (Common Cold)

\medterm{Common Cold In Babies} A viral infection of an infant's upper respiratory tract that commonly causes a runny nose, congestion, and cough.

\medterm{Common Cold Stages And Timeline Of Symptoms} The common cold often begins with sore throat, sneezing, or fatigue, followed by nasal congestion, runny nose, and cough. Symptoms usually peak within several days and improve over one to two weeks, although cough and fatigue may last longer.

\medterm{Common Early Symptoms Of Myasthenia Gravis} Early myasthenia gravis commonly causes fluctuating, fatigable weakness of the eyelids, eye movements, facial muscles, chewing, swallowing, speech, neck, or limbs. Symptoms worsen with use and improve with rest; breathing or swallowing difficulty is an emergency.

\medterm{Common Iliac Arteries} Terminal abdominal aortic branches at L4, each approximately 4 cm long, dividing into
internal and external iliac arteries. They supply the pelvis and lower extremities. The
right common iliac artery crosses anterior to the right common iliac vein.

\medterm{Common Medical Abbreviations And Terms} Medical abbreviations and terms summarize diagnoses, tests, medicines, anatomy, and treatment instructions. Because abbreviations can be ambiguous or unsafe, patients should ask clinicians or pharmacists to explain unfamiliar wording rather than infer its meaning.

\synonyms
Terms Medical Abbreviations (Common Medical Abbreviations And Terms)

\medterm{Common Migraine} A migraine attack without a preceding aura. It usually causes moderate or severe throbbing pain lasting four to 72 hours, often with nausea or sensitivity to light or sound and worsening with routine activity.

\synonyms
Migraine without aura; migraine without neurologic aura.

\medterm{Common Peroneal Nerve} A major nerve running from the back of the thigh around the outer knee into the lower leg and foot. It controls lifting and turning the foot and provides sensation to much of the lower leg and foot.

\medterm{Common Peroneal Nerve Dysfunction} Impairment of the common peroneal nerve, often from compression near the fibular head, trauma, surgery, or systemic nerve disease. It may cause foot drop, weakness of ankle eversion or dorsiflexion, and sensory loss over the lateral leg or dorsum of the foot.

\medterm{Common Peroneal Neuropathy} A peripheral nerve disorder causing weakness of ankle dorsiflexion and eversion, sensory
loss over the anterolateral leg and dorsum of the foot, and possible foot drop.

\medterm{Common Symptoms During Pregnancy} Frequent pregnancy-related symptoms such as nausea, fatigue, breast tenderness, heartburn, constipation, and urinary frequency; severe, sudden, or unusual symptoms may signal complications.

\medterm{Common Variable Immunodeficiency} An immune disorder in which the body makes too few effective antibodies. It causes repeated sinus, ear, or lung infections and may also cause digestive problems, autoimmune disease, or enlarged lymph nodes.

\medterm{Common Warts (Verruca Vulgaris)} Benign, rough-surfaced skin growths caused by infection with human papillomavirus, commonly affecting the hands and fingers.

\synonyms
Verruca Vulgaris

\medterm{Commotio Cordis} Sudden cardiac arrest caused by a blunt, nonpenetrating impact to the chest at a
vulnerable moment in the cardiac cycle, without structural heart injury.

\medterm{Communicable Disease} Any disease caused by bacteria, viruses, or other pathogens that is spread from person
to person.

\medterm{Communicating Hydrocephalus} A type of hydrocephalus in which there is free flow of cerebrospinal fluid between
the ventricles, but impaired absorption of CSF at the arachnoid granulations or
increased CSF production. The ventricles are uniformly dilated, including the fourth
ventricle.

\medterm{Communicating With Patients} The clinical exchange of information, concerns, preferences, and decisions between healthcare professionals and patients, using clear language, listening, empathy, and confirmation of understanding.

\medterm{Communicating With Someone With Aphasia} Communication strategies for a person with impaired language from brain injury or disease, including short sentences, extra time, visual cues, and supported conversation rather than assuming impaired intelligence.

\medterm{Communicating With Someone With Dysarthria} Communication strategies for a person whose speech muscles are weak or poorly coordinated, including reducing background noise, confirming messages, and using writing or augmentative aids when needed.

\medterm{Community Acquired Infection} An infection present or incubating at hospital admission, or developing within 48 hours
of admission without recent healthcare exposure. Common examples include pneumonia,
urinary tract infection, and meningitis caused by organisms such as Streptococcus
pneumoniae and Escherichia coli, generally susceptible to narrow-spectrum antibiotics.

\medterm{Community Health Worker} A trained member of a community who helps connect people with health services, education, prevention, and social support.

\medterm{Community Rehabilitation} Therapy and support delivered at home, in clinics, or in community settings to improve movement, communication, independence, participation, and daily functioning after illness or injury.
\medterm{Community Reintegration} The process of returning to everyday life and taking part in work, education, family, social, and leisure activities after illness or injury. Support may include rehabilitation, changes to the home or workplace, transport help, and practical or emotional support.

\medterm{Community-Acquired} Acquired outside a healthcare facility or before healthcare exposure. For example, community-acquired pneumonia begins outside a hospital and may involve different organisms and resistance patterns from hospital-acquired disease.
\medterm{Community-Acquired Pneumonia} Pneumonia acquired outside a hospital or long-term-care facility. It may cause cough, fever, chest pain, breathlessness, and low oxygen; diagnosis and treatment depend on examination, imaging, severity, likely organisms, comorbidities, and local antibiotic guidance. CURB-65 or PSI can help assess the need for hospital care but does not replace clinical judgment.

\medterm{Community-Acquired Pneumonia In Adults} Pneumonia acquired outside a hospital or healthcare facility; assessment considers severity, oxygenation, comorbidities, likely pathogens, and the need for outpatient or inpatient treatment.

\medterm{Comorbidity} Presence of two or more chronic conditions in same individual. Common in elderly:
hypertension, diabetes, arthritis, COPD, depression. Increases complexity of care,
polypharmacy risk, functional decline. Comprehensive assessment needed for each
condition and interactions.

\medterm{Computational Biomedicine} An interdisciplinary field that uses computers, mathematical models, simulations, and biomedical data to study health and disease. It may support research in areas such as genomics, drug development, medical imaging, disease prediction, diagnosis, and personalized treatment.

\synonyms
Computational Biomedical Science

\medterm{Compact Bone} Dense outer bone tissue arranged in strong cylindrical units called osteons. It provides most of the skeleton’s strength and protects inner marrow while supporting movement and weight-bearing.

\synonyms
Cortical bone; dense bone.

\medterm{Comparison} A reference in an imaging report to earlier studies used to assess whether a finding has changed.

\medterm{Compartment Pressure} The pressure inside a closed group of muscles, nerves, and blood vessels. If it rises too much, blood flow can fall and cause permanent tissue damage; measurements support but do not replace examination.

\medterm{Compartment Syndrome} Dangerous pressure buildup inside a confined muscle group that reduces blood flow and can permanently damage muscle and nerves. Severe pain, swelling, numbness, or weakness after injury requires emergency assessment and often surgery.

\synonyms
Syndrome Compartment (Compartment Syndrome)

\medterm{Competence} A legal determination that a person can make a particular decision or manage a matter. In clinical care, decision-making capacity is assessed for each specific choice.

\medterm{Competitive Antagonist} A medicine or substance that occupies the same receptor as another substance but does not activate it. By competing for that site, it reduces the other substance's effect; increasing the other substance may overcome the competition.

\medterm{Competitive Inhibition} Reduced enzyme activity caused when an inhibitor competes with the normal substance for the enzyme's active site. Because the binding is usually reversible, a higher concentration of the normal substance may lessen the inhibition.

\medterm{Complement} Complement is a system of blood proteins that enhances antibody and immune-cell defense by marking, inflaming, or directly damaging targets.

\medterm{Complement C1Q Deficiency} A rare inherited lack of a complement protein needed for immune defense and removal of antibody-bound material. It can cause repeated infections and autoimmune illness resembling lupus, including rash, joint, kidney, or blood problems.

\medterm{Complement C2/C4 Deficiency} Inherited or acquired lack of complement C2 or C4, proteins needed especially for the classical complement pathway. It can increase susceptibility to infections with encapsulated bacteria and autoimmune disease resembling lupus. Testing may show low CH50 with low C2 or C4; care includes vaccination, prompt infection treatment, and specialist management of autoimmune disease or recurrent infections.

\medterm{Complement C3} A major blood protein in the complement immune system. When activated, it helps mark microbes for removal and promotes inflammation; severe deficiency can cause repeated serious infections.

\medterm{Complement Cascade} A chain reaction involving many immune proteins in blood and tissues. It marks microbes for destruction, attracts immune cells, increases inflammation, and can directly break holes in some microbes.

\medterm{Complement Component 3 (C3)} A blood test measuring complement component C3, a major protein of the complement immune system; abnormal levels can support evaluation of immune-complex disease, infection, or complement deficiency.

\medterm{Complement Component 4} A blood test measuring complement component C4, part of the classical complement pathway; low levels may occur with immune-complex disease, hereditary angioedema, or complement consumption.

\medterm{Complement Deficiencies} Rare primary immunodeficiencies caused by genetic defects in complement pathway
proteins, increasing susceptibility to specific infections and autoimmune diseases.

\medterm{Complement Deficiency} Reduced or absent activity of one or more complement proteins, increasing susceptibility to some infections or autoimmune disease.

\medterm{Complement Fixation Test For C Burnetii} A serologic test detecting complement-fixing antibodies to Coxiella burnetii, the bacterium that causes Q fever; antibody phase patterns and titers help distinguish recent from persistent infection.

\medterm{Complement System} A group of blood proteins that supports immune defense by destroying microbes and promoting inflammation and phagocytosis.

\medterm{Complementary Sequence} A DNA or RNA sequence whose building blocks pair with those of another sequence. This matching allows genetic strands to bind and is used in copying, testing, and regulating genetic information.

\medterm{Complete Blood Count} A blood test that measures red blood cells, white blood cells, platelets, hemoglobin, and hematocrit. It may also measure red-cell size and hemoglobin content and list the different types of white blood cells. It helps assess anemia, infection, inflammation, bleeding, and some bone-marrow disorders.

\medterm{Complete Breech} A fetal position in which both hips and knees are bent, so the baby is sitting cross-legged with the buttocks or feet closest to the birth canal. Delivery planning depends on pregnancy, fetal, and maternal factors.

\synonyms
Flexed Breech; Complete Breech Presentation

\medterm{Complete Cystectomy} Surgery to remove the entire urinary bladder, usually for invasive bladder cancer. Because urine can no longer leave through the bladder, the surgeon creates a urinary diversion such as an ileal conduit, a continent reservoir, or a new bladder from bowel.

\medterm{Complete Denture} A removable artificial set of teeth replacing all teeth in the upper jaw, lower jaw, or both. It rests on the gums and may improve chewing, speech, and appearance but can require adjustment for comfort and stability.
\medterm{Complete Miscarriage} Pregnancy loss in which all pregnancy tissue has passed from the uterus, usually followed by easing bleeding and cramps. A clinician may use examination, hormone testing, or ultrasound to confirm that no tissue remains and to exclude an ectopic pregnancy; heavy bleeding, severe pain, dizziness, or fever needs urgent care.

\synonyms
Complete Spontaneous Abortion

\medterm{Complete Response} No detectable evidence of cancer after treatment on the tests used. It does not prove that every cancer cell has been eliminated, so continued follow-up is needed.

\medterm{Complex Carbohydrates} Carbohydrates made up of longer chains of sugars, chiefly starches and dietary fiber,
found in foods such as whole grains, legumes, vegetables, and potatoes. Fiber is not
fully digested and supports bowel function, satiety, and metabolic health, while starch
is digested into glucose at a rate that depends on food structure and processing.

\medterm{Complex Febrile Seizure} A seizure triggered by fever in a young child that lasts more than 15 minutes, affects only part of the body, or happens more than once during the same illness. It needs medical assessment to exclude brain infection or another cause. Most children recover fully, but the risk of another seizure is higher than with a simple febrile seizure.

\synonyms
Atypical Febrile Seizure; Complicated Febrile Seizure

\medterm{Complex Regional Pain Syndrome} Persistent severe pain and sensitivity in a limb after injury, surgery, or sometimes without an obvious trigger. Swelling, skin-color or temperature changes, sweating, stiffness, and weakness may also occur.

\medterm{Compliance} The extent to which a person follows a treatment recommendation; adherence is generally preferred because it emphasizes an agreed plan.

\medterm{Complicated Bereavement} Alternate index label; see \textit{Complicated grief}.

\medterm{Complicated Grief} Persistent, intense grief that substantially interferes with daily functioning and lasts longer than expected for the person's
cultural and social context.

\synonyms
Complicated Bereavement
Traumatic Grief

\medterm{Complicated Migraine} An older, nonspecific term for migraine accompanied by prolonged or prominent neurologic symptoms; the specific syndrome should be identified when possible.

\synonyms
Migraine with prolonged neurologic symptoms; complex migraine, historical usage.

\medterm{Complicated Urinary Tract Infection} A urinary infection associated with functional or structural abnormality, obstruction,
catheterization, renal impairment, male anatomy, pregnancy, or another factor increasing
treatment failure or complication risk. Management requires culture-directed therapy and
correction of the underlying problem when possible.

\medterm{Complication} Any adverse event arising during or after a surgical procedure, graded by the
Clavien-Dindo classification from minor deviation (Grade I) through organ failure (Grade
IV) to death (Grade V). Common types include surgical site infection, hemorrhage, wound
dehiscence, anastomotic leak, ileus, and venous thromboembolism.

\medterm{Complications (Vascular)} Problems involving blood vessels after disease, injury, or a procedure, such as bleeding, hematoma, thrombosis, embolism, vessel injury, or reduced blood flow to a limb. Symptoms and urgency depend on the vessel and tissue affected.

\synonyms
Vascular

\medterm{Composite Resin} A tooth-colored material used to repair cavities, chipped teeth, or other defects. It is bonded to the tooth and shaped to match it, but may wear or stain sooner than some harder restorations.

\medterm{Compound Fracture (Open Fracture)} A fracture associated with a wound that communicates between the fracture site and the external environment. It carries increased risks of infection and tissue damage.

\synonyms
Open Fracture

\medterm{Compound Heterozygote} A person with two different disease-causing variants in the same gene, one in each copy. Together, the variants may cause an inherited disorder even though neither copy has the same change.

\medterm{Compound Presentation} A birth position in which a hand, arm, foot, or another limb lies beside the baby's head or buttocks. The limb may move back during labor, but the position can obstruct delivery or increase injury risk.

\medterm{Compounding} Preparation, mixing, or altering of a medicine to meet an individual patient's needs when an appropriate approved product is unavailable.

\medterm{Compounding Pharmacy} A pharmacy that prepares customized formulations for an individual patient when a suitable commercially manufactured product is unavailable or unsuitable.

\synonyms
Compounding pharmacy practice; extemporaneous pharmacy.

\medterm{Comprehensive Geriatric Assessment} A multidimensional assessment of an older person's medical conditions, medicines, function, cognition, mood, and social
support, used to develop an individualized care plan.

\medterm{Comprehensive Metabolic Panel} A group of blood tests measuring electrolytes, glucose, kidney function, liver-related chemicals, proteins, and calcium. It helps assess general health and monitor or investigate metabolic, liver, and kidney disorders.

\medterm{Compress} To press together or apply pressure, as with a bandage or device, or to reduce the volume of something.

\medterm{Compressed Nerve} Alternate index label; see \textit{Pinched nerve}.

\medterm{Compression} Pressure applied to tissue or a structure, which may reduce swelling or bleeding but can also impair nerves, blood vessels, or organs when excessive.

\medterm{Compression Fracture} Collapse of a vertebra or other bone because of compression, often related to osteoporosis or trauma.

\medterm{Compression Fractures Of The Back} Collapse of one or more vertebral bodies caused by weakened bone, often from osteoporosis, or by substantial trauma. It can cause sudden back pain, height loss, kyphosis, and rarely nerve or spinal-cord compression.

\medterm{Compression Stockings} Elastic garments that apply graduated pressure to the legs, supporting venous return and reducing swelling or selected risks of venous pooling; correct fit and arterial circulation matter.

\medterm{Compression Therapy} Use of external pressure, usually from garments or bandages, to support venous or lymphatic flow and reduce swelling in selected
conditions.

\medterm{Compression Ultrasound} An ultrasound test in which gentle pressure is applied to a vein. A healthy vein usually collapses; failure to collapse may indicate a blood clot, especially in the deep veins of a limb.

\medterm{Compressive Atelectasis} Collapse of part of a lung because fluid, air, a tumor, or another process presses on it from outside. The lung may re-expand when the pressure is relieved.

\medterm{Compton Scattering} An interaction in which a photon transfers part of its energy to an electron and changes direction. It contributes to scattered
radiation in imaging.

\medterm{Compulsions} Repetitive behaviors or mental acts that the individual feels driven to perform in
response to an obsession or according to rigid rules. They are aimed at preventing or
reducing anxiety or preventing some dreaded event, but are not connected in a realistic
way with what they are designed to neutralize.

\medterm{Compulsive Gambling} Alternate, nonpreferred label; see \textit{Gambling Disorder}.

\medterm{Compulsive Obsessive Disorder (Obsessive Compulsive Disorder)} Alternate index label for Obsessive Compulsive Disorder.

\medterm{Compulsive Overeating} Alternate index label; see \textit{Binge-eating disorder}.

\medterm{Compulsive Sexual Behavior} Recurrent sexual thoughts or behaviors that are difficult to control and cause significant distress or impairment.

\synonyms
Sexual Obsession

\medterm{Compulsive Stealing} Alternate index label; see \textit{Kleptomania}.

\medterm{Computed Tomography} An imaging test that uses rotating X-rays and computer processing to create detailed cross-sectional pictures. It is fast and can show bones, bleeding, organs, blood vessels, tumors, infection, and internal injury.

\medterm{Computed Tomography Pulmonary Angiography} A CT scan performed after injected contrast highlights the lung arteries. It is commonly used to look for a pulmonary embolism and may also show strain on the heart or other chest abnormalities.

\synonyms
CTPA; CT Pulmonary Angiogram; Pulmonary CT Angiography

\medterm{Computed Tomography Scan} An X-ray scan processed by a computer to produce cross-sectional images of the body. It can rapidly detect bleeding, fractures, infection, blocked bowel, blood clots, tumors, and many other urgent conditions.

\medterm{Computerized Axial Tomography} Computed tomography, an imaging method that uses X-rays and computer processing to produce cross-sectional images of the body; it is also called a CT or CAT scan.

\medterm{Concato'S Disease} A historical term for primary polycythemia, now generally called polycythemia vera, a
myeloproliferative neoplasm causing increased red-cell mass and often elevated blood
viscosity.

\medterm{Concealed Conduction} A heart rhythm event in which an electrical impulse enters part of the heart's conduction system but stops before producing a visible heartbeat. The affected tissue may then delay or block the next impulse and alter the heartbeat pattern.

\medterm{Concealed Pregnancy} Pregnancy that is hidden or not recognized until delivery, often by denial (psychogenic)
or deliberate concealment; a medicolegal concept relevant to infanticide.

\medterm{Concentration} The amount of a substance in a given volume, such as the level of glucose or medicine in blood. It is reported in units such as milligrams per deciliter or millimoles per liter.

\medterm{Conception} The beginning of pregnancy, commonly referring to fertilization followed by implantation and establishment of a developing embryo in the uterus.

\medterm{Conception (Fertilization)} The joining of a sperm and egg, usually in a fallopian tube, to form the first cell of an embryo. The resulting cells divide as they travel toward the uterus, where implantation may occur several days later.

\synonyms
Fertilization

\medterm{Conchae} Curved bony structures covered by mucosa inside the nasal cavity that warm, humidify, and filter inhaled air.

\medterm{Conclusion} The part of a report that summarizes the main findings and their clinical significance or recommended next steps.

\medterm{Concomitant} Existing or occurring at the same time as another condition, treatment, or event. Concomitant medicines or illnesses may change diagnosis, treatment choices, expected effects, side effects, or monitoring needs.

\medterm{Concordance} Probability both twins have same trait. Monozygotic (identical) twins > dizygotic
(fraternal) for genetic traits. High concordance suggests genetic influence.

\medterm{Concrescence} Union of adjacent teeth by cementum only. Usually posterior teeth.

\medterm{Concretio Cordis} Abnormal adhesion or fusion of structures within or around the heart, including adhesion
of the pericardial surfaces or fusion of cardiac components.

\medterm{Concussion} A form of mild traumatic brain injury caused by a biomechanical force (direct blow to
the head, face, neck, or elsewhere on the body with impulsive force transmitted to the
head), resulting in a transient disturbance of brain function without structural damage
visible on conventional neuroimaging.





\medterm{Condition} A general term for a state of bodily health or, in medical usage, a specific deviation
from normal in structure or function that may or may not constitute a disease, including
complaints lacking a discrete pathologic basis and overt diagnoses.

\medterm{Conditioning} Learning in which repeated association or reinforcement changes a behavior or response; in medicine, conditioning can also mean preparing a patient or tissue for a treatment.

\medterm{Conditions Caused By Deletion Mutations} Deletion mutations remove DNA segments and can disrupt one gene or several neighboring genes. The clinical effect depends on size, location, inheritance, and whether the remaining copy compensates; genetic counseling and testing may clarify diagnosis and recurrence risk.

\medterm{Condom} A thin barrier sheath worn on the penis or inserted into the vagina that helps prevent pregnancy and reduces transmission of sexually transmitted infections when used correctly.

\medterm{Condoms - Male} A sheath placed over the penis during sexual activity to reduce exchange of semen and genital secretions and lower pregnancy and sexually transmitted infection risk when used correctly.

\medterm{Conduct Disorder} A childhood or adolescent disorder involving a persistent pattern of behavior that violates
the rights of others or major age-appropriate rules. Examples include aggression,
property destruction, deceitfulness, theft, and serious rule violations.

\medterm{Conduction Deafness} Hearing loss caused by impaired transmission of sound through the external or middle
ear, with the inner ear and neural auditory pathways relatively preserved.

\medterm{Conductive Deafness} Hearing loss resulting from a problem in the external auditory canal, tympanic membrane,
or middle ear that reduces conduction of sound to the inner ear.

\medterm{Conductive Hearing Loss} Hearing loss caused by a problem in the outer or middle ear that prevents sound from reaching the inner ear normally. Earwax, fluid, a damaged eardrum, bone changes, or ear-bone disease may be responsible.

\medterm{Condylar Process} The rounded articular projection of a bone, especially the mandibular process that
articulates with the temporal bone to form the temporomandibular joint.

\medterm{Condyle} A rounded projection at the end of a bone that forms part of a joint and helps guide movement against another bone.

\medterm{Condyloma} A wart-like growth, often in the genital or anal area and commonly caused by human papillomavirus; condyloma lata are lesions of secondary syphilis.

\medterm{Condyloma Acuminata} Soft, wart-like growths around the genitals or anus caused by certain types of human papillomavirus. They spread through close skin-to-skin sexual contact; treatment removes visible warts but may not remove the virus, and vaccination helps prevent many infections.

\medterm{Condyloma Acuminatum} A soft, wart-like growth around the genitals or anus caused by certain human papillomavirus types. It spreads through intimate skin contact; treatment removes visible warts, but the virus may remain and recur.

\medterm{Condyloma Latum} A broad, flat, moist, highly infectious lesion of secondary syphilis, usually around the genitals or anus. It differs from a genital wart and indicates syphilis requiring testing and antibiotic treatment.
\medterm{Condylomata Acuminata} Alternate index label; see \textit{Genital warts}.

\medterm{Condylomata Lata} Broad, flat, moist, grayish-white papules occurring in secondary syphilis, typically
found in the anogenital region and other intertriginous areas. They are highly
infectious and contain abundant Treponema pallidum organisms. Diagnosis is by darkfield
microscopy or serological testing for syphilis.

\medterm{Cone Biopsy} A surgical procedure that removes a cone-shaped portion of the cervix for diagnosis or treatment of high-grade cervical precancer and to determine whether invasive cancer is present.

\medterm{Cone-Beam Computed Tomography} A dental or facial scan that uses a cone-shaped X-ray beam to create three-dimensional images of teeth, jawbone, and nearby structures. It should be used when the extra detail will affect care.

\medterm{Cones} Light-sensitive retinal cells responsible for color vision and sharp detail, working best in daylight or other relatively bright conditions.

\synonyms
Cone photoreceptors; cone cells.

\medterm{Confabulation} Unintentional production of inaccurate memories or stories without an attempt to deceive. It may occur with memory disorders or injury affecting the front part of the brain.

\medterm{Confidentiality} The ethical and legal duty of healthcare professionals to protect patient information
from unauthorized disclosure. Founded on the principle of respect for patient autonomy
and trust, confidentiality is a core requirement of medical ethics (GMC guidelines) and
law (GDPR, Data Protection Act, Caldicott principles).

\medterm{Confirmatory Test} A test used to verify a preliminary or screening result, often with a different method
or greater specificity. Confirmation should be selected according to the condition,
specimen, timing, and accepted clinical guidelines.

\medterm{Conformal Radiation Therapy} Radiation treatment planned so that the high-dose region conforms closely to the tumor's three-dimensional shape while limiting exposure of nearby normal tissue.

\medterm{Confounding} A distortion that occurs when another factor is related to both an exposure and an outcome, making their relationship appear stronger, weaker, or different from reality. For example, age may affect both an exposure and disease risk. Study design and statistical adjustment can reduce confounding.

\medterm{Confusion} Impaired ability to think clearly, orient, attend, remember, or respond appropriately; sudden confusion is delirium and may indicate infection, medication effect, metabolic illness, stroke, or another emergency.

\medterm{Confusion (Delirium)} A sudden change in attention and awareness that fluctuates during the day. Infection, medicines, pain, dehydration, metabolic problems, constipation, or urinary retention may cause it and require prompt assessment.

\synonyms
Delirium

\medterm{Confusion Assessment Method} A bedside assessment for delirium. It checks for a sudden or fluctuating change, poor attention, disorganized thinking, and an altered level of alertness; the pattern helps identify delirium but does not replace finding its cause.

\medterm{Congenital} Present at birth, whether inherited, caused by an early developmental difference, or acquired before birth through infection, exposure, or another condition.

\medterm{Congenital Adrenal Hyperplasia} A group of inherited disorders in which the adrenal glands cannot make hormones normally. The most common form can cause salt loss, excess androgen effects, early puberty, or genital differences and may be life-threatening in newborns.

\medterm{Congenital Antithrombin III Deficiency} An inherited deficiency or dysfunction of antithrombin, a natural anticoagulant, causing increased risk of venous thromboembolism.

\medterm{Congenital Biliary Atresia} Alternate index label; see \textit{Biliary atresia}.

\medterm{Congenital Cataract} Clouding of the eye lens present at birth or developing soon afterward. It can interfere with visual development and may require early eye assessment and treatment to prevent lasting reduced vision.

\medterm{Congenital CMV} Cytomegalovirus infection acquired before birth. Some babies have small head size, hearing or vision loss, seizures, jaundice, an enlarged liver or spleen, or developmental problems; others appear well initially but develop later hearing loss.

\medterm{Congenital CNS Malformation} A structural abnormality of the brain, spinal cord, or related developmental structures present from birth, caused by disturbed neural development or other prenatal factors; severity ranges from incidental to life-threatening.

\medterm{Congenital Contractural Arachnodactyly} An inherited connective-tissue disorder causing unusually long, slender fingers and toes, tight joints present from birth, thin muscles, and sometimes a curved spine. Joint stiffness may improve with age.

\medterm{Congenital Cytomegalovirus} Cytomegalovirus infection present at birth after transmission during pregnancy. Some infants have hearing loss, vision problems, small head size, seizures, or developmental disability; others have no symptoms initially.

\medterm{Congenital Cytomegalovirus Infection} Cytomegalovirus infection acquired before birth. It may cause hearing loss,
developmental problems, microcephaly, or no symptoms.

\medterm{Congenital Diaphragmatic Hernia} A defect in the diaphragm allowing abdominal organs to herniate into the thoracic
cavity, causing pulmonary hypoplasia and respiratory distress at birth. Most commonly a
posterolateral (Bochdalek) defect on the left side. Presents with cyanosis, scaphoid
abdomen, and bowel sounds in the chest.

\medterm{Congenital Diaphragmatic Hernia Repair} Congenital diaphragmatic hernia repair is a procedure for congenital diaphragmatic hernia repair; indication, preparation, risks, technical details, and recovery depend on the patient and treatment goal.
\medterm{Congenital Disorder} A medical condition or structural abnormality present in an individual from birth, which
may be caused by genetic factors or environmental influences in the womb.

\medterm{Congenital Fibrinogen Deficiency} An inherited absence or abnormality of fibrinogen that can cause bleeding, thrombosis, or both depending on the specific molecular defect.

\medterm{Congenital Glaucoma} A rare form of glaucoma (1 in 10,000-20,000 births) caused by developmental anomalies of
the trabecular meshwork and anterior chamber angle, impairing aqueous humor outflow.
Most cases are sporadic; familial forms are autosomal recessive (CYP1B1 gene mutations).

\medterm{Congenital Heart Defect - Corrective Surgery} Surgery or catheter-based intervention to repair a structural heart abnormality present at birth, with the specific procedure depending on the defect, its severity, and the patient’s anatomy and condition.

\medterm{Congenital Heart Defects} Problems in the heart or major blood vessels
that are present at birth. They can be mild or serious and may affect how blood
flows through the heart and body. Examples include holes in the heart, abnormal
valves, and problems with the major blood vessels.

\medterm{Congenital Heart Defects In Children} Alternate index label; see \textit{Congenital Heart Defects}.

\medterm{Congenital Heart Disease} A heart or major blood-vessel problem present from birth. It may change blood flow, cause no symptoms, or lead to breathing difficulty, blue color, poor growth, heart failure, or abnormal rhythms. Treatment ranges from monitoring and medicines to catheter procedures or surgery.

\medterm{Congenital Heart Disease In Adults} A heart defect present from birth that continues into adulthood; see \textit{Congenital Heart Defects}.

\medterm{Congenital Heart Murmur} A heart murmur heard in a newborn or infant. It may be innocent or may indicate a
structural heart defect; assessment considers the murmur's characteristics and associated symptoms or examination
findings.

\medterm{Compensated Heart Failure} Heart failure in which symptoms and fluid buildup
are stable or controlled with treatment. The heart failure is still present and can
worsen, especially if a new illness or other problem develops.

\medterm{Congenital Hip Dislocation} Abnormal development or displacement of the hip joint present at birth or developing in
infancy, now generally called developmental dysplasia of the hip.

\medterm{Congenital Hyperinsulinism} An inherited or developmental disorder in which a baby's pancreas releases too much insulin. This can cause dangerously low blood sugar, poor feeding, seizures, brain injury, or death unless recognized and treated promptly.

\medterm{Congenital Hyperthyroidism} Excess thyroid hormone in a newborn, most often caused by thyroid-stimulating antibodies passing from a mother with Graves disease. It may cause a fast heartbeat, irritability, poor feeding, poor weight gain, or an enlarged thyroid.

\medterm{Congenital Hypothyroidism} Too little thyroid hormone from birth, usually because the thyroid is missing, misplaced, or underdeveloped. Newborn screening and early treatment are essential because untreated deficiency can impair growth and brain development.

\medterm{Congenital Malformation} A structural abnormality that develops before birth and is present at birth, caused by genetic, environmental, developmental, or unknown factors.

\medterm{Congenital Megacolon} A birth condition in which part of the intestine lacks the nerve cells needed to move stool. It causes blockage and enlargement of the bowel above it, with delayed first stool, swelling, vomiting, or severe constipation.

\medterm{Congenital Megacolon (Hirschsprung Disease)} Alternate index label for Hirschsprung Disease.

\medterm{Congenital Melanocytic Nevus} A dark mole present at birth or appearing soon afterward, caused by a collection of pigment-producing cells. Most remain harmless, but larger lesions have a higher melanoma risk and need specialist monitoring.

\medterm{Congenital Mitral Valve Anomalies} Structural abnormalities of the mitral valve that are present at birth and may impair blood flow.

\medterm{Congenital Muscular Dystrophy} A group of inherited muscle disorders causing weakness or low muscle tone at birth or during early infancy. Some types also affect the brain, eyes, breathing, heart, or joints; severity varies widely.

\synonyms
CMD; Congenital Dystrophy

\medterm{Congenital Myasthenic Syndromes} A group of inherited disorders that disrupt communication between nerves and muscles. They cause weakness that worsens with activity and may affect the eyes, swallowing, breathing, or limb movement.
\medterm{Congenital Nephrotic Syndrome} Severe kidney disease beginning before birth or during the first three months of life, causing large amounts of protein loss in urine, swelling, low blood protein, poor growth, infections, and increased clotting risk.

\medterm{Congenital Nevus} A melanocytic mole present at birth or developing soon afterward. It is composed of melanocytes in
the epidermis and dermis. Size, location, and changes over time determine the need for
clinical assessment and surveillance.

\medterm{Congenital Penile Curvature} Curvature of the penis present from birth, usually caused by unequal development of the erectile tissues. Mild curvature may need no treatment; significant curvature causing functional difficulty may require specialist assessment.

\medterm{Congenital Plagiocephaly} Alternate index label; see \textit{Craniosynostosis}.

\medterm{Congenital Platelet Function Defects} Inherited disorders in which platelets do not work normally even though their number may be normal. They can cause easy bruising, nosebleeds, heavy menstrual bleeding, or excessive bleeding after injury or surgery.

\medterm{Congenital Pneumonia} A lung infection present at birth, acquired before or during delivery. A newborn may have rapid or difficult breathing, temperature instability, poor feeding, or sepsis and needs urgent evaluation and treatment.

\medterm{Congenital Protein C Or S Deficiency} An inherited lack or abnormality of protein C or protein S, natural substances that limit clotting. It increases the risk of blood clots in veins, sometimes beginning in infancy or childhood.
\medterm{Congenital Rubella} Fetal infection with rubella virus acquired during pregnancy, potentially causing hearing loss, cataracts, congenital heart disease, developmental impairment, and other findings.

\medterm{Congenital Rubella Syndrome} Birth defects caused by rubella infection during pregnancy, especially early pregnancy. Affected babies may have hearing loss, cataracts, heart defects, small head size, developmental problems, enlarged organs, or a low platelet count.

\medterm{Congenital Scoliosis} A lateral curvature of the spine caused by congenital vertebral anomalies that arise
during embryogenesis (failure of formation or segmentation). Failure of formation:
hemivertebra (wedge-shaped vertebra causing progressive curvature). Failure of
segmentation: unilateral unsegmented bar (block of bone on one side).

\medterm{Congenital Syphilis} Treponema pallidum infection acquired transplacentally. Can cause stillbirth,
prematurity, hepatosplenomegaly, rash, snuffles, osteochondritis, and Hutchinson teeth.
Diagnosed by darkfield microscopy, VDRL/RPR in neonatal serum, or PCR.

\medterm{Congenital Toxoplasmosis} Toxoplasma gondii infection acquired in utero, potentially causing chorioretinitis, intracranial calcifications, hydrocephalus, seizures, or developmental impairment.

\medterm{Congenital Zika Syndrome} Fetal brain infection with Zika virus, causing microcephaly, brain abnormalities, eye
defects, hearing loss, and limb contractures. Transmission occurs via mosquito bite or
sexual contact. Diagnosis by Zika virus PCR in maternal serum or amniotic fluid. No
specific treatment. Prevention by avoiding endemic areas during pregnancy.

\medterm{Congestion} A passive vascular phenomenon of engorgement of a tissue or organ with blood from
impaired venous or lymphatic drainage. Chronic venous congestion produces stasis,
hypoxia, dilated vessels, edema, hemorrhage, and hemosiderin-laden macrophages,
exemplified by the lungs (heart failure cells) and liver (nutmeg liver).

\medterm{Congestive Heart Failure} A historical term for heart failure with clinical fluid congestion. Current terminology
is heart failure, which may occur with or without congestion and is classified by
left-ventricular ejection fraction, symptoms, and disease stage.

\medterm{Congestive Heart Failure Overview} Alternate index label; see \textit{Heart failure}.
\medterm{Conization} Surgical removal of a cone-shaped portion of the cervix for diagnosis or treatment of
cervical dysplasia or selected early cervical cancers.

\medterm{Conjoined Twins} Monozygotic, monochorionic, monoamniotic twins who are physically connected at birth,
resulting from incomplete fission of the fertilized ovum after day 13 of fertilization.
Incidence: 1 in 50,000-100,000 births.

\medterm{Conjugated Equine Estrogens} A mixture of estrogenic compounds historically derived from the urine of pregnant mares and used in some hormone-therapy formulations.

\synonyms
CEE; conjugated estrogens.

\medterm{Conjunctiva} A thin mucous membrane that lines the inner eyelids (palpebral conjunctiva) and covers the front surface of the sclera (bulbar conjunctiva). It protects and lubricates the eye and permits smooth movement of the eyelids.

\medterm{Conjunctival} Relating to the conjunctiva, the thin mucous membrane covering the front of the eye and lining the inner eyelids.

\medterm{Conjunctival Squamous-Cell Carcinoma} A malignant tumor of the conjunctival epithelium, typically arising from the limbal
area. It is associated with UV radiation exposure and can invade locally if not treated
promptly.

\medterm{Conjunctivitis} Inflammation of the conjunctiva, the thin, transparent mucous membrane lining the inner
surface of the eyelids and the anterior surface of the sclera. Conjunctivitis is the
most common cause of red eye and is classified by etiology.

\medterm{Conjunctivitis (Pink Eye)} Alternate index label for Pink Eye.

\medterm{Conjunctivitis Or Pink Eye} Inflammation of the conjunctiva causing redness, discharge, irritation, itching, or tearing. Causes include viral or bacterial infection, allergy, and chemical irritation; severe pain, photophobia, vision loss, or contact-lens use requires prompt eye assessment.

\medterm{Conjunctivoplasty} Plastic or reconstructive surgery of the conjunctiva, performed to repair defects,
scarring, adhesions, or other abnormalities of the ocular surface.

\medterm{Conn'S Syndrome} Primary hyperaldosteronism caused by autonomous secretion of aldosterone, leading to
hypertension and often low serum potassium.

\medterm{Connective Tissue} A diverse type of tissue that supports, binds, protects, and insulates organs and
structures throughout the body.

\medterm{Connective Tissue Disease} A disorder affecting tissues that support and connect organs, often because the immune system attacks the body’s own tissues. It may involve skin, joints, muscles, blood vessels, lungs, kidneys, or other organs. Examples include lupus, systemic sclerosis, rheumatoid arthritis, and Sjögren syndrome.

\medterm{Connective Tissue Disease-Related Interstitial Lung Disease} Interstitial lung disease associated with autoimmune connective-tissue disease such as systemic sclerosis, rheumatoid arthritis, polymyositis, dermatomyositis, or Sjogren syndrome; inflammation and fibrosis can impair gas exchange.

\medterm{Consanguinity} A biological relationship between people who share a recent common ancestor; it can increase the chance that both parents carry the same recessive disease variant.

\medterm{Conscious Sedation For Surgical Procedures} Use of medicines to make a person relaxed and sleepy during a procedure while allowing them to respond to voice or touch and usually breathe independently. Heart rate, breathing, oxygen, and alertness require continuous monitoring.
\medterm{Consolidation} A part of the lung whose air spaces have filled with fluid, pus, blood, cells, or another substance. It appears denser on imaging and may result from pneumonia, fluid overload, bleeding, or a tumor.

\medterm{Consolidation Therapy} Additional treatment given after an initial response or remission to destroy remaining disease and lower the chance of relapse. It may include medicines, radiation, surgery, or a transplant depending on the illness.
\medterm{Constipation} A condition involving fewer than three bowel movements per week, hard or lumpy stools, difficult or painful passage, straining, or a feeling that the bowel has not emptied completely. Causes include diet, medicines, dehydration, functional bowel disorders, and structural or systemic disease. Rectal bleeding, vomiting, constant abdominal pain, inability to pass gas, fever, or unexplained weight loss requires medical evaluation.

\medterm{Constipation (Laxatives For Constipation)} Alternate index label for Laxatives For Constipation.






\medterm{Constipation In Infants And Children} Infrequent, hard, painful, or difficult passage of stool in a child, caused by stool withholding, diet, dehydration, illness, medicines, or structural and neurologic disorders.

\medterm{Constitutional Delay In Puberty} A common pattern in which puberty starts later than average but eventually progresses normally. Growth and bone development may also be delayed, often with a family history of late puberty.

\medterm{Constraint-Induced Movement Therapy} Rehabilitation that limits use of a less-affected limb while providing structured practice with an affected limb to
improve functional use in selected people.

\medterm{Constrictive Pericarditis} Thickening and stiffening of the sac around the heart, preventing the heart from filling normally. It can cause breathlessness, leg or abdominal swelling, enlarged neck veins, and liver congestion; causes include infection, surgery, radiation, or unknown disease.

\medterm{Consultation-Liaison Psychiatry} The subspecialty of psychiatry focused on the interface between psychiatry and general
medical illness. Psychiatrists provide consultation on medical-surgical wards for
patients with psychiatric symptoms arising from or coexisting with medical conditions.


\medterm{Consumption} An older term for tuberculosis, especially wasting pulmonary tuberculosis; in general usage it also means the act or amount of using something.


\medterm{Contac Overdose} Overdose of a combination cold medicine that may cause acetaminophen toxicity, antihistamine or decongestant effects, sedation, seizures, or cardiac complications depending on ingredients.

\medterm{Contact Allergen Testing} Patch testing checks whether a substance causes allergic skin inflammation. Small amounts of suspected allergens are placed on the skin, usually the back, and the area is examined after one or more days for a delayed reaction.

\medterm{Contact Dermatitis} Red, itchy, burning, or blistered skin caused by contact with an irritating or allergy-triggering substance. Common triggers include soaps, chemicals, metals, fragrances, plants, and medicines.

\synonyms
Dermatitis, Contact

\medterm{Contact Precautions} Infection control measures used for patients colonised or infected with
multidrug-resistant organisms. Precautions include gown and glove use, patient placement
in single rooms, and dedicated equipment. Contact precautions are used for MRSA, VRE,
CRE, and other resistant organisms.

\medterm{Contact Tracing} Identifying and monitoring persons exposed to infectious disease. Critical for outbreak
control. Examples: COVID-19, Ebola, TB.




\medterm{Contagious Disease} An infectious disease that can be transmitted directly or indirectly from one person,
animal, or source to another.














\medterm{Contamination} Unwanted microorganisms or other material entering a specimen, surface, medicine, or normally sterile area without causing infection in the person. It can produce misleading test results and must be distinguished from true disease.

\medterm{Continuing Medical Education} Education after professional training that updates a healthcare worker's knowledge, skills, judgment, and ability to provide safe, effective patient care.

\medterm{Continuity Of Care} Consistent and coordinated care over time and across clinicians or settings, integrating
the patient’s history, current needs, treatment plan, and preferences. It supports
safety, adherence, prevention, and chronic-disease outcomes.

\medterm{Continuous Bladder Irrigation} Continuous flow of sterile fluid through a urinary catheter into and out of the bladder. It helps prevent blockage and wash out blood clots or debris, often after urinary-tract surgery.

\medterm{Continuous Combined Hormone Therapy} A hormone treatment regimen taking estrogen and progestogen continuously, often to manage menopausal symptoms while reducing regular withdrawal bleeding.

\synonyms
Continuous combined HRT; continuous estrogen-progestogen therapy.

\medterm{Continuous Electronic Fetal Monitoring} The continuous recording of fetal heart rate and uterine contractions using
external transducers (Doppler ultrasound and tocodynamometer) or internal devices (fetal
scalp electrode and intrauterine pressure catheter). Used intrapartum to detect fetal
hypoxia.

\medterm{Continuous Glucose Monitoring} A wearable sensor that measures glucose in fluid beneath the skin at regular intervals and shows current levels, trends, and alerts. It helps manage diabetes but may need confirmation with a blood-glucose test when readings do not fit symptoms.

\medterm{Continuous Kidney Replacement Therapy} A slow, continuous form of dialysis used mainly in intensive care for severe kidney failure. It gradually removes extra fluid and waste and is useful when a person is too unstable for conventional dialysis.

\medterm{Continuous Murmurs} Heart sounds heard throughout both contraction and relaxation because blood flows continuously through an abnormal opening or connection. A persistent murmur requires assessment to identify its cause and significance.
throughout the cardiac cycle.

\medterm{Continuous Positive Airway Pressure} Noninvasive respiratory support that delivers a constant positive pressure throughout
the breathing cycle to keep the upper airway or alveoli open.

\medterm{Continuous Subcutaneous Insulin Infusion} Delivery of rapid-acting insulin through a pump and a subcutaneous infusion set to provide basal insulin and mealtime doses.

\medterm{Contraception} Any method used to prevent pregnancy, including intrauterine devices, implants, pills, injections, condoms, fertility-awareness methods, and sterilization. Methods differ in effectiveness, duration, side effects, cost, and protection against sexually transmitted infections. The best choice depends on health, preferences, reproductive plans, and access to care.

\medterm{Contraception And Birth Control} Contraception encompasses methods used to prevent pregnancy, including hormonal options
like oral contraceptives, patches, rings, injections, and implants, as well as
non-hormonal methods like copper IUDs, barrier methods, diaphragms, and permanent
sterilization.

\medterm{Contraception For Adolescents} Methods to prevent unintended pregnancy. Options: combined oral contraceptives,
progestin-only pills, depot medroxyprogesterone acetate, contraceptive implant
(Nexplanon), intrauterine devices (IUDs), condoms, and abstinence.

\medterm{Contraceptive Device, Intrauterine} A small device placed inside the uterus to prevent pregnancy; copper IUDs impair sperm function, while hormonal IUDs release levonorgestrel and thicken cervical mucus.

\medterm{Contraceptive, Pill} An oral hormonal contraceptive, usually a combined estrogen--progestin pill or a progestin-only pill, that prevents ovulation and/or changes cervical mucus and the uterine lining.

\medterm{Contracted Pelvis} A pelvis whose size or shape is reduced or abnormal enough to interfere with passage of
the fetus during labor.

\medterm{Contractile Proteins} Proteins, especially actin and myosin, that interact inside muscle cells to generate force, shorten fibers, and produce body movement.

\synonyms
Muscle contractile proteins; actomyosin proteins.

\medterm{Contraction} Shortening or tightening of a muscle; uterine contractions during labor thin and dilate the cervix and help deliver the fetus.

\medterm{Contraction Stress Test} A pregnancy test that records the baby's heart rate while contractions occur, sometimes after gentle stimulation or medicine. Repeated slowing of the heart rate may suggest that the placenta cannot supply enough oxygen during labor.

\medterm{Contractions} Rhythmic tightening of the uterus during labor that thins and opens the cervix and helps move the baby through the birth canal. Before labor, Braxton Hicks contractions are often irregular and do not progressively open the cervix.

\medterm{Contractions Braxton-Hicks (Braxton Hicks Contractions)} Alternate index label for Braxton Hicks Contractions.

\medterm{Contracture} A lasting inability to fully move a joint because nearby muscles, tendons, ligaments, skin, or other tissues have shortened or stiffened. It may follow injury, burns, paralysis, inflammation, or prolonged immobility.

\medterm{Contracture Deformity} A fixed limitation of movement caused by shortening or tightening of muscle, tendon, ligament, joint capsule, skin, or other soft tissue. It may follow neurologic disease, prolonged immobility, injury, burns, or inflammation.

\medterm{Contracture Prevention} Measures such as positioning, stretching, exercise, splinting, orthoses, and serial casting that help preserve joint movement in people at risk of fixed shortening.

\medterm{Contraindication} A condition or circumstance in which a treatment, procedure, or medication should not be used because the potential harm outweighs the expected benefit.

\synonyms
Contraindications


\medterm{Contralateral} On or affecting the side opposite a stated body part or injury. For example, a stroke on the left side of the brain may cause weakness on the right side of the body.

\medterm{Contrast} A difference that makes structures distinguishable; in imaging, contrast material is administered to improve visualization of organs, vessels, or tissues.

\medterm{Contrast Agent} A substance given or placed in the body to make organs, tissues, blood vessels, or body spaces easier to see during imaging. Different agents are used for X-ray, CT, MRI, ultrasound, or digestive-tract studies.

\medterm{Contrast Echocardiography} Use of microbubble contrast agents to enhance endocardial border definition and assess
myocardial perfusion. Indicated for suboptimal non-contrast echo images and assessment
of myocardial viability. Microbubbles remain entirely within the intravascular space,
making them ideal for perfusion assessment.

\medterm{Contrast Medium} A substance administered or placed in the body to make tissues, organs, or body spaces more visible during imaging.

\synonyms
Contrast agent; radiographic contrast material.

\medterm{Contrast Radiology} Imaging of the gastrointestinal tract using contrast material, such as barium, to outline the esophagus, stomach, intestine, or colon on radiographs or fluoroscopy.

\medterm{Contrast Reaction} An unwanted response after an imaging contrast agent, ranging from nausea, itching, or a rash to breathing difficulty, low blood pressure, or a severe allergic-like reaction. Imaging staff should be told about previous reactions.
\medterm{Contrast-Associated Acute Kidney Injury} Acute worsening of kidney function occurring after administration of contrast, whether
or not contrast is causal. Risk assessment includes baseline kidney function,
hemodynamic status, diabetes, nephrotoxins, urgency, and the need to avoid unnecessary
contrast exposure.

\medterm{Contrast-Enhanced Ultrasound} An ultrasound examination using tiny gas-filled bubbles in the bloodstream to improve pictures of blood flow and help characterize selected liver, heart, blood-vessel, or other lesions.

\medterm{Contrast-Induced Nephropathy} A sudden decline in kidney function after iodinated contrast, although the contrast is not always the cause. Risk is higher with severe kidney disease, dehydration, low blood pressure, or other kidney stresses.

\medterm{Contrecoup} An injury occurring on the side of an organ opposite the site of impact, especially a
brain injury produced when the brain moves within the skull after trauma.

\medterm{Control (Research)} A comparison participant or group in a study that does not receive the experimental intervention, or receives standard care or a placebo, allowing the intervention's effects to be estimated.

\medterm{Control Group} The comparison group in a study, which may receive usual care, a placebo, or no intervention according to the study design.

\medterm{Controlled Cord Traction} A technique for active management of third-stage labor: gentle, steady traction on the
umbilical cord with one hand while the other hand applies suprapubic counter-pressure
(to prevent uterine inversion). Performed after signs of placental separation. The
Brandt-Andrews maneuver and the Finocchietti technique are CCT methods.

\medterm{Controlled Ovarian Hyperstimulation} Use of hormone medicines to make several ovarian follicles mature at the same time for fertility treatment. Ultrasound and blood tests guide dosing and help reduce the risk of ovarian hyperstimulation, a potentially serious complication.

\synonyms
COH; Ovarian Stimulation

\medterm{Controllers} Long-term asthma medicines taken regularly to reduce airway inflammation, prevent symptoms and attacks, and lower the need for quick-relief inhalers.

\synonyms
Controller medicines; maintenance asthma therapy.


\medterm{Contusion} A bruise caused by a blunt blow that damages small blood vessels without breaking the skin. It may cause pain, swelling, and discoloration; severe injuries can damage muscle, organs, or deeper tissues.

\medterm{Conus Medullaris Syndrome} A group of symptoms caused by damage to the lower end of the spinal cord. It may cause back pain, weakness or numbness in the legs and pelvic area, and loss of bladder, bowel, or sexual function.

\medterm{Convergence} The inward movement of both eyes so that they are directed toward the same nearby object.

\medterm{Convergence Insufficiency} Difficulty turning both eyes inward together to focus on something nearby. It can cause eyestrain, headaches, blurred or double vision, and reading difficulty, especially in children and young adults.

\medterm{Conversion Disorder} Alternate index label; see \textit{Functional neurologic disorder/conversion disorder}.

\medterm{Conversion Disorder (Functional Neurological Symptom Disorder)} A condition in which neurological symptoms, such as weakness, abnormal movements, sensory loss, or seizures, are not explained by a recognized structural neurological disease and cause significant impairment.

\synonyms
Functional Neurological Symptom Disorder

\medterm{Convulsion} A sudden episode of involuntary, often rhythmic muscle contractions that may occur with a seizure or other disturbance of brain function.

\synonyms
Convulsive episode; seizure-related motor activity.

\medterm{Cooking Without Salt} Preparing food with little or no added sodium to reduce dietary sodium intake, often useful in blood-pressure or fluid-retention management when combined with an individualized eating plan.

\medterm{Cooley'S Anemia} Alternate index label; see \textit{Thalassemia}.

\medterm{Coombs Test} A laboratory test that detects antibodies or complement attached to red blood cells, or
antibodies in serum that can react with red cells. It is used in evaluating immune
hemolysis and blood compatibility.

\medterm{Cooper'S Ligament} Fibrous connective-tissue septa extending from the deep surface of the breast toward the
skin and pectoral fascia. Tumor invasion or fibrosis may cause skin retraction.

\medterm{Coordination} The ability to control movements smoothly, accurately, and in the correct order. It depends on the brain, nerves, muscles, vision, and balance and may be tested with finger-to-nose, heel-to-shin, or rapid alternating movements.
movements.

\medterm{COPD} Chronic obstructive pulmonary disease is a progressive, largely irreversible airflow
limitation caused by chronic bronchitis, emphysema, or both, most commonly due to
cigarette smoking. It presents with chronic cough, sputum production, and exertional
dyspnea. Spirometry shows a reduced FEV1/FVC ratio.

\medterm{COPD - Control Drugs} Long-term medicines used to reduce COPD symptoms and exacerbations, such as long-acting bronchodilators and, for selected patients, inhaled corticosteroids.


\medterm{COPD - Quick-Relief Drugs} Fast-acting inhaled bronchodilators used for short-term relief of COPD breathlessness or bronchospasm; frequent need may indicate inadequate control or an exacerbation.


\medterm{COPD And Other Health Problems} The interaction between chronic obstructive pulmonary disease and conditions such as cardiovascular disease, osteoporosis, depression, diabetes, and lung cancer, which can complicate symptoms and care.

\medterm{COPD Exacerbation} Acute worsening of respiratory symptoms beyond normal day-to-day variation in a patient
with COPD that leads to additional therapy. Common triggers include viral or bacterial
infection, air pollution, and cardiopulmonary comorbidity.

\medterm{COPD Flare-Ups} Acute worsening of breathlessness, cough, or sputum beyond usual daily variation in chronic obstructive pulmonary disease, often triggered by infection, pollution, or other irritants.

\medterm{COPD Rehabilitation} A supervised program combining exercise, breathing strategies, education, nutrition, and support to help people with COPD breathe and function better, remain active, and reduce the effects of flare-ups.







\medterm{Copper} An essential trace element required for iron transport, connective-tissue cross-linking,
energy metabolism, neurotransmitter synthesis, and antioxidant enzymes. Deficiency can
cause anemia, neutropenia, impaired neurologic function, and bone abnormalities; excess
is toxic and accumulates in Wilson disease.

\medterm{Copper In Diet} Copper obtained from foods such as shellfish, nuts, seeds, beans, whole grains, and organ meats. The body needs small amounts, while excessive supplements or exposure can be harmful; intake should be tailored for deficiency or disease.

\medterm{Copper Poisoning} Toxic exposure to copper or copper salts that may cause vomiting, abdominal pain, hemolysis, liver injury, kidney injury, or shock.

\medterm{Coprolalia} Involuntary utterance of obscene, taboo, or socially inappropriate words, most often described as a complex vocal tic in Tourette syndrome.

\medterm{Copy Number Variant} A deletion or duplication that changes the number of copies of a DNA segment; its clinical significance depends on the region, size, genes involved, inheritance, and available evidence.

\medterm{Copy Number Variation} A deletion or duplication of a section of DNA that changes how many copies are present. Some variations are harmless, while others alter gene activity and cause developmental, neurologic, or other disorders.

\medterm{CoQ10} A fat-soluble substance made by the body that helps mitochondria produce energy and also acts as an antioxidant. It is sold as a supplement, but benefit varies by condition and it can interact with medicines.
\medterm{Cor} A Latin anatomical word meaning heart, appearing in medical terms that describe the heart, heart-related structures, or diseases affecting circulation.

\medterm{Cor Pulmonale} Enlargement and weakening of the right side of the heart caused by high pressure in the lung arteries, usually from long-term lung disease. It may cause breathlessness, leg swelling, enlarged neck veins, or liver congestion.

\medterm{Cor Triatriatum} A rare birth defect in which a membrane divides the left upper heart chamber into two spaces and may obstruct blood returning from the lungs. Symptoms range from none to breathlessness, poor growth, or heart failure.

\medterm{Coracoid Process} A hook-shaped bony projection from the shoulder blade that points forward. Several muscles and ligaments attach to it, and injury may occur with a direct blow or shoulder dislocation.

\medterm{Cord} A long, narrow structure such as the spinal cord, which contains the central nervous tissue, or the umbilical cord, which connects a fetus to the placenta.

\medterm{Cord Blood Testing} Laboratory testing of blood collected from the umbilical cord at birth, which may include blood type, complete blood count, blood gases, glucose, or tests for infection or inherited disease when clinically indicated.

\medterm{Cord Blood Stem Cells} Blood-forming stem cells collected from the umbilical cord and placenta after birth. They can make different types of blood and immune cells and may be used for stem-cell transplantation.

\medterm{Cord Presentation} A position in which the umbilical cord lies between the baby's presenting part and the cervix while the membranes are intact. If the membranes break, the cord may slip down and become compressed, threatening the baby's oxygen supply.

\medterm{Cordocentesis} A prenatal procedure using ultrasound guidance to take a small blood sample from the umbilical cord. It can assess fetal anemia, infection, or selected genetic conditions but carries risks including bleeding and pregnancy loss.

\medterm{Cordotomy} A procedure that cuts selected pain-carrying nerve fibers in the spinal cord to relieve severe pain that has not responded to other treatment. It can cause weakness, numbness, or other neurologic complications.

\medterm{Core Biopsy} Removal of one or more cylindrical tissue samples with a hollow cutting needle for microscopic examination. It preserves more tissue architecture than fine-needle aspiration and is used for breast, liver, prostate, and other lesions.

\medterm{Core Needle Biopsy} Removal of small cylinders of tissue with a hollow cutting needle for laboratory examination. Imaging may guide the needle, and the samples preserve tissue structure better than fluid or cell aspiration alone.

\medterm{Cori Cycle} A process in which muscles send lactate to the liver, where it is changed into glucose. The glucose returns to the blood and can be used again by muscles for energy.
\medterm{Cori Disease} A rare inherited disorder in which the body cannot break down stored glycogen normally. It can cause low blood sugar, an enlarged liver, muscle weakness, and sometimes heart disease.
\synonyms
Glycogen Storage Disease Type III (GSD III); Forbes Disease; Debrancher Deficiency

\medterm{Corn} A localized thickening of skin caused by repeated pressure or friction, often occurring on the toes or feet.

\synonyms
Clavus; heloma.

\medterm{Cornea} The clear, curved front surface of the eye. It protects the eye and bends light so images focus on the retina; scratches, infection, swelling, or scarring can cause pain and blurred vision.

\medterm{Corneal Abrasion} A scratch or loss of the outer layer of the cornea, often from a fingernail, foreign body, contact lens, makeup tool, or other injury. It can cause pain, tearing, redness, light sensitivity, and blurred vision.

\medterm{Corneal Cross-Linking} A treatment that uses riboflavin and ultraviolet light to strengthen the cornea. It can slow or stop worsening thinning and bulging, especially in keratoconus, but may cause pain during healing.

\medterm{Corneal Disease} Corneal disease affects the transparent front surface of the eye through infection, inflammation, trauma, dystrophy, degeneration, dryness, or abnormal shape. Pain, redness, light sensitivity, discharge, or reduced vision can signal urgent disease because scarring may impair sight.

\medterm{Corneal Edema} Swelling of the cornea from excess fluid, making it cloudy and causing blurred vision, glare, halos, or pain. Causes include damage to its inner cells, glaucoma, inflammation, injury, or eye surgery.

\medterm{Corneal Graft (Corneal Transplant)} Surgical replacement of part or all of a diseased or damaged cornea with donor corneal tissue to restore vision or relieve pain.

\synonyms
Corneal Transplant

\medterm{Corneal Implant} An implanted device or tissue substitute placed in or on the cornea to restore vision,
replace diseased tissue, or alter corneal shape.

\medterm{Corneal Inflammation (Keratitis)} Alternate index label for Keratitis.

\medterm{Corneal Injury} Damage to the clear front surface of the eye, including scratches, foreign bodies, chemical burns, infections, or cuts that may threaten vision.

\medterm{Corneal Transplant} Surgery that replaces part or all of a damaged or diseased cornea with healthy donor tissue. It may restore vision or relieve pain; risks include rejection, infection, increased eye pressure, and need for further treatment.


\medterm{Corneal Ulcer} A corneal ulcer is a defect in the corneal epithelium with associated stromal
inflammation and necrosis, typically caused by infectious or non-infectious mechanisms,
and it is a sight-threatening condition that requires prompt diagnosis and treatment.

\medterm{Corneal Ulcers And Infections} Infection or inflammation of the cornea that may produce an epithelial defect or ulcer, pain, redness, tearing, discharge, and light sensitivity. Bacterial, fungal, protozoal, and viral causes can threaten vision and require urgent ophthalmic assessment.

\medterm{Corneitis} Inflammation of the cornea, often producing pain, photophobia, tearing, and impaired
vision. It may be infectious, immune-mediated, traumatic, or related to exposure.

\medterm{Cornelia De Lange Syndrome} A genetic disorder caused by mutations in genes of the cohesin complex (NIPBL 60\%,
SMC1A, SMC3, RAD21, HDAC8).

\medterm{Corns} A corn is a localized thickening of skin caused by repeated pressure or friction, commonly on toes or other weight-bearing areas. Proper footwear, pressure relief, and treatment of the underlying deformity help; diabetes or poor circulation makes self-removal risky.

\medterm{Corns And Calluses} Corns are small, localized areas of thickened skin caused by repeated pressure or friction,
usually on the toes. Calluses are broader areas of thickened, hardened skin, usually on the
soles of the feet or palms of the hands. Both are usually harmless but may become painful.

\medterm{Coronal} Relating to a plane that divides the body into front and back portions, or to the crown or top of a structure.

\medterm{Coronal Caries} Tooth decay affecting the crown, the visible part of a tooth above the gum. It can begin in pits, grooves, or smooth surfaces when bacterial acids dissolve enamel and dentin.

\medterm{Coronal Plane} A vertical plane that divides the body or an organ into front (anterior) and back
(posterior) portions.

\medterm{Coronary} Relating to the arteries and blood flow that supply the heart muscle. Narrowing or blockage of these arteries can reduce oxygen to the heart and cause chest pain or a heart attack.

\synonyms
Cardiac blood-vessel-related; heart-coronary.

\medterm{Coronary Aneurysm} An abnormal widening or ballooning of a coronary artery. It may result from atherosclerosis, inflammation, infection, injury, or a congenital condition and can cause clotting, blockage, or rarely rupture.

\medterm{Coronary Angiography} An X-ray examination of the heart arteries after contrast is injected through a thin catheter. It shows narrowed or blocked areas and can guide treatment, but the procedure can cause bleeding, stroke, kidney injury, or other complications.

\medterm{Coronary Arteries} The blood vessels that arise from the aorta and supply oxygen to the heart muscle. The right and left arteries divide into branches, and blockage in one can cause angina or a heart attack.

\medterm{Coronary Artery} One of the blood vessels branching from the aorta that supplies oxygen-rich blood
directly to the heart muscle.

\medterm{Coronary Artery Aneurysm} A localized dilatation of a coronary artery exceeding the normal vessel diameter by
approximately 1.5 times. It may be caused by atherosclerosis, Kawasaki disease,
connective-tissue disorders, infection, trauma, or an iatrogenic procedure, and can lead
to thrombosis, distal embolization, myocardial ischemia, or rupture.

\medterm{Coronary Artery Anomaly} A congenital abnormality in the origin, course, structure, or termination of a coronary artery. Some anomalies are incidental, whereas others can impair myocardial perfusion or increase the risk of ischemia, arrhythmia, or sudden cardiac death.

\medterm{Coronary Artery Bypass} A surgical procedure that creates a new route for blood flow around a narrowed or
blocked coronary artery using a vascular graft.

\medterm{Coronary Artery Bypass Grafting} Surgery that uses a blood-vessel graft to create a route around severely narrowed or blocked heart arteries. It can improve blood flow and symptoms and may improve survival in selected people with extensive disease.

\medterm{Coronary Artery Bypass Surgery} An operation that uses a blood vessel from another part of the body to create a new route around narrowed or blocked heart arteries and improve blood flow to heart muscle.
\synonyms
CABG; coronary bypass surgery.

\medterm{Coronary Artery Disease} Narrowing or blockage of the arteries supplying the heart, usually from fatty plaque buildup. It may cause exertional chest pressure, heart attack, heart failure, abnormal rhythms, or sudden death.

\medterm{Coronary Artery Fistula} An abnormal connection between a coronary artery and a heart chamber or vessel that may cause ischemia, failure, arrhythmia, or endocarditis.

\medterm{Coronary Artery Spasm} Temporary tightening of a heart artery that reduces blood flow and may cause chest pain, often at rest. Smoking, stimulant drugs, stress, or other vessel problems can trigger it.

\medterm{Coronary Calcium Score} A CT measurement of calcium in the arteries supplying the heart. More calcium usually means more plaque and higher future cardiovascular risk, but a zero score does not exclude soft plaque or eliminate all risk.

\medterm{Coronary Calcium Scoring} A computed tomography scan measuring calcium-containing plaque in heart arteries. The score helps estimate future cardiovascular risk in selected people but does not show all plaque or replace clinical assessment.

\medterm{Coronary Care Unit} A hospital unit for people with serious or unstable heart conditions. Staff continuously monitor heart rhythm, blood pressure, oxygen, and other vital functions and provide urgent treatment for problems such as heart attack, dangerous rhythms, or heart failure.
\synonyms
CCU

\medterm{Coronary Circulation} The network of arteries and veins that supplies the heart muscle with oxygen and removes waste. Blood flow increases when the heart works harder and falls when an artery is narrowed or blocked.

\medterm{Coronary Computed Tomography Angiography} A noninvasive CT scan performed after injected contrast highlights the heart arteries. It can show plaque, narrowing, or blockage and is useful for selected people with suspected coronary artery disease.

\synonyms
CCTA; Coronary CT Angiography; Cardiac CTA

\medterm{Coronary Dominance} The pattern describing which heart artery gives rise to the vessel supplying the back part of the heart. Most people have right dominance; left or shared patterns are normal variations.

\medterm{Coronary Endarterectomy} Surgery removing hardened plaque from the inner lining of a heart artery. It may improve blood flow in selected complex blockages but is uncommon and carries risks of heart attack, bleeding, or vessel injury.

\synonyms
Coronary plaque endarterectomy; coronary endarterectomy procedure.

\medterm{Coronary Heart Disease} Disease caused by reduced blood flow through the coronary arteries, usually due to
atherosclerotic plaque. It may cause angina, myocardial infarction, heart failure, or
sudden cardiac death. Also called coronary artery disease.

\synonyms
coronary artery disease.

\medterm{Coronary Insufficiency} Inadequate blood flow through the coronary circulation to meet myocardial oxygen demand,
usually because of coronary artery disease.

\medterm{Coronary Microvascular Angina} Angina caused by impaired regulation or constriction of the small coronary vessels, often producing ischemic symptoms despite no flow-limiting epicardial coronary obstruction. Assessment may include evaluation for coronary microvascular dysfunction; management is individualized according to symptoms, cardiovascular risk, and the underlying mechanism.

\medterm{Coronary Microvascular Disease} Alternate index label; see \textit{Small vessel disease}.

\medterm{Coronary Microvascular Dysfunction} Abnormal function of the small heart vessels that limits blood flow even when the major arteries are not severely narrowed. It may cause chest pain, especially in women and people with diabetes or high blood pressure.

\medterm{Coronary Occlusion} Obstruction of a coronary artery, commonly by atherosclerotic plaque and thrombus, that
can cause myocardial ischemia or infarction.

\medterm{Coronary Sinus} A large vein on the back of the heart that collects deoxygenated blood from the heart
muscle and drains it into the right atrium.

\medterm{Coronary Slow Flow Phenomenon} Delayed movement of blood through the small heart vessels despite no major blockage on angiography. It may cause repeated chest pain at rest or with activity and can occasionally lead to abnormal rhythms or heart injury.

\synonyms
CSFP; Slow Coronary Flow; Primary Coronary Slow Flow

\medterm{Coronary Spasm} Temporary contraction of a coronary artery that reduces blood flow to the heart muscle and may cause angina or infarction.

\synonyms
Vasospastic angina mechanism; coronary artery spasm.

\medterm{Coronary Steal Phenomenon} Reduced blood flow to an already poorly supplied part of the heart when nearby healthy vessels widen during certain stress tests or treatments. This can worsen chest pain or reveal hidden coronary disease.

\synonyms
Coronary Steal Syndrome; Myocardial Steal

\medterm{Coronary Stent} A small mesh tube placed in a narrowed heart artery to help keep it open after angioplasty. Medicines are usually needed afterward to reduce the risk of clotting.
\medterm{Coronary Thrombosis} Formation of a blood clot in a coronary artery, potentially obstructing blood flow and causing
myocardial infarction.

\medterm{Coronavirus} A group of viruses that infect people and animals. Human coronaviruses can cause respiratory illness ranging from a mild cold to severe pneumonia, depending on the virus and the person.
\medterm{Coronavirus Disease 2019} A respiratory and systemic infectious disease caused by SARS-CoV-2. Illness ranges from
asymptomatic infection to pneumonia, respiratory failure, and multisystem complications.

\medterm{Corporeal} Relating to the physical body rather than thoughts, feelings, or the spirit. In medical writing, the word may describe a bodily symptom, physical examination, or treatment directed at the body.

\medterm{Corpse} The dead body of a human being; the subject of medicolegal examination. Forensic
assessment of a corpse includes identification, determination of death and time of death
(rigor mortis, livor mortis, algor mortis, vitreous potassium, gastric emptying),
external and internal examination, and collection of evidence.

\medterm{Corpus} Latin for body; in anatomy it refers to the main body of an organ or structure, such as the corpus callosum or corpus uteri.

\medterm{Corpus Callosum} Brain's largest white matter tract with approximately 200 million axons connecting left
and right hemispheres. Four parts: rostrum, genu, body, splenium. Facilitates
interhemispheric communication for sensory, motor, and cognitive integration.

\medterm{Corpus Cavernosum} One of two columns of erectile tissue in the penis that fill with blood during erection.

\synonyms
Cavernous body; penile corpus cavernosum.

\medterm{Corpus Luteum} The temporary endocrine structure formed from an ovarian follicle after ovulation; it produces progesterone during the luteal phase and early pregnancy.

\synonyms
Yellow body; luteal body.

\medterm{Corpus Luteum Cyst} A usually harmless ovarian cyst that forms after ovulation when the normal hormone-producing structure fills with fluid or blood. It often disappears within a few menstrual cycles but may cause pelvic pain, delayed menstruation, or bleeding if it ruptures.

\medterm{Corpus Spongiosum} A column of erectile tissue surrounding the penile urethra and expanding distally to form the glans.

\synonyms
Spongy body; corpus spongiosum penis.


\medterm{Corrected Age} The developmental age of a premature baby after allowing for the weeks born before 40 weeks of pregnancy. It helps interpret growth and milestones, usually during the first two years.

\medterm{Corrigan Pulse} A strong pulse that rises quickly and then collapses rapidly. It is classically associated with aortic-valve leakage but can occur with other high-flow states and requires clinical evaluation.

\medterm{Corrosive} A chemical that destroys tissue on contact, such as a strong acid or alkali. Exposure can burn skin, eyes, airways, or the digestive tract; prompt prolonged rinsing and urgent medical advice may be needed.

\medterm{Corrugator Muscle} A small facial muscle that pulls the eyebrows downward and inward, helping create frowning and other concentrated facial expressions.

\synonyms
Corrugator supercilii; glabellar muscle.

\medterm{Cortex} The outer layer of an organ or structure. Examples include the brain's cerebral cortex, the outer part of the adrenal gland, and the outer layer of a hair.

\synonyms
Cortical layer; outer cortex.

\medterm{Corti'S Ganglion} A group of nerve-cell bodies inside the cochlea that carries hearing signals from the inner ear to the brain.

\medterm{Cortical} Relating to the cortex, the outer layer of an organ such as the brain, kidney, or adrenal gland.

\medterm{Cortical Bone} Dense bone forming the outer shell of most bones and providing much of the skeleton's mechanical strength.

\synonyms
Compact bone; cortical tissue.

\medterm{Cortical Cataract} A cataract that begins in the outer part of the eye's lens and forms spoke-like cloudy areas extending inward. It can cause glare, halos, and blurred vision, especially in bright light.

\medterm{Corticobasal Degeneration} A rare, progressive brain disorder affecting regions involved in movement and higher brain functions. It usually affects one side more than the other and may cause stiffness and slow movement, difficulty carrying out learned movements (apraxia), involuntary jerks, abnormal muscle contractions (dystonia), reduced sensation, or an ``alien limb'' that feels or acts outside the person’s control.

\medterm{Corticosteroid} A steroid hormone from the adrenal cortex or a synthetic analogue; glucocorticoids reduce inflammation and immune activity, while mineralocorticoids regulate salt and water balance.

\medterm{Corticosteroids} Medicines similar to hormones made by the adrenal glands that reduce inflammation and suppress immune activity. They can help many conditions but may cause infection, high blood sugar, mood changes, bone loss, or adrenal suppression.

\medterm{Corticosteroids Mechanism} These medicines enter cells and bind hormone receptors that change gene activity. This reduces production of inflammatory signals and limits the activity and movement of several immune cells.

\medterm{Corticosteroids Overdose} Overdose of corticosteroids that may cause gastrointestinal symptoms, mood changes, hyperglycemia, fluid retention, infection risk, or adrenal suppression.

\medterm{Corticotropin-Releasing Factor} A hypothalamic hormone that stimulates the pituitary gland to release adrenocorticotropic hormone and helps coordinate the stress response.

\synonyms
Corticotropin-releasing hormone; CRH.

\medterm{Cortisol} A hormone made by the adrenal glands that helps the body respond to stress, maintain blood pressure and blood sugar, regulate metabolism, and control inflammation. Too much or too little can cause serious illness.

\medterm{Cortisol Blood Test} A blood test measuring cortisol, a hormone produced by the adrenal cortex. Results help evaluate adrenal insufficiency or cortisol excess, but vary substantially with time of day, stress, illness, medicines, and the assay used.

\medterm{Cortisol Urine Test} A test measuring cortisol excreted in a timed urine collection, commonly 24-hour urine free cortisol, used to evaluate suspected excess cortisol production such as Cushing syndrome.

\medterm{Cortisone} A corticosteroid that is converted in the body to cortisol and can be used to reduce inflammation or replace glucocorticoid activity.

\medterm{Coryza} Acute inflammation of the nasal mucous membranes with nasal discharge and obstruction,
commonly occurring in the common cold or other upper-respiratory infection.

\medterm{Cosmeceutical} A cosmetic product marketed as having biologically active, drug-like benefits; the term has no uniform formal regulatory category in many jurisdictions.

\medterm{Cosmetic Allergies} Cosmetic products can cause irritant dermatitis or allergic contact dermatitis, producing itching, redness, swelling, scaling, or blistering where applied. Stopping the suspected product and medical assessment, sometimes with patch testing, help identify the responsible ingredient.


\medterm{Cosmetic Dermatology} Dermatologic care intended to change appearance rather than treat disease, including procedures for wrinkles, scars, pigmentation, hair, or facial volume. Risks include infection, scarring, pigment changes, allergic reactions, and injury from improperly performed treatment.

\medterm{Cosmetic Ear Surgery} Surgery to change the size, shape, or position of the outer ear, often to correct prominent ears or a structural difference. Risks include bleeding, infection, scarring, altered sensation, and an unsatisfactory result.


\medterm{Costello Syndrome} A rare genetic condition, usually caused by an HRAS gene change, affecting growth, development, skin, heart, and muscles. It may cause feeding difficulty, short stature, distinctive facial features, heart disease, and increased risk of certain tumors.

\medterm{Costochondritis} Costochondritis is inflammation or irritation where ribs join the breastbone, causing localized chest pain often worse with movement or pressing the area.

\synonyms
Chest Wall Pain
Costosternal Chondrodynia
Costosternal Syndrome

\medterm{Costochondritis (Costochondritis And Tietze Syndrome)} Alternate index label for Costochondritis and Tietze Syndrome.

\medterm{Costochondritis And Tietze Syndrome} Costochondritis is inflammation or irritation at the junction of ribs and the breastbone, causing reproducible chest-wall pain. Tietze syndrome is a less common related condition with visible swelling at one or a few joints; heart and lung emergencies must still be excluded.

\synonyms
Costochondritis (Costochondritis And Tietze Syndrome)
Tietze Syndrome (Costochondritis And Tietze Syndrome)

\medterm{Costosternal Chondrodynia} Alternate index label; see \textit{Costochondritis}.

\medterm{Costosternal Syndrome} Alternate index label; see \textit{Costochondritis}.

\medterm{Cot Death} Sudden infant death syndrome, the sudden unexplained death of an apparently healthy
infant, usually during sleep, that remains unexplained after investigation.

\medterm{Cough} Cough is a protective reflex that clears the airways; its duration, triggers, sputum, associated symptoms, and examination findings help distinguish infectious, inflammatory, cardiac, and other causes.
foreign material from the tracheobronchial tree.

\medterm{Cough Augmentation} Techniques to improve cough effectiveness in patients with neuromuscular disease or
chest wall weakness. Methods include manually assisted cough (abdominal thrust timed
with cough), mechanical insufflation-exsufflation (CoughAssist device), and breath
stacking with a resuscitation bag.

\medterm{Cough Chronic (Chronic Cough)} Alternate index label for Chronic Cough.

\medterm{Cough Headaches} Head pain triggered by coughing, sneezing, straining, laughing, or another sudden pressure increase. New, severe, or persistent attacks need assessment because some are caused by a structural brain problem.
\medterm{Cough In Adults} A cough in an adult may be acute, subacute, or chronic and may result from infection, asthma, chronic obstructive pulmonary disease, reflux, upper-airway cough syndrome, medication, or other disease. Coughing blood, breathing difficulty, chest pain, low oxygen, or persistent unexplained cough requires medical assessment.

\medterm{Cough Suppressant} An antitussive medicine or measure that reduces the urge or reflex to cough; treatment should address the underlying cause and avoid suppressing a productive cough when clearance is needed.

\medterm{Coughing Up Blood} Expectoration of blood or blood-streaked sputum from the lower respiratory tract, called hemoptysis; heavy bleeding or breathing difficulty is an emergency.

\medterm{Could A Stiff Neck Be A Sign Of Something Serious} Neck stiffness commonly follows muscle strain, but stiffness with fever, severe headache, light sensitivity, confusion, rash, weakness, trauma, or neurologic change can signal meningitis, bleeding, spinal injury, or another emergency.


\medterm{Count Rate} The number of detected events each second during measurement. In medical imaging, very high rates can cause missed events because equipment needs recovery time, reducing accuracy unless corrected.

\medterm{Countercurrent Multiplication} A kidney process that creates a concentrated salt gradient deep in the kidney. This gradient allows the collecting ducts to reclaim more water when antidiuretic hormone is present.

\medterm{Counterpulsation} Timed inflation and deflation of a balloon or pressure cuffs with each heartbeat to improve blood flow to the heart or reduce the heart's workload. It is used only in selected serious cardiac conditions.

\medterm{Countertransference} The therapist's emotional reactions to the patient, which may be influenced by the
therapist's own unresolved conflicts or may be a response to the patient's transference.
Awareness of countertransference is essential for effective therapy and can provide
valuable clinical information about the patient's interpersonal style.

\medterm{Counting Carbohydrates} Estimating the grams of carbohydrate in foods and drinks to help plan meals, match insulin, and manage blood glucose, especially in diabetes.

\medterm{Coup-Contrecoup Injury} Brain injury in which the brain is damaged both beneath the point of impact and on the opposite side as it moves inside the skull. Falls and vehicle crashes are common causes.

\medterm{Courvoisier Sign} A painless, enlarged gallbladder in a person with jaundice, traditionally suggesting blockage of the bile duct by a tumor rather than a gallstone. It is a clue, not a diagnosis.

\medterm{Covered Stent (Vascular)} A vascular stent lined with a graft material that supports the vessel and excludes a leak or
aneurysm from the bloodstream.

\synonyms
Vascular

\medterm{COVID Toe} A chilblain-like skin lesion reported in association with COVID-19, usually affecting
the toes or fingers and appearing red, purple, or swollen.

\medterm{COVID-19} Coronavirus disease caused by infection with SARS-CoV-2. Illness ranges from asymptomatic or mild disease to pneumonia, respiratory failure, thrombosis, or death; common symptoms include fever, cough, shortness of breath, fatigue, muscle aches, and altered taste or smell. Some people develop persistent symptoms after the acute infection.

\medterm{COVID-19 Antibody Test} A blood test that detects antibodies produced in response to SARS-CoV-2 infection or vaccination; it does not reliably diagnose current infection or establish complete immunity.

\medterm{COVID-19 In Pregnancy} COVID-19 during pregnancy can be more severe than in nonpregnant adults and increases risks such as hospitalization, preterm birth, and stillbirth. Fetal infection before birth is uncommon but possible; vaccination and prompt care reduce risk.


\medterm{COVID-19 Symptoms} Symptoms of SARS-CoV-2 infection, which may include fever, cough, sore throat, congestion, fatigue, muscle aches, headache, loss of smell or taste, vomiting, or diarrhea.

\medterm{COVID-19 Test} A test for current SARS-CoV-2 infection using a respiratory sample. Molecular tests detect viral genetic material and antigen tests detect viral proteins; results depend on timing, sample quality, and the amount of virus present.


\medterm{COVID-19 Vaccines} Vaccines that prime immune responses against SARS-CoV-2 and reduce severe disease, hospitalization, and death; products, schedules, effectiveness, and recommendations vary by age and risk.

\medterm{COVID-19 Vaccines For Children Ages 6 Months And Older} COVID-19 vaccination guidance for children at least six months old, with product and dose schedules determined by age, prior doses, immune status, and current public-health recommendations.

\medterm{COVID-19 Virus Test} A test for current SARS-CoV-2 infection, usually using a nucleic acid amplification test or antigen test on a respiratory specimen; results must be interpreted with symptoms, timing, and test performance.



\medterm{Cow'S Milk - Infants} Guidance on cow’s milk in infancy: it should not replace breast milk or iron-fortified formula before 12 months because of low iron content and risks of renal solute load and intestinal blood loss.

\medterm{Cow'S Milk And Children} Nutritional guidance on cow’s milk for children, balancing protein, calcium, vitamin D, fat, iron, allergy, and total intake so milk does not displace a varied diet.

\medterm{Cowden Syndrome} An inherited condition, usually caused by a PTEN gene change, leading to multiple noncancerous growths and increased risks of breast, thyroid, uterine, kidney, and colorectal cancers. Regular specialist screening is important.

\medterm{Cowper'S Gland} One of a pair of small male accessory sex glands below the prostate. Its secretion
contributes mucus and fluid to the urethra before ejaculation.

\medterm{Cowpox} A zoonotic viral infection caused by cowpox virus or related orthopoxviruses, producing
localized pustular or ulcerative skin lesions and sometimes systemic illness.

\medterm{Cox-2} Cyclooxygenase-2, an inducible enzyme that produces inflammatory prostaglandins; selective COX-2 inhibitors reduce pain and inflammation but may have cardiovascular, renal, and gastrointestinal risks.

\medterm{COX-2 Inhibitors} Medicines that reduce the action of an enzyme involved in pain and inflammation. They may cause less stomach irritation than some older anti-inflammatory medicines, but can still cause bleeding, kidney problems, high blood pressure, or heart risks.

\synonyms
Cyclooxygenase-2 inhibitors; selective COX-2 inhibitors.

\medterm{Coxa Vara} A deformity in which the angle between the femoral neck and shaft is abnormally reduced,
potentially causing a limp, shortened limb, or altered hip mechanics.

\medterm{Coxalgia} Pain arising from the hip or nearby tissues, often felt in the groin and sometimes spreading to the thigh or knee. Causes vary with age and include infection, injury, growth problems, arthritis, cartilage damage, stress fractures, and reduced blood supply.

\medterm{Coxiella Burnetii} A bacterium that causes Q fever, usually spread by breathing dust contaminated by infected farm animals. It can cause fever, headache, hepatitis, or pneumonia and may persist as a serious infection of heart valves.

\medterm{Coxsackie (Coxsackie Virus)} Alternate index label for Coxsackie Virus.

\medterm{Coxsackie Virus Type A16} A virus that commonly causes hand, foot, and mouth disease, especially in children. Fever, painful mouth sores, and a rash or blisters on the hands and feet usually resolve without serious complications.

\medterm{Coxsackievirus} A group of enteroviruses that cause a range of human diseases. Group A viruses commonly
cause hand-foot-and-mouth disease and herpangina, whereas group B viruses are associated
with myocarditis, pericarditis, pleurodynia, and aseptic meningitis.

\medterm{CPAP} Continuous positive airway pressure delivered through a mask to splint open the upper airway during sleep. It is
commonly used for obstructive sleep apnea and can reduce apneas and related symptoms
when used consistently.

\medterm{CPAP Machine} A device that sends a steady stream of air through a mask to keep the upper airway open during sleep. It is commonly used for obstructive sleep apnea; benefit depends on correct mask fit and regular use.

\medterm{CpG} A DNA sequence in which cytosine is followed by guanine on the same strand, often occurring in regions whose methylation helps regulate gene expression.

\medterm{CPK Isoenzymes Test} A blood test that separates forms of creatine kinase to help identify whether increased enzyme levels come from skeletal muscle, heart muscle, or brain. Troponin testing is generally preferred for suspected heart attack.

\medterm{CPR} Cardiopulmonary resuscitation is an emergency procedure using chest compressions and, when appropriate, rescue breaths or an automated external defibrillator to support circulation and breathing during cardiac arrest.

\medterm{CPR - Infant} Cardiopulmonary resuscitation for an unresponsive infant who is not breathing normally, using emergency activation, infant chest compressions, rescue breaths when trained, and an AED according to local guidance.

\medterm{CPR - Young Child (Age 1 Year To Onset Of Puberty)} Cardiopulmonary resuscitation for a child from age one year until puberty who is unresponsive and not breathing normally, using emergency activation, compressions, breaths, and an AED as available.

\medterm{CPT Code For Foreign Body Removal From The Ear} Removal of an ear foreign body is coded according to the documented method, location, anesthesia, complications, and whether a separate evaluation was performed. Objects causing pain, bleeding, hearing loss, or suspected eardrum injury should be removed by a trained clinician.

\medterm{CPVT} Alternate name; see \textit{Catecholaminergic Polymorphic Ventricular Tachycardia}.

\medterm{Crabs} Alternate index label; see \textit{Pubic lice (crabs)}.

\medterm{Cracked Tooth Syndrome} A partial crack in a tooth that may cause sharp pain when biting or when pressure is released, and sensitivity to hot or cold. The crack can be difficult to locate and may worsen over time. Dental treatment depends on its depth and whether the inner pulp is involved.

\medterm{Crackles (Rales)} Discontinuous, non-musical adventitious lung sounds produced by the sudden explosive
opening of collapsed small airways and alveoli during inspiration or the bubbling of air
through fluid in the bronchioles.

\synonyms
Rales

\medterm{Cradle Cap} A common form of seborrheic dermatitis in infants, causing greasy or flaky scales on the
scalp with mild redness. It is usually harmless and improves with gentle scalp care;
persistent inflammation or signs of infection require assessment.

\medterm{Cramp, Muscle} Alternate index label; see \textit{Muscle cramp}.

\medterm{Cramping Leg Pain} Pain caused by involuntary, often painful muscle contraction in the leg; contributors include exertion, dehydration, electrolyte abnormalities, medications, nerve disease, and impaired circulation.

\medterm{Cramps But No Period} Pelvic cramps without menstruation can result from ovulation, pregnancy, implantation, constipation, urinary disease, infection, endometriosis, cysts, or other pelvic conditions. Severe one-sided pain, faintness, fever, vomiting, or possible pregnancy requires urgent assessment.

\medterm{Cramps Heat (Heat Cramps)} Alternate index label for Heat Cramps.

\medterm{Cramps Menstrual (Menstrual Cramps)} Alternate index label for Menstrual Cramps.

\medterm{Cramps, Menstrual} Alternate index label; see \textit{Menstrual cramps}.

\medterm{Cranial} Relating to the skull or toward the head, as opposed to caudal or toward the lower end of the body.

\medterm{Cranial Arteritis} Inflammation of arteries, usually in the scalp and head, that can cause a new headache, scalp tenderness, jaw pain while chewing, fever, or vision loss. Sudden visual symptoms require emergency treatment.

\medterm{Cranial Cavity} The hollow space inside the skull that contains and protects the brain, its coverings, blood vessels, and cerebrospinal fluid. A rise in pressure inside this space can compress the brain and become life-threatening.

\medterm{Cranial Fasciitis} A usually noncancerous, rapidly growing lump in the scalp or nearby soft tissue, most often in infants and young children. It can resemble cancer on examination or imaging and may require biopsy.

\medterm{Cranial Implant} An artificial material or device placed in or on the skull to repair a defect, protect the brain, drain fluid, stimulate nerves, or support treatment. Risks include infection, bleeding, and implant failure.

\medterm{Cranial Mononeuropathy III} Damage to the third nerve controlling several eye muscles and the eyelid. It can cause double vision, a drooping eyelid, an eye that points abnormally, or an enlarged pupil and needs evaluation.

\medterm{Cranial Mononeuropathy III - Diabetic Type} A third-nerve palsy associated with diabetes, often causing painful double vision and a drooping eyelid while the pupil remains normal. Other causes must be excluded, especially when the pupil is enlarged or symptoms are sudden.
\medterm{Cranial Mononeuropathy VI} Damage to the sixth cranial nerve, which controls the muscle that moves an eye outward. It can cause inward turning of the eye and horizontal double vision, especially when looking toward the affected side.

\medterm{Cranial Nerve} One of 12 pairs of nerves arising from the brain or brainstem. They control smell, vision, eye movement, facial sensation and movement, hearing, balance, swallowing, voice, and other head and neck functions.

\medterm{Cranial Nerve Palsy} Weakness or paralysis of one or more cranial nerves, causing deficits such as impaired eye
movement, facial weakness, hearing loss, swallowing difficulty, or tongue weakness. The
cause may be vascular, traumatic, compressive, inflammatory, infectious, or metabolic.

\medterm{Cranial Nerve X} The vagus nerve, the tenth cranial nerve, which carries parasympathetic and sensory fibers and contributes to swallowing, voice, and regulation of thoracic and abdominal organs.

\medterm{Cranial Nerves} The 12 pairs of nerves arising from the brain or brainstem that provide sensory, motor, and autonomic functions to the head, neck, and parts of the trunk.

\medterm{Cranial Sutures} Fibrous joints between skull bones that permit molding in infancy and accommodate brain growth before progressively fusing; premature fusion causes craniosynostosis.

\medterm{Craniectomy} Surgery removing part of the skull, sometimes leaving it off temporarily, to relieve dangerous brain pressure or access the brain. A later operation may replace the removed bone or use an implant.

\medterm{Craniocerebral Injury} An injury involving the scalp, skull, brain, or several of these structures, ranging from minor wounds to life-threatening traumatic brain injury.

\medterm{Craniofacial Abnormalities} Abnormalities of the skull, face, or related tissues that may be present at birth or develop later. Causes include genetic conditions, disrupted development, trauma, infection, and tumors.

\medterm{Craniofacial Dysostosis} A congenital disorder involving premature fusion of cranial sutures and characteristic
abnormalities of the face and skull, such as those seen in Crouzon syndrome.

\medterm{Craniopharyngioma} A usually slow-growing tumor near the pituitary gland and hypothalamus. It may cause headache, vision loss, excessive thirst or urination, hormone problems, growth changes, or increased pressure in the skull.

\medterm{Cranioplasty} Surgical repair or reconstruction of a skull defect using bone, an implant, or another
suitable material.

\medterm{Craniosynostosis} Craniosynostosis is premature fusion of one or more skull sutures, altering head shape and sometimes increasing pressure on the developing brain.

\synonyms
Congenital Plagiocephaly
Craniostenosis

\medterm{Craniosynostosis Repair} Surgery to separate prematurely fused skull bones and reshape the skull when needed. It aims to improve head shape and make room for brain growth; timing and technique depend on the child's age and anatomy.


\medterm{Craniotabes} Soft, thin, or pliable areas of an infant’s skull that indent under gentle pressure, commonly associated with normal newborn molding but also with rickets or other bone disease.

\medterm{Craniotomy} Surgery that temporarily removes a section of skull to reach the brain or its coverings. It may be used to remove a tumor, drain bleeding, repair a blood vessel, treat epilepsy, or relieve pressure.

\medterm{Cranium} The part of the skull that encloses and protects the brain, consisting of the cranial bones and their joints.

\medterm{Crapulence} An unpleasant state of sickness or discomfort caused by excessive eating or drinking, especially excessive alcohol use.

\medterm{Craving} A strong, often difficult-to-control desire for a substance, behavior, or experience, commonly seen in substance-use disorders.

\synonyms
Strong urge; compulsive desire.

\medterm{Crazy Paving Pattern} A lung-imaging pattern of hazy areas crossed by thick lines, resembling paving stones. It can occur with several conditions, including infection, bleeding, inflammation, or protein buildup, so it is not a diagnosis by itself.

\medterm{CRC} Colorectal cancer, malignant disease arising in the colon or rectum; screening and evaluation depend on age, risk, symptoms, and test results.

\medterm{CRE} Enterobacterales bacteria resistant to carbapenem antibiotics, important medicines for serious infections. They may colonize without illness or cause difficult-to-treat urinary, wound, abdominal, lung, or bloodstream infections.

\medterm{CRE Infection} Carbapenem-resistant Enterobacterales are bacteria resistant to important antibiotics and can cause urinary, bloodstream, abdominal, wound, or lung infections. Infection-control precautions, culture-based treatment, and specialist advice are important because therapeutic options may be limited.

\synonyms
Carbapenem Resistant Enterobacteriaceae (CRE Infection)

\medterm{Cream} A spreadable semisolid preparation applied to the skin, often containing a medicine. Creams are usually less greasy than ointments and may be useful on moist or broad areas, but the correct product depends on the skin condition.

\medterm{Creatine} A nitrogen-containing compound stored mainly in skeletal muscle that helps supply energy for short, high-intensity activity.

\medterm{Creatine Kinase} An enzyme found mainly in muscle and released into blood when muscle cells are damaged. Levels may rise after strenuous exercise, injury, seizures, muscle disease, or some medicines; other tests help identify the source.

\medterm{Creatine Phosphokinase Test} A blood test measuring creatine kinase, an enzyme released with muscle injury; elevated levels may occur in skeletal-muscle disease, strenuous exercise, rhabdomyolysis, or myocardial injury.

\medterm{Creating A Family Health History} Collecting a record of relatives’ diseases, ages at diagnosis, ancestry, and causes of death to identify inherited risk and guide screening, prevention, and genetic counseling.

\medterm{Creatinine} A breakdown product of muscle creatine phosphate produced at a constant rate, filtered
by the glomerulus and minimally reabsorbed; serum creatinine is the most used marker of
kidney function.

\medterm{Creatinine Blood Test} A blood test measuring serum creatinine, a waste product generated largely by skeletal muscle. It is used with age, body size, and other factors to estimate kidney filtration, but interpretation is affected by muscle mass, diet, medicines, and acute illness.

\medterm{Creatinine Clearance Test} A test estimating kidney filtration from creatinine in a timed urine collection and blood sample. It can be affected by incomplete urine collection and creatinine secretion, so estimated glomerular filtration rate is often preferred for routine assessment.

\medterm{Creatinine Test} A blood or urine test measuring creatinine to help assess kidney filtration and muscle-related creatinine production.

\synonyms
Creatinine measurement; renal function creatinine test.

\medterm{Creatinine Urine Test} A urine test measuring creatinine concentration or excretion, used to assess kidney handling, interpret other urine measurements, or judge the completeness of a timed collection. Results depend on muscle mass, diet, collection quality, and kidney function.

\medterm{Creeping Eruption} An intensely itchy winding skin track caused by larvae, usually hookworms from contaminated soil or sand, migrating within the superficial skin; it is also called cutaneous larva migrans.

\medterm{Crepitation (Crepitus)} A crackling, grating, or popping sensation or sound produced by movement of tissues, air, or bone. Causes include fractures, joint disease, subcutaneous air, and lung disease.

\synonyms
Crepitus

\medterm{Crepitus} A crackling, grating, or popping sensation or sound caused by movement of tissue, bone, joint surfaces, or air in soft tissue.

\synonyms
Crepitation; crackling sensation.

\medterm{Crescentic Glomerulonephritis} A rapidly worsening inflammation of the kidney's filtering units. It can cause blood or protein in urine, swelling, high blood pressure, and kidney failure and usually requires urgent specialist treatment.

\medterm{CREST Syndrome} A form of systemic sclerosis involving mainly the skin and certain organs. The name refers to calcium deposits, cold-sensitive fingers, swallowing difficulty, tight fingers, and small widened blood vessels; lung, heart, or digestive complications may occur.

\medterm{Creutzfeld-Jakob Disease} A rapidly progressive neurodegenerative disease caused by abnormal prion proteins, leading to
cognitive decline, ataxia, myoclonus, and death.

\medterm{Creutzfeldt Jakob Disease} Creutzfeldt-Jakob disease is a rare, fatal, transmissible spongiform
encephalopathy caused by the accumulation of misfolded prion protein (PrPSc) in the
brain, leading to rapid, progressive dementia and myoclonus.

\medterm{Creutzfeldt-Jakob Disease} A rare, rapidly progressive brain disease caused by abnormal prion proteins. It leads to worsening memory and thinking, poor coordination, involuntary jerks, and loss of function, and is ultimately fatal.
\synonyms
CJD
Mad Cow Disease
Variant CJD
VCJD

\medterm{Cri Du Chat Syndrome} A congenital genetic disorder caused by deletion of part of the short arm of chromosome
5, characterized in infancy by a high-pitched cry, developmental delay, and
characteristic facial features.

\medterm{Cri-Du-Chat Syndrome} A genetic disorder caused by a deletion of the short arm of chromosome 5 (5p-), with the
size of the deletion correlating with severity. Incidence 1 in 15,000-50,000.

\medterm{Crib Death} Alternate index label; see \textit{Sudden infant death syndrome}.

\medterm{Cricoid Cartilage} A complete ring of firm tissue at the lower part of the voice box, just above the windpipe. It supports the airway and connects with other structures involved in breathing and voice.

\medterm{Cricoid Pressure} Pressure applied to the cricoid cartilage during some emergency airway procedures. Its use is intended to reduce stomach contents entering the airway, but it can make ventilation or intubation more difficult and is clinician-dependent.

\medterm{Cricothyrotomy} An emergency operation creating an opening through the front of the neck into the airway when a person cannot be intubated or oxygenated by other means. It provides temporary access for breathing and has serious risks.

\medterm{Crigler-Najjar Syndrome} A rare inherited disorder in which the liver cannot properly process bilirubin. It causes persistent yellowing of the skin and eyes; severe forms can damage the brain unless treated urgently.

\medterm{Crime Scene Investigation} The systematic documentation, collection, preservation, and analysis of physical evidence from a suspected crime scene, often supporting forensic medical and legal conclusions.

\medterm{Criminal Abortion} Abortion performed illegally, outside permitted legal indications and procedures (e.g.,
unsafe methods, unqualified practitioners, or non-medical settings), carrying high
maternal risk.

\medterm{CRISPR-Cas9} A laboratory gene-editing tool that uses a guide molecule to direct an enzyme to a chosen DNA sequence. The enzyme cuts the DNA so the cell can remove, disable, or replace genetic material.

\medterm{CRISPR Therapy} A treatment that uses CRISPR-based genome-editing tools to make a targeted change to DNA in a patient’s cells. Cells may be edited outside the body and returned to the patient, or edited directly inside the body. It is used or studied for selected genetic and other diseases; possible concerns include unintended genetic changes, immune reactions, and effects that require long-term monitoring.

\medterm{Critical Care} Specialized care for people with life-threatening illness or injury who need close monitoring and support for failing organs.

\medterm{Critical Care Nursing} Specialized nursing care for critically ill patients requiring close monitoring, organ support, rapid assessment, complex therapies, and coordination with multidisciplinary intensive-care teams.
\medterm{Critical Findings} Findings requiring immediate communication. Examples: pneumothorax, brain herniation,
aortic dissection.

\medterm{Critical Illness Myopathy} A muscle weakness disorder that develops during severe illness, especially sepsis, organ failure, prolonged intensive-care treatment, or use of certain medicines. It can affect arms, legs, and breathing muscles, making movement and removal from a ventilator difficult. Rehabilitation supports recovery.

\synonyms
CIM; Acute Quadriplegic Myopathy; Critical Care Myopathy

\medterm{Critical Limb Ischemia} Severely reduced blood flow to an arm or leg causing pain even at rest, a wound that will not heal, or gangrene. It threatens the limb and requires urgent vascular assessment.


\medterm{Critical Results} Test findings suggesting a potentially life-threatening problem that require prompt communication, assessment, treatment, and documentation by the responsible healthcare team.

\medterm{Critical Value} A laboratory result indicating a potentially life-threatening condition that requires
prompt communication to a responsible clinician or care team. The threshold is test- and
laboratory-specific and should be defined by institutional policy.

\medterm{Crocodile Tears Syndrome} Gustatory lacrimation, in which eating or smelling food triggers excessive tearing, usually after facial-nerve injury or aberrant regeneration following Bell palsy.

\medterm{Crohn Disease} Chronic transmural inflammatory condition that can affect any segment of the
gastrointestinal tract from the oral cavity to the perianal region, though the terminal
ileum and ileocolonic region are most commonly involved.

\medterm{Crohn Disease (Crohn'S Disease)} Alternate index label for Crohn's Disease.



\medterm{Crohn Disease Cobblestoning} A bumpy bowel-lining appearance seen during endoscopy in some people with Crohn disease. Deep lengthwise ulcers surround raised areas of less-affected tissue, creating a pattern like cobblestones.

\medterm{Crohn'S Disease} A chronic inflammatory bowel disease that can cause patchy, transmural inflammation anywhere from the mouth to the anus, most often involving the terminal ileum or colon. Common features include abdominal cramping or pain, persistent diarrhoea, weight loss, fatigue, fever, anaemia, and extraintestinal inflammation of the joints, eyes, or skin.

\synonyms
Crohn Disease (Crohn's Disease)
Ileitis (Crohn's Disease)


\medterm{Crohn’S Disease Cause Dark Circles Around Eyes} Crohn disease does not directly cause dark circles in every patient, but anemia, fatigue, dehydration, poor sleep, nutritional deficiency, steroid effects, or atopic skin may contribute. Persistent pallor, fatigue, weight loss, bleeding, or active bowel symptoms deserve clinical review.

\medterm{Cross-Matching} Laboratory testing that checks whether donor blood is compatible with a recipient before transfusion. It helps prevent dangerous immune reactions by detecting antibodies that could attack transfused red cells.

\medterm{Cross-Section} A slice or plane made through a structure, or an image showing that slice, such as a cross-sectional CT image.

\medterm{Cross-Sectional Study} A study that measures possible causes and health outcomes at one point in time. It can show how common a condition is and whether factors are linked, but it usually cannot show which came first or prove cause and effect.

\medterm{Cross-Tolerance} Reduced response to one substance because the body has adapted to another substance with a similar effect. It can increase overdose risk when a person changes substances or combines them.

\medterm{Crossmatch} Laboratory test mixing donor and recipient blood to detect preformed antibodies that could cause transfusion or transplant reactions.

\medterm{Crotch} The region where the inner thighs meet the trunk, including the groin and nearby perineal structures.

\medterm{Croup} An acute viral infection causing inflammation and narrowing of the upper airway, typically producing a barking cough, hoarseness, and inspiratory stridor in young children.

\synonyms
Laryngotracheobronchitis (Croup)



\medterm{Crouzon Syndrome} An autosomal dominant craniosynostosis syndrome caused by mutations in the FGFR2 or
FGFR3 genes, leading to premature fusion of multiple cranial sutures. Features:
brachycephaly (short, broad head from bilateral coronal synostosis), proptosis (shallow
orbits), hypertelorism, midface hypoplasia, beaked nose, and Class III malocclusion.

\medterm{Crown} The visible part of a tooth above the gum line, or a dental restoration that covers and protects a damaged tooth.

\synonyms
Dental crown; anatomic tooth crown.

\medterm{Crown (Dental)} A custom-made cap placed over a weakened or damaged tooth to restore its shape, strength, and function. It may be made from metal, ceramic, or a combination; the tooth is reshaped first, and the crown is cemented in place.

\synonyms
Dental

\medterm{Crown Lengthening} A dental procedure that removes or reshapes gum tissue, and sometimes a small amount of bone, to expose more tooth. It may allow a damaged tooth to be restored or change a gummy smile.

\medterm{Crown-Rump Length} The ultrasound measurement from the top of an embryo's head to its bottom. It is the most accurate way to estimate pregnancy age during early pregnancy and helps monitor early growth.

\medterm{Crowning} The stage of childbirth when the baby's head remains visible at the vaginal opening between contractions because it has stretched the surrounding tissues. It occurs shortly before the head is born.

\medterm{CRP} C-reactive protein, a substance made by the liver that rises during inflammation or infection. A CRP test can help assess inflammation and response to treatment, but it does not identify the cause by itself.

\medterm{Crush Injury} Damage caused when a body part is squeezed or trapped under pressure. It can cause broken bones, muscle death, nerve and vessel injury, dangerous chemical changes in blood, shock, or kidney failure.

\medterm{Crush Syndrome} A life-threatening illness after prolonged muscle compression, often when someone is trapped. Damaged muscle releases substances that can cause shock, dangerous heart rhythms, and sudden kidney failure, especially after pressure is removed.

\medterm{Crushing Injury} Damage caused when a body part is compressed between heavy objects. It can produce fractures,
nerve and blood-vessel injury, tissue death, and rhabdomyolysis. Severe limb compression
or sudden release after prolonged entrapment requires emergency treatment.

\medterm{Crust} Dried serum, blood, pus, or other exudate on the skin surface, often covering an erosion, ulcer, vesicle, pustule, or inflamed lesion.

\medterm{Crusted Scabies} A severe form of scabies causing thick crusts full of very large numbers of mites. It is extremely contagious and occurs mainly in people with weakened immunity, older adults, or people with reduced sensation.

\medterm{Crutches} Walking aids that transfer some weight from the legs to the arms while a person stands or walks. Underarm and forearm types are adjusted to fit and require instruction to reduce falls, nerve pressure, and other injuries.
through loop.


\medterm{Crutches And Children - Sitting And Getting Up From A Chair} Instruction for a child using crutches to approach a chair, keep the injured limb protected as prescribed, reach for the armrests, and sit or stand safely without losing balance.

\medterm{Crutches And Children - Stairs} Instruction for safe stair use by a child with crutches, including handrail use, crutch placement, step sequencing, supervision, and adherence to weight-bearing restrictions.

\medterm{Crutches And Children - Standing And Walking} Instruction for fitting and using crutches in children, including posture, handgrip support, gait pattern, weight-bearing limits, turning, and fall prevention.

\medterm{Cry For Help} A verbal or behavioral expression of distress or need for assistance; in a mental-health setting it should be taken seriously and assessed for immediate safety concerns.

\medterm{Crying In Childhood} Vocal expression of distress or emotion in a child; persistent, inconsolable, unusual, or pain-associated crying may indicate illness, injury, abuse, or unmet developmental needs.

\medterm{Crying In Infancy} Crying is an infant’s normal communication of hunger, discomfort, fatigue, or need for contact, but sudden inconsolable crying with lethargy, fever, vomiting, or injury needs urgent assessment.

\medterm{Cryoablation} A treatment that uses controlled freezing to destroy selected abnormal tissue, such as a tumor or abnormal heart pathway. Imaging or specialized equipment guides the treatment, which may cause bleeding, injury, or recurrence.
\medterm{Cryobiopsy} A biopsy method that uses a freezing instrument to remove a larger tissue sample, often from the lung. It can improve diagnosis in selected cases but carries risks such as bleeding and a collapsed lung.

\medterm{Cryoglobulinemia} A condition in which abnormal blood proteins thicken or clump when cold. They can inflame small blood vessels and cause purplish skin spots, joint pain, nerve problems, kidney injury, or poor circulation.

\medterm{Cryoglobulinemic Vasculitis} Systemic vasculitis caused by cryoglobulin-containing immune complexes that deposit in
small and medium vessels, classically presenting with purpura, arthralgia, and
glomerulonephritis. Most cases are associated with hepatitis C virus infection.

\medterm{Cryoglobulins} Immunoglobulins that precipitate at low temperatures and dissolve on warming; circulating cryoglobulins can cause vasculitis, purpura, neuropathy, kidney disease, or hyperviscosity.

\medterm{Cryopexy} A procedure that uses controlled freezing to destroy or seal abnormal tissue, especially
retinal breaks or small retinal detachments.

\medterm{Cryoprecipitate} A blood-product concentrate containing mainly fibrinogen and some clotting proteins. It is given in selected serious bleeding disorders when a person's fibrinogen level is too low.

\medterm{Cryopreservation} Preserving cells, tissues, sperm, eggs, or embryos by cooling them to very low temperatures with protective chemicals. The material can later be thawed for treatment, research, or reproduction.

\medterm{Cryopyrin-Associated Autoinflammatory Syndromes} Alternate index label; see \textit{Cryopyrin-Associated Periodic Syndromes}; the group is also called CAPS.

\medterm{Cryopyrin-Associated Periodic Syndromes} A group of inherited inflammatory disorders causing repeated fever, rash, joint symptoms, eye inflammation, or hearing loss because immune activity is overactive.

\synonyms
Cryopyrin-Associated Autoinflammatory Syndromes

\medterm{Cryostat} An instrument that maintains very low temperatures and cuts thin frozen sections of tissue for rapid microscopic examination.

\medterm{Cryosurgery} Surgery that destroys abnormal tissue by controlled freezing, used in selected skin
lesions, tumors, and other conditions.

\medterm{Cryotherapy} Treatment using controlled cold to destroy or reduce abnormal tissue. It may treat selected skin lesions, tumors, eye problems, or prostate disease; pain, blistering, scarring, and damage to nearby tissue are possible.

\medterm{Cryotherapy For Prostate Cancer} A treatment that destroys prostate tissue by repeated freezing and thawing, used selectively for localized or recurrent prostate cancer.

\medterm{Cryotherapy For The Skin} Destruction of selected skin lesions by controlled freezing, usually with liquid nitrogen, for conditions such as warts, actinic keratoses, or some superficial tumors.

\medterm{Crypt} A small cavity, gland, or recessed space in body tissue, such as an intestinal crypt or tonsillar crypt.

\medterm{Cryptitis} Inflammation of a crypt or small anatomical recess, especially inflammation of the anal
glands and crypts.

\medterm{Crypto} Clinical shorthand that may refer to cryptococcosis or Cryptosporidium infection; the intended meaning depends on the clinical context.

\medterm{Cryptococcal Meningitis} Infection of the meninges by Cryptococcus, usually in people with impaired immunity, causing
headache, fever, altered mental status, or other neurologic symptoms.

\medterm{Cryptococcosis} Cryptococcosis is fungal infection caused by Cryptococcus, often affecting the lungs or brain in people with immune suppression.

\synonyms
Infection Cryptococcus (Cryptococcosis)

\medterm{Cryptococcus Neoformans} An encapsulated yeast that can cause pulmonary or central nervous system infection, especially in
people with impaired immunity.

\medterm{Cryptogenic Organizing Pneumonia} An inflammatory lung illness in which healing tissue blocks small airways and air sacs, without a known cause. It may cause cough, fever, and breathlessness and often improves with specialist-directed treatment.

\medterm{Cryptorchidism} Failure of one or both testicles to move into the scrotum before birth. An undescended testicle has higher risks of infertility, twisting, injury, and cancer and should be assessed during infancy.

\medterm{Cryptosporidiosis} An intestinal infection caused by Cryptosporidium parasites, spread through contaminated water, food, animals, or people. It usually causes watery diarrhea and cramps for days to weeks and can be severe or prolonged with weakened immunity.

\medterm{Cryptosporidium Enteritis} Diarrheal intestinal infection caused by Cryptosporidium parasites, transmitted through contaminated water, food, or contact; illness can be prolonged or severe in immunocompromised people.

\medterm{Crystal Arthropathy} Joint disease caused by crystal deposition. Gout (monosodium urate) and pseudogout
(calcium pyrophosphate). Diagnosis by joint aspiration and polarized microscopy.

\medterm{Crystal-Induced Synovitis} Acute or chronic joint inflammation caused by deposition of crystals such as monosodium urate, calcium pyrophosphate, or basic calcium phosphate in or around the joint.

\medterm{Crystalline Lens} The clear, flexible structure inside the eye that focuses light onto the retina. It changes shape for near vision and becomes cloudy in cataracts, causing blurred or reduced vision.

\medterm{Crystalloids} Aqueous solutions of mineral salts that distribute between intravascular and
interstitial compartments. Normal saline (0.9 percent NaCl) has 154 mEq/L each of sodium
and chloride, pH 5.4, and is associated with hyperchloremic metabolic acidosis at large
volumes.

\medterm{Crystalluria} The presence of crystals in urine, which may reflect precipitation of salts, medication
effects, metabolic disease, or a tendency to form urinary stones.

\medterm{CSF Analysis} Alternate index label; see \textit{Cerebrospinal fluid analysis}.

\medterm{CSF Cell Count} Laboratory measurement of red and white blood cells in cerebrospinal fluid, used to evaluate meningitis, encephalitis, subarachnoid hemorrhage, inflammation, malignancy, and other central nervous system disorders.

\medterm{CSF Coccidioides Complement Fixation Test} A cerebrospinal-fluid test for complement-fixing antibodies to Coccidioides, used with other findings to evaluate suspected coccidioidal meningitis.

\medterm{CSF Glucose Test} A laboratory measurement of glucose in cerebrospinal fluid, usually interpreted with a simultaneous blood-glucose value. Low cerebrospinal-fluid glucose may occur in bacterial, fungal, or tuberculous meningitis and other disorders, but results require the complete clinical and cerebrospinal-fluid profile.

\medterm{CSF Leak} Escape of cerebrospinal fluid through a defect in the skull base or spinal dura, causing positional headache, clear nasal or ear drainage, wound leakage, meningitis risk, or intracranial hypotension.

\medterm{CSF Leak (Cerebrospinal Fluid Leak)} Escape of cerebrospinal fluid through a tear in the membranes surrounding the brain or spinal cord, often causing positional headache.

\synonyms
Spinal Fluid Leak

\medterm{CSF Myelin Basic Protein} A test measuring a protein released when the insulating covering of nerve fibers is damaged. An increased level may support evaluation of multiple sclerosis or other nervous-system injury but is not specific by itself.

\medterm{CSF Oligoclonal Banding} An electrophoretic test for restricted immunoglobulin bands in cerebrospinal fluid, interpreted with serum bands to detect intrathecal antibody production such as in multiple sclerosis.

\medterm{CSF Rhinorrhea} Cerebrospinal fluid rhinorrhea is the leakage of cerebrospinal fluid from the
subarachnoid space through a defect in the skull base into the nasal cavity or paranasal
sinuses, resulting in a clear, watery nasal discharge that may be mistaken for allergic
rhinitis or sinusitis.

\medterm{CSF Smear} Microscopic examination of cerebrospinal fluid for cells, organisms, or abnormal material, supporting evaluation of meningitis, hemorrhage, malignancy, or inflammation.

\medterm{CSF Total Protein} A measurement of protein in cerebrospinal fluid, usually obtained by lumbar puncture. High or low levels may support evaluation of infection, bleeding, inflammation, nerve disorders, or other brain and spinal-cord disease.

\medterm{CSF-VDRL Test} A cerebrospinal-fluid test for reagin antibodies associated with Treponema pallidum infection; a reactive result is highly specific for neurosyphilis, but a nonreactive result does not exclude it.

\medterm{CT} Computed tomography, an imaging method that uses multiple X-ray measurements and computer reconstruction to produce cross-sectional images.

\synonyms
Computed tomography; CAT scan.

\medterm{CT Angiography} A CT scan performed after injected contrast highlights blood vessels. It can show narrowing, blockage, aneurysm, dissection, bleeding, or abnormal vessel connections in areas such as the brain, heart, chest, abdomen, or limbs.

\medterm{CT Angiography - Abdomen And Pelvis} Computed tomography angiography of the abdominal and pelvic blood vessels after intravenous contrast, used to assess aneurysm, stenosis, dissection, occlusion, bleeding, and vascular anatomy.

\medterm{CT Angiography - Arms And Legs} Computed tomography angiography of the peripheral arteries after intravenous contrast, used to evaluate arterial narrowing, blockage, aneurysm, injury, or other vascular abnormalities in the limbs.

\medterm{CT Angiography - Head And Neck} Computed tomography angiography of the blood vessels supplying the brain and neck after intravenous contrast, used to evaluate aneurysm, stenosis, dissection, occlusion, and vascular malformation.

\medterm{CT Angiography – Chest} Computed tomography angiography of the chest uses intravenous contrast and computed tomography to visualize the thoracic aorta, pulmonary arteries, or other chest vessels. It can evaluate embolism, aneurysm, dissection, bleeding, and vascular malformation.

\medterm{CT Chest} A CT scan of the chest that shows the lungs, heart, blood vessels, mediastinum, pleura, ribs, and chest wall. It can evaluate infection, tumors, injury, fluid, blood clots, and other abnormalities.
\medterm{CT Colonography} A CT examination of the colon after bowel cleansing and gentle inflation with air or carbon dioxide. It can screen for polyps or cancer when colonoscopy is incomplete or unsuitable, but abnormal findings usually require colonoscopy.

\medterm{CT Colonography (Virtual Colonoscopy)} A CT scan of the colon after bowel cleansing and gentle inflation with air or carbon dioxide. It can look for polyps or cancer when colonoscopy is incomplete or unsuitable, but an abnormal result usually requires colonoscopy for tissue removal or biopsy.

\synonyms
Virtual Colonoscopy

\medterm{CT Contrast} A substance administered during computed tomography to make selected tissues, vessels, or body spaces more visible.

\medterm{CT Coronary Angiography} A CT scan taken after injected contrast outlines the arteries supplying the heart. It can show plaque or narrowing and is a non-invasive option for selected people with possible coronary artery disease.

\medterm{CT Dose-Length Product} A measurement of the X-ray output from a CT examination that combines scan dose with scan length. It helps compare and optimize imaging protocols but does not directly predict an individual's radiation risk.

\medterm{CT Enterography} A CT scan of the small intestine after drinking contrast to expand the bowel and receiving injected contrast. It helps evaluate Crohn disease, inflammation, tumors, obstruction, and unexplained intestinal bleeding.

\medterm{CT Perfusion} A CT technique that estimates blood flow and related perfusion parameters in tissue, commonly to assess selected cases of acute stroke or tumors.

\medterm{CT Pulmonary Angiography} CT imaging performed after intravenous contrast to show the pulmonary arteries and evaluate suspected pulmonary embolism or other vascular abnormalities.

\medterm{CT Scan} Computed tomography, an imaging method that uses X-rays and computer processing to produce cross-sectional
images of the body.

\medterm{CT Scan (Computed Tomography)} An X-ray examination processed by a computer to create cross-sectional body images. Scans may be done without contrast or with contrast through a vein or mouth, depending on what needs to be assessed.

\synonyms
Computed Tomography

\medterm{Ctenocephalides FelFlea Bite (Flea Bites In Humans)} Alternate index label for Flea Bites in Humans.

\medterm{CTL} Cytotoxic T lymphocyte, an immune cell that recognizes and kills infected, abnormal, or target cells through antigen-specific mechanisms.

\medterm{CTPA} Alternate name; see \textit{Computed Tomography Pulmonary Angiography}.

\medterm{CTS} Commonly, carpal tunnel syndrome: compression of the median nerve at the wrist causing numbness, tingling, pain, or weakness in the hand.

\medterm{Cubital} Relating to the elbow or nearby forearm. Cubital structures include the joint, bones, muscles, nerves, and blood vessels, and problems there may cause pain, restricted movement, numbness, or weakness.

\medterm{Cubital Tunnel Syndrome} Pressure or stretching of the ulnar nerve at the elbow, causing tingling or numbness in the little and ring fingers, hand weakness, or pain. Symptoms often worsen when the elbow stays bent.

\medterm{Cubitus} The elbow or the forearm; the term is also used in names of structures and conditions around the elbow.

\medterm{Cubitus Valgus} An outward angle of the forearm at the elbow that is larger than usual when the arm is straight. It may be present from birth or develop after an elbow injury. In some people, it can put pressure on the ulnar nerve and cause tingling, numbness, or weakness in the ring and little fingers.

\medterm{Cubitus Varus} An abnormal inward angulation of the forearm relative to the upper arm when the elbow is
extended, often developing after a childhood elbow fracture.

\medterm{Cuboid} A cube-shaped bone on the outer side of the midfoot. It helps form the foot arch and connects several foot bones; injury or displacement can cause outer-foot pain, especially during walking.

\medterm{Cul-De-Sac} A blind-ending pouch or recess; in pelvic anatomy, the rectouterine pouch and rectovesical pouch are examples.

\medterm{Culdocentesis} A procedure that places a needle through the back of the vagina to sample or drain fluid from the lowest part of the pelvis. Ultrasound and less-invasive tests have replaced it in many situations.

\medterm{Cullen'S Sign} Bluish or violaceous periumbilical ecchymosis caused by blood tracking through the
fascial planes to the umbilicus. It is a sign of intra-abdominal or retroperitoneal
haemorrhage, classically severe haemorrhagic pancreatitis, and is associated with
serious underlying disease.

\medterm{Cultural Psychiatry} The study and clinical practice of how cultural factors influence mental health, illness
expression, diagnosis, and treatment. Cultural concepts of distress, idioms of distress,
and cultural syndromes must be considered. The DSM-5 includes a Cultural Formulation
Interview.

\medterm{Culture} A laboratory test in which a clinical specimen is placed in conditions that allow microorganisms to grow and be identified; susceptibility testing may help guide treatment.

\medterm{Culture - Colonic Tissue} A laboratory culture of a colon tissue specimen to detect bacteria, fungi, or other microorganisms in suspected infection; results are interpreted with histopathology and clinical findings.

\medterm{Culture - Duodenal Tissue} A laboratory culture of a duodenal tissue specimen to identify microorganisms causing infection or inflammation, interpreted with endoscopic findings and histopathology.

\medterm{Culture And Sensitivity} Microbiological culture of a clinical specimen (blood, urine, sputum, wound swab, CSF)
to identify pathogens and determine antimicrobial susceptibility. Culture media: blood
agar (most bacteria), MacConkey (Gram-negative), chocolate agar (fastidious: N.
gonorrhoeae, H. influenzae), and selective media.

\medterm{Culture Media} Nutrient materials prepared to grow microorganisms in a laboratory, with different formulas chosen for specific organisms, tests, or identification purposes.

\medterm{Culture Swab} A sterile swab used to collect material from skin, wounds, the throat, genital tract, or another site for laboratory culture. The test may identify bacteria or fungi and guide antibiotic choice; viral testing usually uses separate molecular or antigen tests.

\medterm{Culture-Negative Endocarditis} Infective endocarditis in which routine blood cultures do not identify an organism, often because of prior antibiotics, fastidious bacteria, fungi, or noninfectious mimics; diagnosis requires integrated clinical, imaging, and laboratory assessment.

\medterm{Cumulative Dose} The total amount of a substance or radiation received over time. Cumulative exposure matters for medicines, toxins, radiation, carcinogens, and pregnancy-related risks because effects may build gradually.

\medterm{Cumulative Trauma Disorder} Pain or injury that develops after repeated forceful, awkward, or prolonged movements. It may affect muscles, tendons, nerves, or other soft tissues in the hands, arms, neck, shoulders, or back.

\medterm{Cup-To-Disc Ratio} A comparison between the central hollow of the optic nerve and the whole visible optic nerve head. An enlarged or changing ratio may suggest glaucoma but must be interpreted with eye pressure and other findings.

\medterm{Cupping} Enlargement or deepening of the central depression in the optic disc, assessed during glaucoma evaluation and interpreted with eye pressure and other findings.

\synonyms
Optic-disc cupping; glaucomatous cupping.

\medterm{CURB-65 Score} A five-point tool estimating the severity of community-acquired pneumonia using confusion, blood urea, breathing rate, blood pressure, and age 65 or older. It helps guide hospital decisions but does not replace clinical judgment.

\medterm{Cure} Elimination or durable control of a disease so that it no longer causes active illness; the meaning and evidence for cure vary by condition.

\medterm{Curettage} Removal of tissue from a body cavity or surface with a curette, used for diagnosis,
treatment, or removal of abnormal tissue.

\medterm{Curettage Specimen} Tissue scraped or scooped from a body surface, cavity, or bone lesion for laboratory examination. Because the pieces are often small and fragmented, they may not show the full structure or depth of disease.

\medterm{Curette} A spoon-shaped or ring-shaped surgical instrument used to scrape, remove, or clean tissue from a cavity, bone, skin lesion, or tooth-supporting area. Its use and risks depend on the procedure and body site.

\medterm{Curling Ulcer} An acute ulcer in the stomach or duodenum that can develop after severe burns. Reduced blood flow and major physical stress weaken the digestive lining and may cause dangerous bleeding.

\medterm{Curling'S Ulcer} An acute peptic ulcer associated with severe burns, often occurring in the duodenum or
stomach during physiologic stress.

\medterm{Curvature Of The Penis} A bend of the erect or flaccid penis caused by congenital tissue asymmetry, scar tissue such as Peyronie disease, trauma, or other conditions; it may impair intercourse or cause pain.

\synonyms
Peyronie's Disease (Curvature Of The Penis)

\medterm{Curvularia Alcornii} A dark-pigmented environmental fungus that rarely causes infection. It may affect the skin, sinuses, or brain, mainly in people with weakened immunity, and requires specialist laboratory identification and treatment.

\medterm{Cushing Disease} Cushing syndrome caused by an adrenocorticotropic hormone-secreting pituitary tumor.
Excess cortisol can cause central weight gain, rounded facial features, skin changes, muscle weakness, high blood
pressure, high blood glucose, and menstrual or mood changes.

\medterm{Cushing Syndrome} Illness caused by prolonged exposure to too much cortisol or steroid medicine. It can cause central weight gain, easy bruising, muscle weakness, high blood pressure, diabetes, bone loss, infections, and mood or menstrual changes.

\medterm{Cushing Syndrome Complications} Untreated Cushing syndrome carries significant morbidity and mortality from
cardiovascular disease, infections, thromboembolic events, and metabolic derangements.
Cardiovascular complications include accelerated atherosclerosis, hypertension, and
dyslipidemia. Infection risk is increased from immunosuppression.

\medterm{Cushing Syndrome Due To Adrenal Tumor} Excess cortisol caused by an adrenal tumor, producing Cushing syndrome with features such as central weight gain, hypertension, diabetes, muscle weakness, and skin changes.

\medterm{Cushing Ulcer} An acute stomach or duodenal ulcer associated with severe brain disease, head injury, or neurosurgery. Excess nerve stimulation can increase acid production and cause abdominal pain or gastrointestinal bleeding.

\medterm{Cushing'S Disease} A pituitary tumor that makes too much adrenocorticotropic hormone (ACTH), causing the adrenal glands to produce too much cortisol. Features may include weight gain around the trunk, a round face, easy bruising, purple stretch marks, muscle weakness, high blood pressure, high blood glucose, and menstrual or mood changes.

\medterm{Cushing'S Syndrome} A disorder caused by prolonged exposure to too much cortisol or steroid medicine. It may cause weight gain around the trunk, a round face, fat at the back of the neck, thin easily bruised skin, purple stretch marks, muscle weakness, high blood pressure, and high blood glucose. Untreated disease can lead to infections, blood clots, bone loss, and heart or blood-vessel disease.

\synonyms
Hypercortisolism; cortisol excess syndrome.

\medterm{Cushingoid} Resembling the physical features of excess glucocorticoid exposure, such as a round face, central weight gain, thin skin, easy bruising, and muscle weakness.

\medterm{Cusp} A pointed projection, such as a projection of a tooth or one of the leaflets of a heart valve.

\medterm{Cusp Fracture} A break in one of the raised points of a tooth, often around a large filling or weakened tooth. It may cause pain when biting or sensitivity and may require a bonded repair, crown, or other dental treatment.
\medterm{Cut} A break or incision made in tissue, usually by a sharp object; the depth, contamination, and location determine the required care.

\medterm{Cutaneous} Relating to the skin. The skin forms a protective barrier, helps control temperature, senses touch and pain, and supports vitamin D production. Skin findings can provide clues to infections, immune diseases, hormone disorders, and skin cancers.

\medterm{Cutaneous Angiosarcoma} A fast-growing cancer of blood-vessel cells in the skin, usually on the scalp or face of older adults. It may look like a bruise, ulcerate, and spread.

\medterm{Cutaneous Anthrax} A skin infection caused by Bacillus anthracis, usually producing a painless black sore surrounded by swelling.

\medterm{Cutaneous B-Cell Lymphoma} A non-Hodgkin lymphoma mainly affecting B lymphocytes in the skin. It may cause persistent patches, plaques, or lumps and is treated according to subtype, extent, symptoms, and spread.

\medterm{Cutaneous Diphtheria} A skin infection caused by toxigenic or nontoxigenic Corynebacterium diphtheriae, producing chronic
ulcers or a gray membrane.

\medterm{Cutaneous Langerhans Cell Histiocytosis} An uncommon disorder in which abnormal immune cells build up in the skin. It may cause scaly or crusted spots, rash, or purple marks; testing checks whether other organs are involved.

\medterm{Cutaneous Larva Migrans} An intensely itchy, winding skin rash caused by hookworm larvae from contaminated soil or sand entering bare skin. The larvae move beneath the skin and usually die without reaching internal organs.

\medterm{Cutaneous Leiomyoma} A usually noncancerous smooth-muscle tumor in the skin, often causing firm, painful bumps. Multiple lesions may be associated with an inherited condition that also increases kidney-cancer risk and needs specialist assessment.

\medterm{Cutaneous Leishmaniasis} A skin infection caused by Leishmania parasites transmitted by sandflies, usually producing one
or more chronic ulcers.

\medterm{Cutaneous Lupus} A form of lupus erythematosus limited to the skin, manifesting as discoid lesions,
subacute cutaneous plaques, or acute malar erythema. It may occur in isolation or as a
feature of systemic lupus erythematosus.

\medterm{Cutaneous Lupus Erythematosus} An immune disorder causing inflamed skin patches or sores that may worsen with sunlight, scar, or change skin color. It may remain limited to skin or occur with systemic lupus and requires appropriate evaluation.

\medterm{Cutaneous Lymphoma} Cutaneous lymphomas are heterogeneous group of T-cell and B-cell malignancies primarily
involving the skin. Primary cutaneous B-cell lymphomas include marginal zone lymphoma,
follicle center lymphoma, and diffuse large B-cell lymphoma leg type. Primary cutaneous
T-cell lymphomas include mycosis fungoides and Sézary syndrome.

\medterm{Cutaneous Metastases} Secondary deposits of cancer in the skin, originating from a primary tumor elsewhere in
the body. They may present as nodules, plaques, or ulcers, and often indicate advanced
or recurrent disease.

\medterm{Cutaneous Metastasis} A deposit of malignant cells in the skin from a primary tumor elsewhere or, less often, another cutaneous site; it may present as a firm nodule, plaque, or inflammatory-appearing lesion.

\medterm{Cutaneous Polyarteritis Nodosa} A skin-limited necrotizing vasculitis of small- and medium-sized arteries, typically causing tender subcutaneous nodules, livedo reticularis, and ulcers on the lower legs without systemic organ involvement.

\medterm{Cutaneous Skin Tag} A small, soft, usually harmless growth of extra skin, often on a narrow stalk in the neck, armpit, groin, or another skin fold. It may catch on clothing or become irritated.
\medterm{Cutaneous Small-Vessel Vasculitis} Inflammation and destruction of small dermal blood vessels, commonly presenting as
palpable purpura on dependent areas. Causes include infections, medications, autoimmune
disease, and systemic vasculitis; skin biopsy with appropriate specimen handling is
central to diagnosis.

\medterm{Cutaneous T-Cell Lymphoma} A malignancy of skin-homing T lymphocytes. Mycosis fungoides is the most common subtype,
presenting as patches, plaques, and tumors in a clothing distribution. Sézary syndrome
is the leukemic variant with erythroderma, lymphadenopathy, and circulating malignant
cells.

\medterm{Cuticle} The thin outermost layer of a hair shaft formed by overlapping cells that protect the inner cortex.

\synonyms
Hair cuticle; cuticular layer.

\medterm{Cuticle Remover Poisoning} Exposure to cuticle-remover products, which may contain alkaline or acidic chemicals that can injure the skin, eyes, mouth, or gastrointestinal tract.

\medterm{Cutis Anserina} Gooseflesh or piloerection, in which contraction of tiny muscles around hair follicles makes the skin appear rough or dimpled.

\medterm{Cutis Laxa} A group of inherited or acquired connective-tissue disorders characterized by loose, inelastic skin caused by abnormal elastic fibers. Some forms also affect the lungs, blood vessels, heart, gastrointestinal tract, or skeleton.

\medterm{Cuts And Puncture Wounds} Breaks in the skin caused by sharp, blunt, or penetrating trauma; evaluation considers depth, contamination, retained foreign bodies, tendon or nerve injury, vascular damage, and tetanus protection.

\medterm{Cuts Wounds (Cuts, Scrapes And Puncture Wounds)} Alternate index label for Cuts, Scrapes and Puncture Wounds.

\medterm{Cuts, Scrapes And Puncture Wounds} Cuts, scrapes, and puncture wounds damage the skin to varying depths and can introduce bacteria or foreign material. Clean minor wounds with running water, control bleeding, cover them, and seek care for deep or contaminated wounds, bites, impaired circulation, infection, or uncertain tetanus protection.

\synonyms
Cuts Wounds (Cuts, Scrapes And Puncture Wounds)

\medterm{Cutting Balloon (Vascular)} A balloon catheter with tiny blades or scoring surfaces that make controlled cuts in hardened plaque while the balloon widens a narrowed blood vessel. It is used in selected vascular procedures.

\synonyms
Vascular

\medterm{Cutting For The Stone} An older term for surgically opening the urinary tract to remove a stone, historically called lithotomy.

\medterm{CVC} Central venous catheter, a catheter whose tip lies in a large central vein and provides intravenous access for medications, fluids, nutrition, or monitoring.

\medterm{Cyanide} A highly poisonous chemical that prevents cells from using oxygen. Exposure can rapidly cause headache, confusion, seizures, breathing failure, cardiovascular collapse, and death.

\medterm{Cyanide Poisoning} A medical emergency in which cyanide stops cells from using oxygen. It may follow smoke inhalation or industrial, accidental, or intentional exposure and can cause headache, confusion, seizures, collapse, and death within minutes.

\medterm{Cyanoacrylates Poisoning} Exposure to cyanoacrylate adhesives, commonly known as instant glue. Contact may bond tissues or irritate the skin, eyes, or airways; ingestion usually causes local adhesion or irritation.

\medterm{Cyanosis} A blue or gray-blue color of the skin, lips, or moist linings caused by too little oxygen in blood or abnormal blood flow. Sudden cyanosis, especially with breathing difficulty, requires emergency care.

\medterm{Cyanotic} Having a bluish or slate-gray discoloration of skin or mucous membranes, usually reflecting increased deoxygenated hemoglobin or abnormal blood flow.

\medterm{Cyanotic Heart Disease} A congenital or acquired heart disorder that reduces systemic oxygen saturation and causes cyanosis, often through right-to-left shunting or inadequate pulmonary blood flow.

\medterm{Cycle Time} Time for complete image acquisition. PET: minutes. SPECT: 20-30 minutes.

\medterm{Cyclic Citrullinated Peptide} A substance used in a blood test for antibodies associated with rheumatoid arthritis. A positive result supports the diagnosis but must be interpreted with symptoms, examination, and other tests.

\medterm{Cyclic Guanosine Monophosphate} A chemical messenger inside cells that helps regulate blood-vessel widening, smooth-muscle relaxation, vision, and other functions. Its level can be changed by some medicines.

\synonyms
cGMP; cyclic GMP.

\medterm{Cyclic Hormone Therapy} Hormone therapy in which estrogen is given continuously or for part of a cycle and progestogen is added for a defined number of days.

\synonyms
Sequential hormone therapy; cyclic estrogen-progestogen therapy.

\medterm{Cyclic Neutropenia} Repeated episodes of low neutrophils, the white blood cells that fight many infections. Episodes often recur at regular intervals and may cause fever, mouth ulcers, sore throat, or infections.

\medterm{Cyclic Vomiting Syndrome} Repeated, similar attacks of severe nausea and vomiting separated by periods of normal health. It is often related to migraine biology; dehydration and other causes must be considered during evaluation.

\medterm{Cyclooxygenase} An enzyme family that helps convert arachidonic acid into prostaglandins and related mediators involved in pain, inflammation, and platelet function.

\synonyms
COX enzyme; prostaglandin-endoperoxide synthase.

\medterm{Cyclooxygenase-2} An inducible cyclooxygenase enzyme that generates prostaglandins involved in inflammation, pain, and fever and is inhibited by COX-2-selective medicines.

\medterm{Cyclops} A term that may refer to cyclopia, a rare severe congenital malformation with a single central eye, or to a fibrous nodule causing loss of extension after knee surgery; context determines the meaning.


\medterm{Cyclospora Infection} An intestinal infection caused by Cyclospora parasites, usually acquired from contaminated food or water. It can cause prolonged or recurring watery diarrhea, cramps, nausea, tiredness, and weight loss.

\medterm{Cyclospora Infection (Cyclosporiasis)} Alternate index label for Cyclosporiasis.

\medterm{Cyclosporiasis} An intestinal infection caused by the parasite Cyclospora cayetanensis, usually acquired by consuming contaminated food or water. It commonly causes prolonged or relapsing watery diarrhea, abdominal cramps, nausea, fatigue, loss of appetite, and weight loss; diagnosis requires specialized stool testing, and trimethoprim-sulfamethoxazole is the usual effective treatment when appropriate.

\synonyms
Cyclospora Infection (Cyclosporiasis)

\medterm{Cyclosporiasis (Cyclospora)} An intestinal infection caused by Cyclospora parasites, typically producing prolonged watery diarrhea, cramps, nausea, and fatigue.

\medterm{Cyclothymia} A long-term mood disorder with repeated periods of hypomanic and depressive symptoms that do not meet criteria for full episodes.

\medterm{Cyclothymia (Cyclothymic Disorder)} A chronic mood disorder involving repeated periods of hypomanic and depressive symptoms that do not meet criteria for full episodes.

\medterm{Cyclothymic Disorder} A long-term mood disorder with repeated periods of elevated energy or activity and periods of depressive symptoms that are milder than full mania or major depression. Symptoms persist for years and can disrupt relationships or work.

\medterm{Cylindroma} A usually benign adnexal tumor composed of basaloid cells, often presenting as a firm pink or skin-coloured nodule on the scalp or face; multiple lesions may occur in Brooke-Spiegler syndrome.

\medterm{CYP2C19 Polymorphism} CYP2C19 is a cytochrome P450 enzyme involved in the metabolism of approximately 10
percent of drugs, including the antiplatelet agent clopidogrel, proton pump inhibitors,
some antidepressants, and the antiepileptic drug phenytoin.

\medterm{CYP2D6 Polymorphism} CYP2D6 is a cytochrome P450 enzyme responsible for the metabolism of approximately 25
percent of commonly prescribed drugs, including codeine, tamoxifen, many
antidepressants, and antipsychotics.

\medterm{CYP450 System} A family of enzymes, found mainly in the liver, that breaks down many medicines and foreign chemicals. Genetic differences and other medicines can speed up or slow down this process and alter drug effects.

\medterm{Cyproheptadine Overdose} Overdose of cyproheptadine, an antihistamine, that may cause sedation, agitation, anticholinergic effects, seizures, respiratory depression, or cardiac effects.

\medterm{Cyst} A closed sac containing fluid, air, cells, or semisolid material. Cysts can occur in the skin, organs, glands, or other tissues and may be harmless, infected, or related to another disease.

\medterm{Cyst Eyelid (Giardia Lamblia)} Alternate index label for Giardia Lamblia.

\medterm{Cyst, Bartholin'S} Alternate index label; see \textit{Bartholin's cyst}.

\medterm{Cyst, Ganglion} Alternate index label; see \textit{Ganglion cyst}.

\medterm{Cyst, Kidney} Alternate index label; see \textit{Kidney cysts}.

\medterm{Cyst, Ovarian} Alternate index label; see \textit{Ovarian cysts}.

\medterm{Cyst, Pancreatic} Alternate index label; see \textit{Pancreatic cysts}.

\medterm{Cyst, Pilonidal} Alternate index label; see \textit{Pilonidal cyst}.

\medterm{Cyst, Spermatic} Alternate index label; see \textit{Spermatocele}.

\medterm{Cystadenoma} A usually noncancerous gland-forming tumor containing fluid-filled spaces or cyst-like structures. It can occur in several organs, may enlarge or cause symptoms, and sometimes requires removal or monitoring.

\medterm{Cystic Acne} Severe acne with deep, painful lumps that can leave permanent scars. It often needs clinician-directed treatment; squeezing or picking the lesions increases scarring and infection risk.

\medterm{Cystic Echinococcosis} A parasitic infection in which slowly growing cysts form, most often in the liver or lungs, after swallowing Echinococcus eggs from infected dogs or contaminated food, water, or soil. Cysts may cause pain, cough, blockage, or serious allergic reactions if they rupture.

\medterm{Cystic Fibrosis} An inherited condition that makes mucus unusually thick and sticky, affecting the lungs and digestive system and sometimes the liver and reproductive organs. It can cause repeated lung infections, cough, poor digestion, poor growth, or infertility; treatment includes airway clearance, nutrition, enzymes, antibiotics, and selected targeted medicines.

\synonyms
CF

\medterm{Cystic Fibrosis - Nutrition} Nutritional care for cystic fibrosis, often requiring extra calories, salt, fat-soluble vitamins, and pancreatic enzymes when digestion is impaired. Growth, weight, stool changes, and blood tests guide the plan.
\medterm{Cystic Hygroma} A birth defect involving an abnormal collection of lymph vessels, usually forming a soft fluid-filled swelling in the neck. It may grow, become infected, or press on the airway and requires specialist assessment.

\medterm{Cystic Medial Necrosis} Weakening of the middle layer of the aorta caused by loss of supportive cells and elastic fibers. It can contribute to an aortic aneurysm or dissection and may occur with inherited connective-tissue disorders or aging.

\medterm{Cysticercosis} Cysticercosis is tissue infection by larval Taenia solium after swallowing eggs, potentially causing cysts in muscle, skin, eyes, or brain.

\medterm{Cystine And Triple Phosphate (Struvite) Crystals} Crystalline substances that may be found in urine sediment. Hexagonal cystine crystals
support cystinuria in the appropriate setting, while struvite crystals may occur in
infection-related alkaline urine; microscopy is interpreted with clinical and laboratory
findings.

\synonyms
struvite

\medterm{Cystine Stones} Kidney stones formed when inherited cystinuria causes excess cystine in urine. They often recur and may be difficult to dissolve; abundant fluids, medical treatment, and sometimes procedures help prevent or remove them.

\medterm{Cystinuria} An inherited disorder in which the kidneys release too much cystine into urine, allowing recurrent cystine kidney stones to form. It requires long-term fluid intake, monitoring, and sometimes medicines to prevent stones.
\medterm{Cystitis} See \textit{Bladder Infection}.

\synonyms
Bladder Infection

\medterm{Cystitis (Bladder Infection)} See \textit{Bladder Infection}.

\medterm{Cystitis - Acute} Acute inflammation of the urinary bladder, most often caused by bacterial infection and producing dysuria, urinary frequency, urgency, and suprapubic discomfort.

\medterm{Cystitis - Noninfectious} Inflammation of the urinary bladder without a proven bacterial infection. Causes include interstitial cystitis or bladder pain syndrome, medicines, radiation, chemical irritation, autoimmune disease, and injury; symptoms may include frequency, urgency, and painful urination.

\medterm{Cystitis, Interstitial} Alternate index label; see \textit{Interstitial cystitis}.

\medterm{Cystocele} Bulging of the bladder into the front wall of the vagina because supporting tissues have weakened. It may cause pressure, a vaginal bulge, urinary leakage, difficulty emptying the bladder, or no symptoms.

\medterm{Cystography} X-ray or other imaging of the urinary bladder after contrast is placed in it. It can show bladder injury, leaks, fistulas, structural abnormalities, or urine flowing backward toward the kidneys.

\medterm{Cystoid Macular Edema} Swelling at the center of the retina caused by fluid collecting in small cyst-like spaces. It can blur or distort central vision and may follow eye surgery, inflammation, diabetes, or retinal blood-vessel disease.
\medterm{Cystometric Study} A bladder test that measures pressure, sensation, and urine storage while the bladder fills and empties. It helps investigate urgency, leakage, difficulty emptying, weak bladder muscles, or nerve-related urinary problems.

\medterm{Cystometrogram} A recording of bladder pressure and sensation during filling and emptying. It is part of urodynamic testing used to investigate leakage, urgency, blockage, weak contractions, or nerve-related bladder problems.

\medterm{Cystometry} A bladder-function test measuring pressure and sensation while the bladder fills and empties. It helps assess urgency, leakage, obstruction, weak contractions, nerve problems, or other urinary difficulties.

\medterm{Cystoscope} A thin, lighted instrument used to examine the urethra and bladder and, when needed, perform procedures.

\medterm{Cystoscopy} Cystoscopy uses a camera passed through the urethra to inspect the bladder and urethra and may allow biopsy or treatment.

\medterm{Cystostomy} Surgical creation of an opening from the urinary bladder to the skin or another
structure to permit drainage of urine.

\medterm{Cystotome} A surgical instrument used to make an incision in the capsule of the lens during selected eye procedures.

\medterm{Cystourethrography} Radiographic imaging of the urinary bladder and urethra, usually during filling and
voiding after contrast has been introduced.

\medterm{Cysts} Closed sac-like structures containing fluid, air, semisolid material, or other contents; their causes range from benign developmental or obstructive processes to infection or neoplasia.

\medterm{Cysts Choledochal (Choledochal Cysts)} Alternate index label for Choledochal Cysts.

\medterm{Cysts Ovary (Ovarian Cysts)} Alternate index label for Ovarian Cysts.

\medterm{Cysts Pancreatic Inflammatory (Pancreatic Cysts)} Alternate index label for Pancreatic Cysts.

\medterm{Cysts True Pancreas (Pancreatic Cysts)} Alternate index label for Pancreatic Cysts.

\medterm{Cysts, Epidermoid} Alternate index label; see \textit{Epidermoid cysts}.

\medterm{Cysts, Sebaceous} Alternate index label; see \textit{Epidermoid cysts}.

\medterm{Cytochrome B5 Reductase} An enzyme that maintains hemoglobin iron in the ferrous state needed for oxygen binding; inherited deficiency can cause congenital methemoglobinemia, usually presenting with cyanosis and impaired oxygen delivery.

\medterm{Cytochrome C} A protein inside mitochondria that carries electrons during energy production. When released from damaged mitochondria, it can also help start the controlled death of a cell.

\medterm{Cytochrome P450} A family of enzymes, found mainly in liver cells, that chemically changes many medicines, toxins, and hormones so the body can remove them. Genetic differences and drug interactions can alter their activity.

\medterm{Cytogenetics} The study of chromosomes and their number, structure, and arrangement. Tests such as karyotyping, fluorescence in situ hybridization, and chromosome microarray can help identify inherited conditions, leukemia, other cancers, and some fetal abnormalities.

\medterm{Cytogenetics (Karyotype)} Visualization of chromosomes under light microscopy from dividing cells (bone marrow,
blood, amniotic fluid, tissue). Staining: G-banding (Giemsa, light and dark bands).
Resolution: 400-550 bands per haploid genome.

\synonyms
Karyotype

\medterm{Cytokine} A small protein messenger that allows cells to communicate during immunity, inflammation, infection, healing, and blood-cell production. Cytokines can protect tissues, but excessive or poorly controlled activity can cause disease.

\medterm{Cytokine Release Syndrome} A potentially serious whole-body inflammatory reaction caused by a sudden release of immune signaling proteins, sometimes after immune-based cancer treatment or infection. Fever, low blood pressure, breathing difficulty, and organ problems may occur.

\medterm{Cytokine Storm} A severe, dysregulated systemic inflammatory response involving excessive cytokine
release that can cause tissue injury, organ dysfunction, and shock.

\medterm{Cytokines} Small signaling proteins released by immune and other cells that regulate inflammation, immunity, blood-cell production, and tissue responses.

\synonyms
Immune signaling proteins; cytokine mediators.

\medterm{Cytologic Evaluation} Microscopic examination of individual cells or small cell groups from a fluid, brushing, scraping, or fine-needle sample to detect infection, inflammation, or malignancy.

\medterm{Cytology} The microscopic study of individual cells, distinct from histology in that tissue
architecture is not examined, used for screening and diagnosis of neoplastic and
infectious disease. Specimens are exfoliative (cervical smears, sputum) or aspirated
(fine needle aspiration, body fluids), interpreted with Papanicolaou or special stains.

\medterm{Cytology Exam Of Pleural Fluid} Microscopic examination of fluid around the lung for malignant cells and other cellular changes, performed with chemical and microbiologic testing when evaluating an effusion.

\medterm{Cytology Exam Of Urine} Microscopic examination of shed urinary-tract cells, used mainly to detect high-grade urothelial malignancy and selected abnormal cells; a negative result does not exclude cancer.

\medterm{Cytomegalovirus} A herpesvirus that causes significant disease in immunocompromised patients including
transplant recipients and HIV-infected individuals. CMV infects a wide range of cell
types and establishes latent infections in myeloid progenitor cells. Primary infection
is usually asymptomatic in immunocompetent individuals.

\medterm{Cytomegalovirus Infection} A widespread herpesvirus infection. Primary infection in immunocompetent individuals is
usually asymptomatic or causes mononucleosis-like syndrome. Congenital CMV: most common
congenital infection, causes sensorineural hearing loss, microcephaly, and developmental
delay.

\medterm{Cytopathology} The examination of individual cells under a microscope to help diagnose cancer, infection, inflammation, or other disease. Cells may come from fluid, a brushing, a smear, or a needle sample.

\medterm{Cytopenia} A lower-than-normal number of one or more blood-cell types: red blood cells, white blood cells, or platelets. The effects depend on which cells are low and how severe the reduction is.

\medterm{Cytoplasm} The material inside a cell between the cell membrane and nucleus, consisting of cytosol, organelles, and other cellular components.

\medterm{Cytoreductive Surgery} Surgery that removes as much visible cancer as safely possible to reduce the amount of disease. It may be combined with other treatment, and its benefits and risks depend on the cancer's site and spread.

\medterm{Cytoskeleton} A network of protein fibers inside a cell that maintains its shape, moves materials, enables movement, and helps the cell divide. It constantly changes in response to the cell's needs.

\medterm{Cytotoxic} Toxic to cells or capable of killing cells, as with some chemotherapy drugs, toxins, or cytotoxic immune cells.

\medterm{Cytotoxic Alopecia} Hair loss caused by exposure to a cytotoxic treatment or other toxic injury to actively growing hair follicles.

\synonyms
Drug-induced alopecia; chemotherapy-related hair loss.

\medterm{Cytotoxic T Cell} An immune cell that recognizes and kills infected, abnormal, or cancerous cells. It identifies small protein signals displayed on the target cell and can trigger that cell to die.

\medterm{Cytotoxicity} The capacity of an immune effector cell or immune mediator to injure or kill a target
cell. T lymphocytes and natural killer cells commonly use perforin-granzyme pathways or
death-receptor signaling to induce target-cell apoptosis.
\medterm{Common Carotid Artery} One of two major arteries carrying blood from the heart toward the head and neck. Each divides into an internal branch supplying the brain and an external branch supplying the face and scalp.

\medterm{Camelid} Relating to camels, llamas, or alpacas. Camelid antibodies and antibody fragments are used in research and some diagnostic or therapeutic applications; contact with camelids can also be relevant to certain infections or allergies.

\medterm{Cancer of the Ear} Cancer arising in the outer ear, ear canal, middle ear, or nearby tissues. Persistent sores, bleeding, pain, discharge, hearing loss, or a growing lump require examination; treatment depends on the site and tumor type.

\medterm{Cancer Treatment} Medical, surgical, radiation, systemic, or supportive care used to control cancer, remove it, slow its growth, prevent recurrence, or relieve symptoms. Treatment is selected according to cancer type, stage, biomarkers, overall health, and patient goals.

\medterm{Cantu Syndrome} A rare genetic condition, usually caused by an ABCC9 or KCNJ8 gene change. It may cause excessive hair growth, distinctive facial features, loose skin, enlarged bones, skeletal differences, and an enlarged or abnormal heart.

\medterm{Cardiometabolic} Relating to the interaction between heart and blood-vessel disease and metabolic conditions such as obesity, insulin resistance, diabetes, and abnormal blood lipids.

\medterm{Carotid Artery Stenosis} Narrowing of an artery in the neck that supplies the brain, usually from fatty plaque. It may reduce blood flow or release a clot and increase the risk of a transient ischemic attack or stroke.

\medterm{Carotid Occlusion} Complete blockage of an artery supplying the brain, usually from plaque or a blood clot. Some people have no symptoms because other vessels compensate, but sudden weakness, speech trouble, or vision loss is an emergency.

\medterm{Cat Allergy} An immune reaction to proteins in cat saliva, skin flakes, or other secretions. It can cause sneezing, itchy eyes, wheezing, hives, or asthma symptoms; reducing exposure and using appropriate allergy treatment may help.

\medterm{Cat Scratch Disease} An infection caused by \textit{Bartonella henselae}, usually after a scratch or bite from a kitten or cat. It commonly causes a small skin lesion and nearby swollen lymph nodes and is usually self-limited, though serious complications can occur.

\medterm{Catamenial Epilepsy} A pattern in which seizures become more frequent or severe at certain times in the menstrual cycle. Hormone changes may contribute, and a seizure diary can help guide neurologic treatment.

\medterm{Celebrex} A brand name for celecoxib, a selective cyclooxygenase-2 nonsteroidal anti-inflammatory drug used for pain and inflammation. It can cause stomach, kidney, cardiovascular, allergic, or bleeding complications and should be used as directed.

\medterm{Cell Signaling} The process by which cells detect information through receptors and transmit signals inside the cell to change functions such as growth, movement, metabolism, or immune activity.

\medterm{Central Pain Syndrome} Chronic pain caused by damage or dysfunction in the brain or spinal cord. Pain may be burning, aching, electric, or unusually sensitive to touch and may occur after stroke, spinal-cord injury, multiple sclerosis, or other lesions.

\medterm{Cerebrovascular Disease} A group of disorders affecting blood vessels and blood flow in the brain, including stroke, transient ischemic attack, aneurysm, and vascular malformations.

\medterm{Cervical Ectropion} A common condition in which cells normally lining the cervical canal extend onto the outer cervix. It may cause increased discharge or bleeding after sex, but it is usually benign and is not the same as cervical cancer.

\medterm{Char Syndrome} A rare genetic disorder, usually caused by variants in \textit{TFAP2B}, characterized by distinctive facial features, a heart defect such as patent ductus arteriosus, and abnormalities of the fifth fingers or toes.

\medterm{Charles Bonnet Syndrome} Visual hallucinations occurring in a person with significant vision loss who otherwise understands that the images are not real. The images may be simple patterns or detailed scenes; new hallucinations still require medical assessment.

\medterm{CHILD Syndrome} Congenital hemidysplasia with ichthyosiform nevus and limb defects, a rare disorder usually caused by an \textit{NSDHL} variant. It causes one-sided skin abnormalities with limb, skeletal, and sometimes neurologic or internal-organ differences.

\medterm{Chimeric Antigen Receptor} A synthetic receptor engineered to recognize a selected target on a cell surface. It is placed into immune cells, especially T cells, to direct them against cancer cells in some cellular immunotherapies.

\medterm{Chinese Herbal Medicine} The use of plant, mineral, or animal-derived preparations within traditional Chinese medicine. Products vary in composition and may interact with medicines or contain contaminants, so reliable sourcing and medical guidance are important.

\medterm{CHIP} Clonal hematopoiesis of indeterminate potential: an age-related finding in which blood-forming cells acquire a gene change and expand without meeting criteria for a blood cancer. It can increase the future risk of myeloid cancer and cardiovascular disease and may require monitoring.

\medterm{Choroidal Melanoma} A cancer arising from pigment-producing melanocytes in the choroid, the vascular layer beneath the retina. It may cause visual changes or be found incidentally and can spread, especially to the liver; treatment depends on size and location.

\medterm{Chromosome 1} The largest human chromosome, containing many genes involved in development, cell function, and disease. Extra, missing, or rearranged material can cause genetic disorders.

\medterm{Chromosome 13} A human chromosome containing genes important for development and cell function. Abnormal chromosome 13 number or structure can cause disorders such as trisomy 13.

\medterm{Chromosome 2} A large human chromosome containing many genes involved in growth, development, and normal body function. Changes in its number or structure can cause genetic disease.

\medterm{Chromosome 3} A human chromosome containing genes involved in development, immune function, and cell regulation. Deletions, duplications, or other changes may cause genetic disorders or cancer.

\medterm{Chromosome 4} A human chromosome containing many genes involved in development and cellular function. Structural or numerical abnormalities can cause congenital or inherited conditions.

\medterm{Chromosome 5} A human chromosome containing genes involved in development, immunity, and cell regulation. Loss of part of its short arm causes the genetic condition cri-du-chat syndrome.

\medterm{Chromosome X} One of the human sex chromosomes. Most females have two X chromosomes and most males have one X and one Y; genes on the X chromosome can cause X-linked disorders when altered.

\medterm{Chromosome Y} The sex chromosome typically present in most males and absent in most females. It contains the sex-determining \textit{SRY} gene and other genes important for testis development and sperm production.

\medterm{Chronic Idiopathic Constipation} Long-lasting difficult or infrequent bowel movements without an identified structural or metabolic cause. Symptoms may include hard stools, straining, incomplete emptying, or abdominal discomfort; treatment includes diet, fluids, activity, and selected medicines.

\medterm{Chronic Obstructive Pulmonary Disorder (COPD)} A chronic lung disease with persistent airflow limitation, usually involving emphysema, chronic bronchitis, or both. Smoking and other inhaled exposures are important causes; symptoms include breathlessness, cough, and sputum.

\medterm{Circulating Tumor Cell} A cancer cell that has detached from a tumor and entered the bloodstream. Testing may help study cancer spread or treatment response, but clinical use depends on the specific validated test and cancer.

\medterm{Climate Change} Long-term changes in temperature, precipitation, and weather patterns, largely driven by human-produced greenhouse gases. Health effects include heat illness, worsened air quality, infectious-disease changes, injuries, and effects on food, water, and mental health.

\medterm{Clinical Diagnostics} The use of history, examination, laboratory tests, imaging, pathology, and other methods to identify or evaluate a health condition and guide care.

\medterm{Clinical Imaging} Medical imaging used to visualize anatomy, physiology, or disease, including radiography, ultrasound, computed tomography, magnetic resonance imaging, and nuclear medicine.

\medterm{Club Foot} A congenital foot deformity in which the foot is turned inward and downward and the calf may be smaller. It is usually treated early with serial casting, bracing, and sometimes surgery.

\medterm{Coffee} A beverage containing caffeine and other compounds. Moderate intake is tolerated by many adults, but effects vary; excessive caffeine can cause insomnia, anxiety, palpitations, or gastrointestinal symptoms.

\medterm{Cognitive Behavioural Therapy} A structured psychological treatment that helps a person identify and change unhelpful thoughts and behaviors. It is used for conditions such as anxiety, depression, insomnia, pain, and some eating disorders.

\medterm{Colony-Stimulating Factors} Growth factors that stimulate bone marrow to produce specific blood cells. Medicines such as granulocyte colony-stimulating factor can prevent or treat some forms of chemotherapy-related neutropenia.

\medterm{Colorectal Polyp} A growth arising from the lining of the colon or rectum. Many are benign, but some adenomatous or serrated polyps can become cancerous over time and are removed during colonoscopy.

\medterm{Comet Assay} A laboratory technique that measures DNA strand breaks in individual cells by electrophoresis, producing a comet-like pattern. It is mainly used in research and toxicology rather than as a routine clinical test.

\medterm{Compression Bandage} A bandage that applies controlled pressure to a limb or body part. It may reduce swelling, support injured tissue, or help manage venous disease, but improper use can impair circulation or damage skin.

\medterm{Congenital Birth Defect} A structural or functional abnormality present before or at birth. Causes may include genetic changes, infections, nutritional deficiencies, medicines, environmental exposures, or unknown factors.

\medterm{Congenital Heart Defect} A structural abnormality of the heart or major vessels present at birth. It may alter blood flow and range from a small defect needing monitoring to a serious condition requiring urgent treatment or surgery.

\medterm{Conjugated Linoleic Acid} A group of related fats found mainly in meat and dairy products and sold as supplements. Claims about weight loss or muscle improvement are inconsistent, and supplements may adversely affect blood sugar or liver function.

\medterm{Contact Lenses} Thin lenses placed directly on the eye to correct vision or deliver selected treatment. Proper fitting, handwashing, cleaning, and replacement are essential because misuse can cause corneal infection, inflammation, or vision-threatening injury.

\medterm{Coronavirus Disease COVID-19} Illness caused by SARS-CoV-2 infection. It may cause no symptoms, a mild respiratory illness, pneumonia, or severe disease, and some people develop symptoms that persist after the initial infection.

\medterm{Cosmetic Surgery} Surgery intended primarily to change or improve appearance rather than treat disease or restore function. It still carries risks such as infection, bleeding, anesthesia complications, scarring, and unsatisfactory results.

\medterm{Counterfeit Drug} A product falsely labeled as an approved medicine or made to imitate one. It may contain the wrong ingredient, too little or too much active drug, contaminants, or no active drug and can cause treatment failure or poisoning.

\medterm{Cow's Milk Allergy} An immune reaction to proteins in cow's milk, most common in infants and young children. It may cause hives, vomiting, wheezing, eczema, diarrhea, or anaphylaxis and is different from lactose intolerance.

\medterm{Cracked Heels} Dry, thickened skin at the heel that develops fissures, sometimes deep enough to bleed or become infected. Moisturizers, gentle filing, and treatment of contributing skin disease can help; painful or infected cracks need assessment.

\medterm{Cramp} A sudden, involuntary, often painful muscle contraction. Cramps may result from exertion, dehydration, electrolyte changes, pregnancy, medicines, nerve disease, or reduced blood flow.

\medterm{Cranberries} Red fruits containing fiber and plant compounds. Cranberry products may modestly reduce recurrent urinary infections in some people but do not treat an established infection; they can interact with warfarin and other medicines.

\medterm{Craniofacial} Relating to the skull and face. Craniofacial conditions may involve bones, soft tissues, teeth, eyes, ears, or airway and may be congenital, traumatic, developmental, or acquired.

\medterm{Craniofacial Microsomia} A birth condition in which one side of the face develops less fully than the other. It may affect the jaw, ear, eye, facial muscles, and nearby tissues, with severity ranging from mild to complex.

\medterm{CRIA Syndrome} A rare inherited inflammatory disorder caused by a change in the RIPK1 gene. Repeated attacks of fever, swollen lymph nodes, muscle or joint pain, and inflammation often begin during childhood.

\medterm{CRISPR} A gene-editing method that uses a guide molecule and an enzyme to find and change a chosen DNA sequence. It is used in research and some treatments but can make unintended changes and requires careful monitoring.

\medterm{Cruciate Ligament} One of two strong bands crossing inside the knee that stabilize forward, backward, and twisting movement. The anterior cruciate ligament is commonly torn during sports or a sudden change in direction.

\medterm{CTLA-4} A protein on some immune cells that acts as a brake on T-cell activity. Blocking this brake can strengthen an immune attack against cancer but may also cause inflammation in healthy organs.

\medterm{Cupping Therapy} A complementary practice that places cups on the skin to create suction, sometimes after small skin cuts. It may leave bruises or marks and can cause burns, infection, or blood loss; evidence of benefit varies by condition.

\medterm{Curcumin} A yellow compound in turmeric that has anti-inflammatory effects in laboratory studies. Evidence that supplements treat disease is limited or variable, and concentrated products can interact with medicines or cause stomach or liver problems.

\medterm{Cutaneous Vasculitis} Inflammation of small blood vessels in the skin, causing red or purple spots, raised purpura, blisters, sores, or pain. Medicines, infections, autoimmune disease, or inflammation elsewhere in the body may be responsible.

\medterm{Cyclin-Dependent Kinase} An enzyme that helps control when a cell grows and divides after joining with a cyclin protein. Abnormal activity can promote cancer and is blocked by some targeted treatments.

\medterm{Cystinosis} A rare inherited disorder in which cystine builds up inside cells and damages organs. It often affects the kidneys and eyes and may also affect growth, thyroid, muscles, or other tissues; lifelong specialist care is needed.
